% Generated mechanically from the frozen flat.tex target.
% Do not edit this generated file; edit the bound locale source instead.

% OK, start here.
%

\setcounter{chapter}{37}
\chapter{平坦性续篇}


\phantomsection
\label{flat-section-phantom}




\section{引言}
\label{flat-section-introduction}

\noindent
本章讨论平坦模与概形平坦态射的一些高阶结果及其应用。关于平坦性的多数结果
可见 Raynaud 与 Gruson 的论文 \cite{GruRay}。

\medskip\noindent
在阅读本章之前，建议读者先查看下列结果（此表也可作为先前结果的索引）：
\begin{enumerate}
\item 关于平坦模的一般讨论，见《代数》第 \ref{algebra-section-flat} 节。
\item $\text{Tor}$ 群与平坦性之间的关系，见《代数》第
\ref{algebra-section-tor} 节。
\item 平坦性判据，见《代数》第 \ref{algebra-section-criteria-flatness} 节
（诺特情形）、《代数》第 \ref{algebra-section-flatness-artinian} 节
（阿廷情形）、《代数》第 \ref{algebra-section-more-flatness-criteria} 节
（非诺特情形），最后见《态射进阶》第
\ref{more-morphisms-section-criterion-flat-fibres} 节。
\item 一般平坦性，见《代数》第 \ref{algebra-section-generic-flatness} 节与
《态射》第 \ref{morphisms-section-generic-flatness} 节。
\item 平坦轨迹的开性，见《代数》第 \ref{algebra-section-open-flat} 节与
《态射进阶》第 \ref{more-morphisms-section-open-flat} 节。
\item 平坦化，见《代数进阶》诸节
\ref{more-algebra-section-flattening},
\ref{more-algebra-section-flattening-artinian},
\ref{more-algebra-section-flattening-local-base},
\ref{more-algebra-section-flattening-local-source-base} 以及
\ref{more-algebra-section-flattening-Noetherian-complete-local}.
\item 其他结果见《代数进阶》诸节
\ref{more-algebra-section-descent-flatness-integral}、
\ref{more-algebra-section-torsion-flat}、
\ref{more-algebra-section-flat-finite} 以及
\ref{more-algebra-section-blowup-flat}.
\end{enumerate}
作为平坦性材料的应用，我们讨论以下主题：Grothendieck 存在定理的非诺特版本、
爆破与平坦性、Nagata 的紧化定理、h 拓扑、爆破方块与下降、弱正规化、
正特征下向量丛的下降，以及 h 拓扑中完美复形的局部结构。






\section{关于 \'etale 局部化的引理}
\label{flat-section-etale-localization}

\noindent
本节列出一些关于 \'etale 局部化的引理，它们将在本章后文中使用。
初次阅读时可以跳过本节。

\begin{lemma}
\label{flat-lemma-lift-etale}
设 $i : Z \to X$ 是仿射概形的闭浸入，$Z' \to Z$ 是 \'etale 态射，
且 $Z'$ 仿射。则存在 \'etale 态射 $X' \to X$，其中 $X'$ 仿射，
使得作为 $Z$ 上概形有 $Z' \cong Z \times_X X'$。
\end{lemma}

\begin{proof}
见《代数》引理 \ref{algebra-lemma-lift-etale}。
\end{proof}

\begin{lemma}
\label{flat-lemma-etale-at-point}
设
$$
\xymatrix{
X \ar[d] & X' \ar[l] \ar[d] \\
S & S' \ar[l]
}
$$
是概形的交换图，其中 $X' \to X$ 与 $S' \to S$ 均为 \'etale 态射。
设 $s' \in S'$ 为一点。则
$$
X' \times_{S'} \Spec(\mathcal{O}_{S', s'})
\longrightarrow
X \times_S \Spec(\mathcal{O}_{S', s'})
$$
是 \'etale 态射。
\end{lemma}

\begin{proof}
这是因为，作为 $X$ 上 \'etale 概形之间的态射，$X' \to X_{S'}$
是 \'etale 的；见《态射》引理 \ref{morphisms-lemma-etale-permanence}。
而且 \'etale 态射的基变换仍为 \'etale；见《态射》引理
\ref{morphisms-lemma-base-change-etale}。
\end{proof}

\begin{lemma}
\label{flat-lemma-etale-flat-up-down}
设 $X \to T \to S$ 是概形态射，且 $T \to S$ 为 \'etale 态射。
设 $\mathcal{F}$ 是拟凝聚 $\mathcal{O}_X$-模，$x \in X$ 为一点。则
$$
\mathcal{F}\text{ 在 }S\text{ 上于 }x\text{ 处平坦}
\Leftrightarrow
\mathcal{F}\text{ 在 }T\text{ 上于 }x\text{ 处平坦}
$$
特别地，$\mathcal{F}$ 在 $S$ 上平坦，当且仅当 $\mathcal{F}$ 在 $T$ 上平坦。
\end{lemma}

\begin{proof}
\'Etale 态射是平坦态射（见《态射》引理
\ref{morphisms-lemma-etale-flat}），故蕴含 ``$\Leftarrow$'' 由《代数》引理
\ref{algebra-lemma-composition-flat} 得出。反之，假设 $\mathcal{F}$
在点 $x$ 处于 $S$ 上平坦。记 $\tilde x \in X \times_S T$ 为这样一点：
它在 $X$ 中位于 $x$ 上方，在 $T$ 中位于 $x$ 在 $T$ 中的像上方。于是经由
$\text{pr}_2 : X \times_S T \to T$，
$(X \times_S T \to X)^*\mathcal{F}$ 在点 $\tilde x$ 处于 $T$ 上平坦；
见《态射》引理 \ref{morphisms-lemma-base-change-module-flat}。
对角态射 $\Delta_{T/S} : T \to T \times_S T$ 是开浸入；合并《态射》引理
\ref{morphisms-lemma-diagonal-unramified-morphism} 与
\ref{morphisms-lemma-etale-smooth-unramified} 即得。因此可将 $X$
等同于 $X \times_S T$ 的开子概形；$\text{pr}_2$ 在此开子概形上的限制
就是给定态射 $X \to T$，点 $\tilde x$ 对应于其中的点 $x$，并且
$(X \times_S T \to X)^*\mathcal{F}$ 在该开子概形上的限制就是
$\mathcal{F}$。由此可见，$\mathcal{F}$ 在点 $x$ 处于 $T$ 上平坦。
\end{proof}

\begin{lemma}
\label{flat-lemma-etale-flat-up-down-local-ring}
设 $T \to S$ 是 \'etale 态射，$t \in T$ 的像为 $s \in S$。
设 $M$ 是 $\mathcal{O}_{T, t}$-模。则
$$
M\text{ 在 }\mathcal{O}_{S, s}\text{ 上平坦}
\Leftrightarrow
M\text{ 在 }\mathcal{O}_{T, t}\text{ 上平坦}.
$$
\end{lemma}

\begin{proof}
可先用 $s$ 的仿射邻域替换 $S$，再用 $t$ 的仿射邻域替换 $T$。
置 $\mathcal{F} = (\Spec(\mathcal{O}_{T, t}) \to T)_*\widetilde M$。
这是 $T$ 上的拟凝聚层（见概形章引理
\ref{schemes-lemma-push-forward-quasi-coherent}，或直接论证），
其在 $t$ 处的茎为 $M$（细节从略）。应用引理
\ref{flat-lemma-etale-flat-up-down} 即得。
\end{proof}

\begin{lemma}
\label{flat-lemma-flat-up-down-henselization}
设 $S$ 是概形，$s \in S$ 为一点。记 $\mathcal{O}_{S, s}^h$
（分别地，$\mathcal{O}_{S, s}^{sh}$）为亨泽尔化（分别地，严格亨泽尔化）；
见《代数》定义 \ref{algebra-definition-henselization}。
设 $M^{sh}$ 是 $\mathcal{O}_{S, s}^{sh}$-模。下列条件等价：
\begin{enumerate}
\item $M^{sh}$ 在 $\mathcal{O}_{S, s}$ 上平坦；
\item $M^{sh}$ 在 $\mathcal{O}_{S, s}^h$ 上平坦；
\item $M^{sh}$ 在 $\mathcal{O}_{S, s}^{sh}$ 上平坦。
\end{enumerate}
若 $M^{sh} = M^h \otimes_{\mathcal{O}_{S, s}^h} \mathcal{O}_{S, s}^{sh}$，
则上述条件还等价于：
\begin{enumerate}
\item[(4)] $M^h$ 在 $\mathcal{O}_{S, s}$ 上平坦；
\item[(5)] $M^h$ 在 $\mathcal{O}_{S, s}^h$ 上平坦。
\end{enumerate}
若 $M^h = M \otimes_{\mathcal{O}_{S, s}} \mathcal{O}_{S, s}^h$，
则上述条件还等价于：
\begin{enumerate}
\item[(6)] $M$ 在 $\mathcal{O}_{S, s}$ 上平坦。
\end{enumerate}
\end{lemma}

\begin{proof}
由《代数进阶》引理 \ref{more-algebra-lemma-dumb-properties-henselization}，局部环同态
$\mathcal{O}_{S, s} \to \mathcal{O}_{S, s}^h \to \mathcal{O}_{S, s}^{sh}$
是忠实平坦的。因此 (3) $\Rightarrow$ (2) $\Rightarrow$ (1) 与
(5) $\Rightarrow$ (4) 由《代数》引理
\ref{algebra-lemma-composition-flat} 得出。由忠实平坦性，等价关系
(6) $\Leftrightarrow$ (5) 与 (5) $\Leftrightarrow$ (3) 由《代数》引理
\ref{algebra-lemma-flatness-descends} 得出。因此只需证明
(1) $\Rightarrow$ (2) $\Rightarrow$ (3) 与 (4) $\Rightarrow$ (5)。
为此可假设 $S$ 是仿射概形。

\medskip\noindent
假设 (1) 成立。由引理 \ref{flat-lemma-etale-flat-up-down-local-ring}，
对任意 \'etale 邻域 $(T, t) \to (S, s)$，$M^{sh}$ 在
$\mathcal{O}_{T, t}$ 上平坦。由于 $\mathcal{O}_{S, s}^h$ 与
$\mathcal{O}_{S, s}^{sh}$ 是形如 $\mathcal{O}_{T, t}$ 的局部环的有向余极限
（见《代数》引理 \ref{algebra-lemma-henselization-different} 与
\ref{algebra-lemma-strict-henselization-different}），由《代数》引理
\ref{algebra-lemma-colimit-rings-flat}，$M^{sh}$ 在
$\mathcal{O}_{S, s}^h$ 与 $\mathcal{O}_{S, s}^{sh}$ 上平坦。
故 (1) 蕴含 (2) 与 (3)。将 $\mathcal{O}_{S, s}$ 替换为
$\mathcal{O}_{S, s}^h$，当然也得到 (2) $\Rightarrow$ (3)。
同样的论证证明 (4) $\Rightarrow$ (5)。
\end{proof}

\begin{lemma}
\label{flat-lemma-tor-amplitude-up-down-henselization}
设 $S$ 是概形，$s \in S$ 为一点。记 $\mathcal{O}_{S, s}^h$
（分别地，$\mathcal{O}_{S, s}^{sh}$）为亨泽尔化（分别地，严格亨泽尔化）；
见《代数》定义 \ref{algebra-definition-henselization}。
设 $M^{sh}$ 是 $D(\mathcal{O}_{S, s}^{sh})$ 的对象，
$a, b \in \mathbf{Z}$。下列条件等价：
\begin{enumerate}
\item $M^{sh}$ 在 $\mathcal{O}_{S, s}$ 上的 Tor 振幅包含于 $[a, b]$；
\item $M^{sh}$ 在 $\mathcal{O}_{S, s}^h$ 上的 Tor 振幅包含于 $[a, b]$；
\item $M^{sh}$ 在 $\mathcal{O}_{S, s}^{sh}$ 上的 Tor 振幅包含于 $[a, b]$。
\end{enumerate}
若 $M^{sh} =
M^h \otimes_{\mathcal{O}_{S, s}^h}^\mathbf{L} \mathcal{O}_{S, s}^{sh}$
其中 $M^h \in D(\mathcal{O}_{S, s}^h)$，则上述条件还等价于：
\begin{enumerate}
\item[(4)] $M^h$ 在 $\mathcal{O}_{S, s}$ 上的 Tor 振幅包含于 $[a, b]$；
\item[(5)] $M^h$ 在 $\mathcal{O}_{S, s}^h$ 上的 Tor 振幅包含于 $[a, b]$。
\end{enumerate}
若 $M^h = M \otimes_{\mathcal{O}_{S, s}}^\mathbf{L} \mathcal{O}_{S, s}^h$
其中 $M \in D(\mathcal{O}_{S, s})$，则上述条件还等价于：
\begin{enumerate}
\item[(6)] $M$ 在 $\mathcal{O}_{S, s}$ 上的 Tor 振幅包含于 $[a, b]$。
\end{enumerate}
\end{lemma}

\begin{proof}
由《代数进阶》引理 \ref{more-algebra-lemma-dumb-properties-henselization}，局部环同态
$\mathcal{O}_{S, s} \to \mathcal{O}_{S, s}^h \to \mathcal{O}_{S, s}^{sh}$
是忠实平坦的。因此 (3) $\Rightarrow$ (2) $\Rightarrow$ (1) 与
(5) $\Rightarrow$ (4) 由《代数进阶》引理
\ref{more-algebra-lemma-flat-push-tor-amplitude} 得出。由忠实平坦性，
等价关系 (6) $\Leftrightarrow$ (5) 与 (5) $\Leftrightarrow$ (3) 由
《代数进阶》引理 \ref{more-algebra-lemma-flat-descent-tor-amplitude} 得出。
因此只需证明 (1) $\Rightarrow$ (3)、(2) $\Rightarrow$ (3) 与
(4) $\Rightarrow$ (5)。

\medskip\noindent
假设 (1) 成立。特别地，$M^{sh}$ 在次数 $< a$ 与 $> b$ 上的上同调为零。
故可用自由 $\mathcal{O}_{X, x}^{sh}$-模复形 $P^\bullet$ 表示
$M^{sh}$，且当 $i > b$ 时 $P^i = 0$（例如见一般性很强的导出范畴章引理
\ref{derived-lemma-subcategory-left-resolution}）。注意，对所有 $n$，
$P^n$ 在 $\mathcal{O}_{S, s}$ 上平坦。考虑 $\Coker(d_P^{a - 1})$。
由《代数进阶》引理 \ref{more-algebra-lemma-last-one-flat}，这是平坦
$\mathcal{O}_{S, s}$-模。故由引理 \ref{flat-lemma-flat-up-down-henselization}，
它也是平坦 $\mathcal{O}_{S, s}^{sh}$-模。因此
$\tau_{\geq a}P^\bullet$ 是平坦 $\mathcal{O}_{S, s}^{sh}$-模复形，
它在 $D(\mathcal{O}_{S, s}^{sh}$ 中表示 $M^{sh}$，于是可知
$M^{sh}$ 的 Tor 振幅包含于 $[a, b]$；见《代数进阶》引理
\ref{more-algebra-lemma-tor-amplitude}。故 (1) 蕴含 (3)。将
$\mathcal{O}_{S, s}$ 替换为 $\mathcal{O}_{S, s}^h$，当然也得到
(2) $\Rightarrow$ (3)。同样的论证证明 (4) $\Rightarrow$ (5)。
\end{proof}

\begin{lemma}
\label{flat-lemma-finite-flat-weak-assassin-up-down}
设 $g : T \to S$ 是概形的有限平坦态射，$\mathcal{G}$ 是拟凝聚
$\mathcal{O}_S$-模。设 $t \in T$ 的像为 $s \in S$。则
$$
t \in \text{WeakAss}(g^*\mathcal{G})
\Leftrightarrow
s \in \text{WeakAss}(\mathcal{G})
$$
\end{lemma}

\begin{proof}
蕴含 ``$\Leftarrow$'' 立即由《因子》引理
\ref{divisors-lemma-weakly-ass-pullback} 得出。假设
$t \in \text{WeakAss}(g^*\mathcal{G})$。设
$\Spec(A) \subset S$ 是 $s$ 的仿射开邻域，且 $\mathcal{G}$ 是与
$A$-模 $M$ 相伴的拟凝聚层。设 $\mathfrak p \subset A$ 是与 $s$
对应的素理想。由于 $g$ 有限平坦，对某个有限平坦 $A$-代数 $B$，有
$g^{-1}(\Spec(A)) = \Spec(B)$。注意，$g^*\mathcal{G}$ 是与
$B$-模 $M \otimes_A B$ 相伴的拟凝聚 $\mathcal{O}_{\Spec(B)}$-模，
而 $g_*g^*\mathcal{G}$ 是与 $A$-模 $M \otimes_A B$ 相伴的拟凝聚
$\mathcal{O}_{\Spec(A)}$-模。由《代数》引理
\ref{algebra-lemma-finite-flat-local}，对某个整数 $n \geq 0$，有
$B_{\mathfrak p} \cong A_{\mathfrak p}^{\oplus n}$。由于假设至少存在
一个位于 $s$ 上方的 $T$ 中点，故 $n \geq 1$。于是考察茎可知
$$
s \in \text{WeakAss}(\mathcal{G})
\Leftrightarrow
s \in \text{WeakAss}(g_*g^*\mathcal{G})
$$
现在，由《因子》引理 \ref{divisors-lemma-weakly-associated-finite}，
假设 $t \in \text{WeakAss}(g^*\mathcal{G})$ 蕴含
$s \in \text{WeakAss}(g_*g^*\mathcal{G})$，故由上式有
$s \in  \text{WeakAss}(\mathcal{G})$。
\end{proof}

\begin{lemma}
\label{flat-lemma-etale-weak-assassin-up-down}
设 $h : U \to S$ 是概形的 \'etale 态射，$\mathcal{G}$ 是拟凝聚
$\mathcal{O}_S$-模。设 $u \in U$ 的像为 $s \in S$。则
$$
u \in \text{WeakAss}(h^*\mathcal{G})
\Leftrightarrow
s \in \text{WeakAss}(\mathcal{G})
$$
\end{lemma}

\begin{proof}
用 $s$ 与 $u$ 的仿射邻域分别替换 $S$ 与 $U$ 后，可假设 $g$
是仿射概形的标准 \'etale 态射；见《态射》引理
\ref{morphisms-lemma-etale-locally-standard-etale}。
% P05-EN-CH38-0003: frozen proof names g and X although the lemma's morphism and source are h and U.
因此可假设 $S = \Spec(A)$ 且 $X = \Spec(A[x, 1/g]/(f))$，其中 $f$
首一，且 $f'$ 在 $A[x, 1/g]$ 中可逆。注意，
$A[x, 1/g]/(f) = (A[x]/(f))_g$ 也是有限自由 $A$-代数
$A[x]/(f)$ 的局部化。因此可将 $U$ 看作概形
$T = \Spec(A[x]/(f))$ 的开子概形，而后者在 $S$ 上有限局部自由。
这就归结到上面的引理 \ref{flat-lemma-finite-flat-weak-assassin-up-down}。
\end{proof}

\begin{lemma}
\label{flat-lemma-weakly-associated-henselization}
设 $S$ 是概形，$s \in S$ 为一点。记 $\mathcal{O}_{S, s}^h$
（分别地，$\mathcal{O}_{S, s}^{sh}$）为亨泽尔化（分别地，严格亨泽尔化）；
见《代数》定义 \ref{algebra-definition-henselization}。
设 $\mathcal{F}$ 是拟凝聚 $\mathcal{O}_S$-模。下列条件等价：
\begin{enumerate}
\item $s$ 是 $\mathcal{F}$ 的弱伴随点；
\item $\mathfrak m_s$ 是 $\mathcal{F}_s$ 的弱伴随素理想；
\item $\mathfrak m_s^h$ 是
$\mathcal{F}_s \otimes_{\mathcal{O}_{S, s}} \mathcal{O}_{S, s}^h$
的弱伴随素理想；
\item $\mathfrak m_s^{sh}$ 是
$\mathcal{F}_s \otimes_{\mathcal{O}_{S, s}} \mathcal{O}_{S, s}^{sh}$
的弱伴随素理想。
\end{enumerate}
\end{lemma}

\begin{proof}
(1) 与 (2) 的等价性就是定义；见《因子》定义
\ref{divisors-definition-weakly-associated}。将《因子》引理
\ref{divisors-lemma-weakly-ass-pullback} 应用于下列平坦态射
（见《代数进阶》引理 \ref{more-algebra-lemma-dumb-properties-henselization}）：
$$
\Spec(\mathcal{O}_{S, s}^{sh}) \to
\Spec(\mathcal{O}_{S, s}^h) \to
\Spec(\mathcal{O}_{S, s})
$$
以及闭点，即得 (2) $\Rightarrow$ (3) $\Rightarrow$ (4)。
为证明 (4) $\Rightarrow$ (2)，可用一个仿射邻域替换 $S$。假设
$x \in \mathcal{F}_s \otimes_{\mathcal{O}_{S, s}} \mathcal{O}_{S, s}^{sh}$
是这样一个元素：其零化子之根基等于 $\mathfrak m_s^{sh}$
（见《代数》引理 \ref{algebra-lemma-weakly-ass-local}）。由《代数》引理
\ref{algebra-lemma-strict-henselization-different}，
$\mathcal{O}_{S, s}^{sh}$ 等于当 \'etale 邻域
$f : (U, u) \to (S, s)$ 变化时 $\mathcal{O}_{U, u}$ 的极限，
故可假设 $x$ 是某个
$x' \in \mathcal{F}_s \otimes_{\mathcal{O}_{S, s}} \mathcal{O}_{U, u}$
的像。局部环同态 $\mathcal{O}_{U, u} \to \mathcal{O}_{S, s}^{sh}$
忠实平坦（因为这是严格亨泽尔化），因而普遍单射（《代数》引理
\ref{algebra-lemma-faithfully-flat-universally-injective}）。于是 $x'$
的零化子是 $x$ 的零化子的逆像，故该零化子之根基等于
$\mathfrak m_u$。因此 $u$ 是 $f^*\mathcal{F}$ 的弱伴随点。
由引理 \ref{flat-lemma-etale-weak-assassin-up-down}，$s$ 是
$\mathcal{F}$ 的弱伴随点。
\end{proof}





\section{有限型模的局部结构}
\label{flat-section-local-structure-module}

\noindent
使本章许多论证得以成立的关键技术结果，是几何引理
\ref{flat-lemma-elementary-devissage}。

\begin{lemma}
\label{flat-lemma-sheaf-lives-on-subscheme}
设 $f : X \to S$ 是仿射概形的有限型态射，$\mathcal{F}$ 是有限型
拟凝聚 $\mathcal{O}_X$-模。设 $x \in X$ 在 $S$ 中的像为
$s = f(x)$，并置 $\mathcal{F}_s = \mathcal{F}|_{X_s}$。
则存在有限表示闭浸入 $i : Z \to X$ 与有限型拟凝聚
$\mathcal{O}_Z$-模 $\mathcal{G}$，使得 $i_*\mathcal{G} = \mathcal{F}$ 且
$Z_s = \text{Supp}(\mathcal{F}_s)$。
\end{lemma}

\begin{proof}
设态射 $f : X \to S$ 由环同态 $A \to B$ 给出，并设 $\mathcal{F}$
是与 $B$-模 $M$ 相伴的拟凝聚层。由《态射》引理
\ref{morphisms-lemma-locally-finite-type-characterize}，$A \to B$
是有限型环同态；由《性质》引理 \ref{properties-lemma-finite-type-module}，
$M$ 是有限 $B$-模。特别地，$\mathcal{F}$ 的支集是 $\Spec(B)$ 中由
$M$ 的零化子 $I = \{x \in B \mid xm = 0\ \forall m \in M\}$
截出的闭子概形；见《代数》引理 \ref{algebra-lemma-support-closed}。
设 $\mathfrak q \subset B$ 是与 $x$ 对应的素理想，
$\mathfrak p \subset A$ 是与 $s$ 对应的素理想。注意，
$X_s = \Spec(B \otimes_A \kappa(\mathfrak p))$，而 $\mathcal{F}_s$
是与 $B \otimes_A \kappa(\mathfrak p)$-模
$M \otimes_A \kappa(\mathfrak p)$ 相伴的拟凝聚层。由《态射》引理
\ref{morphisms-lemma-support-finite-type}，$\mathcal{F}_s$ 的支集等于
$V(I(B \otimes_A \kappa(\mathfrak p)))$。由于
$B \otimes_A \kappa(\mathfrak p)$ 在 $\kappa(\mathfrak p)$ 上有限型，
存在有限多个元素 $f_1, \ldots, f_m \in I$，使得
$$
I(B \otimes_A \kappa(\mathfrak p)) =
(f_1, \ldots, f_n)(B \otimes_A \kappa(\mathfrak p)).
$$
% P05-EN-CH38-0004: frozen display uses f_1,...,f_n after choosing f_1,...,f_m.
记 $i : Z \to X$ 为由 $(f_1, \ldots, f_m)$ 截出的闭子概形，即
$Z = \Spec(B/(f_1, \ldots, f_m))$。由于 $I$ 零化 $M$，可将 $M$
视为 $B/(f_1, \ldots, f_m)$-模。换言之，$\mathcal{F}$ 是 $Z$
上某个有限型模的推前。按构造有
$Z_s = \text{Supp}(\mathcal{F}_s)$，故引理得证。
\end{proof}

\begin{lemma}
\label{flat-lemma-elementary-devissage}
设 $f : X \to S$ 是局部有限型概形态射，$\mathcal{F}$ 是有限型
拟凝聚 $\mathcal{O}_X$-模。设 $x \in X$ 在 $S$ 中的像为
$s = f(x)$。置 $\mathcal{F}_s = \mathcal{F}|_{X_s}$ 且
$n = \dim_x(\text{Supp}(\mathcal{F}_s))$。则可构造：
\begin{enumerate}
\item 初等 \'etale 邻域 $g : (X', x') \to (X, x)$ 与
$e : (S', s') \to (S, s)$；
\item 交换图
$$
\xymatrix{
X \ar[dd]_f & X' \ar[dd] \ar[l]^g & Z' \ar[l]^i \ar[d]^\pi \\
& & Y' \ar[d]^h \\
S & S' \ar[l]_e & S' \ar@{=}[l]
}
$$
\item 点 $z' \in Z'$，满足 $i(z') = x'$、$y' = \pi(z')$、
$h(y') = s'$；
\item 有限型拟凝聚 $\mathcal{O}_{Z'}$-模 $\mathcal{G}$；
\end{enumerate}
使得下列性质成立：
\begin{enumerate}
\item $X'$、$Z'$、$Y'$、$S'$ 都是仿射概形；
\item $i$ 是有限表示闭浸入；
\item $i_*(\mathcal{G}) \cong g^*\mathcal{F}$；
\item $\pi$ 有限，且 $\pi^{-1}(\{y'\}) = \{z'\}$；
\item 扩张 $\kappa(y')/\kappa(s')$ 是纯超越的；
\item $h$ 是相对维数为 $n$ 且纤维几何整的光滑态射。
\end{enumerate}
\end{lemma}

\begin{proof}
设 $V \subset S$ 是 $s$ 的仿射邻域，$U \subset f^{-1}(V)$ 是 $x$
的仿射邻域。只需对 $f|_U : U \to V$ 与 $\mathcal{F}|_U$ 证明引理。
因此在余下证明中假设 $X$ 与 $S$ 均仿射。

\medskip\noindent
首先假设 $X_s = \text{Supp}(\mathcal{F}_s)$；特别地，
$n = \dim_x(X_s)$。应用《态射进阶》引理
\ref{more-morphisms-lemma-local-local-structure-finite-type} 与
\ref{more-morphisms-lemma-local-local-structure-finite-type-affine}，
得到交换图
$$
\xymatrix{
X \ar[dd] & X' \ar[l]^g \ar[d]^\pi \\
& Y' \ar[d]^h  \\
S & S' \ar[l]_e
}
$$
以及点 $x' \in X'$。置 $Z' = X'$、$i = \text{id}$ 且
$\mathcal{G} = g^*\mathcal{F}$，便得到此情形的解。

\medskip\noindent
一般地，依引理 \ref{flat-lemma-sheaf-lives-on-subscheme} 选取闭浸入
$Z \to X$ 与 $Z$ 上的层 $\mathcal{G}$。将上一段的结果应用于
$Z \to S$ 与 $\mathcal{G}$，得到图表
$$
\xymatrix{
X \ar[dd]_f & Z \ar[l] \ar[dd]_{f|_Z} & Z' \ar[l]^g \ar[d]^\pi \\
& & Y' \ar[d]^h \\
S & S \ar@{=}[l] & S' \ar[l]_e
}
$$
以及满足全部所需性质的点 $z' \in Z'$。我们将用引理
\ref{flat-lemma-lift-etale} 把 $Z'$ 嵌入一个在 $X$ 上 \'etale 的概形中。
由于要求 $X'$ 是 $S'$ 上的概形，不能直接应用该引理。改为考虑态射
$$
\xymatrix{
Z' \ar[r] & Z \times_S S' \ar[r] & X \times_S S'
}
$$
由《态射》引理 \ref{morphisms-lemma-etale-permanence}，第一个态射为
\'etale；第二个态射是闭浸入的基变换，因而为闭浸入。最后，由于
$X$、$S$、$S'$、$Z$、$Z'$ 均仿射，可应用引理
\ref{flat-lemma-lift-etale}，得到仿射概形的 \'etale 态射
$X' \to X \times_S S'$，使得
$$
Z' = (Z \times_S S') \times_{(X \times_S S')} X' = Z \times_X X'.
$$
由于 $Z \to X$ 是有限表示闭浸入，$Z' \to X'$ 亦然。设
$x' \in X'$ 为与 $z' \in Z'$ 对应的点。于是补全后的图表
$$
\xymatrix{
X \ar[dd] & X' \ar[dd] \ar[l] & Z' \ar[l]^i \ar[d]^\pi \\
& & Y' \ar[d]^h \\
S & S' \ar[l]_e & S' \ar@{=}[l]
}
$$
就是原问题的一个解。
\end{proof}

\begin{lemma}
\label{flat-lemma-devissage-finite-presentation}
假设与记号同引理 \ref{flat-lemma-elementary-devissage}。若 $f$ 局部有限表示，
则 $\pi$ 有限表示。在此情形，下列条件等价：
\begin{enumerate}
\item $\mathcal{F}$ 在 $x$ 的某个邻域上是有限表示
$\mathcal{O}_X$-模；
\item $\mathcal{G}$ 在 $z'$ 的某个邻域上是有限表示
$\mathcal{O}_{Z'}$-模；
\item $\pi_*\mathcal{G}$ 在 $y'$ 的某个邻域上是有限表示
$\mathcal{O}_{Y'}$-模。
\end{enumerate}
仍假设 $f$ 局部有限表示，则下列条件也彼此等价：
\begin{enumerate}
\item[(a)] $\mathcal{F}_x$ 是有限表示 $\mathcal{O}_{X, x}$-模；
\item[(b)] $\mathcal{G}_{z'}$ 是有限表示 $\mathcal{O}_{Z', z'}$-模；
\item[(c)] $(\pi_*\mathcal{G})_{y'}$ 是有限表示
$\mathcal{O}_{Y', y'}$-模。
\end{enumerate}
\end{lemma}

\begin{proof}
假设 $f$ 局部有限表示。由于 $Z' \to S$ 是这类态射的复合，它也局部
有限表示；见《态射》引理
\ref{morphisms-lemma-composition-finite-presentation}。注意，
$Y' \to S$ 是光滑态射与 \'etale 态射的复合，因而也局部有限表示。
故《态射》引理 \ref{morphisms-lemma-finite-presentation-permanence}
蕴含 $\pi$ 局部有限表示。由于 $\pi$ 有限，它也分离且拟紧，
因而 $\pi$ 实际上有限表示。

\medskip\noindent
为证明 (1)、(2)、(3) 等价，还考虑：
(4) $g^*\mathcal{F}$ 在 $x'$ 的某个邻域上是有限表示
$\mathcal{O}_{X'}$-模。有限表示模的拉回仍有限表示；见模章引理
\ref{modules-lemma-pullback-finite-presentation}。故
(1) $\Rightarrow$ (4)。\'Etale 态射 $g$ 是开的；见《态射》引理
\ref{morphisms-lemma-etale-open}。因此，对 $x'$ 的任意开邻域
$U' \subset X'$，像 $g(U')$ 是 $x$ 的开邻域，且
$\{U' \to g(U')\}$ 是 \'etale 覆盖。于是由下降章引理
\ref{descent-lemma-finite-presentation-descends}，有
(4) $\Rightarrow$ (1)。利用下降章引理
\ref{descent-lemma-finite-finitely-presented-module} 与一些简单拓扑论证
（见《态射进阶》引理
\ref{more-morphisms-lemma-finite-morphism-single-point-in-fibre}），可知
(4) $\Leftrightarrow$ (2) $\Leftrightarrow$ (3)。

\medskip\noindent
为证明 (a)、(b)、(c) 等价，考虑环同态
$$
\mathcal{O}_{X, x} \to
\mathcal{O}_{X', x'} \to
\mathcal{O}_{Z', z'} \leftarrow
\mathcal{O}_{Y', y'}
$$
第一个环同态忠实平坦。因此 $\mathcal{F}_x$ 在
$\mathcal{O}_{X, x}$ 上有限表示，当且仅当 $g^*\mathcal{F}_{x'}$
在 $\mathcal{O}_{X', x'}$ 上有限表示；见《代数》引理
\ref{algebra-lemma-descend-properties-modules}。第二个环同态满射
（故有限），且按假设有限表示，因此 $g^*\mathcal{F}_{x'}$
在 $\mathcal{O}_{X', x'}$ 上有限表示，当且仅当 $\mathcal{G}_{z'}$
在 $\mathcal{O}_{Z', z'}$ 上有限表示；见《代数》引理
\ref{algebra-lemma-finite-finitely-presented-extension}。由于 $\pi$ 有限、
有限表示且 $\pi^{-1}(\{y'\}) = \{x'\}$，环同态
$\mathcal{O}_{Y', y'} \leftarrow \mathcal{O}_{Z', z'}$ 有限且有限表示；
见《态射进阶》引理
\ref{more-morphisms-lemma-finite-morphism-single-point-in-fibre}。
% P05-EN-CH38-0005: frozen proof writes fibre {x'} although pi has domain Z' and the distinguished point is z'.
因此 $\mathcal{G}_{z'}$ 在 $\mathcal{O}_{Z', z'}$ 上有限表示，
当且仅当 $\pi_*\mathcal{G}_{y'}$ 在 $\mathcal{O}_{Y', y'}$ 上有限表示；
见《代数》引理 \ref{algebra-lemma-finite-finitely-presented-extension}。
\end{proof}

\begin{lemma}
\label{flat-lemma-devissage-flat}
假设与记号同引理 \ref{flat-lemma-elementary-devissage}。下列条件等价：
\begin{enumerate}
\item $\mathcal{F}$ 在 $x$ 的某个邻域上于 $S$ 上平坦；
\item $\mathcal{G}$ 在 $z'$ 的某个邻域上于 $S'$ 上平坦；
\item $\pi_*\mathcal{G}$ 在 $y'$ 的某个邻域上于 $S'$ 上平坦。
\end{enumerate}
下列条件也等价：
\begin{enumerate}
\item[(a)] $\mathcal{F}_x$ 在 $\mathcal{O}_{S, s}$ 上平坦；
\item[(b)] $\mathcal{G}_{z'}$ 在 $\mathcal{O}_{S', s'}$ 上平坦；
\item[(c)] $(\pi_*\mathcal{G})_{y'}$ 在 $\mathcal{O}_{S', s'}$ 上平坦。
\end{enumerate}
\end{lemma}

\begin{proof}
为证明 (1)、(2)、(3) 等价，还考虑：
(4) $g^*\mathcal{F}$ 在 $x'$ 的某个邻域上于 $S$ 上平坦。
我们将使用引理 \ref{flat-lemma-etale-flat-up-down}，不再说明地把在 $S$
上平坦与在 $S'$ 上平坦等同起来。\'Etale 态射 $g$ 平坦且为开映射；
见《态射》引理 \ref{morphisms-lemma-etale-open}。因此，对 $x'$ 的任意
开邻域 $U' \subset X'$，像 $g(U')$ 是 $x$ 的开邻域，且态射
$U' \to g(U')$ 满且平坦。故由《态射》引理
\ref{morphisms-lemma-flat-permanence}，有 (4) $\Leftrightarrow$ (1)。注意
$$
\Gamma(X', g^*\mathcal{F}) =
\Gamma(Z', \mathcal{G}) =
\Gamma(Y', \pi_*\mathcal{G})
$$
因此，$g^*\mathcal{F}$、$\mathcal{G}$ 与 $\pi_*\mathcal{G}$ 在
$S'$ 上平坦三者彼此等价（这里使用了 $X'$、$Z'$、$Y'$、$S'$
均仿射）。再用关于仿射邻域的若干略去的拓扑论证（比较《态射进阶》引理
\ref{more-morphisms-lemma-finite-morphism-single-point-in-fibre}），便得到
(4) $\Leftrightarrow$ (2) $\Leftrightarrow$ (3)。

\medskip\noindent
为证明 (a)、(b)、(c) 等价，考虑局部环同态的交换图
$$
\xymatrix{
\mathcal{O}_{X', x'} \ar[r]_\iota &
\mathcal{O}_{Z', z'} &
\mathcal{O}_{Y', y'} \ar[l]^\alpha &
\mathcal{O}_{S', s'} \ar[l]^\beta \\
\mathcal{O}_{X, x} \ar[u]^\gamma & & &
\mathcal{O}_{S, s} \ar[lll]_\varphi \ar[u]_\epsilon
}
$$
我们将使用引理 \ref{flat-lemma-etale-flat-up-down-local-ring}，不再说明地把
在 $\mathcal{O}_{S, s}$ 上平坦与在 $\mathcal{O}_{S', s'}$ 上平坦
等同起来。同态 $\gamma$ 忠实平坦。因此 $\mathcal{F}_x$ 在
$\mathcal{O}_{S, s}$ 上平坦，当且仅当 $g^*\mathcal{F}_{x'}$
在 $\mathcal{O}_{S', s'}$ 上平坦；见《代数》引理
\ref{algebra-lemma-flatness-descends-more-general}。作为
$\mathcal{O}_{S', s'}$-模，$g^*\mathcal{F}_{x'}$、
$\mathcal{G}_{z'}$ 与 $\pi_*\mathcal{G}_{y'}$ 全都同构；见《态射进阶》引理
\ref{more-morphisms-lemma-finite-morphism-single-point-in-fibre}。
证明完毕。
\end{proof}









\section{一步 d\'evissage}
\label{flat-section-one-step-devissage}

\noindent
本节解释模的一步 d\'evissage 是什么。一步 d\'evissage 在底与靶上
\'etale 局部存在。我们讨论一步 d\'evissage 的基变换、Zariski 缩小与
\'etale 局部化。

\begin{definition}
\label{flat-definition-one-step-devissage}
设 $S$ 是概形，$X$ 在 $S$ 上局部有限型，$\mathcal{F}$ 是有限型
拟凝聚 $\mathcal{O}_X$-模，且 $s \in S$ 为一点。所谓
{\it $s$ 上 $\mathcal{F}/X/S$ 的一步 d\'evissage}，是指给定
$S$ 上概形的态射
$$
\xymatrix{
X & Z \ar[l]_i \ar[r]^\pi & Y
}
$$
以及有限型拟凝聚 $\mathcal{O}_Z$-模 $\mathcal{G}$，满足：
\begin{enumerate}
\item $X$、$S$、$Z$、$Y$ 均仿射；
\item $i$ 是有限表示闭浸入；
\item $\mathcal{F} \cong i_*\mathcal{G}$；
\item $\pi$ 有限；
\item 结构态射 $Y \to S$ 光滑，且其纤维几何不可约、维数为
$\dim(\text{Supp}(\mathcal{F}_s))$。
\end{enumerate}
此时称 $(Z, Y, i, \pi, \mathcal{G})$ 为 $s$ 上
$\mathcal{F}/X/S$ 的一步 d\'evissage。
\end{definition}

\noindent
注意，这样的一步 d\'evissage 只有在 $X$ 与 $S$ 仿射时才可能存在。
上述定义只要求 $X$ 在 $S$ 上（局部）有限型，下文仍在这一设定中工作。
在 \cite{GruRay} 中，作者始终采用如下设定：$X \to S$ 局部有限表示，
而 $\mathcal{F}$ 是有限型拟凝聚 $\mathcal{O}_X$-模。这样选择的好处是，
询问 $\mathcal{F}$ 作为 $\mathcal{O}_X$-模是否有限表示是“有意义”的，
而在我们的设定中则“没有意义”。有关讨论见《态射进阶》第
\ref{more-morphisms-section-finite-type-finite-presentation} 节；
其中的观察表明，在我们的设定下，可考虑 $\mathcal{F}$ “相对于 $S$
局部有限表示”这一条件，并可始终使用该概念。不过，我们将依靠引理
\ref{flat-lemma-devissage-finite-presentation} 的结果与注记
\ref{flat-remark-finite-presentation} 中的观察，在问题出现时逐案处理。

\begin{definition}
\label{flat-definition-one-step-devissage-at-x}
设 $S$ 是概形，$X$ 在 $S$ 上局部有限型，$\mathcal{F}$ 是有限型
拟凝聚 $\mathcal{O}_X$-模。设 $x \in X$ 在 $S$ 中的像为 $s$。
所谓 {\it $x$ 处 $\mathcal{F}/X/S$ 的一步 d\'evissage}，是指系统
$(Z, Y, i, \pi, \mathcal{G}, z, y)$，其中
$(Z, Y, i, \pi, \mathcal{G})$ 是 $s$ 上 $\mathcal{F}/X/S$
的一步 d\'evissage，并且
\begin{enumerate}
\item $\dim_x(\text{Supp}(\mathcal{F}_s)) = \dim(\text{Supp}(\mathcal{F}_s))$,
\item $z \in Z$ 是满足 $i(z) = x$ 与 $\pi(z) = y$ 的点；
\item $\pi^{-1}(\{y\}) = \{z\}$；
\item 扩张 $\kappa(y)/\kappa(s)$ 是纯超越的。
\end{enumerate}
\end{definition}

\noindent
$x$ 处 $\mathcal{F}/X/S$ 的一步 d\'evissage 只有在 $X$ 与 $S$
仿射时才可能存在。结合定义 \ref{flat-definition-one-step-devissage} 的条件
(5)，条件 (1) 保证 $Y \to S$ 的相对维数等于
$\dim_x(\text{Supp}(\mathcal{F}_s))$。

\begin{lemma}
\label{flat-lemma-elementary-devissage-variant}
设 $f : X \to S$ 是局部有限型概形态射，$\mathcal{F}$ 是有限型
拟凝聚 $\mathcal{O}_X$-模，且 $x \in X$ 在 $S$ 中的像为
$s = f(x)$。则存在带点概形的交换图
$$
\xymatrix{
(X, x) \ar[d]_f & (X', x') \ar[l]^g \ar[d] \\
(S, s) & (S', s') \ar[l] \\
}
$$
使得 $(S', s') \to (S, s)$ 与 $(X', x') \to (X, x)$ 是初等
\'etale 邻域，且 $g^*\mathcal{F}/X'/S'$ 在 $x'$ 处有一步
d\'evissage。
\end{lemma}

\begin{proof}
这立即由定义 \ref{flat-definition-one-step-devissage-at-x} 与引理
\ref{flat-lemma-elementary-devissage} 得出。
\end{proof}

\begin{lemma}
\label{flat-lemma-base-change-one-step}
设 $S$、$X$、$\mathcal{F}$、$s$ 如定义
\ref{flat-definition-one-step-devissage}，且
$(Z, Y, i, \pi, \mathcal{G})$ 是 $s$ 上 $\mathcal{F}/X/S$
的一步 d\'evissage。设 $(S', s') \to (S, s)$ 是任意带点概形态射。
由这些数据，令 $X', Z', Y', i', \pi'$ 分别为
$X, Z, Y, i, \pi$ 沿 $S' \to S$ 的基变换；令 $\mathcal{F}'$
为 $\mathcal{F}$ 到 $X'$ 的拉回，$\mathcal{G}'$ 为
$\mathcal{G}$ 到 $Z'$ 的拉回。若 $S'$ 仿射，则
$(Z', Y', i', \pi', \mathcal{G}')$ 是 $s'$ 上
$\mathcal{F}'/X'/S'$ 的一步 d\'evissage。
\end{lemma}

\begin{proof}
仿射概形的纤维积仿射；见概形章引理
\ref{schemes-lemma-fibre-product-affines}。基变换保持闭浸入、有限表示态射、
有限态射、光滑态射、纤维几何不可约的态射与相对维数为 $n$ 的态射；
见《态射》诸引理 \ref{morphisms-lemma-base-change-closed-immersion}、
\ref{morphisms-lemma-base-change-finite-presentation},
\ref{morphisms-lemma-base-change-finite},
\ref{morphisms-lemma-base-change-smooth},
\ref{morphisms-lemma-base-change-relative-dimension-d} 以及《态射进阶》引理
\ref{more-morphisms-lemma-base-change-fibres-geometrically-irreducible}.
由于沿有限态射 $i$ 的推前与基变换交换，有
$i'_*\mathcal{G}' \cong \mathcal{F}'$；见概形上同调章引理
\ref{coherent-lemma-affine-base-change}。又有
$\dim(\text{Supp}(\mathcal{F}_s)) = \dim(\text{Supp}(\mathcal{F}'_{s'}))$
这是因为下式成立，并可应用《态射》引理
\ref{morphisms-lemma-dimension-fibre-after-base-change}：
$$
\text{Supp}(\mathcal{F}_s) \times_s s' = \text{Supp}(\mathcal{F}'_{s'}).
$$
引理得证。
\end{proof}

\begin{lemma}
\label{flat-lemma-base-change-one-step-at-x}
设 $S$、$X$、$\mathcal{F}$、$x$、$s$ 如定义
\ref{flat-definition-one-step-devissage-at-x}，且
$(Z, Y, i, \pi, \mathcal{G}, z, y)$ 是 $x$ 处
$\mathcal{F}/X/S$ 的一步 d\'evissage。设带点概形态射
$(S', s') \to (S, s)$ 诱导同构 $\kappa(s) = \kappa(s')$。
令 $(Z', Y', i', \pi', \mathcal{G}')$ 如引理
\ref{flat-lemma-base-change-one-step} 中所构造，并令 $x' \in X'$
（分别地，$z' \in Z'$、$y' \in Y'$）为同时映到
$x \in X$（分别地，$z \in Z$、$y \in Y$）与 $s' \in S'$
的唯一一点。若 $S'$ 仿射，则
$(Z', Y', i', \pi', \mathcal{G}', z', y')$ 是 $x'$ 处
$\mathcal{F}'/X'/S'$ 的一步 d\'evissage。
\end{lemma}

\begin{proof}
由引理 \ref{flat-lemma-base-change-one-step}，
$(Z', Y', i', \pi', \mathcal{G}')$ 是 $s'$ 上
$\mathcal{F}'/X'/S'$ 的一步 d\'evissage。由于假设
$\kappa(s) = \kappa(s')$ 保证纤维 $X'_{s'}$、$Z'_{s'}$、$Y'_{s'}$
分别同构于 $X_s$、$Z_s$、$Y_s$，定义
\ref{flat-definition-one-step-devissage-at-x} 的性质 (1) -- (4) 对
$(Z', Y', i', \pi', \mathcal{G}', z', y')$ 成立。
\end{proof}

\begin{definition}
\label{flat-definition-shrink}
设 $S$、$X$、$\mathcal{F}$、$x$、$s$ 如定义
\ref{flat-definition-one-step-devissage-at-x}，且
$(Z, Y, i, \pi, \mathcal{G}, z, y)$ 是 $x$ 处
$\mathcal{F}/X/S$ 的一步 d\'evissage。此情形的一个
{\it 标准缩小}，是指给定标准开集 $S' \subset S$、$X' \subset X$、
$Z' \subset Z$、$Y' \subset Y$，满足 $s \in S'$、$x \in X'$、
$z \in Z'$、$y \in Y'$，并且
$$
(Z', Y', i|_{Z'}, \pi|_{Z'}, \mathcal{G}|_{Z'}, z, y)
$$
是 $x$ 处 $\mathcal{F}|_{X'}/X'/S'$ 的一步 d\'evissage。
\end{definition}

\begin{lemma}
\label{flat-lemma-shrink}
假设与记号同定义 \ref{flat-definition-shrink}。则：
\begin{enumerate}
\item
\label{flat-item-shrink-base}
若 $S' \subset S$ 是 $s$ 的标准开邻域，则置
$X' = X_{S'}$、$Z' = Z_{S'}$、$Y' = Y_{S'}$，得到一个标准缩小。
\item
\label{flat-item-shrink-on-Y}
设 $W \subset Y$ 是 $y$ 的标准开邻域。则存在满足
$Y' = W \times_S S'$ 的标准缩小。
\item
\label{flat-item-shrink-on-X}
设 $U \subset X$ 是 $x$ 的开邻域。则存在满足 $X' \subset U$
的标准缩小。
\end{enumerate}
\end{lemma}

\begin{proof}
(1) 立即由引理 \ref{flat-lemma-base-change-one-step-at-x} 与如下事实得出：
仿射概形态射下标准开集的逆像仍为标准开集；见《代数》引理
\ref{algebra-lemma-spec-functorial}。

\medskip\noindent
设 $W \subset Y$ 如 (2)。由于 $Y \to S$ 光滑，它是开映射；见《态射》引理
\ref{morphisms-lemma-smooth-open}。故可找到包含于 $W$ 之像中的
$s$ 的标准开邻域 $S'$。于是 $W_{S'} \to S'$ 的纤维是 $Y \to S$
在 $S'$ 上各纤维的非空开子概形，因而也几何不可约。置
$Y' = W_{S'}$ 与 $Z' = \pi^{-1}(Y')$，可见 $Z' \subset Z$
是 $z$ 的标准开邻域。取函数 $\overline{h} \in \Gamma(Z, \mathcal{O}_Z)$
使 $Z' = D(\overline{h})$。由于 $i : Z \to X$ 是闭浸入，可找到
$h \in \Gamma(X, \mathcal{O}_X)$ 使得
$i^\sharp(h) = \overline{h}$。取 $X' = D(h) \subset X$，
便得到 (2) 所述的标准缩小。

\medskip\noindent
设 $U \subset X$ 如 (3)。缩小 $U$ 后，可假设 $U$ 是标准开集。
由《态射进阶》引理
\ref{more-morphisms-lemma-finite-morphism-single-point-in-fibre}，
存在 $y$ 的标准开邻域 $W \subset Y$，使得
$\pi^{-1}(W) \subset i^{-1}(U)$。应用 (2)，得到标准缩小
$X', S', Z', Y'$，其中 $Y' = W_{S'}$。由于
$Z' \subset \pi^{-1}(W) \subset i^{-1}(U)$，可用 $X' \cap U$
替换 $X'$（因 $U$ 也是标准开集，交集仍为标准开集），而不破坏
标准缩小定义中的任何条件。结论得证。
\end{proof}

\begin{lemma}
\label{flat-lemma-elementary-etale-neighbourhood}
设 $S$、$X$、$\mathcal{F}$、$x$、$s$ 如定义
\ref{flat-definition-one-step-devissage-at-x}，且
$(Z, Y, i, \pi, \mathcal{G}, z, y)$ 是 $x$ 处
$\mathcal{F}/X/S$ 的一步 d\'evissage。设
$$
\xymatrix{
(Y, y) \ar[d] & (Y', y') \ar[l] \ar[d] \\
(S, s) & (S', s') \ar[l]
}
$$
是带点概形的交换图，且水平箭头为初等 \'etale 邻域。则存在交换图
$$
\xymatrix{
& & (X'', x'') \ar[lld] \ar[d] & (Z'', z'') \ar[l] \ar[lld] \ar[d] \\
(X, x) \ar[d] & (Z, z) \ar[l] \ar[d] &
(S'', s'') \ar[lld] & (Y'', y'') \ar[lld] \ar[l] \\
(S, s) & (Y, y) \ar[l]
}
$$
满足下列性质：
\begin{enumerate}
\item $(S'', s'') \to (S', s')$ 是初等 \'etale 邻域，且态射
$S'' \to S$ 是复合 $S'' \to S' \to S$；
\item $Y''$ 是 $Y' \times_{S'} S''$ 的开子概形；
\item $Z'' = Z \times_Y Y''$,
\item $(X'', x'') \to (X, x)$ 是初等 \'etale 邻域；
\item $(Z'', Y'', i'', \pi'', \mathcal{G}'', z'', y'')$ 是层
$\mathcal{F}''$ 在 $x''$ 处的一步 d\'evissage。
\end{enumerate}
这里 $\mathcal{F}''$（分别地，$\mathcal{G}''$）是 $\mathcal{F}$
（分别地，$\mathcal{G}$）沿态射 $X'' \to X$（分别地，$Z'' \to Z$）
的拉回，而 $i'' : Z'' \to X''$ 与 $\pi'' : Z'' \to Y''$ 如图所示。
\end{lemma}

\begin{proof}
任取初等 \'etale 邻域 $(S'', s'') \to (S', s')$，其中 $S''$ 仿射。
任取包含点 $y'' = (y', s'')$ 的仿射开邻域
$Y'' \subset Y' \times_{S'} S''$。由 (3) 得到仿射的
$(Z'', z'')$。此外，$Z_{S''} \to X_{S''}$ 是闭浸入，而
$Z'' \to Z_{S''}$ 是 \'etale 态射。因此可应用引理
\ref{flat-lemma-lift-etale}，找到仿射概形的 \'etale 态射
$X'' \to X_{S'}$，使 $Z'' \cong X'' \times_{X_{S'}} Z_{S'}$。
% P05-EN-CH38-0006: frozen proof constructs over S'' but then writes X_{S'} and Z_{S'}.
记 $i'' : Z'' \to X''$ 为相应的闭浸入。置 $x'' = i''(z'')$，
得到引理所述的交换图。性质 (1)、(2)、(3)、(4) 按构造成立。
故只需为适当选择的 $(S'', s'') \to (S', s')$ 与 $Y''$ 证明 (5)。

\medskip\noindent
先列出对第一段中任意选择的 $(S'', s'') \to (S', s')$ 与 $Y''$
都成立的性质。按构造有 $Z'' = X'' \times_X Z$，故
$i''_*\mathcal{G}'' = \mathcal{F}''$（记号同引理陈述）；见概形上同调章
引理 \ref{coherent-lemma-affine-base-change}。置
$n = \dim(\text{Supp}(\mathcal{F}_s)) =
\dim_x(\text{Supp}(\mathcal{F}_s))$。态射 $Y'' \to S''$ 光滑且相对
维数为 $n$：事实上，$Y' \to S'$ 是相对维数为 $n$ 的 \'etale 态射与
相对维数为 $n$ 的光滑态射之复合 $Y' \to Y_{S'} \to S'$，且基变换保持相对维数为
$n$ 的光滑态射。又有 $\kappa(y'') = \kappa(y)$ 与
$\kappa(s) = \kappa(s'')$，故 $\kappa(y'')$ 是
$\kappa(s'')$ 的纯超越扩张。纤维态射 $X''_{s''} \to X_s$ 是
$\kappa(s) = \kappa(s'')$ 上仿射概形的 \'etale 态射，把点 $x''$
映到点 $x$，并把 $\mathcal{F}_s$ 拉回为 $\mathcal{F}''_{s''}$。因此
$$
\dim(\text{Supp}(\mathcal{F}''_{s''})) =
\dim(\text{Supp}(\mathcal{F}_s)) = n =
\dim_x(\text{Supp}(\mathcal{F}_s)) =
\dim_{x''}(\text{Supp}(\mathcal{F}''_{s''}))
$$
这是因为维数在 \'etale 局部化下不变；见下降章引理
\ref{descent-lemma-dimension-at-point-local}。由于
$\pi'' : Z'' \to Y''$ 是 $\pi$ 的基变换，$\pi''$ 有限；又因
$\kappa(y) = \kappa(y'')$，有 $\pi^{-1}(\{y''\}) = \{z''\}$。
% P05-EN-CH38-0007: frozen proof uses pi^{-1} on a point of Y''; the base-changed map is pi''.

\medskip\noindent
至此，定义 \ref{flat-definition-one-step-devissage} 的所有条件均已验证，
唯有尚未验证 $Y'' \to S''$ 的纤维几何不可约。一般而言这当然未必成立，
此处将使用 $\kappa(s) \subset \kappa(y)$ 为纯超越扩张。具体地，
设 $T \subset Y'_{s'}$ 为包含
$y' = (y, s')$ 的不可约分支。由于 $Y'_{s'}$ 在 $\kappa(s')$ 上光滑，
$T$ 是 $Y'_{s'}$ 的开子概形。又因 $\kappa(s') = \kappa(s)$
在 $\kappa(y') = \kappa(y)$ 中代数闭，由簇章引理
\ref{varieties-lemma-geometrically-connected-if-connected-and-point}，
$T$ 几何连通。由于 $T$ 光滑，$T$ 几何不可约。于是可应用《态射进阶》引理
\ref{more-morphisms-lemma-normal-morphism-irreducible}，找到初等
\'etale 态射 $(S'', s'') \to (S', s')$ 与仿射开集
$Y'' \subset Y'_{S''}$，使得 $Y'' \to S''$ 的全部纤维几何不可约，
且 $T = Y''_{s''}$。缩小（先缩小 $Y''$，再缩小 $S''$）后，可假设
$Y''$ 与 $S''$ 均仿射。引理得证。
\end{proof}

\begin{lemma}
\label{flat-lemma-existence-alpha}
设 $S$、$X$、$\mathcal{F}$、$s$ 如定义
\ref{flat-definition-one-step-devissage}，且
$(Z, Y, i, \pi, \mathcal{G})$ 是 $s$ 上 $\mathcal{F}/X/S$
的一步 d\'evissage。设 $\xi \in Y_s$ 为（唯一的）一般点。
则存在整数 $r > 0$ 与 $\mathcal{O}_Y$-模同态
$$
\alpha : \mathcal{O}_Y^{\oplus r} \longrightarrow \pi_*\mathcal{G}
$$
使得
$$
\alpha :
\kappa(\xi)^{\oplus r}
\longrightarrow
(\pi_*\mathcal{G})_\xi \otimes_{\mathcal{O}_{Y, \xi}} \kappa(\xi)
$$
为同构。此外，此时有
$$
\dim(\text{Supp}(\Coker(\alpha)_s)) < \dim(\text{Supp}(\mathcal{F}_s)).
$$
\end{lemma}

\begin{proof}
按假设，概形 $S$ 与 $Y$ 仿射。写作 $S = \Spec(A)$ 与
$Y = \Spec(B)$。由于 $\pi$ 有限，$\mathcal{O}_Y$-模
$\pi_*\mathcal{G}$ 是有限型拟凝聚 $\mathcal{O}_Y$-模。因此对某个有限
$B$-模 $N$，有 $\pi_*\mathcal{G} = \widetilde{N}$。设
$\mathfrak p \subset B$ 为与 $\xi$ 对应的素理想。为得到 $\alpha$，
置 $r = \dim_{\kappa(\mathfrak p)} N \otimes_B \kappa(\mathfrak p)$，
并选取 $x_1, \ldots, x_r \in N$，使其构成
$N \otimes_B \kappa(\mathfrak p)$ 的一组基。取
$\alpha : B^{\oplus r} \to N$ 为公式
$\alpha(b_1, \ldots, b_r) = \sum b_ix_i$ 给出的同态。显然，
$\alpha : \kappa(\mathfrak p)^{\oplus r} \to
N \otimes_B \kappa(\mathfrak p)$ 如所需为同构。最后，设 $\alpha$
是任意具有此性质的同态。则 $N' = \Coker(\alpha)$ 是有限 $B$-模，
且 $N' \otimes \kappa(\mathfrak p) = 0$。由 Nakayama 引理
（《代数》引理 \ref{algebra-lemma-NAK}），有
$N'_{\mathfrak p} = 0$。纤维 $Y_s$ 几何不可约、维数为 $n$，其一般点为
$\xi$；而刚刚已经看到 $\xi$ 不在 $\Coker(\alpha)$ 的支集中，
故引理最后一个断言成立。
\end{proof}


\section{完全 d\'evissage}
\label{flat-section-complete-devissage}

\noindent
本节解释模的完全 d\'evissage 是什么，并证明它存在。本节内容主要是记账工作。

\begin{definition}
\label{flat-definition-complete-devissage}
设 $S$ 是概形，$X$ 在 $S$ 上局部有限型，$\mathcal{F}$ 是有限型
拟凝聚 $\mathcal{O}_X$-模，且 $s \in S$ 为一点。所谓
{\it $s$ 上 $\mathcal{F}/X/S$ 的完全 d\'evissage}，是指给定图表
$$
\xymatrix{
X & Z_1 \ar[l]^{i_1} \ar[d]^{\pi_1} \\
& Y_1 & Z_2 \ar[l]^{i_2} \ar[d]^{\pi_2} \\
& & Y_2 & Z_3 \ar[l] \ar[d] \\
& & & ... & ... \ar[l] \ar[d] \\
& & & & Y_n
}
$$
其中各概形均在 $S$ 上，并给定有限型拟凝聚
$\mathcal{O}_{Z_k}$-模 $\mathcal{G}_k$ 与 $\mathcal{O}_{Y_k}$-模同态
$$
\alpha_k :
\mathcal{O}_{Y_k}^{\oplus r_k}
\longrightarrow
\pi_{k, *}\mathcal{G}_k,
\quad
k = 1, \ldots, n
$$
满足下列性质：
\begin{enumerate}
\item $(Z_1, Y_1, i_1, \pi_1, \mathcal{G}_1)$ 是 $s$ 上
$\mathcal{F}/X/S$ 的一步 d\'evissage；
\item 同态 $\alpha_k$ 诱导同构
$$
\kappa(\xi_k)^{\oplus r_k} \longrightarrow
(\pi_{k, *}\mathcal{G}_k)_{\xi_k}
\otimes_{\mathcal{O}_{Y_k, \xi_k}} \kappa(\xi_k)
$$
其中 $\xi_k \in (Y_k)_s$ 是唯一的一般点；
\item 对 $k = 2, \ldots, n$，系统
$(Z_k, Y_k, i_k, \pi_k, \mathcal{G}_k)$
是 $s$ 上 $\Coker(\alpha_{k - 1})/Y_{k - 1}/S$ 的一步 d\'evissage；
\item $\Coker(\alpha_n) = 0$.
\end{enumerate}
此时称
$(Z_k, Y_k, i_k, \pi_k, \mathcal{G}_k, \alpha_k)_{k = 1, \ldots, n}$
为 $s$ 上 $\mathcal{F}/X/S$ 的完全 d\'evissage。
\end{definition}

\begin{definition}
\label{flat-definition-complete-devissage-at-x}
设 $S$ 是概形，$X$ 在 $S$ 上局部有限型，$\mathcal{F}$ 是有限型
拟凝聚 $\mathcal{O}_X$-模。设 $x \in X$ 的像为 $s \in S$。
所谓 {\it $x$ 处 $\mathcal{F}/X/S$ 的完全 d\'evissage}，是指给定系统
$$
(Z_k, Y_k, i_k, \pi_k, \mathcal{G}_k, \alpha_k, z_k, y_k)_{k = 1, \ldots, n}
$$
使得 $(Z_k, Y_k, i_k, \pi_k, \mathcal{G}_k, \alpha_k)$ 是 $s$ 上
$\mathcal{F}/X/S$ 的完全 d\'evissage，并且
\begin{enumerate}
\item $(Z_1, Y_1, i_1, \pi_1, \mathcal{G}_1, z_1, y_1)$ 是 $x$ 处
$\mathcal{F}/X/S$ 的一步 d\'evissage；
\item 对 $k = 2, \ldots, n$，系统
$(Z_k, Y_k, i_k, \pi_k, \mathcal{G}_k, z_k, y_k)$
是 $y_{k - 1}$ 处 $\Coker(\alpha_{k - 1})/Y_{k - 1}/S$
的一步 d\'evissage。
\end{enumerate}
\end{definition}

\noindent
再次注意，完全 d\'evissage 只有在 $X$ 与 $S$ 仿射时才可能存在。

\begin{lemma}
\label{flat-lemma-base-change-complete}
设 $S$、$X$、$\mathcal{F}$、$s$ 如定义
\ref{flat-definition-complete-devissage}，且
$(S', s') \to (S, s)$ 是任意带点概形态射。设
$(Z_k, Y_k, i_k, \pi_k, \mathcal{G}_k, \alpha_k)_{k = 1, \ldots, n}$
是 $s$ 上 $\mathcal{F}/X/S$ 的完全 d\'evissage。由这些数据，
令 $X', Z'_k, Y'_k, i'_k, \pi'_k$ 分别为
$X, Z_k, Y_k, i_k, \pi_k$ 沿 $S' \to S$ 的基变换；令
$\mathcal{F}'$ 为 $\mathcal{F}$ 到 $X'$ 的拉回，$\mathcal{G}'_k$
为 $\mathcal{G}_k$ 到 $Z'_k$ 的拉回，$\alpha'_k$ 为
$\alpha_k$ 到 $Y'_k$ 的拉回。若 $S'$ 仿射，则
$(Z'_k, Y'_k, i'_k, \pi'_k, \mathcal{G}'_k, \alpha'_k)_{k = 1, \ldots, n}$
是 $s'$ 上 $\mathcal{F}'/X'/S'$ 的完全 d\'evissage。
\end{lemma}

\begin{proof}
由引理 \ref{flat-lemma-base-change-one-step}，一步 d\'evissage 的基变换仍为
一步 d\'evissage。因此只需证明 $\Coker(\alpha_k)$ 的形成与基变换交换，
且定义 \ref{flat-definition-complete-devissage} 的条件 (2) 在基变换下保持。
第一个断言成立，因为 $\pi'_{k, *}\mathcal{G}'_k$ 是
$\pi_{k, *}\mathcal{G}_k$ 的拉回（由概形上同调章引理
\ref{coherent-lemma-affine-base-change}），且 $\otimes$ 右正合。
第二个断言成立，因为由同一理由有
$$
(\pi_{k, *}\mathcal{G}_k)_{\xi_k}
\otimes_{\mathcal{O}_{Y_k, \xi_k}} \kappa(\xi_k)
\otimes_{\kappa(\xi_k)} \kappa(\xi'_k)
\cong
(\pi'_{k, *}\mathcal{G}'_k)_{\xi'_k}
\otimes_{\mathcal{O}_{Y'_k, \xi'_k}} \kappa(\xi'_k)
$$
其中记号含义显然。
\end{proof}

\begin{lemma}
\label{flat-lemma-base-change-complete-at-x}
设 $S$、$X$、$\mathcal{F}$、$x$、$s$ 如定义
\ref{flat-definition-complete-devissage-at-x}。设带点概形态射
$(S', s') \to (S, s)$ 诱导同构 $\kappa(s) = \kappa(s')$。设
$(Z_k, Y_k, i_k, \pi_k, \mathcal{G}_k, \alpha_k, z_k, y_k)_{k = 1, \ldots, n}$
是 $x$ 处 $\mathcal{F}/X/S$ 的完全 d\'evissage。令
$(Z'_k, Y'_k, i'_k, \pi'_k, \mathcal{G}'_k, \alpha'_k)_{k = 1, \ldots, n}$
如引理 \ref{flat-lemma-base-change-complete} 中所构造；令 $x' \in X'$
（分别地，$z'_k \in Z'$、$y'_k \in Y'$）为同时映到
$x \in X$（分别地，$z_k \in Z_k$、$y_k \in Y_k$）与
$s' \in S'$ 的唯一一点。若 $S'$ 仿射，则
$(Z'_k, Y'_k, i'_k, \pi'_k, \mathcal{G}'_k, \alpha'_k,
z'_k, y'_k)_{k = 1, \ldots, n}$
是 $x'$ 处 $\mathcal{F}'/X'/S'$ 的完全 d\'evissage。
\end{lemma}

\begin{proof}
合并引理 \ref{flat-lemma-base-change-complete} 与
\ref{flat-lemma-base-change-one-step-at-x} 即得。
\end{proof}

\begin{definition}
\label{flat-definition-shrink-complete}
设 $S$、$X$、$\mathcal{F}$、$x$、$s$ 如定义
\ref{flat-definition-complete-devissage-at-x}。考虑 $x$ 处
$(Z_k, Y_k, i_k, \pi_k, \mathcal{G}_k, \alpha_k, z_k, y_k)_{k = 1, \ldots, n}$
的 $\mathcal{F}/X/S$ 完全 d\'evissage。此情形的一个
{\it 标准缩小}，是指给定标准开集 $S' \subset S$、$X' \subset X$、
$Z'_k \subset Z_k$、$Y'_k \subset Y_k$，使得 $s_k \in S'$、
$x_k \in X'$、$z_k \in Z'$、$y_k \in Y'$，并且
% P05-EN-CH38-0008: frozen definition uses undefined s_k,x_k and unindexed Z',Y' in its membership conditions.
$$
(Z'_k, Y'_k, i'_k, \pi'_k,
\mathcal{G}'_k, \alpha'_k, z_k, y_k)_{k = 1, \ldots, n}
$$
是 $x$ 处 $\mathcal{F}'/X'/S'$ 的一步 d\'evissage，其中
% P05-EN-CH38-0009: frozen definition calls the whole indexed system a one step rather than complete devissage.
$\mathcal{G}'_k = \mathcal{G}_k|_{Z'_k}$，且
$\mathcal{F}' = \mathcal{F}|_{X'}$.
\end{definition}

\begin{lemma}
\label{flat-lemma-shrink-complete}
假设与记号同定义 \ref{flat-definition-shrink-complete}。则：
\begin{enumerate}
\item
\label{flat-item-shrink-base-complete}
若 $S' \subset S$ 是 $s$ 的标准开邻域，则置
$X' = X_{S'}$、$Z'_k = Z_{S'}$、$Y'_k = Y_{S'}$，得到标准缩小。
% P05-EN-CH38-0010: frozen formula drops the k indices from Z_k and Y_k on the right.
\item
\label{flat-item-shrink-on-Y-complete}
设 $W \subset Y_n$ 是 $y$ 的标准开邻域。则存在满足
$Y'_n = W \times_S S'$ 的标准缩小。
\item
\label{flat-item-shrink-on-X-complete}
设 $U \subset X$ 是 $x$ 的开邻域。则存在满足 $X' \subset U$
的标准缩小。
\end{enumerate}
\end{lemma}

\begin{proof}
(1) 立即由引理 \ref{flat-lemma-base-change-complete-at-x} 与
\ref{flat-lemma-shrink} 得出。

\medskip\noindent
证明 (2)。为方便起见，记 $X = Y_0$。应用引理
\ref{flat-lemma-shrink} (\ref{flat-item-shrink-on-Y})，找到如下标准缩小：
$S', Y'_{n - 1}, Z'_n, Y'_n$
它是 $y_{n - 1}$ 处 $\Coker(\alpha_{n - 1})/Y_{n - 1}/S$
的一步 d\'evissage 的标准缩小，并满足 $Y'_n = W \times_S S'$。
重复此过程，找到标准缩小
$S'', Y''_{n - 2}, Z''_{n - 1}, Y''_{n - 1}$
它是 $y_{n - 2}$ 处 $\Coker(\alpha_{n - 2})/Y_{n - 2}/S$
的一步 d\'evissage 的标准缩小，并满足
$Y''_{n - 1} = Y'_{n - 1} \times_S S''$。如此继续，直至得到
$S^{(n)}, Y^{(n)}_0, Z^{(n)}_1, Y^{(n)}_1$.
此时显然，取 $S^{(n)}$、$X^{(n)}$、
$Z_k^{(n - k)} \times_S S^{(n)}$ 与
$Y_k^{(n - k)} \times_S S^{(n)}$，便得到所需性质的标准缩小。

\medskip\noindent
证明 (3)。对完全 d\'evissage 的长度作归纳。首先应用引理
\ref{flat-lemma-shrink} (\ref{flat-item-shrink-on-X})，找到标准缩小
$S', X', Z'_1, Y'_1$
它来自 $x$ 处 $\mathcal{F}/X/S$ 的一步 d\'evissage，且
$X' \subset U$。若 $n = 1$，结论已得。若 $n > 1$，则由归纳假设，
可找到完全 d\'evissage 的标准缩小 $S''$、$Y''_1$、$Z''_k$、$Y''_k$：
$(Z_k, Y_k, i_k, \pi_k, \mathcal{G}_k, \alpha_k, z_k, y_k)_{k = 2, \ldots, n}$
它是 $x$ 处 $\Coker(\alpha_1)/Y_1/S$ 的完全 d\'evissage，且
$Y''_1 \subset Y'_1$。利用引理 \ref{flat-lemma-shrink}
(\ref{flat-item-shrink-on-Y})，可找到 $S''' \subset S'$、
$X''' \subset X'$、$Z'''_1$ 与
$Y'''_1 = Y''_1 \times_S S'''$，它们构成标准缩小。问题的解取为
$$
S''', X''', Z'''_1, Y'''_1, Z''_2 \times_S S''',
Y''_2 \times_S S''', \ldots, Z''_n \times_S S''', Y''_n \times_S S'''
$$
引理证明完毕。
\end{proof}

\begin{proposition}
\label{flat-proposition-existence-complete-at-x}
设 $S$ 是概形，$X$ 在 $S$ 上局部有限型，$x \in X$ 的像为
$s \in S$。则存在交换图
$$
\xymatrix{
(X, x) \ar[d] & (X', x') \ar[l]^g \ar[d] \\
(S, s) & (S', s') \ar[l]
}
$$
其中水平箭头为初等 \'etale 邻域，且 $g^*\mathcal{F}/X'/S'$ 在 $x$
处有完全 d\'evissage。
% P05-EN-CH38-0011: frozen conclusion locates the pulled-back sheaf's devissage at x rather than x'.
\end{proposition}

\begin{proof}
对整数 $d = \dim_x(\text{Supp}(\mathcal{F}_s))$ 作归纳。由引理
\ref{flat-lemma-elementary-devissage-variant}，存在图表
$$
\xymatrix{
(X, x) \ar[d] & (X', x') \ar[l]^g \ar[d] \\
(S, s) & (S', s') \ar[l]
}
$$
其中水平箭头为初等 \'etale 邻域，且 $g^*\mathcal{F}/X'/S'$ 在
$x'$ 处有一步 d\'evissage。由问题的局部性，可用
$(X', x') \to (S', s')$ 替换 $(X, x) \to (S, s)$。替换后可假设
存在 $\mathcal{F}/X/S$ 在 $x$ 处的一步 d\'evissage
$(Z_1, Y_1, i_1, \pi_1, \mathcal{G}_1)$。

\medskip\noindent
应用引理 \ref{flat-lemma-existence-alpha}，找到同态
$$
\alpha_1 :
\mathcal{O}_{Y_1}^{\oplus r_1}
\longrightarrow
\pi_{1, *}\mathcal{G}_1
$$
它诱导 $\kappa(\xi_1)$ 上向量空间的同构，其中
$\xi_1 \in Y_1$ 是 $Y_1$ 在 $s$ 上纤维的唯一一般点。此外，
$\dim_{y_1}(\text{Supp}(\Coker(\alpha_1)_s)) < d$。
$\Coker(\alpha_1)_s$ 在 $y_1$ 处的茎可能为零。在此情形，由引理
\ref{flat-lemma-shrink} (\ref{flat-item-shrink-on-Y}) 缩小 $Y_1$ 后，
可假设 $\Coker(\alpha_1) = 0$，从而得到长度为零的完全 d\'evissage。

\medskip\noindent
现在假设 $\Coker(\alpha_1)_s$ 在 $y_1$ 处的茎非零。在此情形，
由归纳假设存在交换图
\begin{equation}
\label{flat-equation-overcome-this}
\vcenter{
\xymatrix{
(Y_1, y_1) \ar[d] & (Y'_1, y'_1) \ar[l]^h \ar[d] \\
(S, s) & (S', s') \ar[l]
}
}
\end{equation}
其中水平箭头为初等 \'etale 邻域，且
$h^*\Coker(\alpha_1)/Y'_1/S'$ 有完全 d\'evissage
$$
(Z_k, Y_k, i_k, \pi_k, \mathcal{G}_k, \alpha_k, z_k, y_k)_{k = 2, \ldots, n}
$$
于 $y'_1$。（特别地，$i_2 : Z_2 \to Y'_1$ 是到 $Y'_2$ 的闭浸入。）
% P05-EN-CH38-0012: frozen parenthetical gives i_2 target Y'_1, then calls it an immersion into Y'_2.
此时，把引理 \ref{flat-lemma-elementary-etale-neighbourhood} 应用于
$S, X, \mathcal{F}, x, s$、系统
$(Z_1, Y_1, i_1, \pi_1, \mathcal{G}_1)$ 与图
(\ref{flat-equation-overcome-this})，得到图表
$$
\xymatrix{
& & (X'', x'') \ar[lld] \ar[d] & (Z''_1, z''_1) \ar[l] \ar[lld] \ar[d] \\
(X, x) \ar[d] & (Z_1, z_1) \ar[l] \ar[d] &
(S'', s'') \ar[lld] & (Y''_1, y''_1) \ar[lld] \ar[l] \\
(S, s) & (Y_1, y_1) \ar[l]
}
$$
它具有所引引理列出的全部性质。特别地，
$Y''_1 \subset Y'_1 \times_{S'} S''$。置
$X_1 = Y'_1 \times_{S'} S''$，并记 $\mathcal{F}_1$ 为
$\Coker(\alpha_1)$ 的拉回。由引理
\ref{flat-lemma-base-change-complete-at-x}，系统
\begin{equation}
\label{flat-equation-shrink-this}
(Z_k \times_{S'} S'',
Y_k \times_{S'} S'', i''_k, \pi''_k, \mathcal{G}''_k,
\alpha''_k, z''_k, y''_k)_{k = 2, \ldots, n}
\end{equation}
是 $\mathcal{F}_1$ 到 $X_1$ 的完全 d\'evissage。再次由问题的性质，
可用 $(X'', x'') \to (S'', s'')$ 替换
$(X, x) \to (S, s)$。由此可假设：
\begin{enumerate}
\item[(a)] 存在 $\mathcal{F}/X/S$ 在 $x$ 处的一步 d\'evissage
$(Z_1, Y_1, i_1, \pi_1, \mathcal{G}_1)$；
\item[(b)] 存在 $\alpha_1 : \mathcal{O}_{Y_1}^{\oplus r_1}
\to \pi_{1, *}\mathcal{G}_1$，使得
$\alpha \otimes \kappa(\xi_1)$ 为同构；
% P05-EN-CH38-0013: frozen condition switches alpha_1 to alpha.
\item[(c)] $Y_1 \subset X_1$ 是开集，$y_1 = x_1$，且
$\mathcal{F}_1|_{Y_1} \cong \Coker(\alpha_1)$；
\item[(d)] 存在完全 d\'evissage
$(Z_k, Y_k, i_k, \pi_k, \mathcal{G}_k, \alpha_k, z_k, y_k)_{k = 2, \ldots, n}$
它是 $\mathcal{F}_1/X_1/S$ 在 $x_1$ 处的完全 d\'evissage。
\end{enumerate}
为完成证明，只需缩小这一步 d\'evissage 与完全 d\'evissage，使二者拼接成
一个完全 d\'evissage。（建议读者用引理 \ref{flat-lemma-shrink} 与
\ref{flat-lemma-shrink-complete} 自行完成，而不必阅读下述证明。）
由于 $Y_1 \subset X_1$ 是 $x_1$ 的开邻域，可应用引理
\ref{flat-lemma-shrink-complete} (\ref{flat-item-shrink-on-X-complete})，
找到数据 (d) 的标准缩小 $S', X'_1, Z'_2, Y'_2, \ldots, Y'_n$，
使 $X'_1 \subset Y_1$。注意，$X'_1$ 也是仿射概形 $Y_1$ 的标准开集。
接着如下缩小数据 (a)：先由引理 \ref{flat-lemma-shrink}
(\ref{flat-item-shrink-base}) 把底 $S$ 缩小到 $S'$，再用引理
\ref{flat-lemma-shrink} (\ref{flat-item-shrink-on-Y}) 把所得结果缩小为
$S''$、$X''$、$Z''_1$、$Y''_1$，最终满足
$Y''_1 = X'_1 \times_S S''$ 与 $S'' \subset S'$。于是
$$
Z''_1, Y''_1, Z'_2 \times_{S'} S'', Y'_2 \times_{S'} S'', \ldots,
Y'_n \times_{S'} S''
$$
给出所求的完全 d\'evissage。
\end{proof}

\noindent
再做一些记账工作，得到下列推论。

\begin{lemma}
\label{flat-lemma-existence-complete}
设 $X \to S$ 是有限型概形态射，$\mathcal{F}$ 是有限型拟凝聚
$\mathcal{O}_X$-模，且 $s \in S$ 为一点。则存在初等 \'etale 邻域
$(S', s') \to (S, s)$ 与 \'etale 态射
$h_i : Y_i \to X_{S'}$，$i = 1, \ldots, n$，使得对每个 $i$，
在 $s'$ 上存在 $\mathcal{F}_i/Y_i/S'$ 的完全 d\'evissage，其中
$\mathcal{F}_i$ 是 $\mathcal{F}$ 到 $Y_i$ 的拉回，并且
$X_s = (X_{S'})_{s'} \subset \bigcup h_i(Y_i)$。
\end{lemma}

\begin{proof}
对每一点 $x \in X_s$，可找到图表
$$
\xymatrix{
(X, x) \ar[d] & (X', x') \ar[l]^g \ar[d] \\
(S, s) & (S', s') \ar[l]
}
$$
其中水平箭头为初等 \'etale 邻域，且 $g^*\mathcal{F}/X'/S'$ 在
$x'$ 处有完全 d\'evissage。由于 $X \to S$ 有限型，纤维 $X_s$
拟紧；又由于上述每个 $g : X' \to X$ 都是开映射，可用有限多个
$g(X'_{s'})$ 覆盖 $X_s$。因此可找到这类图表的有限族
$$
\vcenter{
\xymatrix{
(X, x) \ar[d] & (X'_i, x'_i) \ar[l]^{g_i} \ar[d] \\
(S, s) & (S'_i, s'_i) \ar[l]
}
}
\quad i = 1, \ldots, n
$$
使得 $X_s = \bigcup g_i(X'_i)$。置
$S' = S'_1 \times_S \ldots \times_S S'_n$，并令
$Y_i = X_i \times_{S'_i} S'$ 为 $X'_i$ 到 $S'$ 的基变换。
由引理 \ref{flat-lemma-base-change-complete}，$\mathcal{F}$ 到 $Y_i$
的拉回在 $s$ 上有完全 d\'evissage，结论得证。
% P05-EN-CH38-0014: frozen proof defines Y_i from X_i although only X'_i was introduced, and concludes over s rather than s'.
\end{proof}



\section{代数表述}
\label{flat-section-translation}

\noindent
把具有完全 d\'evissage 的含义用代数语言写清楚或许有用。为此引入下列概念
（它并不十分有用，所以我们给它一个长得不可思议的名称）。

\begin{definition}
\label{flat-definition-elementary-etale-neighbourhood}
设 $R \to S$ 是环同态，$\mathfrak q$ 是 $S$ 中位于 $R$ 的素理想
$\mathfrak p$ 上方的素理想。所谓{\it 环同态 $R \to S$ 在
$\mathfrak q$ 处的初等 \'etale 局部化}，是指给定环及相伴素理想的交换图
$$
\xymatrix{
S \ar[r] & S' \\
R \ar[u] \ar[r] & R' \ar[u]
}
\quad\quad
\xymatrix{
\mathfrak q \ar@{-}[r] & \mathfrak q' \\
\mathfrak p \ar@{-}[u] \ar@{-}[r] & \mathfrak p' \ar@{-}[u]
}
$$
使得 $R \to R'$ 与 $S \to S'$ 是 \'etale 环同态，且
$\kappa(\mathfrak p) = \kappa(\mathfrak p')$ 与
$\kappa(\mathfrak q) = \kappa(\mathfrak q')$。
\end{definition}

\begin{definition}
\label{flat-definition-complete-devissage-algebra}
设 $R \to S$ 是有限型环同态，$\mathfrak r$ 是 $R$ 的素理想，
$N$ 是有限 $S$-模。所谓 {\it $\mathfrak r$ 上 $N/S/R$ 的完全
d\'evissage}，是指给定 $R$-代数同态
$$
\xymatrix{
& A_1 & & A_2 & & ... & & A_n \\
S \ar[ru] & & B_1 \ar[lu] \ar[ru] & & ... \ar[lu] \ar[ru] & &
... \ar[lu] \ar[ru] & & B_n \ar[lu]
}
$$
有限 $A_i$-模 $M_i$ 与 $B_i$-模同态
$\alpha_i : B_i^{\oplus r_i} \to M_i$，满足：
\begin{enumerate}
\item $S \to A_1$ 满且有限表示；
\item $B_i \to A_{i + 1}$ 满且有限表示；
\item $B_i \to A_i$ 有限；
\item $R \to B_i$ 光滑且纤维几何不可约；
\item 作为 $S$-模有 $N \cong M_1$；
\item 作为 $B_i$-模有 $\Coker(\alpha_i) \cong M_{i + 1}$；
\item 同态 $\alpha_i : \kappa(\mathfrak p_i)^{\oplus r_i}
\to M_i \otimes_{B_i} \kappa(\mathfrak p_i)$ 为同构，
其中 $\mathfrak p_i = \mathfrak rB_i$；
\item $\Coker(\alpha_n) = 0$.
\end{enumerate}
% P05-EN-CH38-0015: frozen conditions (2) and (6) lack an i<n range, and condition (7) treats rB_i as a prime without qualification.
在此情形，称
$(A_i, B_i, M_i, \alpha_i)_{i = 1, \ldots, n}$
为 $\mathfrak r$ 上 $N/S/R$ 的完全 d\'evissage。
\end{definition}

\begin{remark}
\label{flat-remark-finite-presentation}
注意，对所有 $i$，$R$-代数 $B_i$ 在 $R$ 上有限表示；当 $i \geq 2$
时，$A_i$ 亦然。若 $S$ 在 $R$ 上有限表示，则 $A_1$ 也在 $R$
上有限表示。在此情形，完全 d\'evissage 中的全部环同态均有限表示。
见《代数》引理 \ref{algebra-lemma-compose-finite-type}。仍假设 $S$
在 $R$ 上有限表示，则下列条件等价：
\begin{enumerate}
\item $M$ 在 $S$ 上有限表示；
\item $M_1$ 在 $A_1$ 上有限表示；
\item $M_1$ 在 $B_1$ 上有限表示；
\item 每个 $M_i$ 作为 $A_i$-模与作为 $B_i$-模均有限表示。
\end{enumerate}
% P05-EN-CH38-0016: frozen remark switches the finite module from N to undefined M.
等价关系 (1) $\Leftrightarrow$ (2) 与 (2) $\Leftrightarrow$ (3)
由《代数》引理 \ref{algebra-lemma-finite-finitely-presented-extension} 得出。
若 $M_1$ 有限表示，则 $\Coker(\alpha_1)$ 也有限表示（见《代数》引理
\ref{algebra-lemma-extension}），因而 $M_2$ 有限表示，依此类推。
\end{remark}

\begin{definition}
\label{flat-definition-complete-devissage-at-x-algebra}
设 $R \to S$ 是有限型环同态，$\mathfrak q$ 是 $S$ 中位于 $R$
的素理想 $\mathfrak r$ 上方的素理想，且 $N$ 是有限 $S$-模。
所谓 {\it $N/S/R$ 在 $\mathfrak q$ 处的完全 d\'evissage}，是指给定
$\mathfrak r$ 上 $N/S/R$ 的完全 d\'evissage
$(A_i, B_i, M_i, \alpha_i)_{i = 1, \ldots, n}$ 与位于
$\mathfrak r$ 上方的素理想 $\mathfrak q_i \subset B_i$，使得
\begin{enumerate}
\item $\kappa(\mathfrak r) \subset \kappa(\mathfrak q_i)$ 是纯超越扩张；
\item 存在唯一的素理想 $\mathfrak q'_i \subset A_i$ 位于
$\mathfrak q_i \subset B_i$ 上方；
\item $\mathfrak q = \mathfrak q'_1 \cap S$，且
$\mathfrak q_i = \mathfrak q'_{i + 1} \cap A_i$,
\item $R \to B_i$ 的相对维数为
$\dim_{\mathfrak q_i}(\text{Supp}(M_i \otimes_R \kappa(\mathfrak r)))$.
\end{enumerate}
\end{definition}

\begin{remark}
\label{flat-remark-same-notion}
设 $A \to B$ 是有限型环同态，$N$ 是有限 $B$-模，$\mathfrak q$
是 $B$ 中位于 $A$ 的素理想 $\mathfrak r$ 上方的素理想。置
$X = \Spec(B)$、$S = \Spec(A)$，并在 $X$ 上置
$\mathcal{F} = \widetilde{N}$。设 $x$ 为与 $\mathfrak q$ 对应的点，
$s \in S$ 为与 $\mathfrak p$ 对应的点。则
% P05-EN-CH38-0017: frozen remark introduces r but later calls the corresponding base prime p.
\begin{enumerate}
\item 若存在 $s$ 上 $\mathcal{F}/X/S$ 的完全 d\'evissage，则存在
$\mathfrak p$ 上 $N/B/A$ 的完全 d\'evissage；
\item 存在 $\mathcal{F}/X/S$ 在 $x$ 处的完全 d\'evissage，当且仅当存在
$N/B/A$ 在 $\mathfrak q$ 处的完全 d\'evissage。
\end{enumerate}
这里有一点细微差别：在“$\mathfrak p$ 上 $N/B/A$ 的完全
d\'evissage”的表述中，我们省略了相对维数条件，所以 (1) 中的蕴含
只有一个方向。在 $\mathfrak q$ 处完全 d\'evissage 的概念中已经包含
此条件。无论如何，我们只会使用如下事实：$\mathcal{F}/X/S$
存在完全 d\'evissage 蕴含 $N/B/A$ 存在完全 d\'evissage。
\end{remark}

\begin{lemma}
\label{flat-lemma-existence-algebra}
设 $R \to S$ 是有限型环同态，$M$ 是有限 $S$-模，
$\mathfrak q$ 是 $S$ 的素理想。则存在环同态 $R \to S$ 在
$\mathfrak q$ 处的初等 \'etale 局部化
$R' \to S', \mathfrak q', \mathfrak p'$，使得
$(M \otimes_S S')/S'/R'$ 在 $\mathfrak q'$ 处存在完全 d\'evissage。
\end{lemma}

\begin{proof}
这是命题 \ref{flat-proposition-existence-complete-at-x} 经由注记
\ref{flat-remark-same-notion} 得到的重述。
\end{proof}




\section{局部化与普遍单射同态}
\label{flat-section-localize-universally-injective}


\begin{lemma}
\label{flat-lemma-homothety-spectrum}
设 $R \to S$ 是环同态，$N$ 是 $S$-模。假设
\begin{enumerate}
\item $R$ 是极大理想为 $\mathfrak m$ 的局部环；
\item $\overline{S} = S/\mathfrak m S$ 是诺特环；并且
% P05-EN-CH38-0018: frozen source writes m_R although only the maximal ideal m is defined.
\item $\overline{N} = N/\mathfrak m_R N$ 是有限 $\overline{S}$-模。
\end{enumerate}
令 $\Sigma \subset S$ 为由在 $\overline{N}$ 上不是零因子的元素组成的乘法子集。
则 $\Sigma^{-1}S$ 是半局部环，其谱由包含于
$\text{Ass}_S(\overline{N})$ 某一元素的素理想 $\mathfrak q \subset S$ 组成。
此外，$\Sigma^{-1}S$ 的任一极大理想都对应于 $\overline{N}$
作为 $\overline{S}$-模的一个伴随素理想。
\end{lemma}

\begin{proof}
注意
$\text{Ass}_S(\overline{N}) = \text{Ass}_{\overline{S}}(\overline{N})$，参见
代数，引理 \ref{algebra-lemma-ass-quotient-ring}。
由代数，引理 \ref{algebra-lemma-finite-ass}，这是有限集。
记 $\{\mathfrak q_1, \ldots, \mathfrak q_r\} = \text{Ass}_S(\overline{N})$。
由代数，引理 \ref{algebra-lemma-ass-zero-divisors}，有
$\Sigma = S \setminus (\bigcup \mathfrak q_i)$。
结合代数，引理 \ref{algebra-lemma-spec-localization} 中对
$\Spec(\Sigma^{-1}S)$ 的描述和代数，引理 \ref{algebra-lemma-silly}，
可知 $\Sigma^{-1}S$ 的素理想对应于 $S$ 中包含于某个
$\mathfrak q_i$ 的素理想。因此，$\Sigma^{-1}S$ 的极大理想与集合
$\{\mathfrak q_1, \ldots, \mathfrak q_r\}$ 中关于包含关系的极大元一一对应。
引理得证。
\end{proof}

\begin{lemma}
\label{flat-lemma-homothety-universally-injective}
采用引理 \ref{flat-lemma-homothety-spectrum} 的假设与记号。再假设
\begin{enumerate}
\item $S$ 是局部环，且 $R \to S$ 是局部同态；
\item $S$ 在 $R$ 上本质有限表示；
\item $N$ 是有限表示 $S$-模；并且
\item $N$ 在 $R$ 上平坦。
\end{enumerate}
则每个 $s \in \Sigma$ 都定义一个普遍单射的 $R$-模同态
$s : N \to N$，并且同态 $N \to \Sigma^{-1}N$ 是 $R$-普遍单射的。
\end{lemma}

\begin{proof}
由代数，引理 \ref{algebra-lemma-mod-injective-general}，序列
$0 \to N \to N \to N/sN \to 0$ 正合，且 $N/sN$ 在 $R$ 上平坦。
这说明 $s : N \to N$ 普遍单射；参见代数，引理
\ref{algebra-lemma-flat-tor-zero}。同态 $N \to \Sigma^{-1}N$
是同态 $s : N \to N$ 的有向余极限，因而普遍单射。
\end{proof}

\begin{lemma}
\label{flat-lemma-base-change-universally-flat-local}
设 $R \to S$ 是环同态，$N$ 是 $S$-模，$S \to S'$ 是环同态。假设
\begin{enumerate}
\item $R \to S$ 是局部环之间的局部同态；
\item $S$ 在 $R$ 上本质有限表示；
\item $N$ 是有限表示 $S$-模；
\item $N$ 在 $R$ 上平坦；
\item $S \to S'$ 平坦；并且
\item $\Spec(S') \to \Spec(S)$ 的像包含 $S$ 中位于
$\mathfrak m_R$ 上方的所有素理想 $\mathfrak q$，其中
$\mathfrak q$ 是 $N/\mathfrak m_R N$ 的伴随素理想。
\end{enumerate}
则 $N \to N \otimes_S S'$ 是 $R$-普遍单射的。
\end{lemma}

\begin{proof}
% P05-EN-CH38-0019: frozen source defines N' using tensor over R although S' is introduced as an S-algebra and the proof later uses tensor over S.
置 $N' = N \otimes_R S'$。考虑交换图
$$
\xymatrix{
N \ar[d] \ar[r] & N' \ar[d] \\
\Sigma^{-1}N \ar[r] & \Sigma^{-1}N'
}
$$
其中 $\Sigma \subset S$ 是由在 $N/\mathfrak m_R N$ 上不是零因子的元素组成的集合。
若能证明同态 $N \to \Sigma^{-1}N'$ 普遍单射，则 $N \to N'$ 也普遍单射
（参见代数，引理 \ref{algebra-lemma-universally-injective-permanence}）。

\medskip\noindent
由引理 \ref{flat-lemma-homothety-spectrum}，环 $\Sigma^{-1}S$ 是半局部环，
其极大理想对应于 $N/\mathfrak m_R N$ 的伴随素理想。因此，由假设，
$\Spec(\Sigma^{-1}S') \to \Spec(\Sigma^{-1}S)$ 的像包含所有这些极大理想。
由代数，引理 \ref{algebra-lemma-ff-rings}，环同态
$\Sigma^{-1}S \to \Sigma^{-1}S'$ 忠实平坦。因而
$\Sigma^{-1}N \to \Sigma^{-1}N'$，即同态
$$
N \otimes_S \Sigma^{-1}S \longrightarrow N \otimes_S \Sigma^{-1}S'
$$
普遍单射；参见代数，引理
\ref{algebra-lemma-faithfully-flat-universally-injective} 和
\ref{algebra-lemma-universally-injective-tensor}。最后应用引理
\ref{flat-lemma-homothety-universally-injective}，可知
$N \to \Sigma^{-1}N$ 普遍单射。普遍单射模同态的复合仍普遍单射
（参见代数，引理 \ref{algebra-lemma-composition-universally-injective}），
故 $N \to \Sigma^{-1}N'$ 普遍单射，得证。
\end{proof}

\begin{lemma}
\label{flat-lemma-base-change-universally-flat}
设 $R \to S$ 是环同态，$N$ 是 $S$-模，$S \to S'$ 是环同态。假设
\begin{enumerate}
\item $R \to S$ 有限表示，且 $N$ 是有限表示 $S$-模；
\item $N$ 在 $R$ 上平坦；
\item $S \to S'$ 平坦；并且
\item $\Spec(S') \to \Spec(S)$ 的像包含所有满足下述条件的素理想
$\mathfrak q$：若 $\mathfrak p$ 是 $\mathfrak q$ 在 $R$ 中的逆像，则
$\mathfrak q$ 是 $N \otimes_R \kappa(\mathfrak p)$ 的伴随素理想。
\end{enumerate}
则 $N \to N \otimes_S S'$ 是 $R$-普遍单射的。
\end{lemma}

\begin{proof}
由代数，引理 \ref{algebra-lemma-universally-injective-check-stalks}，只须证明：
对于 $S$ 中位于 $R$ 的 $\mathfrak p$ 上方的任一素理想 $\mathfrak q$，
% P05-EN-CH38-0020: frozen source tensors over R and localizes S' at q although the map is S to S' and no prime above q is specified.
$N_{\mathfrak q} \to (N \otimes_R S')_{\mathfrak q}$ 是
$R_{\mathfrak p}$-普遍单射的。因此可将引理
\ref{flat-lemma-base-change-universally-flat-local} 应用于环同态
$R_{\mathfrak p} \to S_{\mathfrak q} \to S'_{\mathfrak q}$
以及模 $N_{\mathfrak q}$。
\end{proof}

\noindent
读者不妨将下述引理与代数，引理 \ref{algebra-lemma-mod-injective}、
\ref{algebra-lemma-mod-injective-general} 以及
第 \ref{flat-section-variants-mod-injective} 节的结果比较。
在每种情形中，结论都是同态 $u : M \to N$ 普遍单射且余核平坦。

\begin{lemma}
\label{flat-lemma-universally-injective-local}
设 $(R, \mathfrak m)$ 是局部环，$u : M \to N$ 是 $R$-模同态。
若 $M$ 是投射 $R$-模，$N$ 是平坦 $R$-模，并且
$\overline{u} : M/\mathfrak mM \to N/\mathfrak mN$ 单射，
则 $u$ 普遍单射。
\end{lemma}

\begin{proof}
由代数，定理 \ref{algebra-theorem-projective-free-over-local-ring}，模 $M$ 自由。
只要对 $M$ 的每个有限生成直和因子证明结论，引理便随之成立。因此可设
$M$ 有限自由。将 $N = \colim_i N_i$ 写成有限自由模的有向余极限；
参见代数，定理 \ref{algebra-theorem-lazard}。由于 $M$ 有限自由，
同态 $u : M \to N$ 对某个 $i$ 通过 $N_i$ 分解。将相应的 $R$-模同态记为
$u_i : M \to N_i$。因为 $\overline{u}$ 单射，所以
$\overline{u_i} : M/\mathfrak mM \to N_i/\mathfrak mN_i$ 单射，
并且对每个 $i' \geq i$，它与同态
$N_i/\mathfrak mN_i \to N_{i'}/\mathfrak mN_{i'}$ 的复合仍单射。
由于 $M$ 和 $N_{i'}$ 都是局部环 $R$ 上的有限自由模，这说明对每个
$i' \geq i$，$M \to N_{i'}$ 都是分裂单射。因此，对任一 $R$-模 $Q$，
以及所有 $i' \geq i$，同态 $M \otimes_R Q \to N_{i'} \otimes_R Q$
均单射。又因为 $- \otimes_R Q$ 与余极限交换，故可得
% P05-EN-CH38-0021: frozen conclusion retains N_{i'} after passing to the colimit; the target should be N.
$M \otimes_R Q \to N_{i'} \otimes_R Q$ 单射，如所要求。
\end{proof}

\begin{lemma}
\label{flat-lemma-invert-universally-injective}
采用引理 \ref{flat-lemma-homothety-spectrum} 的假设与记号。
再假设 $N$ 作为 $R$-模是投射的。则每个 $s \in \Sigma$ 都定义一个
普遍单射的 $R$-模同态 $s : N \to N$，并且同态
$N \to \Sigma^{-1}N$ 是 $R$-普遍单射的。
\end{lemma}

\begin{proof}
取 $s \in \Sigma$。由引理 \ref{flat-lemma-universally-injective-local}，
同态 $s : N \to N$ 普遍单射。同态 $N \to \Sigma^{-1}N$
是同态 $s : N \to N$ 的有向余极限，因而普遍单射。
\end{proof}






\section{完备化与 Mittag-Leffler 模}
\label{flat-section-completion-ML}


\begin{lemma}
\label{flat-lemma-completed-direct-sum-ML}
设 $R$ 是环，$I \subset R$ 是理想，$A$ 是集合。
假设 $R$ 是诺特环且关于 $I$ 完备。则完备化
$(\bigoplus\nolimits_{\alpha \in A} R)^\wedge$ 平坦且为 Mittag-Leffler 模。
\end{lemma}

\begin{proof}
由《代数进阶》，引理
\ref{more-algebra-lemma-ui-completion-direct-sum-into-product}，同态
$(\bigoplus\nolimits_{\alpha \in A} R)^\wedge
\to \prod_{\alpha \in A} R$ 普遍单射。因此，由代数，引理
\ref{algebra-lemma-ui-flat-domain} 和 \ref{algebra-lemma-pure-submodule-ML}，
只须证明 $\prod_{\alpha \in A} R$ 平坦且为 Mittag-Leffler 模。
由代数，命题 \ref{algebra-proposition-characterize-coherent}
（以及代数，引理 \ref{algebra-lemma-Noetherian-coherent}），可知
$\prod_{\alpha \in A} R$ 平坦。又因为 $R$ 的诸副本之积是
Mittag-Leffler 模，结论成立；参见代数，引理
\ref{algebra-lemma-product-over-Noetherian-ring}。
\end{proof}

\begin{lemma}
\label{flat-lemma-lift-ML}
设 $R$ 是环，$I \subset R$ 是理想，$M$ 是 $R$-模。假设
\begin{enumerate}
\item $R$ 是诺特环且 $I$-进完备；
\item $M$ 在 $R$ 上平坦；并且
\item $M/IM$ 是投射 $R/I$-模。
\end{enumerate}
则 $I$-进完备化 $M^\wedge$ 是平坦 Mittag-Leffler $R$-模。
\end{lemma}

\begin{proof}
选择满射 $F \to M$，其中 $F$ 是自由 $R$-模。由代数，引理
\ref{algebra-lemma-split-completed-sequence}，模 $M^\wedge$ 是模
$F^\wedge$ 的直和因子。因此只须对 $F$ 证明引理；此时结论由引理
\ref{flat-lemma-completed-direct-sum-ML} 得出。
\end{proof}

\noindent
在引理 \ref{flat-lemma-universally-injective-to-completion} 和
\ref{flat-lemma-universally-injective-to-completion-flat} 中，若 $R \to S$
有限型，则 $S$ 为诺特环这一假设成立；参见代数，引理
\ref{algebra-lemma-Noetherian-permanence}。

\begin{lemma}
\label{flat-lemma-universally-injective-to-completion}
设 $R$ 是环，$I \subset R$ 是理想，$R \to S$ 是环同态，
$N$ 是 $S$-模。假设
\begin{enumerate}
\item $R$ 是诺特环；
\item $S$ 是诺特环；
\item $N$ 是有限 $S$-模；并且
\item 对任一有限 $R$-模 $Q$，每个
$\mathfrak q \in \text{Ass}_S(Q \otimes_R N)$ 均满足
$IS + \mathfrak q \not = S$。
\end{enumerate}
则从 $N$ 到 $N$ 的 $I$-进完备化的同态 $N \to N^\wedge$，
作为 $R$-模同态是普遍单射的。
\end{lemma}

\begin{proof}
须证明对任一有限 $R$-模 $Q$，同态
$Q \otimes_R N \to Q \otimes_R N^\wedge$ 单射；参见代数，定理
\ref{algebra-theorem-universally-exact-criteria}。由于存在典范同态
$Q \otimes_R N^\wedge \to (Q \otimes_R N)^\wedge$，只须证明典范同态
$Q \otimes_R N \to (Q \otimes_R N)^\wedge$ 单射。因此可用
$Q \otimes_R N$ 替换 $N$，并且只须证明同态 $N \to N^\wedge$ 单射。

\medskip\noindent
令 $K = \Ker(N \to N^\wedge)$。由于 $N$ 是
$\prod_{\mathfrak q \in \text{Ass}(N)} N_{\mathfrak q}$ 的子模，
只须对每个 $\mathfrak q \in \text{Ass}(N)$ 证明
$K_{\mathfrak q} = 0$；参见代数，引理
\ref{algebra-lemma-zero-at-ass-zero}。取
$\mathfrak q \in \text{Ass}(N)$。由最后一个假设，存在素理想
$\mathfrak q' \supset IS + \mathfrak q$。由于 $K_{\mathfrak q}$
是 $K_{\mathfrak q'}$ 的局部化，只须证明 $K_{\mathfrak q'}$ 为零。
注意 $K = \bigcap I^nN$，故
$K_{\mathfrak q'} \subset \bigcap I^nN_{\mathfrak q'}$。
由代数，引理 \ref{algebra-lemma-intersect-powers-ideal-module-zero}，
$K_{\mathfrak q'} = 0$。
\end{proof}

\begin{lemma}
\label{flat-lemma-universally-injective-to-completion-flat}
设 $R$ 是环，$I \subset R$ 是理想，$R \to S$ 是环同态，
$N$ 是 $S$-模。假设
\begin{enumerate}
\item $R$ 是诺特环；
\item $S$ 是诺特环；
\item $N$ 是有限 $S$-模；
\item $N$ 在 $R$ 上平坦；并且
\item 对任一素理想 $\mathfrak q \subset S$，若它是
$N \otimes_R \kappa(\mathfrak p)$ 的伴随素理想，其中
$\mathfrak p = R \cap \mathfrak q$，则有 $IS + \mathfrak q \not = S$。
\end{enumerate}
则从 $N$ 到 $N$ 的 $I$-进完备化的同态 $N \to N^\wedge$，
作为 $R$-模同态是普遍单射的。
\end{lemma}

\begin{proof}
结论由引理 \ref{flat-lemma-universally-injective-to-completion} 得出，
因为代数，引理 \ref{algebra-lemma-bourbaki-fibres} 和注记
\ref{algebra-remark-bourbaki} 保证张量积 $N \otimes_R Q$
的伴随素理想集合包含于诸模
$N \otimes_R \kappa(\mathfrak p)$ 的伴随素理想集合之中。
\end{proof}




\section{投射模}
\label{flat-section-projective}

\noindent
下述引理可用于在基底上作诺特归纳以证明投射性；参见引理
\ref{flat-lemma-fibres-irreducible-flat-projective}。

\begin{lemma}
\label{flat-lemma-flat-pure-over-complete-projective}
设 $R$ 是环，$I \subset R$ 是理想，$R \to S$ 是环同态，
$N$ 是 $S$-模。假设
\begin{enumerate}
\item $R$ 是诺特环且 $I$-进完备；
\item $R \to S$ 有限型；
\item $N$ 是有限 $S$-模；
\item $N$ 在 $R$ 上平坦；
\item $N/IN$ 作为 $R/I$-模是投射的；并且
\item 对任一素理想 $\mathfrak q \subset S$，若它是
$N \otimes_R \kappa(\mathfrak p)$ 的伴随素理想，其中
$\mathfrak p = R \cap \mathfrak q$，则有 $IS + \mathfrak q \not = S$。
\end{enumerate}
则 $N$ 作为 $R$-模是投射的。
\end{lemma}

\begin{proof}
由引理 \ref{flat-lemma-universally-injective-to-completion-flat}，同态
$N \to N^\wedge$ 普遍单射。由引理 \ref{flat-lemma-lift-ML}，模
$N^\wedge$ 是 Mittag-Leffler 模。再由代数，引理
\ref{algebra-lemma-pure-submodule-ML}，$N$ 是 Mittag-Leffler 模。
因此 $N$ 作为 $R$-模可数生成、平坦且为 Mittag-Leffler 模，
故由代数，引理 \ref{algebra-lemma-countgen-projective}，它是投射模。
\end{proof}

\begin{lemma}
\label{flat-lemma-fibres-irreducible-flat-projective}
设 $R$ 是环，$R \to S$ 是环同态。假设
\begin{enumerate}
\item $R$ 是诺特环；
\item $R \to S$ 有限型且平坦；并且
\item 每个纤维环 $S \otimes_R \kappa(\mathfrak p)$ 在
$\kappa(\mathfrak p)$ 上几何整。
\end{enumerate}
则 $S$ 作为 $R$-模是投射的。
\end{lemma}

\begin{proof}
考虑集合
$$
\{I \subset R \mid S/IS\text{ 作为 }R/I\text{-模不投射}\}
$$
须证明此集合为空。反设它非空，则它含有极大元 $I$。
令 $J = \sqrt{I}$ 为其根理想。若 $I \not = J$，则 $S/JS$
作为 $R/J$-模是投射的，$S/IS$ 在 $R/I$ 上平坦，并且 $J/I$
是 $R/I$ 中的幂零理想。应用代数，引理
\ref{algebra-lemma-lift-projective}，可知 $S/IS$ 是投射
$R/I$-模，矛盾。因此可设 $I$ 是根理想。换言之，问题归结为在下述情形
证明引理：$R$ 是既约环，并且对 $R$ 的每个非零理想 $I$，
$S/IS$ 都是投射 $R/I$-模。

\medskip\noindent
假设 $R$ 是既约环，并且对 $R$ 的每个非零理想 $I$，$S/IS$
都是投射 $R/I$-模。由一般平坦性，即代数，引理
\ref{algebra-lemma-generic-flatness-Noetherian}
（应用于为整环的局部化 $R_g$），或更一般的代数，引理
\ref{algebra-lemma-generic-flatness-reduced}，存在非零元 $f \in R$，
使得 $S_f$ 是自由 $R_f$-模。将 $R$ 的 $(f)$-进完备化记为
$R^\wedge = \lim R/(f^n)$。注意环同态
$$
R \longrightarrow R_f \times R^\wedge
$$
忠实平坦；参见代数，引理 \ref{algebra-lemma-completion-flat}。
因此，由投射性的忠实平坦下降（参见代数，定理
\ref{algebra-theorem-ffdescent-projectivity}），只须证明
$S \otimes_R R^\wedge$ 是投射 $R^\wedge$-模。为此将使用引理
\ref{flat-lemma-flat-pure-over-complete-projective} 的判据。首先注意，
$S/fS = (S \otimes_R R^\wedge)/f(S \otimes_R R^\wedge)$
是投射 $R/(f)$-模，并且 $S \otimes_R R^\wedge$ 作为相应性质的基变换，
在 $R^\wedge$ 上平坦且有限型。其次，设 $\mathfrak p^\wedge$ 是
$R^\wedge$ 的素理想，令 $\mathfrak p \subset R$ 为 $R$ 中相应的素理想。
由于 $R \to S$ 的纤维环几何整，任一基变换的纤维环也如此。因此
% P05-EN-CH38-0022: frozen source has the number-agreement typo “is a prime ideals lying”.
$\mathfrak q^\wedge = \mathfrak p^\wedge(S \otimes_R R^\wedge)$
是位于 $\mathfrak p^\wedge$ 上方的素理想，并且是
$S \otimes_R \kappa(\mathfrak p^\wedge)$ 的唯一伴随素理想。因此，只须有
$f(S \otimes_R R^\wedge) + \mathfrak q^\wedge \not = S \otimes_R R^\wedge$
便得证。这确实成立：$f$ 位于 $f$-进完备环 $R^\wedge$ 的雅可比根中，
所以 $\mathfrak p^\wedge + fR^\wedge \not = R^\wedge$；又因为各纤维非空
（不可约空间非空），$R^\wedge \to S \otimes_R R^\wedge$ 在谱上满射。
\end{proof}

\begin{lemma}
\label{flat-lemma-fibres-irreducible-flat-projective-nonnoetherian}
设 $R$ 是环，$R \to S$ 是环同态。假设
\begin{enumerate}
\item $R \to S$ 有限表示且平坦；并且
\item 每个纤维环 $S \otimes_R \kappa(\mathfrak p)$ 在
$\kappa(\mathfrak p)$ 上几何整。
\end{enumerate}
则 $S$ 作为 $R$-模是投射的。
\end{lemma}

\begin{proof}
可以找到环的余笛卡儿图
$$
\xymatrix{
S_0 \ar[r] & S \\
R_0 \ar[u] \ar[r] & R \ar[u]
}
$$
其中 $R_0$ 在 $\mathbf{Z}$ 上有限型，且同态 $R_0 \to S_0$
有限型、平坦并具有几何整纤维；参见《态射进阶》，引理
\ref{more-morphisms-lemma-Noetherian-approximation-flat}、
\ref{more-morphisms-lemma-Noetherian-approximation-geometrically-reduced},
\ref{more-morphisms-lemma-Noetherian-approximation-geometrically-irreducible},
以及 \ref{more-morphisms-lemma-Noetherian-approximation-combine}。
由引理 \ref{flat-lemma-fibres-irreducible-flat-projective}，$S_0$ 是投射
$R_0$-模。因此 $S = S_0 \otimes_{R_0} R$ 是投射 $R$-模；
参见代数，引理 \ref{algebra-lemma-ascend-properties-modules}。
\end{proof}

\begin{remark}
\label{flat-remark-how-in-RG}
引理 \ref{flat-lemma-fibres-irreducible-flat-projective-nonnoetherian}
是本章后续结果发展中的关键一步。该引理在 \cite{GruRay} 中的对应结果是
\cite[I 命题 3.3.1]{GruRay}：若 $R \to S$ 光滑且纤维几何整，
则 $S$ 作为 $R$-模是投射的。这是引理
\ref{flat-lemma-fibres-irreducible-flat-projective-nonnoetherian} 的特殊情形；
不过我们稍后终究会改进该引理，所以此时拥有一个更强结果并无太大收益。
下面简述 \cite{GruRay} 中给出的证明。
\begin{enumerate}
\item 先如上归结到 $R$ 为诺特环的情形。
\item 投射性沿忠实平坦环同态下降；参见代数，定理
\ref{algebra-theorem-ffdescent-projectivity}。因此可在 $R$ 的 fppf
拓扑下局部工作，从而可设 $R \to S$ 有截面 $\sigma : S \to R$
（只须使用基变换到 $S$ 的通常技巧）。置 $I = \Ker(S \to R)$。
\item 在 $R$ 上再作一些局部化，可设 $I/I^2$ 是自由 $R$-模，且 $S$
关于 $I$ 的完备化 $S^\wedge$ 同构于 $R[[t_1, \ldots, t_n]]$；
参见《态射》，引理 \ref{morphisms-lemma-section-smooth-morphism}。
这里使用了 $R \to S$ 光滑这一事实。
\item 为证明 $S$ 作为 $R$-模是投射的，只须证明 $S$ 作为 $R$-模
平坦、可数生成且为 Mittag-Leffler 模；参见代数，引理
\ref{algebra-lemma-countgen-projective}。前两个性质显然，故只须证明
$S$ 作为 $R$-模是 Mittag-Leffler 模。由代数，引理
\ref{algebra-lemma-power-series-ML}，模 $R[[t_1, \ldots, t_n]]$
在 $R$ 上为 Mittag-Leffler 模。因此，代数，引理
\ref{algebra-lemma-pure-submodule-ML} 表明，只须证明
% P05-EN-CH38-0023: frozen source has the grammatical surplus article in “show that the S to S^wedge”.
$S \to S^\wedge$ 作为 $R$-模同态普遍单射。
\item 应用引理 \ref{flat-lemma-base-change-universally-flat} 可知
$S \to S^\wedge$ 是 $R$-普遍单射的。事实上，由于 $R \to S$
具有几何整纤维，任一纤维环的每个伴随点都只是该纤维环的一般点，
而这个一般点属于 $\Spec(S^\wedge) \to \Spec(S)$ 的像。
\end{enumerate}
刚才概述的证明与导向引理
\ref{flat-lemma-fibres-irreducible-flat-projective-nonnoetherian} 之证明的论证发展
彼此类似。在二者中，完备化都起到关键作用；在两处，具有几何整纤维这一
假设都保证从 $S$ 到 $S$ 的完备化的同态是 $R$-普遍单射的。
\end{remark}













\section{有限型平坦模，第一部分}
\label{flat-section-finite-type-flat-I}

\noindent
在某些情形下，给定有限表示环同态 $R \to S$ 与有限 $S$-模 $N$，
$N$ 在 $R$ 上平坦蕴含 $N$ 有限表示。本节将证明这一结论“逐点”成立。
需要指出，命题 \ref{flat-proposition-finite-type-flat-at-point} 的第一个证明使用
第 \ref{flat-section-local-structure-module} 节的几何结果，但不使用完全
d\'evissage 的存在性。

\begin{lemma}
\label{flat-lemma-induction-step}
设 $(R, \mathfrak m)$ 是局部环，$R \to S$ 是具有几何整纤维的有限表示
平坦环同态。记 $\mathfrak p = \mathfrak mS$。设
$\mathfrak q \subset S$ 是位于 $\mathfrak m$ 上方的素理想，$N$ 是有限
$S$-模。存在 $r \geq 0$ 以及 $S$-模同态
$$
\alpha : S^{\oplus r} \longrightarrow N
$$
使得
$\alpha : \kappa(\mathfrak p)^{\oplus r} \to N \otimes_S \kappa(\mathfrak p)$
是同构。对任一这样的 $\alpha$，下列条件等价：
\begin{enumerate}
\item $N_{\mathfrak q}$ 在 $R$ 上平坦；
\item $\alpha$ 是 $R$-普遍单射的，且
$\Coker(\alpha)_{\mathfrak q}$ 在 $R$ 上平坦；
\item $\alpha$ 单射，且 $\Coker(\alpha)_{\mathfrak q}$ 在 $R$ 上平坦；
\item $\alpha_{\mathfrak p}$ 是同构，且
$\Coker(\alpha)_{\mathfrak q}$ 在 $R$ 上平坦；并且
\item $\alpha_{\mathfrak q}$ 单射，且
$\Coker(\alpha)_{\mathfrak q}$ 在 $R$ 上平坦。
\end{enumerate}
\end{lemma}

\begin{proof}
为得到 $\alpha$，置
$r = \dim_{\kappa(\mathfrak p)} N \otimes_S \kappa(\mathfrak p)$，并选取
$x_1, \ldots, x_r \in N$，使其构成
$N \otimes_S \kappa(\mathfrak p)$ 的一组基。定义
$\alpha(s_1, \ldots, s_r) = \sum s_i x_i$，便证明了存在性。

\medskip\noindent
固定一个 $\alpha$。最有意义的蕴含是 (1) $\Rightarrow$ (2)，先来证明它。
假设 (1)。因为 $S/\mathfrak mS$ 是分式域为
$\kappa(\mathfrak p)$ 的整环，可知
$(S/\mathfrak mS)^{\oplus r} \to
N_{\mathfrak p}/\mathfrak mN_{\mathfrak p} = N \otimes_S \kappa(\mathfrak p)$
单射。因此，由引理 \ref{flat-lemma-universally-injective-local} 和
\ref{flat-lemma-fibres-irreducible-flat-projective-nonnoetherian}.
同态 $S^{\oplus r} \to N_{\mathfrak p}$ 是 $R$-普遍单射的。
于是 $S^{\oplus r} \to N$ 是 $R$-普遍单射的；参见代数，引理
\ref{algebra-lemma-universally-injective-permanence}。其局部化
$\alpha_{\mathfrak q}$ 也为 $R$-普遍单射；参见代数，引理
\ref{algebra-lemma-universally-injective-localize}。由代数，引理
\ref{algebra-lemma-ui-flat-domain}，可知
$\Coker(\alpha)_{\mathfrak q}$ 在 $R$ 上平坦。

\medskip\noindent
蕴含 (2) $\Rightarrow$ (3) 是立即的。若 (3) 成立，则
$\alpha_{\mathfrak p}$ 作为单射模同态的局部化是单射的。由 Nakayama
引理（代数，引理 \ref{algebra-lemma-NAK}），$\alpha_{\mathfrak p}$
也满射。因此 (3) $\Rightarrow$ (4)。若 (4) 成立，则
$\alpha_{\mathfrak p}$ 是同构，而
$S_{\mathfrak q} \to S_{\mathfrak p}$ 单射，故
% P05-EN-CH38-0024: frozen proof says alpha is injective although this argument establishes the localized map alpha_q required by (5).
$\alpha$ 单射。事实上，结合引理 \ref{flat-lemma-invert-universally-injective}
与 \ref{flat-lemma-fibres-irreducible-flat-projective-nonnoetherian}，可知
$S \setminus \mathfrak p$ 中的元素在 $S$ 上都是非零因子。
因此 (4) $\Rightarrow$ (5)。最后，若 (5) 成立，则
$N_{\mathfrak q}$ 作为平坦模的扩张在 $R$ 上平坦；参见代数，引理
\ref{algebra-lemma-flat-ses}。故 (5) $\Rightarrow$ (1)，证明完成。
\end{proof}

\begin{lemma}
\label{flat-lemma-complete-devissage-flat-finite-type-module}
设 $(R, \mathfrak m)$ 是局部环，$R \to S$ 是有限表示环同态，
$N$ 是有限 $S$-模，$\mathfrak q$ 是 $S$ 中位于 $\mathfrak m$
上方的素理想。假设 $N_{\mathfrak q}$ 在 $R$ 上平坦，并且在
$\mathfrak q$ 处存在 $N/S/R$ 的完全 d\'evissage。则 $N$ 是有限表示
$S$-模，作为 $R$-模自由，并且存在 $R$-模同构
$$
N \cong B_1^{\oplus r_1} \oplus \ldots \oplus B_n^{\oplus r_n}
$$
其中每个 $B_i$ 都是具有几何不可约纤维的光滑 $R$-代数。
\end{lemma}

\begin{proof}
设 $(A_i, B_i, M_i, \alpha_i, \mathfrak q_i)_{i = 1, \ldots, n}$
为给定的完全 d\'evissage。对 $n$ 作归纳以证明引理。注意，$N$ 作为
$S$-模有限表示，当且仅当 $M_1$ 作为 $B_1$-模有限表示；参见注记
\ref{flat-remark-finite-presentation}。
% P05-EN-CH38-0025: frozen source has the article error “an B_1-module”.
还要注意，作为 $R$-模有 $N_{\mathfrak q} \cong (M_1)_{\mathfrak q_1}$，
因为：(a) $N_{\mathfrak q} \cong (M_1)_{\mathfrak q'_1}$，其中
$\mathfrak q'_1$ 是 $A_1$ 中位于 $\mathfrak q_1$ 上方的唯一素理想；
(b) 由代数，引理 \ref{algebra-lemma-unique-prime-over-localize-below}，
$(A_1)_{\mathfrak q'_1} = (A_1)_{\mathfrak q_1}$；从而
(c) $(M_1)_{\mathfrak q'_1} \cong (M_1)_{\mathfrak q_1}$。
因此 $(M_1)_{\mathfrak q_1}$ 是平坦 $R$-模。为证明引理，故可用
$(B_1, M_1)$ 替换 $(S, N)$。由引理 \ref{flat-lemma-induction-step}，同态
$\alpha_1 : B_1^{\oplus r_1} \to M_1$ 是 $R$-普遍单射的，并且
% P05-EN-CH38-0026: frozen proof retains q after replacing the active prime by q_1.
$\Coker(\alpha_1)_{\mathfrak q}$ 在 $R$ 上平坦。注意
$(A_i, B_i, M_i, \alpha_i, \mathfrak q_i)_{i = 2, \ldots, n}$ 是
$\mathfrak q_1$ 处 $\Coker(\alpha_1)/B_1/R$ 的完全 d\'evissage。
因此，归纳假设说明 $\Coker(\alpha_1)$ 是有限表示 $B_1$-模，
作为 $R$-模自由，并且具有引理中的分解。这又说明 $M_1$ 是有限表示
$B_1$-模；参见代数，引理 \ref{algebra-lemma-extension}。进而，作为
$R$-模有
$M_1 \cong B_1^{\oplus r_1} \oplus \Coker(\alpha_1)$，
从而得到引理中的分解。最后，由引理
\ref{flat-lemma-fibres-irreducible-flat-projective-nonnoetherian}，$B_1$
作为 $R$-模是投射的；再由代数，定理
\ref{algebra-theorem-projective-free-over-local-ring}，它作为 $R$-模自由。
证明完成。
\end{proof}

\begin{proposition}
\label{flat-proposition-finite-type-flat-at-point}
设 $f : X \to S$ 是概形态射，$\mathcal{F}$ 是 $X$ 上的拟凝聚层，
$x \in X$ 的像为 $s \in S$。假设
\begin{enumerate}
\item $f$ 局部有限表示；
\item $\mathcal{F}$ 有限型；并且
\item $\mathcal{F}$ 在点 $x$ 处于 $S$ 上平坦。
\end{enumerate}
则存在初等 \'etale 邻域 $(S', s') \to (S, s)$ 以及开子概形
$$
V \subset X \times_S \Spec(\mathcal{O}_{S', s'})
$$
它包含 $X \times_S \Spec(\mathcal{O}_{S', s'})$ 中映到 $x$ 的唯一点，
并且 $\mathcal{F}$ 在 $V$ 上的拉回是有限表示 $\mathcal{O}_V$-模，
且在 $\mathcal{O}_{S', s'}$ 上平坦。
\end{proposition}

\begin{proof}[第一个证明]
这个证明较长，但不使用完全 d\'evissage 的存在性。问题在 $x$ 与 $s$
附近是局部的，故可设 $X$ 和 $S$ 仿射。在证明过程中，我们会有限次地
% P05-EN-CH38-0027: frozen source is missing “a finite number of” in “we will finitely many times replace”.
用 $(S, s)$ 的初等 \'etale 邻域替换 $S$。此时目标是在这样的替换之后，
找到包含 $x$ 的开集
$V \subset X \times_S \Spec(\mathcal{O}_{S, s})$，使得
$\mathcal{F}|_V$ 在 $S$ 上平坦且有限表示。当然，在证明的任意阶段，也可
用 $\Spec(\mathcal{O}_{S, s})$ 替换 $S$；也就是说，可设 $S$ 为局部概形。
对整数 $n = \dim_x(\text{Supp}(\mathcal{F}_s))$ 作归纳以证明命题。

\medskip\noindent
可以选择
\begin{enumerate}
\item 初等 \'etale 邻域 $g : (X', x') \to (X, x)$ 与
$e : (S', s') \to (S, s)$；
\item 交换图
$$
\xymatrix{
X \ar[dd]_f & X' \ar[dd] \ar[l]^g & Z' \ar[l]^i \ar[d]^\pi \\
& & Y' \ar[d]^h \\
S & S' \ar[l]_e & S' \ar@{=}[l]
}
$$
\item 点 $z' \in Z'$，满足 $i(z') = x'$、$y' = \pi(z')$、
$h(y') = s'$；
\item 有限型拟凝聚 $\mathcal{O}_{Z'}$-模 $\mathcal{G}$；
\end{enumerate}
如引理 \ref{flat-lemma-elementary-devissage} 中所述。按照证明第一段的说明，
将用 $\Spec(\mathcal{O}_{S', s'})$ 替换 $S$。考虑图
$$
\xymatrix{
X_{\mathcal{O}_{S', s'}} \ar[ddr]_f &
X'_{\mathcal{O}_{S', s'}} \ar[dd] \ar[l]^g &
Z'_{\mathcal{O}_{S', s'}} \ar[l]^i \ar[d]^\pi \\
& & Y'_{\mathcal{O}_{S', s'}} \ar[dl]^h \\
& \Spec(\mathcal{O}_{S', s'})
}
$$
这里通过 $\Spec(\mathcal{O}_{S', s'}) \to S'$ 对 $S'$ 上的概形
$X', Z', Y'$ 作了基变换，并通过
$\Spec(\mathcal{O}_{S', s'}) \to S$ 对 $S$ 上的概形 $X$ 作了基变换。
$g$ 仍为 \'etale；参见引理 \ref{flat-lemma-etale-at-point}。分别用
$X_{\mathcal{O}_{S', s'}}$、$X'_{\mathcal{O}_{S', s'}}$、
$Z'_{\mathcal{O}_{S', s'}}$、$Y'_{\mathcal{O}_{S', s'}}$
替换 $X$、$X'$、$Z'$、$Y'$ 后，可设已有引理
\ref{flat-lemma-elementary-devissage} 中的图，并且 $S = S'$ 是闭点为 $s$
的局部概形。由引理 \ref{flat-lemma-devissage-finite-presentation} 和
\ref{flat-lemma-devissage-flat}，对 $Y' \to S$、层 $\pi_*\mathcal{G}$ 与点
$y'$ 的结论蕴含对 $X \to S$、$\mathcal{F}$ 与 $x$ 的结论。因此可设
$S$ 局部，且 $X \to S$ 是仿射概形之间的光滑态射，其纤维几何不可约、
维数为 $n$。

\medskip\noindent
归纳的初始情形：$n = 0$。因为 $X \to S$ 光滑，且纤维几何不可约、
维数为 $0$，所以 $X \to S$ 是开浸入；参见下降章，引理
\ref{descent-lemma-universally-injective-etale-open-immersion}。又因为 $S$
局部且闭点属于 $X \to S$ 的像，故 $X = S$。于是 $\mathcal{F}$
对应于有限平坦 $\mathcal{O}_{S, s}$ 模。在此情形，结论由代数，引理
\ref{algebra-lemma-finite-flat-local} 得出；该引理说明 $\mathcal{F}$
事实上有限自由。

\medskip\noindent
归纳步骤。假设只要闭纤维中支撑的维数 $< n$，结论就成立。记
$S = \Spec(A)$、$X = \Spec(B)$，并对某个 $B$-模 $N$ 记
$\mathcal{F} = \widetilde{N}$。注意 $A$ 是局部环，将其极大理想记为
$\mathfrak m$。由于 $X \to S$ 的纤维几何不可约，
$\mathfrak p = \mathfrak mB$ 是位于 $\mathfrak m$ 上方的唯一极小素理想。
最后，令 $\mathfrak q \subset B$ 为与 $x$ 对应的素理想。由引理
\ref{flat-lemma-induction-step}
可选择同态
$$
\alpha : B^{\oplus r} \to N
$$
使得 $\kappa(\mathfrak p)^{\oplus r} \to N \otimes_B \kappa(\mathfrak p)$
是同构。此外，因为 $N_{\mathfrak q}$ 在 $A$ 上平坦，该引理还说明
$\alpha$ 单射，且 $\Coker(\alpha)_{\mathfrak q}$ 在 $A$ 上平坦。
置 $Q = \Coker(\alpha)$。注意 $Q/\mathfrak mQ$ 的支撑不包含
$\mathfrak p$，所以必有
$\dim_{\mathfrak q}(\text{Supp}(Q/\mathfrak mQ)) < n$。综合关于 $Q$
的全部已知信息，可将归纳假设应用于 $Q$。因此存在初等 \'etale 态射
% P05-EN-CH38-0028: frozen source writes (S', s) although the local ring on the next line is at s'.
$(S', s) \to (S, s)$，使得结论对 $B \otimes_A A'$ 上的
$Q \otimes_A A'$ 成立，其中 $A' = \mathcal{O}_{S', s'}$。
用 $A'$ 替换 $A$ 后，有正合序列
$$
0 \to B^{\oplus r} \to N \to Q \to 0
$$
（这里使用了上面提到的 $\alpha$ 单射），其各项均为有限 $B$-模；
还得到元素 $g \in B$、$g \not \in \mathfrak q$，使得 $Q_g$
在 $B_g$ 上有限表示且在 $A$ 上平坦。由于局部化正合，序列
$$
0 \to B_g^{\oplus r} \to N_g \to Q_g \to 0
$$
仍正合。因为 $B_g$ 和 $Q_g$ 在 $A$ 上平坦，所以 $N_g$ 在 $A$
上平坦；参见代数，引理 \ref{algebra-lemma-flat-ses}。又因为 $B_g$
和 $Q_g$ 在 $B_g$ 上有限表示，$N_g$ 也如此；参见代数，引理
\ref{algebra-lemma-extension}。
\end{proof}

\begin{proof}[第二个证明]
应用命题 \ref{flat-proposition-existence-complete-at-x}，得到带点概形的交换图
$$
\xymatrix{
(X, x) \ar[d] & (X', x') \ar[l]^g \ar[d] \\
(S, s) & (S', s') \ar[l]
}
$$
其中水平箭头都是初等 \'etale 邻域，并且
% P05-EN-CH38-0029: frozen source locates the complete devissage of g^*F on X' at x rather than x'.
$g^*\mathcal{F}/X'/S'$ 在 $x$ 处有完全 d\'evissage（特别地，$S'$ 和
$X'$ 仿射）。由《态射》，引理 \ref{morphisms-lemma-flat-permanence}，
$g^*\mathcal{F}$ 在点 $x'$ 处于 $S$ 上平坦；再由引理
\ref{flat-lemma-etale-flat-up-down}，它在点 $x'$ 处于 $S'$ 上平坦。
由注记 \ref{flat-remark-same-notion}，可得
$$
\Gamma(X', g^*\mathcal{F})/
\Gamma(X', \mathcal{O}_{X'})/
\Gamma(S', \mathcal{O}_{S'})
$$
在 $\Gamma(X', \mathcal{O}_{X'})$ 中与 $x'$ 对应的素理想处有完全
d\'evissage。可将这个完全 d\'evissage 基变换到
$\Gamma(S', \mathcal{O}_{S'})$ 在与 $s'$ 对应的素理想处的局部环
$\mathcal{O}_{S', s'}$。因此，引理
\ref{flat-lemma-complete-devissage-flat-finite-type-module} 说明
$$
% P05-EN-CH38-0030: frozen display switches from g^*F to an undefined F'.
\Gamma(X', \mathcal{F}')
\otimes_{\Gamma(S', \mathcal{O}_{S'})}
\mathcal{O}_{S', s'}
$$
在 $\mathcal{O}_{S', s'}$ 上平坦，并且在
$\Gamma(X', \mathcal{O}_{X'})
\otimes_{\Gamma(S', \mathcal{O}_{S'})} \mathcal{O}_{S', s'}$ 上有限表示。
换言之，$\mathcal{F}$ 在
$X' \times_{S'} \Spec(\mathcal{O}_{S', s'})$ 上的限制有限表示，且在
$\mathcal{O}_{S', s'}$ 上平坦。
% P05-EN-CH38-0031: frozen prose names F on X' rather than the previously defined pullback g^*F.
由于态射
$X' \times_{S'} \Spec(\mathcal{O}_{S', s'})
\to X \times_S \Spec(\mathcal{O}_{S', s'})$
是 \'etale 的（引理 \ref{flat-lemma-etale-at-point}），其像
$V \subset X \times_S \Spec(\mathcal{O}_{S', s'})$ 是开子概形；由
\'etale 下降，$\mathcal{F}$ 在 $V$ 上的限制有限表示，且在
$\mathcal{O}_{S', s'}$ 上平坦。这里使用了《态射》，引理
\ref{morphisms-lemma-etale-open}，下降章，引理
\ref{descent-lemma-finite-presentation-descends}，以及《态射》，引理
\ref{morphisms-lemma-flat-permanence}。
\end{proof}

\begin{lemma}
\label{flat-lemma-open-in-fibre-where-flat}
设 $f : X \to S$ 是局部有限型概形态射，$\mathcal{F}$ 是有限型拟凝聚
$\mathcal{O}_X$-模，且 $s \in S$。则集合
$$
\{x \in X_s \mid \mathcal{F}\text{ 在 }x\text{ 处于 }S\text{ 上平坦}\}
$$
在纤维 $X_s$ 中是开集。
\end{lemma}

\begin{proof}
% P05-EN-CH38-0032: frozen proof refers to U although the displayed flat locus is not named U in the statement.
设 $x \in U$。按命题 \ref{flat-proposition-finite-type-flat-at-point}，选择初等
\'etale 邻域 $(S', s') \to (S, s)$ 与开集
$V \subset X \times_S \Spec(\mathcal{O}_{S', s'})$。注意，由于
$\kappa(s) = \kappa(s')$，有 $X_{s'} = X_s$。若
$x' \in V \cap X_{s'}$，则 $\mathcal{F}$ 在 $X \times_S S'$ 上的拉回
在点 $x'$ 处于 $S'$ 上平坦。因此，$\mathcal{F}$ 在点 $x'$ 处于
$S$ 上平坦；参见《态射》，引理 \ref{morphisms-lemma-flat-permanence}。
换言之，$X_s \cap V \subset U$ 是 $U$ 中点 $x$ 的开邻域。
\end{proof}

\begin{lemma}
\label{flat-lemma-finite-type-flat-at-point}
设 $f : X \to S$ 是概形态射，$\mathcal{F}$ 是 $X$ 上的拟凝聚层，
$x \in X$ 的像为 $s \in S$。假设
\begin{enumerate}
\item $f$ 局部有限型；
\item $\mathcal{F}$ 有限型；并且
\item $\mathcal{F}$ 在点 $x$ 处于 $S$ 上平坦。
\end{enumerate}
则存在初等 \'etale 邻域 $(S', s') \to (S, s)$ 以及开子概形
$$
V \subset X \times_S \Spec(\mathcal{O}_{S', s'})
$$
它包含 $X \times_S \Spec(\mathcal{O}_{S', s'})$ 中映到 $x$ 的唯一点，
并且 $\mathcal{F}$ 在 $V$ 上的拉回在 $\mathcal{O}_{S', s'}$ 上平坦。
\end{lemma}

\begin{proof}
（本引理与命题 \ref{flat-proposition-finite-type-flat-at-point} 的唯一区别是，
这里不假设 $f$ 有限表示。）问题在 $X$ 和 $S$ 上是局部的，故可设
$X$、$S$ 仿射。记 $X = \Spec(B)$、$S = \Spec(A)$，并写成
$B = A[x_1, \ldots, x_n]/I$。换言之，得到闭浸入
$i : X \to \mathbf{A}^n_S$。记 $t = i(x) \in \mathbf{A}^n_S$。
将命题 \ref{flat-proposition-finite-type-flat-at-point} 应用于
$\mathbf{A}^n_S \to S$、层 $i_*\mathcal{F}$ 与点 $t$，得到初等
\'etale 邻域 $(S', s') \to (S, s)$ 以及开子概形
$$
W \subset \mathbf{A}^n_{\mathcal{O}_{S', s'}}
$$
使得 $i_*\mathcal{F}$ 在 $W$ 上的拉回在 $\mathcal{O}_{S', s'}$ 上平坦。
这意味着
$V := W \cap \big(X \times_S \Spec(\mathcal{O}_{S', s'})\big)$
就是所需的开子概形。
\end{proof}

\begin{lemma}
\label{flat-lemma-finite-type-flat-along-fibre}
设 $f : X \to S$ 是概形态射，$\mathcal{F}$ 是 $X$ 上的拟凝聚层，
$s \in S$。假设
\begin{enumerate}
\item $f$ 有限表示；
\item $\mathcal{F}$ 有限型；并且
\item $\mathcal{F}$ 在纤维 $X_s$ 的每一点处于 $S$ 上平坦。
\end{enumerate}
则存在初等 \'etale 邻域 $(S', s') \to (S, s)$ 以及开子概形
$$
V \subset X \times_S \Spec(\mathcal{O}_{S', s'})
$$
它包含纤维 $X_s = X \times_S s'$，并且 $\mathcal{F}$ 在 $V$ 上的拉回
是有限表示 $\mathcal{O}_V$-模，且在 $\mathcal{O}_{S', s'}$ 上平坦。
\end{lemma}

\begin{proof}
对每个点 $x \in X_s$，可应用命题
\ref{flat-proposition-finite-type-flat-at-point}，找到初等 \'etale 邻域
$(S_x, s_x) \to (S, s)$ 与开集
$V_x \subset X \times_S \Spec(\mathcal{O}_{S_x, s_x})$，使得
$x \in X_s = X \times_S s_x$ 包含于 $V_x$，并且 $\mathcal{F}$ 在
$V_x$ 上的拉回是有限表示 $\mathcal{O}_{V_x}$-模，且在
$\mathcal{O}_{S_x, s_x}$ 上平坦。特别地，可将纤维 $(V_x)_{s_x}$
视为 $X_s$ 中点 $x$ 的开邻域。因为 $X_s$ 拟紧，可找到有限个点
$x_1, \ldots, x_n \in X_s$，使得 $X_s$ 是各
$(V_{x_i})_{s_{x_i}}$ 的并。选择支配每个邻域
$(S_{x_i}, s_{x_i})$ 的初等 \'etale 邻域
$(S' , s') \to (S, s)$；参见《态射进阶》，引理
\ref{more-morphisms-lemma-elementary-etale-neighbourhoods}。
置 $V = \bigcup V_i$，其中 $V_i$ 是开集 $V_{x_i}$ 在态射
% P05-EN-CH38-0033: frozen source has the number-agreement error “V_i is the inverse images”.
$$
X \times_S \Spec(\mathcal{O}_{S', s'})
\longrightarrow
X \times_S \Spec(\mathcal{O}_{S_{x_i}, s_{x_i}})
$$
下的逆像。按构造，$V$ 包含 $X_s$；同样按构造，$\mathcal{F}$ 在
$V$ 上的拉回是有限表示 $\mathcal{O}_V$-模，且在
$\mathcal{O}_{S', s'}$ 上平坦。
\end{proof}

\begin{lemma}
\label{flat-lemma-finite-type-flat-along-fibre-variant}
设 $f : X \to S$ 是概形态射，$\mathcal{F}$ 是 $X$ 上的拟凝聚层，
$s \in S$。假设
\begin{enumerate}
\item $f$ 有限型；
\item $\mathcal{F}$ 有限型；并且
\item $\mathcal{F}$ 在纤维 $X_s$ 的每一点处于 $S$ 上平坦。
\end{enumerate}
则存在初等 \'etale 邻域 $(S', s') \to (S, s)$ 以及开子概形
$$
V \subset X \times_S \Spec(\mathcal{O}_{S', s'})
$$
它包含纤维 $X_s = X \times_S s'$，并且 $\mathcal{F}$ 在 $V$ 上的拉回
在 $\mathcal{O}_{S', s'}$ 上平坦。
\end{lemma}

\begin{proof}
（本引理与引理 \ref{flat-lemma-finite-type-flat-along-fibre} 的唯一区别是，
这里不假设 $f$ 有限表示。）对每个点 $x \in X_s$，可应用引理
\ref{flat-lemma-finite-type-flat-at-point}，找到初等 \'etale 邻域
$(S_x, s_x) \to (S, s)$ 与开集
$V_x \subset X \times_S \Spec(\mathcal{O}_{S_x, s_x})$，使得
$x \in X_s = X \times_S s_x$ 包含于 $V_x$，并且 $\mathcal{F}$ 在
$V_x$ 上的拉回在 $\mathcal{O}_{S_x, s_x}$ 上平坦。特别地，可将纤维
$(V_x)_{s_x}$ 视为 $X_s$ 中点 $x$ 的开邻域。因为 $X_s$ 拟紧，
可找到有限个点 $x_1, \ldots, x_n \in X_s$，使得 $X_s$ 是各
$(V_{x_i})_{s_{x_i}}$ 的并。选择支配每个邻域
$(S_{x_i}, s_{x_i})$ 的初等 \'etale 邻域
$(S' , s') \to (S, s)$；参见《态射进阶》，引理
\ref{more-morphisms-lemma-elementary-etale-neighbourhoods}。
置 $V = \bigcup V_i$，其中 $V_i$ 是开集 $V_{x_i}$ 在态射
% P05-EN-CH38-0034: frozen source repeats the number-agreement error “V_i is the inverse images”.
$$
X \times_S \Spec(\mathcal{O}_{S', s'})
\longrightarrow
X \times_S \Spec(\mathcal{O}_{S_{x_i}, s_{x_i}})
$$
下的逆像。按构造，$V$ 包含 $X_s$；同样按构造，$\mathcal{F}$ 在
$V$ 上的拉回在 $\mathcal{O}_{S', s'}$ 上平坦。
\end{proof}

\begin{lemma}
\label{flat-lemma-finite-type-flat-at-point-X}
设 $S$ 是概形，$X$ 在 $S$ 上局部有限型，$x \in X$ 的像为
$s \in S$。若 $X$ 在点 $x$ 处于 $S$ 上平坦，则存在初等
\'etale 邻域 $(S', s') \to (S, s)$ 以及开子概形
$$
V \subset X \times_S \Spec(\mathcal{O}_{S', s'})
$$
它包含 $X \times_S \Spec(\mathcal{O}_{S', s'})$ 中映到 $x$ 的唯一点，
并且 $V \to \Spec(\mathcal{O}_{S', s'})$ 平坦且有限表示。
\end{lemma}

\begin{proof}
问题在 $X$ 和 $S$ 上是局部的，故可设 $X$、$S$ 仿射。记
$X = \Spec(B)$、$S = \Spec(A)$，并写成
$B = A[x_1, \ldots, x_n]/I$。换言之，得到闭浸入
$i : X \to \mathbf{A}^n_S$。记 $t = i(x) \in \mathbf{A}^n_S$。
将命题 \ref{flat-proposition-finite-type-flat-at-point} 应用于
$\mathbf{A}^n_S \to S$、层 $\mathcal{F} = i_*\mathcal{O}_X$ 与点 $t$，
得到初等 \'etale 邻域 $(S', s') \to (S, s)$ 以及开子概形
$$
W \subset \mathbf{A}^n_{\mathcal{O}_{S', s'}}
$$
使得 $i_*\mathcal{O}_X$ 的拉回平坦且有限表示。这意味着
$V := W \cap \big(X \times_S \Spec(\mathcal{O}_{S', s'})\big)$
就是所需的开子概形。
\end{proof}

\begin{lemma}
\label{flat-lemma-finite-type-flat-at-point-local}
设 $f : X \to S$ 是局部有限表示态射，$\mathcal{F}$ 是有限型拟凝聚
$\mathcal{O}_X$-模。若 $x \in X$，且 $\mathcal{F}$ 在点 $x$ 处于
$S$ 上平坦，则 $\mathcal{F}_x$ 是有限表示 $\mathcal{O}_{X, x}$-模。
\end{lemma}

\begin{proof}
令 $s = f(x)$。由命题 \ref{flat-proposition-finite-type-flat-at-point}，
存在初等 \'etale 邻域 $(S', s') \to (S, s)$，使得 $\mathcal{F}$ 在
$X \times_S \Spec(\mathcal{O}_{S', s'})$ 上的拉回，在与 $x$ 对应的点
$x' \in X_{s'} = X_s$ 的某个邻域中有限表示。环同态
$$
\mathcal{O}_{X, x} \longrightarrow
\mathcal{O}_{X \times_S \Spec(\mathcal{O}_{S', s'}), x'}
=
\mathcal{O}_{X \times_S S', x'}
$$
作为 \'etale 环同态的局部化，是平坦局部同态。因此由下降，
$\mathcal{F}_x$ 在 $\mathcal{O}_{X, x}$ 上有限表示；参见代数，引理
\ref{algebra-lemma-descend-properties-modules}（还要使用平坦局部环同态
忠实平坦这一事实；参见代数，引理 \ref{algebra-lemma-local-flat-ff}）。
\end{proof}

\begin{lemma}
\label{flat-lemma-finite-type-flat-at-point-local-X}
设 $f : X \to S$ 是局部有限型态射，$x \in X$ 的像为 $s \in S$。
若 $f$ 在点 $x$ 处于 $S$ 上平坦，则 $\mathcal{O}_{X, x}$ 在
$\mathcal{O}_{S, s}$ 上本质有限表示。
\end{lemma}

\begin{proof}
可设 $X$ 和 $S$ 仿射。记 $X = \Spec(B)$、$S = \Spec(A)$，并写成
$B = A[x_1, \ldots, x_n]/I$。换言之，得到闭浸入
$i : X \to \mathbf{A}^n_S$。记 $t = i(x) \in \mathbf{A}^n_S$。
将引理 \ref{flat-lemma-finite-type-flat-at-point-local} 应用于
$\mathbf{A}^n_S \to S$、层 $\mathcal{F} = i_*\mathcal{O}_X$ 与点 $t$，
可得 $\mathcal{O}_{X, x}$ 在 $\mathcal{O}_{\mathbf{A}^n_S, t}$
上有限表示，这蕴含所求结论。
\end{proof}







\section{从开集延拓性质}
\label{flat-section-extending-properties}

\noindent
本节汇集若干下述形式的结果：若 $f : X \to S$ 是平坦概形态射，且
$f$ 在 $S$ 的某个稠密开集上满足某一性质，则 $f$ 在整个 $S$ 上满足同一性质。

\begin{lemma}
\label{flat-lemma-flat-finite-type-finitely-presented-over-dense-open}
\begin{slogan}
在一个（良好）开集上延拓有限表示模所得的 $S$-平坦有限型模，
也是 $X$-有限表示的。
\end{slogan}
设 $f : X \to S$ 是概形态射，$\mathcal{F}$ 是拟凝聚
$\mathcal{O}_X$-模，$U \subset S$ 是开集。假设
\begin{enumerate}
\item $f$ 局部有限表示；
\item $\mathcal{F}$ 有限型且在 $S$ 上平坦；
\item $U \subset S$ 逆紧且概形论稠密；
\item $\mathcal{F}|_{f^{-1}U}$ 有限表示。
\end{enumerate}
则 $\mathcal{F}$ 有限表示。
\end{lemma}

\begin{proof}
问题在 $X$ 和 $S$ 上是局部的，故可设 $X$、$S$ 仿射。记
$S = \Spec(A)$、$X = \Spec(B)$。设 $N$ 是有限 $B$-模，且
$\mathcal{F}$ 是与 $N$ 相伴的拟凝聚层。对某些 $f_i \in A$，有
$U = D(f_1) \cup \ldots \cup D(f_n)$；参见代数，引理
\ref{algebra-lemma-qc-open}。因为 $U$ 概形论稠密，同态
$A \to A_{f_1} \times \ldots \times A_{f_n}$ 单射。取位于
$\mathfrak p \subset A$ 上方的素理想 $\mathfrak q \subset B$；它对应于
映到 $s \in S$ 的点 $x \in X$。由引理
\ref{flat-lemma-finite-type-flat-at-point-local}，模 $N_\mathfrak q$ 在
$B_\mathfrak q$ 上有限表示。选择 $B$-模满射
$\varphi : B^{\oplus m} \to N$。选择
$k_1, \ldots, k_t \in \Ker(\varphi)$，并置
$N' = B^{\oplus m}/\sum Bk_j$。存在典范满射 $N' \to N$，且 $N$
是以这种方式构造的 $B$-模 $N'$ 的滤过余极限。因此可选择
$k_1, \ldots, k_t$，使得 (a) $N'_{f_i} \cong N_{f_i}$，
$i = 1, \ldots, n$；以及 (b) $N'_\mathfrak q \cong N_\mathfrak q$。
特别地，这说明 $N'_\mathfrak q$ 在 $A$ 上平坦。由平坦性的开性
（参见代数，定理 \ref{algebra-theorem-openness-flatness}），存在
$g \in B$、$g \not \in \mathfrak q$，使得 $N'_g$ 在 $A$ 上平坦。
考虑交换图
$$
\xymatrix{
N'_g \ar[r] \ar[d] & N_g \ar[d] \\
\prod N'_{gf_i} \ar[r] & \prod N_{gf_i}
}
$$
由 $k_1, \ldots, k_t$ 的选择，底部箭头是同构。因为
$A \to \prod A_{f_i}$ 单射且 $N'_g$ 在 $A$ 上平坦，左侧竖直箭头单射。
因此顶部水平箭头单射，从而是同构。这证明 $N_g$ 在 $B_g$ 上有限表示。
最后应用代数，引理 \ref{algebra-lemma-cover} 即得结论。
\end{proof}

\begin{lemma}
\label{flat-lemma-flat-finite-type-finitely-presented-over-dense-open-X}
设 $f : X \to S$ 是概形态射，$U \subset S$ 是开集。假设
\begin{enumerate}
\item $f$ 局部有限型且平坦；
\item $U \subset S$ 逆紧且概形论稠密；
\item $f|_{f^{-1}U} : f^{-1}U \to U$ 局部有限表示。
\end{enumerate}
% P05-EN-CH38-0035: frozen conclusion has the duplicated preposition “is of locally of finite presentation”.
则 $f$ 局部有限表示。
\end{lemma}

\begin{proof}
问题在 $X$ 和 $S$ 上是局部的，故可设 $X$、$S$ 仿射。选择闭浸入
$i : X \to \mathbf{A}^n_S$，并将引理
\ref{flat-lemma-flat-finite-type-finitely-presented-over-dense-open} 应用于
$i_*\mathcal{O}_X$。省略若干细节。
\end{proof}

\begin{lemma}
\label{flat-lemma-flat-finite-presentation-dimension-over-dense-open}
设 $f : X \to S$ 是平坦且局部有限型的概形态射，$U \subset S$ 是稠密
开集，并且 $X_U \to U$ 的相对维数 $\leq e$；参见《态射》，定义
\ref{morphisms-definition-relative-dimension-d}。若还满足下列条件之一：
\begin{enumerate}
\item $f$ 局部有限表示；或
\item $U \subset S$ 逆紧；
\end{enumerate}
则 $f$ 的相对维数 $\leq e$。
\end{lemma}

\begin{proof}
先证明情形 (1)。设 $W \subset X$ 是《态射进阶》下列引理中构造并研究的开子概形：
\ref{more-morphisms-lemma-flat-finite-presentation-CM-open} 和
\ref{more-morphisms-lemma-flat-finite-presentation-CM-pieces}.
注意，每个纤维的每个一般点都包含于 $W$，所以只须对 $W$ 证明结论。
由于 $W = \bigcup_{d \geq 0} U_d$，只须证明当 $d > e$ 时
$U_d = \emptyset$。因为 $f$ 平坦且局部有限表示，所以它是开态射，
因而 $f(U_d)$ 是开集（《态射》，引理 \ref{morphisms-lemma-fppf-open}）。
因此，若 $U_d$ 非空，则 $f(U_d) \cap U \not = \emptyset$，如所需。

\medskip\noindent
再证明情形 (2)。可用其既约化替换 $S$，此时 $U$ 概形论稠密。
于是由引理 \ref{flat-lemma-flat-finite-type-finitely-presented-over-dense-open-X}，
$f$ 局部有限表示。这样便归结到情形 (1)。
\end{proof}

\begin{lemma}
\label{flat-lemma-proper-flat-finite-over-dense-open}
设 $f : X \to S$ 是平坦固有概形态射，$U \subset S$ 是稠密开集，
且 $X_U \to U$ 有限。若 $f$ 局部有限表示或 $U \subset S$ 逆紧，
则 $f$ 有限。
\end{lemma}

\begin{proof}
由引理 \ref{flat-lemma-flat-finite-presentation-dimension-over-dense-open}，
$f$ 的纤维维数为零。因此 $f$ 拟有限（《态射》，引理
\ref{morphisms-lemma-locally-quasi-finite-rel-dimension-0}），从而具有有限
纤维（《态射》，引理 \ref{morphisms-lemma-quasi-finite}）。故由《态射进阶》，
引理 \ref{more-morphisms-lemma-characterize-finite}，$f$ 有限。
\end{proof}

\begin{lemma}
\label{flat-lemma-zariski}
设 $f : X \to S$ 是概形态射，$U \subset S$ 是开集。若
\begin{enumerate}
\item $f$ 分离、局部有限型且平坦；
\item $f^{-1}(U) \to U$ 是同构；并且
\item $U \subset S$ 逆紧且概形论稠密；
\end{enumerate}
则 $f$ 是开浸入。
\end{lemma}

\begin{proof}
由引理 \ref{flat-lemma-flat-finite-type-finitely-presented-over-dense-open-X}，
态射 $f$ 局部有限表示。像 $f(X) \subset S$ 是开集（《态射》，引理
\ref{morphisms-lemma-fppf-open}），故可用 $f(X)$ 替换 $S$。于是须证明
$f$ 是同构。可设 $S$ 仿射。还可归结到 $X$ 拟紧的情形，因为只须证明：
像为 $S$ 的任一拟紧开集 $X' \subset X$ 都同构地映到 $S$。因此可设
$f$ 拟紧。由引理
\ref{flat-lemma-flat-finite-presentation-dimension-over-dense-open}，$f$ 的全部纤维
维数都是 $0$。故 $f$ 拟有限；参见《态射》，引理
\ref{morphisms-lemma-locally-quasi-finite-rel-dimension-0}。
取 $s \in S$。选择初等 \'etale 邻域
$g : (T, t) \to (S, s)$，使得 $X \times_S T = V \amalg W$，其中
$V \to T$ 有限且 $W_t = \emptyset$；参见《态射进阶》，引理
\ref{more-morphisms-lemma-etale-splits-off-quasi-finite-part}。将给定态射记为
$\pi : V \amalg W \to T$。因为 $\pi$ 平坦且局部有限表示，
$\pi(V)$ 在 $T$ 中是开集（《态射》，引理 \ref{morphisms-lemma-fppf-open}）。
缩小 $T$ 后，可设 $T = \pi(V)$。由于 $f$ 在 $U$ 上是同构，
$\pi$ 在 $g^{-1}U$ 上是同构。结合 $\pi(V) = T$，这说明
$\pi^{-1}g^{-1}U$ 包含于 $V$。由《态射》，引理
\ref{morphisms-lemma-flat-morphism-scheme-theoretically-dense-open}，
$\pi^{-1}g^{-1}U \subset V \amalg W$ 概形论稠密。因此
$W = \emptyset$。于是 $X \times_S T = V$ 在 $T$ 上有限。这说明在用
$s$ 的某个开邻域替换 $S$ 后，$f$ 有限；例如可用下降章，引理
\ref{descent-lemma-descending-property-finite}。此时 $f$ 有限局部自由
（《态射》，引理 \ref{morphisms-lemma-finite-flat}）；再将 $S$ 缩小为
$s$ 的更小开邻域，可设 $f$ 是某个次数 $d$ 的有限局部自由态射
（《态射》，引理 \ref{morphisms-lemma-finite-locally-free}）。但由于在稠密
开集 $U$ 上次数为 $1$，显然有 $d = 1$。故 $f$ 是同构。
\end{proof}





\section{有限表示平坦模}
\label{flat-section-finitely-presented-flat}

\noindent
在某些情形下，给定一个有限表示的环映射 $R \to S$ 以及一个有限表示的
$S$-模 $N$，$N$ 在 $R$ 上的平坦性蕴含 $N$ 作为 $R$-模是投射的，
至少在用一个 \'etale 扩张替换 $S$ 之后如此。
% P05-EN-CH38-0036: Frozen English has the grammatical typo "collect a some results".
本节汇集若干这类结果。

\begin{lemma}
\label{flat-lemma-induction-step-fp}
设 $R$ 是一个环。设 $R \to S$ 是一个有限表示的平坦环映射，
其纤维几何整。设 $\mathfrak q \subset S$ 是位于素理想
$\mathfrak r \subset R$ 之上的一个素理想。置
$\mathfrak p = \mathfrak r S$。设 $N$ 是一个有限表示的 $S$-模。
则存在 $r \geq 0$ 以及一个 $S$-模映射
$$
\alpha : S^{\oplus r} \longrightarrow N
$$
使得
$\alpha : \kappa(\mathfrak p)^{\oplus r} \to N \otimes_S \kappa(\mathfrak p)$
是同构。对于任意这样的 $\alpha$，以下条件等价：
\begin{enumerate}
\item $N_{\mathfrak q}$ 是 $R$-平坦的；
\item 存在 $f \in R$，$f \not \in \mathfrak r$，使得
$\alpha_f : S_f^{\oplus r} \to N_f$ 是 $R_f$-普遍单射的，并且存在
$g \in S$，$g \not \in \mathfrak q$，使得 $\Coker(\alpha)_g$
是 $R$-平坦的；
\item $\alpha_{\mathfrak r}$ 是 $R_{\mathfrak r}$-普遍单射的，并且
$\Coker(\alpha)_{\mathfrak q}$ 是 $R$-平坦的；
\item $\alpha_{\mathfrak r}$ 是单射，并且
$\Coker(\alpha)_{\mathfrak q}$ 是 $R$-平坦的；
\item $\alpha_{\mathfrak p}$ 是同构，并且
$\Coker(\alpha)_{\mathfrak q}$ 是 $R$-平坦的；
\item $\alpha_{\mathfrak q}$ 是单射，并且
$\Coker(\alpha)_{\mathfrak q}$ 是 $R$-平坦的。
\end{enumerate}
\end{lemma}

\begin{proof}
为得到 $\alpha$，置
$r = \dim_{\kappa(\mathfrak p)} N \otimes_S \kappa(\mathfrak p)$，并选取
$x_1, \ldots, x_r \in N$，使它们构成
$N \otimes_S \kappa(\mathfrak p)$ 的一组基。定义
$\alpha(s_1, \ldots, s_r) = \sum s_i x_i$。这就证明了存在性。

\medskip\noindent
固定 $\alpha$ 的一个选择。我们可以将
引理 \ref{flat-lemma-induction-step}
应用于映射
$\alpha_{\mathfrak r} : S_{\mathfrak r}^{\oplus r} \to N_{\mathfrak r}$.
于是可见 (1)、(3)、(4)、(5) 和 (6) 彼此等价。
又因 (2) 蕴含 (3) 也是显然的，我们只需证明 (1) 蕴含 (2)。

\medskip\noindent
假设 (1) 成立。由平坦性的开放性，参见
《代数》，定理 \ref{algebra-theorem-openness-flatness}，集合
$$
U_1 = \{\mathfrak q' \subset S \mid N_{\mathfrak q'}\text{ 在 }R\text{ 上平坦}\}
$$
在 $\Spec(S)$ 中是开的。由假设，它包含 $\mathfrak q$，因而也包含
$\mathfrak p$。由于 $S^{\oplus r}$ 和 $N$ 都是有限表示的 $S$-模，集合
$$
U_2 = \{\mathfrak q' \subset S \mid
\alpha_{\mathfrak q'}\text{ 是同构}\}
$$
在 $\Spec(S)$ 中是开的，参见
《代数》，引理 \ref{algebra-lemma-map-between-finitely-presented}。
由 (5)，它包含 $\mathfrak p$。因为 $R \to S$ 是有限表示且平坦的，映射
$\Phi : \Spec(S) \to \Spec(R)$ 是开映射，参见
《代数》，命题 \ref{algebra-proposition-fppf-open}。
对于任意素理想 $\mathfrak r' \in \Phi(U_1 \cap U_2)$，存在一个位于
$\mathfrak r'$ 之上的素理想 $\mathfrak q'$，使得 $N_{\mathfrak q'}$
是平坦的且 $\alpha_{\mathfrak q'}$ 是同构。这蕴含
$\alpha \otimes \kappa(\mathfrak p')$ 是同构，其中
$\mathfrak p' = \mathfrak r' S$。因此，由蕴含关系
(1) $\Rightarrow$ (3)，$\alpha_{\mathfrak r'}$ 是
$R_{\mathfrak r'}$-普遍单射的。
故若选取 $f \in R$，$f \not \in \mathfrak r$，使得
$D(f) \subset \Phi(U_1 \cap U_2)$，则可推出 $\alpha_f$ 是
$R_f$-普遍单射的，参见
《代数》，引理 \ref{algebra-lemma-universally-injective-check-stalks}。
同样的论证还表明，对于任意
$\mathfrak q' \in U_1 \cap \Phi^{-1}(\Phi(U_1 \cap U_2))$
模 $\Coker(\alpha)_{\mathfrak q'}$ 都是 $R$-平坦的。注意
$\mathfrak q \in U_1 \cap \Phi^{-1}(\Phi(U_1 \cap U_2))$。
因此可以找到 $g \in S$，$g \not \in \mathfrak q$，使得
$D(g) \subset U_1 \cap \Phi^{-1}(\Phi(U_1 \cap U_2))$，结论得证。
\end{proof}

\begin{lemma}
\label{flat-lemma-complete-devissage-flat-finitely-presented-module}
设 $R \to S$ 是一个有限表示的环映射。设 $N$ 是一个有限表示的
$S$-模，并且在 $R$ 上平坦。设 $\mathfrak r \subset R$ 是一个素理想。
假设在 $\mathfrak r$ 上存在 $N/S/R$ 的一个完全 d\'evissage。
则存在 $f \in R$，$f \not \in \mathfrak r$，使得
$$
N_f \cong B_1^{\oplus r_1} \oplus \ldots \oplus B_n^{\oplus r_n}
$$
% P05-EN-CH38-0037: Frozen English says "as R-modules" after localization;
% the displayed decomposition naturally concerns R_f-modules. Translation follows source.
作为 $R$-模成立，其中每个 $B_i$ 都是一个纤维几何不可约的光滑
$R_f$-代数。此外，$N_f$ 作为 $R_f$-模是投射的。
\end{lemma}

\begin{proof}
设 $(A_i, B_i, M_i, \alpha_i)_{i = 1, \ldots, n}$ 是给定的完全
d\'evissage。我们对 $n$ 作归纳来证明该引理。注意，该引理的断言完全
关乎 $N$ 作为 $R$-模的结构。因此可以用 $M_1$ 替换 $N$，并将 $M_1$
视为一个 $B_1$-模。关于 $M_1$ 为何是有限表示的 $B_1$-模，参见
注记 \ref{flat-remark-finite-presentation}。由
引理 \ref{flat-lemma-induction-step-fp}，在对某个
$f \in R$，$f \not \in \mathfrak r$ 用 $R_f$ 替换 $R$ 之后，可以假设
映射 $\alpha_1 : B_1^{\oplus r_1} \to M_1$ 是 $R$-普遍单射的。
由于 $M_1$ 和 $B_1^{\oplus r_1}$ 都是 $R$-平坦的，并且作为
$B_1$-模有限表示，可见 $\Coker(\alpha_1)$ 是 $R$-平坦的
（《代数》，引理 \ref{algebra-lemma-ui-flat-domain}），
且作为 $B_1$-模有限表示。注意，
$(A_i, B_i, M_i, \alpha_i)_{i = 2, \ldots, n}$ 是
$\Coker(\alpha_1)$ 的一个完全 d\'evissage。因此，归纳假设蕴含：
在对某个 $f \in R$，$f \not \in \mathfrak r$ 用 $R_f$ 替换 $R$ 后，
可以假设 $\Coker(\alpha_1)$ 具有该引理所述的分解，并且是投射的。
特别地，
% P05-EN-CH38-0038: Frozen English uses equality for a splitting obtained up to isomorphism.
$M_1 = B_1^{\oplus r_1} \oplus \Coker(\alpha_1)$.
这证明了关于分解的断言。关于投射性的断言来自以下事实：由
引理 \ref{flat-lemma-fibres-irreducible-flat-projective-nonnoetherian}，
$B_1$ 作为 $R$-模是投射的。
\end{proof}

\begin{remark}
\label{flat-remark-complete-devissage-flat-finitely-presented-module}
引理 \ref{flat-lemma-complete-devissage-flat-finitely-presented-module}
有如下变体：我们弱化平坦性条件，只假设 $N$ 在位于 $\mathfrak r$
之上的某个给定素理想 $\mathfrak q$ 处平坦；但同时加强 d\'evissage
条件，假设存在一个{\it 在 $\mathfrak q$ 处}的完全 d\'evissage。
与引理 \ref{flat-lemma-complete-devissage-flat-finite-type-module} 比较。
\end{remark}

\noindent
以下是本节的主要结果。

\begin{proposition}
\label{flat-proposition-finite-presentation-flat-at-point}
设 $f : X \to S$ 是一个概形态射。设 $\mathcal{F}$ 是 $X$ 上的一个
拟凝聚层。设 $x \in X$ 的像为 $s \in S$。假设
\begin{enumerate}
\item $f$ 局部有限表示；
\item $\mathcal{F}$ 有限表示；
\item $\mathcal{F}$ 在 $x$ 处于 $S$ 上平坦。
\end{enumerate}
则存在带基点概形的一个交换图
$$
\xymatrix{
(X, x) \ar[d] & (X', x') \ar[l]^g \ar[d] \\
(S, s) & (S', s') \ar[l]
}
$$
其中水平箭头都是初等 \'etale 邻域，$X'$、$S'$ 都是仿射的，并且
$\Gamma(X', g^*\mathcal{F})$ 是一个投射的
$\Gamma(S', \mathcal{O}_{S'})$-模。
\end{proposition}

\begin{proof}
由平坦性的开放性，参见
《态射进阶》，定理 \ref{more-morphisms-theorem-openness-flatness}，
可以用 $x$ 的一个开邻域替换 $X$，并假设 $\mathcal{F}$ 在 $S$ 上平坦。
接着应用命题 \ref{flat-proposition-existence-complete-at-x}，得到命题所述的
一个图，使得 $g^*\mathcal{F}/X'/S'$ 在 $s'$ 上有一个完全
d\'evissage。（特别地，$S'$ 和 $X'$ 都是仿射的。）由
《态射》，引理 \ref{morphisms-lemma-flat-permanence}，可见
$g^*\mathcal{F}$ 在 $S$ 上平坦；再由
引理 \ref{flat-lemma-etale-flat-up-down}，可见它在 $S'$ 上平坦。
通过注记 \ref{flat-remark-same-notion}，我们推出
$$
\Gamma(X', g^*\mathcal{F})/
\Gamma(X', \mathcal{O}_{X'})/
\Gamma(S', \mathcal{O}_{S'})
$$
在 $\Gamma(S', \mathcal{O}_{S'})$ 中对应于 $s'$ 的素理想上有一个完全
d\'evissage。因此，引理
\ref{flat-lemma-complete-devissage-flat-finitely-presented-module} 蕴含：
用 $s'$ 的一个标准开邻域替换 $S'$ 后，命题的结论成立。
\end{proof}

\noindent
在本节余下部分，我们证明该结果的若干变体。第一个是一个“整体”版本。

\begin{lemma}
\label{flat-lemma-finite-presentation-flat-along-fibre}
设 $f : X \to S$ 是一个概形态射。设 $\mathcal{F}$ 是 $X$ 上的一个
拟凝聚层。设 $s \in S$。假设
\begin{enumerate}
\item $f$ 有限表示；
\item $\mathcal{F}$ 有限表示；
\item $\mathcal{F}$ 在纤维 $X_s$ 的每一点处于 $S$ 上平坦。
\end{enumerate}
则存在一个初等 \'etale 邻域 $(S', s') \to (S, s)$ 以及概形的一个交换图
$$
\xymatrix{
X \ar[d] & X' \ar[l]^g \ar[d] \\
S & S' \ar[l]
}
$$
使得 $g$ 是 \'etale 的，$X_s \subset g(X')$，概形 $X'$、$S'$ 都是
仿射的，并且 $\Gamma(X', g^*\mathcal{F})$ 是一个投射的
$\Gamma(S', \mathcal{O}_{S'})$-模。
\end{lemma}

\begin{proof}
对于每一点 $x \in X_s$，可以应用
命题 \ref{flat-proposition-finite-presentation-flat-at-point} 找到一个交换图
$$
\xymatrix{
(X, x) \ar[d] & (Y_x, y_x) \ar[l]^{g_x} \ar[d] \\
(S, s) & (S_x, s_x) \ar[l]
}
$$
其水平箭头都是初等 \'etale 邻域，$Y_x$、$S_x$ 都是仿射的，并且
$\Gamma(Y_x, g_x^*\mathcal{F})$ 是一个投射的
$\Gamma(S_x, \mathcal{O}_{S_x})$-模。特别地，
$g_x(Y_x) \cap X_s$ 是 $X_s$ 中 $x$ 的一个开邻域。
由于 $X_s$ 拟紧，可以找到有限多个点 $x_1, \ldots, x_n \in X_s$，
使 $X_s$ 是各 $g_{x_i}(Y_{x_i}) \cap X_s$ 的并。选取一个支配每个邻域
$(S_{x_i}, s_{x_i})$ 的初等 \'etale 邻域
$(S' , s') \to (S, s)$，参见《态射进阶》，
引理 \ref{more-morphisms-lemma-elementary-etale-neighbourhoods}。
还可以假设 $S'$ 是仿射的。置
$X' = \coprod Y_{x_i} \times_{S_{x_i}} S'$，并赋予它显然的态射
$g : X' \to X$。由构造，$g(X')$ 包含 $X_s$，并且
$$
\Gamma(X', g^*\mathcal{F})
=
\bigoplus \Gamma(Y_{x_i}, g_{x_i}^*\mathcal{F})
\otimes_{\Gamma(S_{x_i}, \mathcal{O}_{S_{x_i}})}
\Gamma(S', \mathcal{O}_{S'}).
$$
这是一个投射的 $\Gamma(S', \mathcal{O}_{S'})$-模，参见
《代数》，引理 \ref{algebra-lemma-ascend-properties-modules}。
\end{proof}

\noindent
以下两个引理是上述结果在 $\mathcal{F} = \mathcal{O}_X$ 情形下的改述。

\begin{lemma}
\label{flat-lemma-finite-presentation-flat-at-point-X}
设 $f : X \to S$ 局部有限表示。设 $x \in X$ 的像为 $s \in S$。
若 $f$ 在 $x$ 处于 $S$ 上平坦，则存在带基点概形的一个交换图
$$
\xymatrix{
(X, x) \ar[d] & (X', x') \ar[l]^g \ar[d] \\
(S, s) & (S', s') \ar[l]
}
$$
其水平箭头都是初等 \'etale 邻域，$X'$、$S'$ 都是仿射的，并且
$\Gamma(X', \mathcal{O}_{X'})$ 是一个投射的
$\Gamma(S', \mathcal{O}_{S'})$-模。
\end{lemma}

\begin{proof}
这是命题 \ref{flat-proposition-finite-presentation-flat-at-point} 的一个特例。
\end{proof}

\begin{lemma}
\label{flat-lemma-finite-presentation-flat-along-fibre-X}
设 $f : X \to S$ 有限表示。设 $s \in S$。若 $X$ 在 $X_s$ 的所有点处
于 $S$ 上平坦，则存在一个初等 \'etale 邻域
$(S', s') \to (S, s)$ 以及概形的一个交换图
$$
\xymatrix{
X \ar[d] & X' \ar[l]^g \ar[d] \\
S & S' \ar[l]
}
$$
其中 $g$ 是 \'etale 的，$X_s \subset g(X')$，$X'$、$S'$ 都是仿射的，
并且 $\Gamma(X', \mathcal{O}_{X'})$ 是一个投射的
$\Gamma(S', \mathcal{O}_{S'})$-模。
\end{lemma}

\begin{proof}
这是引理 \ref{flat-lemma-finite-presentation-flat-along-fibre} 的一个特例。
\end{proof}

\noindent
以下各引理说明命题 \ref{flat-proposition-finite-presentation-flat-at-point}
在只假设态射与层为有限型（而不必有限表示）时的推论。

\begin{lemma}
\label{flat-lemma-finite-type-flat-at-point-free}
设 $f : X \to S$ 是一个概形态射。设 $\mathcal{F}$ 是 $X$ 上的一个
拟凝聚层。设 $x \in X$ 的像为 $s \in S$。假设
\begin{enumerate}
\item $f$ 局部有限表示；
\item $\mathcal{F}$ 为有限型；
\item $\mathcal{F}$ 在 $x$ 处于 $S$ 上平坦。
\end{enumerate}
则存在一个初等 \'etale 邻域 $(S', s') \to (S, s)$ 以及带基点概形的
一个交换图
$$
\xymatrix{
(X, x) \ar[d] & (X', x') \ar[l]^g \ar[d] \\
(S, s) & (\Spec(\mathcal{O}_{S', s'}), s') \ar[l]
}
$$
使得 $X' \to X \times_S \Spec(\mathcal{O}_{S', s'})$ 是 \'etale 的，
$\kappa(x) = \kappa(x')$，概形 $X'$ 在 $\mathcal{O}_{S', s'}$ 上仿射且
有限表示，层 $g^*\mathcal{F}$ 在 $\mathcal{O}_{X'}$ 上有限表示，并且
$\Gamma(X', g^*\mathcal{F})$ 是一个自由的
$\mathcal{O}_{S', s'}$-模。
\end{lemma}

\begin{proof}
为证明该引理，可以用任意初等 \'etale 邻域替换 $(S, s)$，也可以用
$\Spec(\mathcal{O}_{S, s})$ 替换 $S$。因此，由
命题 \ref{flat-proposition-finite-type-flat-at-point}，可以假设
$\mathcal{F}$ 在 $x$ 的一个邻域中有限表示且在 $S$ 上平坦。
在此情形下，结论来自
命题 \ref{flat-proposition-finite-presentation-flat-at-point}，因为
《代数》，定理 \ref{algebra-theorem-projective-free-over-local-ring}
保证在局部环上投射 $=$ 自由。
\end{proof}

\begin{lemma}
\label{flat-lemma-finite-type-flat-at-point-free-variant}
设 $f : X \to S$ 是一个概形态射。设 $\mathcal{F}$ 是 $X$ 上的一个
拟凝聚层。设 $x \in X$ 的像为 $s \in S$。假设
\begin{enumerate}
\item $f$ 局部有限型；
\item $\mathcal{F}$ 为有限型；
\item $\mathcal{F}$ 在 $x$ 处于 $S$ 上平坦。
\end{enumerate}
则存在一个初等 \'etale 邻域 $(S', s') \to (S, s)$ 以及带基点概形的
一个交换图
$$
\xymatrix{
(X, x) \ar[d] & (X', x') \ar[l]^g \ar[d] \\
(S, s) & (\Spec(\mathcal{O}_{S', s'}), s') \ar[l]
}
$$
使得 $X' \to X \times_S \Spec(\mathcal{O}_{S', s'})$ 是 \'etale 的，
$\kappa(x) = \kappa(x')$，概形 $X'$ 是仿射的，并且
$\Gamma(X', g^*\mathcal{F})$ 是一个自由的
$\mathcal{O}_{S', s'}$-模。
\end{lemma}

\begin{proof}
（与引理 \ref{flat-lemma-finite-type-flat-at-point-free} 唯一的区别在于，
这里不假设 $f$ 有限表示。）该问题在 $X$ 与 $S$ 上都是局部的。
因此可以假设 $X$ 与 $S$ 都是仿射的，例如
$X = \Spec(B)$ 且 $S = \Spec(A)$。由于 $B$ 是有限型 $A$-代数，
可以找到一个满射 $A[x_1, \ldots, x_n] \to B$。换言之，可以选取
一个闭浸入 $i : X \to \mathbf{A}^n_S$。置 $t = i(x)$ 且
$\mathcal{G} = i_*\mathcal{F}$。注意，作为 $\mathcal{O}_{S, s}$-模，
$\mathcal{G}_t \cong \mathcal{F}_x$。故 $\mathcal{G}$ 在 $t$ 处于
$S$ 上平坦。将引理 \ref{flat-lemma-finite-type-flat-at-point-free}
应用于态射 $\mathbf{A}^n_S \to S$、点 $t$ 以及层 $\mathcal{G}$。
于是可以找到一个初等 \'etale 邻域 $(S', s') \to (S, s)$ 以及带基点
概形的一个交换图
$$
\xymatrix{
(\mathbf{A}^n_S, t) \ar[d] & (Y, y) \ar[l]^h \ar[d] \\
(S, s) & (\Spec(\mathcal{O}_{S', s'}), s') \ar[l]
}
$$
使得 $Y \to \mathbf{A}^n_{\mathcal{O}_{S', s'}}$ 是 \'etale 的，
$\kappa(t) = \kappa(y)$，概形 $Y$ 是仿射的，并且
$\Gamma(Y, h^*\mathcal{G})$ 是一个投射的
$\mathcal{O}_{S', s'}$-模。于是，$Y$ 的闭子概形
$X' = Y \times_{\mathbf{A}^n_S} X$ 给出原问题的一个解。
\end{proof}

\begin{lemma}
\label{flat-lemma-finite-type-flat-along-fibre-free}
设 $f : X \to S$ 是一个概形态射。设 $\mathcal{F}$ 是 $X$ 上的一个
拟凝聚层。设 $s \in S$。假设
\begin{enumerate}
\item $f$ 有限表示；
\item $\mathcal{F}$ 为有限型；
\item $\mathcal{F}$ 在 $X_s$ 的所有点处于 $S$ 上平坦。
\end{enumerate}
则存在一个初等 \'etale 邻域 $(S', s') \to (S, s)$ 以及概形的一个交换图
$$
\xymatrix{
X \ar[d] & X' \ar[l]^g \ar[d] \\
S & \Spec(\mathcal{O}_{S', s'}) \ar[l]
}
$$
使得 $X' \to X \times_S \Spec(\mathcal{O}_{S', s'})$ 是 \'etale 的，
$X_s = g((X')_{s'})$，概形 $X'$ 在 $\mathcal{O}_{S', s'}$ 上仿射且
有限表示，层 $g^*\mathcal{F}$ 在 $\mathcal{O}_{X'}$ 上有限表示，并且
$\Gamma(X', g^*\mathcal{F})$ 是一个自由的
$\mathcal{O}_{S', s'}$-模。
\end{lemma}

\begin{proof}
对于每一点 $x \in X_s$，可以应用
引理 \ref{flat-lemma-finite-type-flat-at-point-free} 找到一个初等 \'etale 邻域
$(S_x , s_x) \to (S, s)$ 以及一个交换图
$$
\xymatrix{
(X, x) \ar[d] & (Y_x, y_x) \ar[l]^{g_x} \ar[d] \\
(S, s) & (\Spec(\mathcal{O}_{S_x, s_x}), s_x) \ar[l]
}
$$
使得 $Y_x \to X \times_S \Spec(\mathcal{O}_{S_x, s_x})$ 是 \'etale 的，
$\kappa(x) = \kappa(y_x)$，概形 $Y_x$ 在 $\mathcal{O}_{S_x, s_x}$ 上
仿射且有限表示，层 $g_x^*\mathcal{F}$ 在 $\mathcal{O}_{Y_x}$ 上有限表示，
并且 $\Gamma(Y_x, g_x^*\mathcal{F})$ 是一个自由的
$\mathcal{O}_{S_x, s_x}$-模。特别地，$g_x((Y_x)_{s_x})$ 是 $X_s$ 中
$x$ 的一个开邻域。由于 $X_s$ 拟紧，可以找到有限多个点
$x_1, \ldots, x_n \in X_s$，使 $X_s$ 是各
$g_{x_i}((Y_{x_i})_{s_{x_i}})$ 的并。选取一个支配每个邻域
$(S_{x_i}, s_{x_i})$ 的初等 \'etale 邻域
$(S' , s') \to (S, s)$，参见《态射进阶》，
引理 \ref{more-morphisms-lemma-elementary-etale-neighbourhoods}。置
$$
X' = \coprod Y_{x_i} \times_{\Spec(\mathcal{O}_{S_{x_i}, s_{x_i}})}
\Spec(\mathcal{O}_{S', s'})
$$
并赋予它显然的态射 $g : X' \to X$。由构造，
$X_s = g(X'_{s'})$，并且
$$
\Gamma(X', g^*\mathcal{F})
=
\bigoplus \Gamma(Y_{x_i}, g_{x_i}^*\mathcal{F})
\otimes_{\mathcal{O}_{S_{x_i}, s_{x_i}}}
\mathcal{O}_{S', s'}.
$$
这是自由模经基变换所得模的直和，因而是一个自由的
$\mathcal{O}_{S', s'}$-模。略去若干次要细节。
\end{proof}

\begin{lemma}
\label{flat-lemma-finite-type-flat-along-fibre-free-variant}
设 $f : X \to S$ 是一个概形态射。设 $\mathcal{F}$ 是 $X$ 上的一个
拟凝聚层。设 $s \in S$。假设
\begin{enumerate}
\item $f$ 有限型；
\item $\mathcal{F}$ 为有限型；
\item $\mathcal{F}$ 在 $X_s$ 的所有点处于 $S$ 上平坦。
\end{enumerate}
则存在一个初等 \'etale 邻域 $(S', s') \to (S, s)$ 以及概形的一个交换图
$$
\xymatrix{
X \ar[d] & X' \ar[l]^g \ar[d] \\
S & \Spec(\mathcal{O}_{S', s'}) \ar[l]
}
$$
使得 $X' \to X \times_S \Spec(\mathcal{O}_{S', s'})$ 是 \'etale 的，
$X_s = g((X')_{s'})$，概形 $X'$ 是仿射的，并且
$\Gamma(X', g^*\mathcal{F})$ 是一个自由的
$\mathcal{O}_{S', s'}$-模。
\end{lemma}

\begin{proof}
（与引理 \ref{flat-lemma-finite-type-flat-along-fibre-free} 唯一的区别在于，
这里不假设 $f$ 有限表示。）对于每一点 $x \in X_s$，可以应用
引理 \ref{flat-lemma-finite-type-flat-at-point-free-variant}
找到一个初等 \'etale 邻域 $(S_x , s_x) \to (S, s)$ 以及一个交换图
$$
\xymatrix{
(X, x) \ar[d] & (Y_x, y_x) \ar[l]^{g_x} \ar[d] \\
(S, s) & (\Spec(\mathcal{O}_{S_x, s_x}), s_x) \ar[l]
}
$$
使得 $Y_x \to X \times_S \Spec(\mathcal{O}_{S_x, s_x})$ 是 \'etale 的，
$\kappa(x) = \kappa(y_x)$，概形 $Y_x$ 是仿射的，并且
$\Gamma(Y_x, g_x^*\mathcal{F})$ 是一个自由的
$\mathcal{O}_{S_x, s_x}$-模。特别地，$g_x((Y_x)_{s_x})$ 是 $X_s$ 中
$x$ 的一个开邻域。由于 $X_s$ 拟紧，可以找到有限多个点
$x_1, \ldots, x_n \in X_s$，使 $X_s$ 是各
$g_{x_i}((Y_{x_i})_{s_{x_i}})$ 的并。选取一个支配每个邻域
$(S_{x_i}, s_{x_i})$ 的初等 \'etale 邻域
$(S' , s') \to (S, s)$，参见《态射进阶》，
引理 \ref{more-morphisms-lemma-elementary-etale-neighbourhoods}。置
$$
X' = \coprod Y_{x_i} \times_{\Spec(\mathcal{O}_{S_{x_i}, s_{x_i}})}
\Spec(\mathcal{O}_{S', s'})
$$
并赋予它显然的态射 $g : X' \to X$。由构造，
$X_s = g(X'_{s'})$，并且
$$
\Gamma(X', g^*\mathcal{F})
=
\bigoplus \Gamma(Y_{x_i}, g_{x_i}^*\mathcal{F})
\otimes_{\mathcal{O}_{S_{x_i}, s_{x_i}}}
\mathcal{O}_{S', s'}.
$$
这是自由模经基变换所得模的直和，因而是一个自由的
$\mathcal{O}_{S', s'}$-模。
\end{proof}




\section{有限型平坦模，第二部分}
\label{flat-section-finite-type-flat-II}

\noindent
我们将需要以下引理。

\begin{lemma}
\label{flat-lemma-weak-bourbaki-pre-pre}
设 $R \to S$ 是一个有限表示的环映射。设 $N$ 是一个有限表示的
$S$-模。设 $\mathfrak q \subset S$ 是一个位于
$\mathfrak p \subset R$ 之上的素理想。置
$\overline{S} = S \otimes_R \kappa(\mathfrak p)$，
$\overline{\mathfrak q} = \mathfrak q \overline{S}$，以及
$\overline{N} = N \otimes_R \kappa(\mathfrak p)$。则可以找到
$g \in S$，$g \not \in \mathfrak q$，使得
$\overline{g} \in \mathfrak r$ 对所有满足
$\mathfrak r \in \text{Ass}_{\overline{S}}(\overline{N})$
且 $\mathfrak r \not \subset \overline{\mathfrak q}$ 的素理想成立。
\end{lemma}

\begin{proof}
事实上，若 $\text{Ass}_{\overline{S}}(\overline{N}) =
\{\mathfrak r_1, \ldots, \mathfrak r_n\}$（有限性来自《代数》，
引理 \ref{algebra-lemma-finite-ass}），则重新编号后可以假设
$$
\mathfrak r_1 \subset \overline{\mathfrak q},
\ldots,
\mathfrak r_r \subset \overline{\mathfrak q}, \quad
\mathfrak r_{r + 1} \not \subset \overline{\mathfrak q},
\ldots,
\mathfrak r_n \not \subset \overline{\mathfrak q}
$$
由于 $\overline{\mathfrak q}$ 是素理想，乘积
$\mathfrak r_{r + 1} \ldots \mathfrak r_n$ 不包含于
$\overline{\mathfrak q}$。因此，可以选取 $\overline{S}$ 中的一个元素
$a$，它属于 $\mathfrak r_{r + 1}, \ldots, \mathfrak r_n$，但不属于
$\overline{\mathfrak q}$。若存在映到 $a$ 的 $g \in S$，则这个 $g$
可用。一般而言，可以找到一个非零元素
$\lambda \in \kappa(\mathfrak p)$，使得 $\lambda a$ 是某个
$g \in S$ 的像。
\end{proof}

\noindent
以下引理在下文有一个稍强的变体，即引理 \ref{flat-lemma-weak-bourbaki}。

\begin{lemma}
\label{flat-lemma-weak-bourbaki-pre}
设 $R \to S$ 是一个有限表示的环映射。设 $N$ 是一个有限表示的
$S$-模，并且作为 $R$-模平坦。设 $M$ 是一个 $R$-模。
设 $\mathfrak q$ 是 $S$ 中位于 $\mathfrak p \subset R$ 之上的一个
素理想。则
$$
\mathfrak q \in \text{WeakAss}_S(M \otimes_R N)
\Leftrightarrow
\Big(
\mathfrak p \in \text{WeakAss}_R(M)
\text{ 且 }
\overline{\mathfrak q} \in \text{Ass}_{\overline{S}}(\overline{N})
\Big)
$$
这里 $\overline{S} = S \otimes_R \kappa(\mathfrak p)$，
$\overline{\mathfrak q} = \mathfrak q \overline{S}$，并且
$\overline{N} = N \otimes_R \kappa(\mathfrak p)$.
\end{lemma}

\begin{proof}
按引理 \ref{flat-lemma-weak-bourbaki-pre-pre} 选取 $g \in S$。
将命题 \ref{flat-proposition-finite-presentation-flat-at-point}
应用于概形态射 $\Spec(S_g) \to \Spec(R)$、与 $N_g$ 对应的拟凝聚模，
以及分别对应于素理想 $\mathfrak qS_g$ 和 $\mathfrak p$ 的点。
翻译成代数语言，我们得到环的一个交换图
$$
\xymatrix{
S \ar[r] & S_g \ar[r] & S' \\
& R \ar[lu] \ar[u] \ar[r] & R' \ar[u]
}
\quad\quad
\xymatrix{
\mathfrak q \ar@{-}[r] \ar@{-}[rd] &
\mathfrak qS_g \ar@{-}[d] \ar@{-}[r] & \mathfrak q' \ar@{-}[d] \\
& \mathfrak p \ar@{-}[r] & \mathfrak p'
}
$$
并带有图中所示的素理想；水平箭头都是 \'etale 的，并且
$N \otimes_S S'$ 作为 $R'$-模是投射的。置
$N' = N \otimes_S S'$, $M' = M \otimes_R R'$,
% P05-EN-CH38-0039: Frozen English tensors S' over R' with kappa(q'),
% although q' is an S'-prime and the next displayed identity uses kappa(p').
$\overline{S}' = S' \otimes_{R'} \kappa(\mathfrak q')$,
$\overline{\mathfrak q}' = \mathfrak q' \overline{S}'$,
以及
$$
\overline{N}' = N' \otimes_{R'} \kappa(\mathfrak p') =
\overline{N} \otimes_{\overline{S}} \overline{S}'
$$
由引理 \ref{flat-lemma-etale-weak-assassin-up-down}，有
\begin{align*}
\text{WeakAss}_{S'}(M' \otimes_{R'} N') & =
(\Spec(S') \to \Spec(S))^{-1}\text{WeakAss}_S(M \otimes_R N) \\
\text{WeakAss}_{R'}(M') & =
(\Spec(R') \to \Spec(R))^{-1}\text{WeakAss}_R(M) \\
\text{Ass}_{\overline{S}'}(\overline{N}') & =
(\Spec(\overline{S}') \to \Spec(\overline{S}))^{-1}
\text{Ass}_{\overline{S}}(\overline{N})
\end{align*}
对 $\overline{N}$ 和 $\overline{N}'$ 应用《代数》，
引理 \ref{algebra-lemma-ass-weakly-ass}。特别地，有
\begin{align*}
\mathfrak q \in \text{WeakAss}_S(M \otimes_R N)
& \Leftrightarrow
\mathfrak q' \in \text{WeakAss}_{S'}(M' \otimes_{R'} N') \\
\mathfrak p \in \text{WeakAss}_R(M)
& \Leftrightarrow
\mathfrak p' \in \text{WeakAss}_{R'}(M') \\
\overline{\mathfrak q} \in \text{Ass}_{\overline{S}}(\overline{N})
& \Leftrightarrow
\overline{\mathfrak q}' \in \text{WeakAss}_{\overline{S}'}(\overline{N}')
\end{align*}
对 $g$ 的审慎选择以及上面关于
$\text{Ass}_{\overline{S}'}(\overline{N}')$ 的公式表明
\begin{equation}
\label{flat-equation-key-observation}
\text{若 }\mathfrak r' \in \text{Ass}_{\overline{S}'}(\overline{N}')
\text{ 位于 }\mathfrak r \subset \overline{S}\text{ 上方，则 }
\mathfrak r \subset \overline{\mathfrak q}
\end{equation}
这将是证明后面部分的一个关键观察。下文将不加说明地使用
《代数》，引理 \ref{algebra-lemma-weakly-ass-local} 给出的弱伴随素理想
刻画。

\medskip\noindent
假设
$\overline{\mathfrak q} \not \in \text{Ass}_{\overline{S}}(\overline{N})$。
那么
$\overline{\mathfrak q}' \not \in \text{Ass}_{\overline{S}'}(\overline{N}')$。
由《代数》，引理 \ref{algebra-lemma-ass-zero-divisors}、
\ref{algebra-lemma-finite-ass} 和
\ref{algebra-lemma-silly}
，存在一个元素 $\overline{a}' \in \overline{\mathfrak q}'$，它不是
$\overline{N}'$ 上的零因子。对某个非零
$\lambda \in \kappa(\mathfrak p)$ 用 $\lambda \overline{a}'$ 替换
$\overline{a}'$ 后，可以找到映到 $\overline{a}'$ 的
$a' \in \mathfrak q'$。由
引理 \ref{flat-lemma-invert-universally-injective}，映射
$a' : N'_{\mathfrak p'} \to N'_{\mathfrak p'}$ 是
$R'_{\mathfrak p'}$-普遍单射的。特别地，映射
$a' : M' \otimes_{R'} N' \to M' \otimes_{R'} N'$ 在
$\mathfrak p'$ 处局部化后是单射，因而在 $\mathfrak q'$ 处局部化后
也是单射。显然，这蕴含
$\mathfrak q' \not \in \text{WeakAss}_{S'}(M' \otimes_{R'} N')$。
我们得出：$\mathfrak q \in \text{WeakAss}_S(M \otimes_R N)$ 蕴含
$\overline{\mathfrak q} \in \text{Ass}_{\overline{S}}(\overline{N})$。

\medskip\noindent
假设 $\mathfrak q \in \text{WeakAss}_S(M \otimes_R N)$。我们要证明
% P05-EN-CH38-0040: Frozen English writes WeakAss_S(M), although p is an R-prime.
$\mathfrak p \in \text{WeakAss}_S(M)$。取元素
$z \in M \otimes_R N$，使得 $\mathfrak q$ 是包含
$J = \text{Ann}_S(z)$ 的极小素理想。设 $f_i \in \mathfrak p$，
$i \in I$ 是理想 $\mathfrak p$ 的一组生成元。由于 $\mathfrak q$
位于 $\mathfrak p$ 之上，对于每个 $i$，可以选取 $n_i \geq 1$ 以及
$g_i \in S$，$g_i \not \in \mathfrak q$，使得
$g_i f_i^{n_i} \in J$，即 $g_i f_i^{n_i} z = 0$。
设 $z' \in (M' \otimes_{R'} N')_{\mathfrak p'}$ 为 $z$ 的像。
注意，$z'$ 非零，因为 $z$ 在 $(M \otimes_R N)_\mathfrak q$ 中的像非零，
并且 $S_\mathfrak q \to S'_{\mathfrak q'}$ 忠实平坦。我们断言
$f_i^{n_i} z' = 0$。

\medskip\noindent
断言的证明：设 $g'_i \in S'$ 为 $g_i$ 的像。由关键观察
(\ref{flat-equation-key-observation})，可见像
$\overline{g}'_i \in \overline{S}'$ 不属于 $\mathfrak r'$，其中
% P05-EN-CH38-0041: Frozen English uses overline N instead of overline N' here.
$\mathfrak r' \in \text{Ass}_{\overline{S}'}(\overline{N})$。
于是由引理 \ref{flat-lemma-invert-universally-injective}，可见
$g'_i : N'_{\mathfrak p'} \to N'_{\mathfrak p'}$ 是
$R'_{\mathfrak p'}$-普遍单射的。特别地，映射
$g'_i : M' \otimes_{R'} N' \to M' \otimes_{R'} N'$ 在
$\mathfrak p'$ 处局部化后是单射。由于 $g_i f_i^{n_i} z' = 0$，
断言得证。

\medskip\noindent
我们的断言表明，$z'$ 在 $R'_{\mathfrak p'}$ 中的零化理想包含各元素
$f_i^{n_i}$。由于 $R \to R'$ 是 \'etale 的，由《代数》，
引理 \ref{algebra-lemma-etale-at-prime}，有
$\mathfrak p'R'_{\mathfrak p'} = \mathfrak pR'_{\mathfrak p'}$
。因此，$z'$ 在 $R'_{\mathfrak p'}$ 中的零化理想的根等于
% P05-EN-CH38-0042: Frozen English localizes R at the R'-prime p'; likely R' was intended.
$\mathfrak p' R_{\mathfrak p'}$（这里使用了 $z'$ 非零）。另一方面，
$$
z' \in (M' \otimes_{R'} N')_{\mathfrak p'} =
M'_{\mathfrak p'} \otimes_{R'_{\mathfrak p'}} N'_{\mathfrak p'}
$$
模 $N'_{\mathfrak p'}$ 在局部环 $R'_{\mathfrak p'}$ 上投射，因而自由
（《代数》，定理 \ref{algebra-theorem-projective-free-over-local-ring}）。
因此，可以找到一个有限自由直和项 $F' \subset N'_{\mathfrak p'}$，
使得 $z' \in M'_{\mathfrak p'} \otimes_{R'_{\mathfrak p'}} F'$。
若 $F'$ 的秩为 $n$，则可推出
$\mathfrak p' R'_{\mathfrak p'} \in
\text{WeakAss}_{R'_{\mathfrak p'}}({M'_{\mathfrak p'}}^{\oplus n})$.
这蕴含
$\mathfrak p'R'_{\mathfrak p'} \in \text{WeakAss}(M'_{\mathfrak p'})$
，例如可由《代数》，引理 \ref{algebra-lemma-weakly-ass} 得到。
于是 $\mathfrak p' \in \text{WeakAss}_{R'}(M')$，进而得到
$\mathfrak p \in \text{WeakAss}_R(M)$。这就完成了引理中等价关系的
蕴含“$\Rightarrow$”的证明。

\medskip\noindent
假设 $\mathfrak p \in \text{WeakAss}_R(M)$ 且
$\overline{\mathfrak q} \in \text{Ass}_{\overline{S}}(\overline{N})$。
我们要证明 $\mathfrak q$ 是 $M \otimes_R N$ 的一个弱伴随素理想。
注意，$\overline{\mathfrak q}'$ 是
$\text{Ass}_{\overline{S}'}(\overline{N}')$ 的一个极大元。
这来自 (\ref{flat-equation-key-observation})，以及如下事实：
$\overline{S}'$ 中位于 $\overline{\mathfrak q}$ 之上的各素理想之间
不存在包含关系（因为 \'etale 态射的纤维是离散的，参见《态射》，
引理 \ref{morphisms-lemma-etale-over-field}）。因此，用
$R', S', \mathfrak p', \mathfrak q', M', N'$ 分别替换
$R, S, \mathfrak p, \mathfrak q, M, N$ 后，除引理的假设外，还可假设
\begin{enumerate}
\item $\mathfrak p \in \text{WeakAss}_R(M)$；
\item $\overline{\mathfrak q} \in \text{Ass}_{\overline{S}}(\overline{N})$；
\item $N$ 作为 $R$-模是投射的；
\item $\overline{\mathfrak q}$ 在
$\text{Ass}_{\overline{S}}(\overline{N})$ 中是极大的。
\end{enumerate}
还有一步约化：可以用 $R, S, M, N$ 在 $\mathfrak p$ 处的局部化分别
替换它们。这又给出一个条件，即
\begin{enumerate}
\item[(5)] $R$ 是极大理想为 $\mathfrak p$ 的局部环。
\end{enumerate}
最后我们将证明 (1) -- (5) 蕴含
$\mathfrak q \in \text{WeakAss}(M \otimes_R N)$.

\medskip\noindent
由于 $R$ 是局部的且 $\mathfrak p \in \text{WeakAss}_R(M)$，可以选取
$y \in M$，使其零化理想 $I$ 的根等于 $\mathfrak p$。
对某些 $\overline{g}_i \in \overline{S}$，写成
$\overline{\mathfrak q} = (\overline{g}_1, \ldots, \overline{g}_n)$。
选取映到 $\overline{g}_i$ 的 $g_i \in S$。那么
$\mathfrak q = \mathfrak pS + g_1S + \ldots + g_nS$。考虑映射
$$
\Psi : N/IN \longrightarrow (N/IN)^{\oplus n}, \quad
z \longmapsto (g_1z, \ldots, g_nz).
$$
这是投射 $R/I$-模之间的一个同态。由于 $\mathfrak p/I$ 局部幂零，
局部环 $R/I$ 是自伴随的（《代数进阶》，
定义 \ref{more-algebra-definition-auto-ass}）。映射
$\Psi \otimes \kappa(\mathfrak p)$ 不是单射，因为
$\overline{\mathfrak q} \in \text{Ass}_{\overline{S}}(\overline{N})$。
因此，《代数进阶》，引理 \ref{more-algebra-lemma-P-fPD-zero}
蕴含 $\Psi$ 不是单射。在 $\Psi$ 的核中选取一个非零元
$z \in N/IN$。由构造，零化理想 $J = \text{Ann}_S(z)$ 包含 $IS$
以及各 $g_i$。因此，$\sqrt{J} \subset S$ 包含 $\mathfrak q$。
设 $\mathfrak s \subset S$ 是包含 $J$ 的一个极小素理想。那么
$\mathfrak q \subset \mathfrak s$，$\mathfrak s$ 位于 $\mathfrak p$
之上，并且 $\mathfrak s \in \text{WeakAss}_S(N/IN)$。
最后一个事实来自弱伴随素理想的定义。将该引理已证明的
“$\Rightarrow$”部分应用于环映射 $R \to S$ 以及模 $R/I$ 和 $N$，
得到
$\overline{\mathfrak s} \in \text{Ass}_{\overline{S}}(\overline{N})$.
由于 $\overline{\mathfrak q} \subset \overline{\mathfrak s}$，
$\overline{\mathfrak q}$ 的极大性（见上面的条件 (4)）蕴含
$\overline{\mathfrak q} = \overline{\mathfrak s}$。这表明
$\mathfrak q = \mathfrak s$，从而得到所需结论。
\end{proof}

\begin{lemma}
\label{flat-lemma-bourbaki-finite-type-general-base-at-point}
设 $S$ 是一个概形。设 $f : X \to S$ 局部有限型。设 $x \in X$
的像为 $s \in S$。设 $\mathcal{F}$ 是 $X$ 上的一个有限型拟凝聚层，
$\mathcal{G}$ 是 $S$ 上的一个拟凝聚层。若 $\mathcal{F}$ 在 $x$ 处于
$S$ 上平坦，则
$$
x \in \text{WeakAss}_X(\mathcal{F} \otimes_{\mathcal{O}_X} f^*\mathcal{G})
\Leftrightarrow
s \in \text{WeakAss}_S(\mathcal{G})
\text{ 且 }
x \in \text{Ass}_{X_s}(\mathcal{F}_s).
$$
\end{lemma}

\begin{proof}
本段将问题约化到 $f$ 有限表示的情形。该问题在 $X$ 与 $S$ 上都是
局部的，故可假设 $X$ 与 $S$ 都是仿射的。写成
$X = \Spec(B)$、$S = \Spec(A)$，并写
$B = A[x_1, \ldots, x_n]/I$。换言之，我们得到 $S$ 上的一个闭浸入
$i : X \to \mathbf{A}^n_S$。记 $t = i(x) \in \mathbf{A}^n_S$。
注意，$i_*\mathcal{F}$ 是 $\mathbf{A}^n_S$ 上的一个有限型拟凝聚层，
它在 $t$ 处于 $S$ 上平坦；并且
$$
i_*(\mathcal{F} \otimes_{\mathcal{O}_X} f^*\mathcal{G}) =
i_*\mathcal{F} \otimes_{\mathcal{O}_{\mathbf{A}^n_S}} p^*\mathcal{G}
$$
其中 $p : \mathbf{A}^n_S \to S$ 是投影。注意，$t$ 是
$i_*(\mathcal{F} \otimes_{\mathcal{O}_X} f^*\mathcal{G})$ 的弱伴随点，
当且仅当 $x$ 是
$\mathcal{F} \otimes_{\mathcal{O}_X} f^*\mathcal{G}$ 的弱伴随点，参见
《除子》，引理 \ref{divisors-lemma-weakly-associated-finite}。
类似地，$x \in \text{Ass}_{X_s}(\mathcal{F}_s)$ 当且仅当
$t \in \text{Ass}_{\mathbf{A}^n_s}((i_*\mathcal{F})_s)$（参见
《代数》，引理 \ref{algebra-lemma-ass-quotient-ring}）。
因此，只需在 $X = \mathbf{A}^n_S$ 的情形下证明该引理。
于是可以假设 $X \to S$ 有限表示。

\medskip\noindent
本段将问题约化到 $\mathcal{F}$ 有限表示且在 $S$ 上平坦的情形。
按命题 \ref{flat-proposition-finite-type-flat-at-point} 选取一个初等
\'etale 邻域 $e : (S', s') \to (S, s)$ 以及一个开集
$V \subset X \times_S \Spec(\mathcal{O}_{S', s'})$。
设 $x' \in X' = X \times_S S'$ 是映到 $x$ 与 $s'$ 的唯一一点。
那么只需对 $X' \to S'$、$x'$、$s'$、
$(X' \to X)^*\mathcal{F}$ 以及 $e^*\mathcal{G}$ 证明该断言，参见
引理 \ref{flat-lemma-etale-weak-assassin-up-down}。
% P05-EN-CH38-0043: Frozen English omits "be" in "Let v in V the unique point".
设 $v \in V$ 是映到 $x'$ 的唯一一点，并设
$s' \in \Spec(\mathcal{O}_{S', s'})$ 为闭点。那么
$\mathcal{O}_{V, v} = \mathcal{O}_{X', x'}$，且
$\mathcal{O}_{\Spec(\mathcal{O}_{S', s'}), s'} =
\mathcal{O}_{S', s'}$；$\mathcal{F}$ 与 $\mathcal{G}$ 的拉回的茎也同样。
此外，$V_{s'} \subset X'_{s'}$ 是一个开子概形。由于一个点为弱伴随点
这一条件只依赖于层的茎，可以将
$X' \to S'$、$x'$、$s'$、$(X' \to X)^*\mathcal{F}$ 以及 $e^*\mathcal{G}$
替换为
$V \to \Spec(\mathcal{O}_{S', s'})$, $v$, $s'$, $(V \to X)^*\mathcal{F}$,
以及 $(\Spec(\mathcal{O}_{S', s'}) \to S)^*\mathcal{G}$。
因此可以假设 $f$ 有限表示，并且 $\mathcal{F}$ 有限表示且在 $S$ 上平坦。

\medskip\noindent
假设 $f$ 有限表示且 $\mathcal{F}$ 有限表示并在 $S$ 上平坦。
将 $X$ 与 $S$ 分别缩小为 $x$ 与 $s$ 的仿射邻域后，此情形由
引理 \ref{flat-lemma-weak-bourbaki-pre} 处理。
\end{proof}

\begin{lemma}
\label{flat-lemma-weak-bourbaki}
设 $R \to S$ 是一个本质有限型的环映射。设 $N$ 是某个有限 $S$-模的
一个局部化，并且在 $R$ 上平坦。设 $M$ 是一个 $R$-模。则
$$
\text{WeakAss}_S(M \otimes_R N)
=
\bigcup\nolimits_{\mathfrak p \in \text{WeakAss}_R(M)}
\text{Ass}_{S \otimes_R \kappa(\mathfrak p)}(N \otimes_R \kappa(\mathfrak p))
$$
\end{lemma}

\begin{proof}
该引理是引理 \ref{flat-lemma-bourbaki-finite-type-general-base-at-point}
用代数语言作出的翻译。略去翻译的细节。
\end{proof}

\begin{lemma}
\label{flat-lemma-bourbaki-finite-type-general-base}
设 $f : X \to S$ 是一个局部有限型的态射。设 $\mathcal{F}$ 是 $X$
上的一个有限型拟凝聚层，并且在 $S$ 上平坦。设 $\mathcal{G}$ 是 $S$
上的一个拟凝聚层。则有
$$
\text{WeakAss}_X(\mathcal{F} \otimes_{\mathcal{O}_X} f^*\mathcal{G}) =
\bigcup\nolimits_{s \in \text{WeakAss}_S(\mathcal{G})}
\text{Ass}_{X_s}(\mathcal{F}_s)
$$
\end{lemma}

\begin{proof}
这是引理 \ref{flat-lemma-bourbaki-finite-type-general-base-at-point}
的直接推论。
\end{proof}

\begin{theorem}
\label{flat-theorem-finite-type-flat}
\begin{slogan}
平坦轨迹是开的（非 Noether 情形）。
\end{slogan}
设 $f : X \to S$ 是一个概形态射。设 $\mathcal{F}$ 是一个拟凝聚
$\mathcal{O}_X$-模。假设
\begin{enumerate}
\item $X \to S$ 局部有限表示；
\item $\mathcal{F}$ 是有限型 $\mathcal{O}_X$-模；
\item $S$ 的弱伴随点集在 $S$ 中局部有限。
\end{enumerate}
则 $U = \{x \in X \mid \mathcal{F}\text{ 在 }x\text{ 处于 }S\text{ 上平坦}\}$
在 $X$ 中是开的，并且 $\mathcal{F}|_U$ 是有限表示的
$\mathcal{O}_U$-模且在 $S$ 上平坦。
\end{theorem}

\begin{proof}
设 $x \in X$，使得 $\mathcal{F}$ 在 $x$ 处于 $S$ 上平坦。
我们必须找到 $x$ 的一个开邻域，使 $\mathcal{F}$ 在该邻域上的限制
是一个有限表示的 $S$-平坦模。该问题在 $X$ 与 $S$ 上都是局部的，
故可假设 $X$ 与 $S$ 都是仿射的。由
引理 \ref{flat-lemma-finite-type-flat-at-point-local}，$\mathcal{F}_x$ 是
一个有限表示的 $\mathcal{O}_{X, x}$-模。因此，由
《代数》，引理 \ref{algebra-lemma-construct-fp-module-from-stalk}，
存在一个有限表示的 $\mathcal{O}_X$-模 $\mathcal{F}'$ 以及一个映射
$\varphi : \mathcal{F}' \to \mathcal{F}$，它诱导同构
$\varphi_x : \mathcal{F}'_x \to \mathcal{F}_x$。特别地，可见
$\mathcal{F}'$ 在 $x$ 处于 $S$ 上平坦。因此，由平坦性的开放性
（《态射进阶》，定理 \ref{more-morphisms-theorem-openness-flatness}），
缩小 $X$ 后，可以假设 $\mathcal{F}'$ 在 $S$ 上平坦。
由于 $\mathcal{F}$ 为有限型，再次缩小 $X$ 后可以假设 $\varphi$ 是
满射，参见《模》，引理 \ref{modules-lemma-finite-type-surjective-on-stalk}；
也可以使用 Nakayama 引理（《代数》，引理 \ref{algebra-lemma-NAK}）。
由引理 \ref{flat-lemma-bourbaki-finite-type-general-base}，有
$$
\text{WeakAss}_X(\mathcal{F}') \subset
\bigcup\nolimits_{s \in \text{WeakAss}(S)} \text{Ass}_{X_s}(\mathcal{F}'_s)
$$
由假设，$\text{WeakAss}(S)$ 是有限的；又由《除子》，
引理 \ref{divisors-lemma-finite-ass}，
$\text{Ass}_{X_s}(\mathcal{F}'_s)$ 是有限的，故可推出
$\text{WeakAss}_X(\mathcal{F}')$ 是有限的。利用《代数》，
引理 \ref{algebra-lemma-silly}，再次缩小 $X$ 后，可以假设
$\text{WeakAss}_X(\mathcal{F}')$ 包含于 $x$ 的一般化点集。
现在考虑 $\mathcal{K} = \Ker(\varphi)$。我们有
$\text{WeakAss}_X(\mathcal{K}) \subset \text{WeakAss}_X(\mathcal{F}')$
（由《除子》，引理 \ref{divisors-lemma-ses-weakly-ass}）。
另一方面，$\varphi_x$ 是同构，而且对于所有 $x' \leadsto x$，
$\varphi_{x'}$ 也是同构。我们推出
$\text{WeakAss}_X(\mathcal{K}) = \emptyset$，从而由《除子》，
引理 \ref{divisors-lemma-weakly-ass-zero}，有 $\mathcal{K} = 0$。
\end{proof}

\begin{lemma}
\label{flat-lemma-finite-type-flat-algebra}
设 $R \to S$ 是一个有限表示的环映射。设 $M$ 是一个有限 $S$-模。
假设 $\text{WeakAss}_R(R)$ 是有限的。则
$$
U = \{\mathfrak q \subset S \mid M_{\mathfrak q}\text{ 在 }R\text{ 上平坦}\}
$$
在 $\Spec(S)$ 中是开的；并且对于每个满足 $D(g) \subset U$ 的
$g \in S$，局部化 $M_g$ 是一个有限表示的 $S_g$-模，且在 $R$ 上平坦。
\end{lemma}

\begin{proof}
这立即来自定理 \ref{flat-theorem-finite-type-flat}。
\end{proof}

\begin{lemma}
\label{flat-lemma-finite-type-flat-X}
设 $f : X \to S$ 是一个局部有限型的概形态射。假设 $S$ 的弱伴随点集
在 $S$ 中局部有限。则 $f$ 平坦的各点 $x \in X$ 构成一个开子概形
$U \subset X$，并且 $U \to S$ 平坦且局部有限表示。
\end{lemma}

\begin{proof}
该问题在 $X$ 与 $S$ 上都是局部的，故可假设 $X$ 与 $S$ 都是仿射的。
那么 $X \to S$ 对应于一个有限型环映射 $A \to B$。选取一个满射
$A[x_1, \ldots, x_n] \to B$，并将 $B$ 视为一个
$A[x_1, \ldots, x_n]$-模。应用
引理 \ref{flat-lemma-finite-type-flat-algebra} 即完成证明。
\end{proof}

\begin{lemma}
\label{flat-lemma-finite-type-flat-over-integral}
设 $f : X \to S$ 是一个局部有限型且平坦的概形态射。若 $S$ 是整的，
则 $f$ 局部有限表示。
\end{lemma}

\begin{proof}
这是引理 \ref{flat-lemma-finite-type-flat-X} 的特例。
\end{proof}

\begin{proposition}
\label{flat-proposition-flat-finite-type-finite-presentation-domain}
设 $R$ 是一个整环。设 $R \to S$ 是一个有限型环映射。
设 $M$ 是一个有限 $S$-模。
\begin{enumerate}
\item 若 $S$ 在 $R$ 上平坦，则 $S$ 是一个有限表示的 $R$-代数。
\item 若 $M$ 作为 $R$-模平坦，则 $M$ 作为 $S$-模有限表示。
\end{enumerate}
\end{proposition}

\begin{proof}
(1) 是引理 \ref{flat-lemma-finite-type-flat-over-integral} 的一个特例。
对于 (2)，选取一个满射 $R[x_1, \ldots, x_n] \to S$。
由引理 \ref{flat-lemma-finite-type-flat-algebra}，可见 $M$ 作为
$R[x_1, \ldots, x_n]$-模有限表示。最后应用《代数》，
引理 \ref{algebra-lemma-finitely-presented-over-subring}。
\end{proof}

\begin{lemma}[定理 \ref{flat-theorem-finite-type-flat} 的有限型版本]
\label{flat-lemma-finite-type-flat}
设 $f : X \to S$ 是一个概形态射。设 $\mathcal{F}$ 是一个拟凝聚
$\mathcal{O}_X$-模。假设
\begin{enumerate}
\item $X \to S$ 局部有限型；
\item $\mathcal{F}$ 是有限型 $\mathcal{O}_X$-模；
\item $S$ 的弱伴随点集在 $S$ 中局部有限。
\end{enumerate}
则 $U = \{x \in X \mid \mathcal{F}\text{ 在 }x\text{ 处于 }S\text{ 上平坦}\}$
在 $X$ 中是开的，并且 $\mathcal{F}|_U$ 在 $S$ 上平坦且局部相对于
$S$ 有限表示（参见《态射进阶》，定义
\ref{more-morphisms-definition-relatively-finitely-presented-sheaf}).
\end{lemma}

\begin{proof}
该问题在 $X$ 与 $S$ 上都是局部的。因此可以假设 $X$ 与 $S$ 都是
仿射的。那么可以选取一个闭浸入 $i : X \to \mathbf{A}^n_S$。
将定理 \ref{flat-theorem-finite-type-flat} 应用于
$X' = \mathbf{A}^n_S \to S$ 以及有限型拟凝聚模
$\mathcal{F}' = i_*\mathcal{F}$，可得
$$
U' = \{x' \in X' \mid \mathcal{F}'\text{ 在 }x'\text{ 处于 }S\text{ 上平坦}\}
$$
在 $X'$ 中是开的，并且 $\mathcal{F}'|_{U'}$ 有限表示。
由于 $\mathcal{F}'$ 在 $X' \setminus i(X)$ 上的限制为零，且对于所有
$x \in X$ 都有 $\mathcal{F}'_{i(x)} \cong \mathcal{F}_x$，可见
$$
U' = i(U) \amalg (X' \setminus i(X))
$$
因此 $U = i^{-1}(U')$ 是开的。此外，显然
$\mathcal{F}'|_{U'} = (i|_U)_*(\mathcal{F}|_U)$.
故由《态射进阶》，引理
\ref{more-morphisms-lemma-relative-finite-presentation} 和
\ref{more-morphisms-lemma-finite-morphism-relative-finite-presentation}.
可推出 $\mathcal{F}|_U$ 相对于 $S$ 有限表示。
\end{proof}

\begin{lemma}
\label{flat-lemma-finite-type-flat-algebra-bis}
设 $R \to S$ 是一个有限型环映射。设 $M$ 是一个有限 $S$-模。
假设 $\text{WeakAss}_R(R)$ 是有限的。则
$$
U = \{\mathfrak q \subset S \mid M_{\mathfrak q}\text{ 在 }R\text{ 上平坦}\}
$$
在 $\Spec(S)$ 中是开的；并且对于每个满足 $D(g) \subset U$ 的
$g \in S$，局部化 $M_g$ 在 $R$ 上平坦，且作为 $S_g$-模相对于
$R$ 有限表示（参见《代数进阶》，定义
\ref{more-algebra-definition-relatively-finitely-presented}).
\end{lemma}

\begin{proof}
这是引理 \ref{flat-lemma-finite-type-flat} 用代数语言作出的翻译。
\end{proof}








\section{相对纯模的例子}
\label{flat-section-examples-pure-modules}

\noindent
在本节中，我们讨论一些结果的例子；它们将为稍后引入的
{\it 相对纯模}概念和 {\it 非纯性}概念提供动机。每个例子都表述为一个
引理。请注意，关于伴随素理想的条件与下列引理中出现的条件相似：
引理 \ref{flat-lemma-base-change-universally-flat}、
\ref{flat-lemma-universally-injective-to-completion},
\ref{flat-lemma-universally-injective-to-completion-flat} 和
\ref{flat-lemma-flat-pure-over-complete-projective}。
另请参见《代数》，引理 \ref{algebra-lemma-compare-relative-assassins}，
其中有相关讨论。

\begin{lemma}
\label{flat-lemma-explain-why-pure}
设 $R$ 是极大理想为 $\mathfrak m$ 的局部环。设 $R \to S$ 是一个环映射。
设 $N$ 是一个 $S$-模。假设
\begin{enumerate}
\item $N$ 作为 $R$-模是投射的；
\item $S/\mathfrak mS$ 是 Noether 环，并且 $N/\mathfrak mN$ 是有限的
$S/\mathfrak mS$-模。
\end{enumerate}
则对于任意素理想 $\mathfrak q \subset S$，若它是
$N \otimes_R \kappa(\mathfrak p)$ 的伴随素理想，其中
$\mathfrak p = R \cap \mathfrak q$，则有
$\mathfrak q + \mathfrak m S \not = S$。
\end{lemma}

\begin{proof}
注意，引理 \ref{flat-lemma-homothety-spectrum} 与
\ref{flat-lemma-invert-universally-injective} 的假设都满足。以下不再另行说明，
直接使用这两个引理的结论。令 $\Sigma \subset S$ 为在 $N/\mathfrak mN$
上不是零因子的元素所成的乘法集。映射
$N \to \Sigma^{-1}N$ 是 $R$-普遍单射的。因此，任何作为
$N \otimes_R \kappa(\mathfrak p)$ 伴随素理想的
$\mathfrak q \subset S$，也是
$\Sigma^{-1}N \otimes_R \kappa(\mathfrak p)$ 的伴随素理想。
显然，这蕴含 $\mathfrak q$ 对应于 $\Sigma^{-1}S$ 的一个素理想。
于是 $\mathfrak q \subset \mathfrak q'$，其中 $\mathfrak q'$ 对应于
$N/\mathfrak mN$ 的一个伴随素理想，结论得证。
\end{proof}

\noindent
以下引理给出这一现象的另一个（稍显无聊的）例子。

\begin{lemma}
\label{flat-lemma-explain-why-pure-complete}
设 $R$ 是一个环。设 $I \subset R$ 是一个理想。设 $R \to S$ 是一个环映射。
设 $N$ 是一个 $S$-模。若 $N$ 是 $I$-进完备的，则对于任意 $R$-模 $M$，
以及任意素理想 $\mathfrak q \subset S$，若它是
$N \otimes_R M$ 的伴随素理想，则有 $\mathfrak q + I S \not = S$。
\end{lemma}

\begin{proof}
令 $S^\wedge$ 表示 $S$ 的 $I$-进完备化。注意，$N$ 是一个
$S^\wedge$-模，因而 $N \otimes_R M$ 也是一个 $S^\wedge$-模。
取 $z \in N \otimes_R M$，使得 $\mathfrak q = \text{Ann}_S(z)$。
由于 $z \not = 0$，可见 $\text{Ann}_{S^\wedge}(z) \not = S^\wedge$，
从而 $\mathfrak q S^\wedge \not = S^\wedge$。因此存在一个极大理想
$\mathfrak m \subset S^\wedge$，使得 $\mathfrak q S^\wedge \subset \mathfrak m$。
由《代数》，引理 \ref{algebra-lemma-radical-completion}，
$IS^\wedge \subset \mathfrak m$，结论得证。
\end{proof}

\noindent
注意，以下引理给出引理 \ref{flat-lemma-explain-why-pure} 的另一种证明，
因为局部环上的投射模是自由模，参见《代数》，定理
\ref{algebra-theorem-projective-free-over-local-ring}。

\begin{lemma}
\label{flat-lemma-explain-why-pure-direct-sum-finite-modules}
设 $R$ 是极大理想为 $\mathfrak m$ 的局部环。设 $R \to S$ 是一个环映射。
设 $N$ 是一个 $S$-模，并假设 $N$ 作为 $R$-模同构于有限 $R$-模的一个直和。
则对于任意 $R$-模 $M$，
以及任意素理想 $\mathfrak q \subset S$，若它是 $N \otimes_R M$ 的伴随素理想，
则有 $\mathfrak q + \mathfrak m S \not = S$。
\end{lemma}

\begin{proof}
写 $N = \bigoplus_{i \in I} M_i$，其中每个 $M_i$ 都是有限的 $R$-模。
设 $M$ 是一个 $R$-模，设 $\mathfrak q \subset S$ 是
$N \otimes_R M$ 的伴随素理想，并且 $\mathfrak q + \mathfrak m S = S$。
取 $z \in N \otimes_R M$，使得 $\mathfrak q = \text{Ann}_S(z)$。
稍微调整直和分解后，可以假设对于某个 $1 \in I$ 有
$z \in M_1 \otimes_R M$。写成
$1 = f + \sum x_j g_j$，其中 $f \in \mathfrak q$、$x_j \in \mathfrak m$、
$g_j \in S$。对于任意 $g \in S$，记 $g'$ 为以下 $R$-线性映射：
$$
M_1 \to N \xrightarrow{g} N \to M_1
$$
其中第一个箭头是包含映射，第二个箭头是乘以 $g$，第三个箭头是投影映射。
由于每个 $x_j \in R$，我们得到等式
$$
f' + \sum x_j g'_j = \text{id}_{M_1} \in \text{End}_R(M_1)
$$
由 Nakayama 引理（《代数》，引理 \ref{algebra-lemma-NAK}），可见 $f'$
是满射；于是由《代数》，引理 \ref{algebra-lemma-fun}，可见 $f'$ 是同构。
特别地，以下映射
$$
M_1 \otimes_R M \to N \otimes_R M \xrightarrow{f} N \otimes_R M
\to M_1 \otimes_R M
$$
是同构。这与 $fz = 0$ 的假设矛盾。
\end{proof}

\begin{lemma}
\label{flat-lemma-explain-why-pure-ML}
设 $R$ 是极大理想为 $\mathfrak m$ 的 Hensel 局部环。设 $R \to S$ 是一个环映射。
设 $N$ 是一个 $S$-模。假设 $N$ 作为 $R$-模是可数生成的且为
Mittag-Leffler 模。则对于任意 $R$-模 $M$ 以及任意素理想
$\mathfrak q \subset S$，若它是 $N \otimes_R M$ 的伴随素理想，则有
$\mathfrak q + \mathfrak m S \not = S$.
\end{lemma}

\begin{proof}
由《代数》，引理 \ref{algebra-lemma-split-ML-henselian}，该引理约化为
引理 \ref{flat-lemma-explain-why-pure-direct-sum-finite-modules}。
\end{proof}

\noindent
设 $f : X \to S$ 是一个概形态射，$\mathcal{F}$ 是 $X$ 上的一个
拟凝聚模。设 $\xi \in \text{Ass}_{X/S}(\mathcal{F})$，并令
$Z = \overline{\{\xi\}}$。图示如下：
$$
\xymatrix{
\xi \ar@{|->}[d] & Z \ar[r] \ar[rd] & X \ar[d]^f \\
f(\xi) & & S
}
$$
注意，$f(Z) \subset \overline{\{f(\xi)\}}$，并且 $f(Z)$ 是闭集
当且仅当等式成立，即 $f(Z) = \overline{\{f(\xi)\}}$。
由引理 \ref{flat-lemma-explain-why-pure} 可知，若 $S$、$X$ 都是仿射的，
纤维 $X_s$ 是 Noether 概形，$\mathcal{F}$ 是有限型的，并且
$\Gamma(X, \mathcal{F})$ 是投射的 $\Gamma(S, \mathcal{O}_S)$-模，
则 $f(Z) = \overline{\{f(\xi)\}}$ 是一个闭子集。
对于下列情形，也有略有不同但类似的陈述：
引理 \ref{flat-lemma-explain-why-pure-complete}、
\ref{flat-lemma-explain-why-pure-direct-sum-finite-modules} 和
\ref{flat-lemma-explain-why-pure-ML}.




\section{非纯性}
\label{flat-section-impure}

\noindent
我们想用 $\text{Ass}_{X/S}(\mathcal{F})$ 中点的特化来形式化
第 \ref{flat-section-examples-pure-modules} 节中举例说明的现象。
我们还希望在 $s \in S$ 的附近局部地工作。为此作出以下定义。

\begin{situation}
\label{flat-situation-pre-pure}
这里 $S$、$X$ 是概形，且 $f : X \to S$ 是有限型态射。
此外，$\mathcal{F}$ 是有限型拟凝聚 $\mathcal{O}_X$-模。
最后，$s$ 是 $S$ 的一个点。
\end{situation}

\noindent
在此情形下，考虑一个态射 $g : T \to S$、一个满足 $g(t) = s$ 的
点 $t \in T$、一个特化 $t' \leadsto t$，以及 $X$ 的基变换中位于
$t'$ 之上的一点 $\xi \in X_T$。图示如下：
\begin{equation}
\label{flat-equation-impurity}
\vcenter{
\xymatrix{
\xi \ar@{|->}[d] & \\
t' \ar@{~>}[r] & t \ar@{|->}[r] & s
}
}
\quad\quad
\vcenter{
\xymatrix{
X_T \ar[d] \ar[r] & X \ar[d] \\
T \ar[r]^g & S
}
}
\end{equation}
此外，记 $\mathcal{F}_T$ 为 $\mathcal{F}$ 到 $X_T$ 的拉回。

\begin{definition}
\label{flat-definition-impurity}
在情形 \ref{flat-situation-pre-pure} 中，如果
图 (\ref{flat-equation-impurity}) 满足
$\xi \in \text{Ass}_{X_T/T}(\mathcal{F}_T)$ 且
$\overline{\{\xi\}} \cap X_t = \emptyset$，则称它定义了
$\mathcal{F}$ 在 $s$ 之上的一个 {\it 非纯性}。
我们将用如下说法表示这一点：“设
$(g : T \to S, t' \leadsto t, \xi)$ 是 $\mathcal{F}$ 在 $s$ 之上的
一个非纯性”。
\end{definition}

\begin{lemma}
\label{flat-lemma-impure-finite-presentation}
在情形 \ref{flat-situation-pre-pure} 中，若存在 $\mathcal{F}$ 在 $s$ 之上的
一个非纯性，则存在一个 $\mathcal{F}$ 在 $s$ 之上的非纯性
$(g : T \to S, t' \leadsto t, \xi)$，使得 $g$ 局部有限表示，且 $t$
是 $g$ 在 $s$ 之上的纤维中的闭点。
\end{lemma}

\begin{proof}
设 $(g : T \to S, t' \leadsto t, \xi)$ 是 $\mathcal{F}$ 在 $s$ 之上的
任意非纯性。将《极限》，引理 \ref{limits-lemma-separate} 应用于
$t \in T$ 与 $Z = \overline{\{\xi\}}$，得到 $t$ 的一个开邻域
$V \subset T$、一个交换图
$$
\xymatrix{
V \ar[d] \ar[r]_a & T' \ar[d]^b \\
T \ar[r]^g & S,
}
$$
以及一个闭子概形 $Z' \subset X_{T'}$，满足
\begin{enumerate}
\item 态射 $b : T' \to S$ 局部有限表示；
\item $Z' \cap X_{a(t)} = \emptyset$；
\item $Z \cap X_V$ 通过态射 $X_V \to X_{T'}$ 映入 $Z'$。
\end{enumerate}
由于 $t'$ 特化到 $t$，可以用 $t$ 的开邻域 $V$ 替换 $T$。
于是有交换图
$$
\xymatrix{
X_T \ar[d] \ar[r] &
X_{T'} \ar[d] \ar[r] &
X \ar[d] \\
T \ar[r]^a & T' \ar[r]^b & S
}
$$
其中 $b \circ a = g$。令 $\xi' \in X_{T'}$ 表示 $\xi$ 的像。由
《除子》，引理 \ref{divisors-lemma-base-change-relative-assassin}，
可见 $\xi' \in \text{Ass}_{X_{T'}/T'}(\mathcal{F}_{T'})$。
此外，按构造，$\overline{\{\xi'\}}$ 的闭包包含于闭子集 Z' 中，
而 $Z'$ 避开纤维 $X_{a(t)}$。因此
$(T' \to S, a(t') \leadsto a(t), \xi')$ 是 $\mathcal{F}$ 在 $s$ 之上的
一个非纯性。

\medskip\noindent
因此可以假设 $g : T \to S$ 局部有限表示。令
$Z = \overline{\{\xi\}}$。由假设 $Z_t = \emptyset$。由《态射进阶》，
引理 \ref{more-morphisms-lemma-empty-generic-fibre}，这意味着对于
$\overline{\{t\}}$ 的某个开子集中的 $t''$，有 $Z_{t''} = \emptyset$。
由于 $T \to S$ 在 $s$ 上的纤维是 Jacobson 概形，参见《态射》，
引理 \ref{morphisms-lemma-ubiquity-Jacobson-schemes}，可以找到
$t'' \in \overline{\{t\}}$ 中的一个闭点，使 $Z_{t''} = \emptyset$。
于是 $(g : T \to S, t' \leadsto t'', \xi)$ 就是所需的非纯性。
\end{proof}

\begin{lemma}
\label{flat-lemma-impure-limit}
在情形 \ref{flat-situation-pre-pure} 中，设
$(g : T \to S, t' \leadsto t, \xi)$ 是 $\mathcal{F}$ 在 $s$ 之上的一个
非纯性。假设 $T = \lim_{i \in I} T_i$ 是 $S$ 上仿射概形的一个有向极限。
则对于某个 $i$，三元组
$(T_i \to S, t'_i \leadsto t_i, \xi_i)$ 是 $\mathcal{F}$ 在 $s$ 之上的
一个非纯性。
\end{lemma}

\begin{proof}
命题中的记号含义如下：令 $p_i : T \to T_i$ 为投影态射，令
$t_i = p_i(t)$ 且 $t'_i = p_i(t')$。最后，$\xi_i \in X_{T_i}$ 是
$\xi$ 的像。由《除子》，引理 \ref{divisors-lemma-base-change-relative-assassin}，
$\xi_i$ 确实是 $\mathcal{F}_{T_i}$ 关于 $T_i$ 的相对伴随点。
因此唯一需要证明的是：对某个 $i$，有
$\overline{\{\xi_i\}} \cap X_{t_i} = \emptyset$。

\medskip\noindent
第一种证明。令 $Z_i = \overline{\{\xi_i\}} \subset X_{T_i}$，
$Z = \overline{\{\xi\}} \subset X_T$，并赋予它们约化的诱导概形结构。
由《极限》，引理 \ref{limits-lemma-inverse-limit-irreducibles}，
$Z = \lim Z_i$。取一个域 $k$ 以及一个像为 $t$ 的态射
$\Spec(k) \to T$。则
$$
\emptyset =
Z \times_T \Spec(k) = (\lim Z_i) \times_{(\lim T_i)} \Spec(k)
= \lim Z_i \times_{T_i} \Spec(k)
$$
这是因为极限与纤维积交换（极限与极限交换）。每个
$Z_i \times_{T_i} \Spec(k)$ 都是拟紧的，因为 $X_{T_i} \to T_i$
有限型，从而 $Z_i \to T_i$ 有限型。由《极限》，引理
\ref{limits-lemma-limit-nonempty}，因此对于某个 $i$，
$Z_i \times_{T_i} \Spec(k)$ 为空。由于复合态射
$\Spec(k) \to T \to T_i$ 的像是 $t_i$，所需结论成立。

\medskip\noindent
第二种证明。置 $Z = \overline{\{\xi\}}$。将《极限》，引理
\ref{limits-lemma-separate} 应用于此情形，得到 $t$ 的一个开邻域
$V \subset T$、一个交换图
$$
\xymatrix{
V \ar[d] \ar[r]_a & T' \ar[d]^b \\
T \ar[r]^g & S,
}
$$
以及一个闭子概形 $Z' \subset X_{T'}$，满足
\begin{enumerate}
\item 态射 $b : T' \to S$ 局部有限表示；
\item $Z' \cap X_{a(t)} = \emptyset$；
\item $Z \cap X_V$ 通过态射 $X_V \to X_{T'}$ 映入 $Z'$。
\end{enumerate}
可以假设 $V$ 是 $T$ 的仿射开子集，因此由《极限》，引理
\ref{limits-lemma-descend-opens} 和 \ref{limits-lemma-limit-affine}，
可以找到某个 $i$ 和一个仿射开子集 $V_i \subset T_i$，使得
$V = f_i^{-1}(V_i)$。由《极限》，命题
\ref{limits-proposition-characterize-locally-finite-presentation}，
在必要时将 $i$ 增大后，可以找到态射 $a_i : V_i \to T'$，使得
$a = a_i \circ f_i|_V$。诱导的态射 $X_{V_i} \to X_{T'}$ 将 $\xi_i$
映入 $Z'$。由于 $Z' \cap X_{a(t)} = \emptyset$，可得
$(T_i \to S, t'_i \leadsto t_i, \xi_i)$ 是 $\mathcal{F}$ 在 $s$ 之上的
一个非纯性。
\end{proof}

\begin{lemma}
\label{flat-lemma-quasi-finite-impurity-elementary}
在情形 \ref{flat-situation-pre-pure} 中，若存在一个 $\mathcal{F}$ 在 $s$ 之上的
非纯性 $(g : T \to S, t' \leadsto t, \xi)$，且 $g$ 在 $t$ 处拟有限，
则存在一个非纯性 $(g : T \to S, t' \leadsto t, \xi)$，使
$(T, t) \to (S, s)$ 是初等 \'etale 邻域。
\end{lemma}

\begin{proof}
设 $(g : T \to S, t' \leadsto t, \xi)$ 是 $\mathcal{F}$ 在 $s$ 之上的一个
非纯性，且 $g$ 在 $t$ 处拟有限。缩小 $T$ 后，可以假设 $g$ 局部有限型。
将《态射进阶》，引理
\ref{more-morphisms-lemma-etale-makes-quasi-finite-finite-at-point}
应用于 $T \to S$ 与 $t \mapsto s$，得到图
$$
\xymatrix{
T \ar[d] & T \times_S U \ar[l] \ar[d] & V \ar[l] \ar[ld] \\
S & U \ar[l]
}
$$
其中 $(U, u) \to (S, s)$ 是一个初等 \'etale 邻域，且
$V \subset T \times_S U$ 是 $v = (t, u)$ 的一个开邻域，满足
$V \to U$ 有限，并且 $v$ 是 $V$ 中位于 $u$ 之上的唯一点。
由于态射 $V \to T$ 是 \'etale 的，因而平坦，可见存在一个特化
$v' \leadsto v$，使得 $v' \mapsto t'$。注意
$\kappa(t') \subset \kappa(v')$ 是有限可分扩张。任取
任取点 $\zeta \in X_{v'}$，它映到 $\xi \in X_{t'}$。由《除子》，引理
\ref{divisors-lemma-base-change-relative-assassin}，可见
$\zeta \in \text{Ass}_{X_V/V}(\mathcal{F}_V)$。
此外，$\overline{\{\zeta\}}$ 不与纤维 $X_v$ 相交，因为按假设
$\overline{\{\xi\}}$ 不与 $X_t$ 相交。换言之，
$(V \to S, v' \leadsto v, \zeta)$ 是 $\mathcal{F}$ 在 $S$ 之上的一个
非纯性。

\medskip\noindent
接着，令 $u' \in U'$ 为 $v'$ 的像，令 $\theta \in X_U$ 为 $\zeta$ 的像。
则 $\theta \mapsto u'$ 且 $u' \leadsto u$。由《除子》，引理
\ref{divisors-lemma-base-change-relative-assassin}，可见
$\theta \in \text{Ass}_{X_U/U}(\mathcal{F})$。
此外，由于 $\pi : X_V \to X_U$ 是有限的，有
$\pi\big(\overline{\{\zeta\}}\big) = \overline{\{\pi(\zeta)\}}$。
由于 $v$ 是 $V$ 中位于 $u$ 之上的唯一点，可见
$X_u \cap \overline{\{\pi(\zeta)\}} = \emptyset$，因为
$X_v \cap \overline{\{\zeta\}} = \emptyset$。因此
$(U \to S, u' \leadsto u, \theta)$ 是 $\mathcal{F}$ 在 $s$ 之上的一个
非纯性，结论得证。
\end{proof}

\begin{lemma}
\label{flat-lemma-Noetherian-impurity-quasi-finite}
在情形 \ref{flat-situation-pre-pure} 中，假设 $S$ 局部 Noether。
若存在 $\mathcal{F}$ 在 $s$ 之上的一个非纯性，则存在一个
$\mathcal{F}$ 在 $s$ 之上的非纯性 $(g : T \to S, t' \leadsto t, \xi)$，
使得 $g$ 在 $t$ 处拟有限。
\end{lemma}

\begin{proof}
可以用 $s$ 的一个仿射邻域替换 $S$。由引理
\ref{flat-lemma-impure-finite-presentation}，可以假设存在一个非纯性
$(g : T \to S, t' \leadsto t, \xi)$，其中 $g$ 局部有限型，且 $t$ 是
$g$ 在 $s$ 之上的纤维中的闭点。可以用
$\overline{\{t'\}}$ 上的约化诱导概形结构替换 $T$。令
$Z = \overline{\{\xi\}} \subset X_T$。由假设 $Z_t = \emptyset$，且
$Z \to T$ 的像包含 $t'$。由《态射进阶》，引理
\ref{more-morphisms-lemma-relative-assassin-in-neighbourhood}，存在
Z 的一个非空开子集 $V \subset Z$，使得对于任意 $w \in f(V)$，
$V_w$ 的任意一般点 $\xi'$ 都属于
$\text{Ass}_{X_T/T}(\mathcal{F}_T)$。由《态射进阶》，引理
\ref{more-morphisms-lemma-nonempty-generic-fibre}，存在一个非空开集
$W \subset T$，满足 $W \subset f(V)$。由《态射进阶》，引理
\ref{more-morphisms-lemma-quasi-finite-quasi-section-meeting-nearby-open-X}，
存在一个闭子概形 $T' \subset T$，使得 $t \in T'$，$T' \to S$ 在 $t$
处拟有限，并且存在 $z \in T' \cap W$，满足 $z \leadsto t$ 且不映到 $s$。
取非空概形 $V_z$ 的任意一般点 $\xi'$。则
$(T' \to S, z \leadsto t, \xi')$ 是所需的非纯性。
\end{proof}

\noindent
下文将使用 $S$ 在 $s$ 处的 Hensel 化
$S^h = \Spec(\mathcal{O}_{S, s}^h)$，参见《\'Etale 上同调》，定义
\ref{etale-cohomology-definition-etale-local-rings}。
由于 $S^h \to S$ 将 $S^h$ 的闭点映到 $s$，并诱导剩余域的同构，
我们也用 $s \in S^h$ 表示这个闭点。因此
$(S^h, s) \to (S, s)$ 是带基点概形的一个态射。

\begin{lemma}
\label{flat-lemma-impurity-on-henselization}
在情形 \ref{flat-situation-pre-pure} 中，若存在 $\mathcal{F}$ 在 $s$ 之上的
非纯性 $(S^h \to S, s' \leadsto s, \xi)$，则存在一个
$\mathcal{F}$ 在 $s$ 之上的非纯性 $(T \to S, t' \leadsto t, \xi)$，
其中 $(T, t) \to (S, s)$ 是初等 \'etale 邻域。
\end{lemma}

\begin{proof}
可以用 $s$ 的一个仿射邻域替换 $S$。设 $S = \Spec(A)$，且 $s$ 对应于
$\mathfrak p \subset A$。则
$\mathcal{O}_{S, s}^h = \colim_{(T, t)} \Gamma(T, \mathcal{O}_T)$，
其中极限遍历 $(S, s)$ 的仿射初等 \'etale 邻域 $(T, t)$ 所成的
余滤范畴的相反范畴，参见《态射进阶》，引理
\ref{more-morphisms-lemma-describe-henselization} 及其证明。因此
$S^h = \lim_i T_i$，由引理 \ref{flat-lemma-impure-limit} 即得结论。
\end{proof}

\begin{lemma}
\label{flat-lemma-pure-along-X-s}
在情形 \ref{flat-situation-pre-pure} 中，以下各条件等价：
\begin{enumerate}
\item 存在 $\mathcal{F}$ 在 $s$ 之上的非纯性
$(S^h \to S, s' \leadsto s, \xi)$，其中 $S^h$ 是 $S$ 在 $s$ 处的
Hensel 化；
\item 存在 $\mathcal{F}$ 在 $s$ 之上的非纯性
$(T \to S, t' \leadsto t, \xi)$，使 $(T, t) \to (S, s)$ 是初等
\'etale 邻域；
\item 存在 $\mathcal{F}$ 在 $s$ 之上的非纯性
$(T \to S, t' \leadsto t, \xi)$，使 $T \to S$ 在 $t$ 处拟有限。
\end{enumerate}
\end{lemma}

\begin{proof}
由于 \'etale 态射局部拟有限，(2) 蕴含 (3) 是显然的。在引理
\ref{flat-lemma-quasi-finite-impurity-elementary} 中已经证明 (3) 蕴含 (2)。
在引理 \ref{flat-lemma-impurity-on-henselization} 中已经证明 (1) 蕴含 (2)。
最后，若 $(T \to S, t' \leadsto t, \xi)$ 是 $\mathcal{F}$ 在 $s$ 之上的
一个非纯性，且 $(T, t) \to (S, s)$ 是初等 \'etale 邻域，则可以选取
结构态射 $S^h \to T \to S$ 的分解 $S^h \to S$。
选取任意点 $s' \in S^h$，它映到 $t'$，并选取任意
$\xi' \in X_{s'}$，它映到 $\xi \in X_{t'}$。则
$(S^h \to S, s' \leadsto s, \xi')$ 是 $\mathcal{F}$ 在 $s$ 之上的
一个非纯性。细节略去。
\end{proof}





\section{相对纯模}
\label{flat-section-pure}

\noindent
相对于底空间纯的模这一概念由 \cite{GruRay} 引入。

\begin{definition}
\label{flat-definition-pure}
设 $f : X \to S$ 是有限型概形态射。
设 $\mathcal{F}$ 是有限型拟凝聚 $\mathcal{O}_X$-模。
\begin{enumerate}
\item 设 $s \in S$。若不存在 $\mathcal{F}$ 在 $s$ 之上的非纯性
$(g : T \to S, t' \leadsto t, \xi)$，且 $(T, t) \to (S, s)$ 是一个
初等 \'etale 邻域，则称 $\mathcal{F}$ {\it 沿 $X_s$ 纯}。
\item 若不存在 $\mathcal{F}$ 在 $s$ 之上的任何非纯性，则称
$\mathcal{F}$ {\it 沿 $X_s$ 普遍纯}。
\item 若 $\mathcal{O}_X$ 沿 $X_s$ 纯，则称 $X$ {\it 沿 $X_s$ 纯}。
\item 若对每个 $s \in S$，$\mathcal{F}$ 都沿 $X_s$ 普遍纯，则称
$\mathcal{F}$ {\it 普遍 $S$-纯}，或称其 {\it 相对于 $S$ 普遍纯}。
\item 若对每个 $s \in S$，$\mathcal{F}$ 都沿 $X_s$ 纯，则称
$\mathcal{F}$ {\it $S$-纯}，或称其 {\it 相对于 $S$ 纯}。
\item 若 $\mathcal{O}_X$ 相对于 $S$ 纯，则称 $X$ {\it $S$-纯}，或称其
{\it 相对于 $S$ 纯}。
\end{enumerate}
\end{definition}

\noindent
这里我们有意将范围限制为有限型态射，而不只是局部有限型态射；讨论见
注记 \ref{flat-remark-discuss-finite-type}。在该定义的情形下，引理
\ref{flat-lemma-pure-along-X-s} 告诉我们，下列条件等价：
\begin{enumerate}
\item $\mathcal{F}$ 沿 $X_s$ 纯；
\item 不存在非纯性 $(g : T \to S, t' \leadsto t, \xi)$，且 $g$
在 $t$ 处拟有限；
\item 不存在形如 $(S^h \to S, s' \leadsto s, \xi)$ 的非纯性，其中
$S^h$ 是 $S$ 在 $s$ 处的 Hensel 化。
\end{enumerate}
记 $X^h = X \times_S S^h$，并令 $\mathcal{F}^h$ 为 $\mathcal{F}$
到 $X^h$ 的拉回，则最后一个条件可以用下面更积极的方式表述：
\begin{enumerate}
\item[(4)] $\text{Ass}_{X^h/S^h}(\mathcal{F}^h)$ 的所有点都专化到
$X_s$ 中的点。
\end{enumerate}
特别地，显然 $\mathcal{F}$ 沿 $X_s$ 纯，当且仅当 $\mathcal{F}$ 拉回到
$X \times_S \Spec(\mathcal{O}_{S, s})$ 后沿 $X_s$ 纯。

\begin{remark}
\label{flat-remark-discuss-finite-type}
设 $f : X \to S$ 是局部有限型态射，且 $\mathcal{F}$ 是拟凝聚有限型
$\mathcal{O}_X$-模。在此情形，(1) 与 (2) 仍然等价，因为引理
\ref{flat-lemma-quasi-finite-impurity-elementary} 的证明并未使用 $f$ 是拟紧的
事实。(3) 与 (4) 等价也很清楚。不过，我们不知道 (1) 与 (3) 是否等价。
此时有时使用等价条件 (3) 和 (4) 来定义纯性会更方便，正如
\cite{GruRay} 中所做的那样。另一方面，对于许多应用，正确的概念似乎
确实应当是普遍纯。
\end{remark}

\noindent
一个自然问题是，相对于底空间纯这一性质是否在基变换下保持；也就是说，
纯是否等同于普遍纯。事实证明，在 Noether 环概形上（见引理
\ref{flat-lemma-Noetherian-base-change}），或者层是平坦的情况下（见引理
\ref{flat-lemma-finite-type-flat-pure-along-fibre-is-universal} 与
\ref{flat-lemma-finite-type-flat-pure-is-universal}），答案是肯定的。
一般而言这并不成立，即使态射和层都是有限表示的也不成立；反例见
例子一节 \ref{examples-section-pure-not-universally}。首先，我们将
这里对“普遍”的用法与通常的概念对应起来。

\begin{lemma}
\label{flat-lemma-base-change-universally}
设 $f : X \to S$ 是有限型概形态射。设 $\mathcal{F}$ 是有限型拟凝聚
$\mathcal{O}_X$-模。设 $s \in S$。下列条件等价：
\begin{enumerate}
\item $\mathcal{F}$ 沿 $X_s$ 普遍纯；
\item 对每个带点概形态射 $(S', s') \to (S, s)$，拉回
$\mathcal{F}_{S'}$ 沿 $X_{s'}$ 纯。
\end{enumerate}
特别地，$\mathcal{F}$ 相对于 $S$ 普遍纯，当且仅当 $\mathcal{F}$ 的每个
基变换 $\mathcal{F}_{S'}$ 都相对于 $S'$ 纯。
\end{lemma}

\begin{proof}
这是形式性的。
\end{proof}

\begin{lemma}
\label{flat-lemma-quasi-finite-base-change}
设 $f : X \to S$ 是有限型概形态射。设 $\mathcal{F}$ 是有限型拟凝聚
$\mathcal{O}_X$-模。设 $s \in S$。设 $(S', s') \to (S, s)$ 是带点概形态射。
若 $S' \to S$ 在 $s'$ 处拟有限，且 $\mathcal{F}$ 沿 $X_s$ 纯，则
$\mathcal{F}_{S'}$ 沿 $X_{s'}$ 纯。
\end{lemma}

\begin{proof}
% P05-EN-CH38-0044: Frozen English line 4917 reads "It" where the conditional sentence requires "If"; the canonical target renders the intended conditional and records the source defect.
若 $(T \to S', t' \leadsto t, \xi)$ 是 $\mathcal{F}_{S'}$ 在 $s'$ 之上的
一个非纯性，且 $T \to S'$ 在 $t$ 处拟有限，则
$(T \to S, t' \to t, \xi)$ 是 $\mathcal{F}$ 在 $s$ 之上的非纯性，且
$T \to S$ 在 $t$ 处拟有限；见《态射》，引理
\ref{morphisms-lemma-composition-quasi-finite}。因此，立即由
\ref{flat-definition-pure} 之后给出的纯性刻画 (2) 得到本引理。
\end{proof}

\begin{lemma}
\label{flat-lemma-Noetherian-base-change}
设 $f : X \to S$ 是有限型概形态射。设 $\mathcal{F}$ 是有限型拟凝聚
$\mathcal{O}_X$-模。设 $s \in S$。若 $\mathcal{O}_{S, s}$ 是 Noether 环，
则 $\mathcal{F}$ 沿 $X_s$ 纯，当且仅当 $\mathcal{F}$ 沿 $X_s$ 普遍纯。
\end{lemma}

\begin{proof}
首先可以将 $S$ 替换为 $\Spec(\mathcal{O}_{S, s})$，即可以假设 $S$ 是
Noether 环。然后使用引理 \ref{flat-lemma-Noetherian-impurity-quasi-finite} 和
\ref{flat-definition-pure} 之后讨论中给出的纯性刻画 (2)，即可得出结论。
\end{proof}

\noindent
纯性满足平坦下降。

\begin{lemma}
\label{flat-lemma-flat-descend-pure}
设 $f : X \to S$ 是有限型概形态射。设 $\mathcal{F}$ 是有限型拟凝聚
$\mathcal{O}_X$-模。设 $s \in S$。设 $(S', s') \to (S, s)$ 是带点概形态射。
假设 $S' \to S$ 在 $s'$ 处平坦。
\begin{enumerate}
\item 若 $\mathcal{F}_{S'}$ 沿 $X_{s'}$ 纯，则 $\mathcal{F}$ 沿 $X_s$ 纯。
\item 若 $\mathcal{F}_{S'}$ 沿 $X_{s'}$ 普遍纯，则 $\mathcal{F}$ 沿 $X_s$
普遍纯。
\end{enumerate}
\end{lemma}

\begin{proof}
设 $(T \to S, t' \leadsto t, \xi)$ 是 $\mathcal{F}$ 在 $s$ 之上的非纯性。
置 $T_1 = T \times_S S'$，并令 $t_1$ 为 $T_1$ 中映到 $t$ 和 $s'$ 的唯一点。
由于 $T_1 \to T$ 在 $t_1$ 处平坦，见《态射》，引理
\ref{morphisms-lemma-base-change-flat}，存在位于 $t' \leadsto t$ 之上的
专化 $t'_1 \leadsto t_1$，见《代数》，一节
\ref{algebra-section-going-up}。选取 $\xi_1 \in X_{t'_1}$，使其对应于
$\Spec(\kappa(t'_1) \otimes_{\kappa(t')} \kappa(\xi))$ 的一个泛点，见
《概形》，引理 \ref{schemes-lemma-points-fibre-product}。由《除子》，引理
\ref{divisors-lemma-base-change-relative-assassin}，有
$\xi_1 \in \text{Ass}_{X_{T_1}/T_1}(\mathcal{F}_{T_1})$。
由于 $\{\xi_1\}$ 在 $X_{T_1}$ 中的 Zariski 闭包映入 $\{\xi\}$ 在 $X_T$
中的 Zariski 闭包，我们得到该闭包与 $X_{t_1}$ 不相交。因此
$(T_1 \to S', t'_1 \leadsto t_1, \xi_1)$ 是 $\mathcal{F}_{S'}$ 在 $s'$
之上的非纯性。换言之，我们证明了本引理第 (2) 部分的逆否命题。最后，若
$(T, t) \to (S, s)$ 是初等 \'etale 邻域，则 $(T_1, t_1) \to (S', s')$
也是初等 \'etale 邻域，由此可见 (1) 成立。
\end{proof}

\begin{lemma}
\label{flat-lemma-supported-on-closed}
设 $i : Z \to X$ 是在概形 $S$ 上有限型概形的闭浸入。设 $s \in S$。
设 $\mathcal{F}$ 是 $Z$ 上的有限型拟凝聚层。则 $\mathcal{F}$ 沿 $Z_s$
（普遍）纯，当且仅当 $i_*\mathcal{F}$ 沿 $X_s$（普遍）纯。
\end{lemma}

\begin{proof}
这由《除子》，引理
\ref{divisors-lemma-relative-weak-assassin-finite} 得出。
\end{proof}








\section{相对纯层的例子}
\label{flat-section-examples-pure-sheaves}

\noindent
下面给出一些可以看出纯性含义的例子情形。

\begin{lemma}
\label{flat-lemma-proper-pure}
设 $f : X \to S$ 是有限型概形态射。
设 $\mathcal{F}$ 是有限型拟凝聚 $\mathcal{O}_X$-模。
\begin{enumerate}
\item 若 $\mathcal{F}$ 的支集在 $S$ 上适当，则 $\mathcal{F}$ 相对于 $S$
普遍纯。
\item 若 $f$ 适当，则 $\mathcal{F}$ 相对于 $S$ 普遍纯。
\item 若 $f$ 适当，则 $X$ 相对于 $S$ 普遍纯。
\end{enumerate}
\end{lemma}

\begin{proof}
首先将 (1) 归约为 (2)。具体地，令 $Z \subset X$ 为 $\mathcal{F}$ 的
概形论支集。令 $i : Z \to X$ 为相应的闭浸入，并写成
$\mathcal{F} = i_*\mathcal{G}$，其中 $\mathcal{G}$ 是有限型拟凝聚
$\mathcal{O}_Z$-模，见《态射》，一节 \ref{morphisms-section-support}。
在情形 (1) 中，按假设 $Z \to S$ 适当。因此由引理
\ref{flat-lemma-supported-on-closed}，情形 (1) 归约为情形 (2)。

\medskip\noindent
假设 $f$ 适当。设 $(g : T \to S, t' \leadsto t, \xi)$ 是 $\mathcal{F}$
在 $s \in S$ 之上的非纯性。由于 $f$ 适当，它是普遍闭的。因此
$f_T : X_T \to T$ 是闭态射。由于 $f_T(\xi) = t'$，这蕴含
$t \in f(\overline{\{\xi\}})$，矛盾。
\end{proof}

\begin{lemma}
\label{flat-lemma-quasi-finite-pure}
设 $f : X \to S$ 是概形的分离有限型态射。设 $\mathcal{F}$ 是有限型
拟凝聚 $\mathcal{O}_X$-模。假设对每个 $s \in S$，$\text{Supp}(\mathcal{F}_s)$
都是有限的。则下列条件等价：
\begin{enumerate}
\item $\mathcal{F}$ 相对于 $S$ 纯；
\item $\mathcal{F}$ 的概形论支集在 $S$ 上有限；
\item $\mathcal{F}$ 相对于 $S$ 普遍纯。
\end{enumerate}
特别地，给定拟有限分离态射 $X \to S$，可知 $X$ 相对于 $S$ 纯，当且仅当
$X \to S$ 有限。
\end{lemma}

\begin{proof}
令 $Z \subset X$ 为 $\mathcal{F}$ 的概形论支集，见《态射》，定义
\ref{morphisms-definition-scheme-theoretic-support}。则 $Z \to S$ 是带有
有限纤维的分离有限型概形态射。因此它是分离且拟有限的，见《态射》，引理
\ref{morphisms-lemma-quasi-finite}。由引理 \ref{flat-lemma-supported-on-closed}，
只需对 $Z \to S$ 以及视为 $Z$ 上有限型拟凝聚模的层 $\mathcal{F}$ 证明本引理。
于是可以假设 $X \to S$ 分离且拟有限，并且 $\text{Supp}(\mathcal{F}) = X$。

\medskip\noindent
由引理 \ref{flat-lemma-proper-pure} 以及《态射》，引理
\ref{morphisms-lemma-finite-proper}，(2) 蕴含 (3)。显然 (3) 蕴含 (1)。
假设 (1) 成立。我们证明 (2) 成立。显然可以假设 $S$ 仿射。由《态射进阶》，
引理 \ref{more-morphisms-lemma-quasi-finite-separated-pass-through-finite}，
可以找到图表
$$
\xymatrix{
X \ar[rd]_f \ar[rr]_j & & T \ar[ld]^\pi \\
& S &
}
$$
其中 $\pi$ 有限，$j$ 是拟紧开浸入。若证明 $j$ 是闭的，则 $j$ 是闭浸入，
从而 $f = \pi \circ j$ 有限。要证明 $j$ 闭，只需证明专化沿 $j$ 可提升，见
《概形》，引理 \ref{schemes-lemma-quasi-compact-closed}。令 $x \in X$，置
$t' = j(x)$，并设 $t' \leadsto t$ 是一个专化。我们要证明 $t \in j(X)$。
置 $s' = f(x)$、$s = \pi(t)$，于是 $s' \leadsto s$。由《态射进阶》，引理
\ref{more-morphisms-lemma-etale-splits-off-quasi-finite-part-technical}，
可以找到初等 \'etale 邻域 $(U, u) \to (S, s)$ 以及分解
$$
T_U = T \times_S U = V \amalg W
$$
为开闭子概形，使得 $V \to U$ 有限，并且 $V$ 中存在唯一映到 $u$ 的点 $v$，
且 $v$ 在 $T$ 中映到 $t$。由于 $V \to T$ 是 \'etale，可以提升一般化，见
《态射》，引理 \ref{morphisms-lemma-generalizations-lift-flat} 与
\ref{morphisms-lemma-etale-flat}。因此存在专化 $v' \leadsto v$，使得 $v'$
映到 $t' \in T$。特别地，$v' \in X_U \subset T_U$。记 $u' \in U$ 为 $t'$ 的像。
注意，由于 $X_{u'}$ 是有限离散集且
$X_{u'} = \text{Supp}(\mathcal{F}_{u'})$，有
$v' \in \text{Ass}_{X_U/U}(\mathcal{F})$。由于 $\mathcal{F}$ 相对于 $S$ 纯，
可知 $v'$ 必须专化到 $X_u$ 中的一个点。由于 $v$ 是 $V$ 中位于 $u$ 之上的唯一点
（且 $W$ 中没有点能成为 $v'$ 的专化），可知 $v \in X_u$。故 $t \in X$。
\end{proof}

\begin{lemma}
\label{flat-lemma-flat-geometrically-integral-fibres-pure}
设 $f : X \to S$ 是有限型平坦概形态射，且纤维几何整。则 $X$ 在 $S$ 上
普遍纯。
\end{lemma}

\begin{proof}
令 $\xi \in X$，其中 $s' = f(\xi)$，并设 $s' \leadsto s$ 是 $S$ 中的专化。
若 $\xi$ 是 $X_{s'}$ 的伴随点，则 $\xi$ 是唯一的泛点，因为 $X_{s'}$ 是整概形。
令 $\xi_0$ 为 $X_s$ 的唯一泛点。由于 $X \to S$ 平坦，可以在 $X$ 中将
$s' \leadsto s$ 提升为专化 $\xi' \leadsto \xi_0$，见《态射》，引理
\ref{morphisms-lemma-generalizations-lift-flat}。由于 $\xi$ 是 $X_{s'}$ 的泛点，
% P05-EN-CH38-0045: Frozen English line 5138 begins "The $\xi \leadsto \xi'$" and omits the copula; the canonical target preserves the mathematical relation and records the grammatical omission.
有 $\xi \leadsto \xi'$，从而 $\xi \leadsto \xi_0$。这意味着
$(\text{id}_S, s' \to s, \xi)$ 不是 $\mathcal{O}_X$ 在 $s$ 之上的非纯性。
由于“$f$ 有限型、平坦且纤维几何整”这一假设在基变换下保持，可知任意基变换后
都不存在非纯性。因此 $X$ 是普遍 $S$-纯。
\end{proof}

\begin{lemma}
\label{flat-lemma-affine-locally-projective-pure}
设 $f : X \to S$ 是有限型仿射概形态射。设 $\mathcal{F}$ 是有限型拟凝聚
$\mathcal{O}_X$-模，且 $f_*\mathcal{F}$ 在 $S$ 上局部投射，见《性质》，定义
\ref{properties-definition-locally-projective}。则 $\mathcal{F}$ 在 $S$ 上普遍纯。
\end{lemma}

\begin{proof}
将情形归约到 $S$ 是 Hensel 局部环的谱后，这由引理
\ref{flat-lemma-explain-why-pure} 得出。
\end{proof}




\section{纯性的一个判据}
\label{flat-section-criterion-purity}

\noindent
我们首先证明，对于有限型拟凝聚层的一个平坦族，相对伴随点会专化到邻近
纤维的相对伴随点（如果它们确实发生专化）。

\begin{lemma}
\label{flat-lemma-associated-point-specializes}
设 $f : X \to S$ 是有限型概形态射。设 $\mathcal{F}$ 是有限型拟凝聚
$\mathcal{O}_X$-模。设 $s \in S$。假设 $\mathcal{F}$ 在 $X_s$ 的所有点处
都在 $S$ 上平坦。设 $x' \in \text{Ass}_{X/S}(\mathcal{F})$，满足
$f(x') = s'$，且 $s' \leadsto s$ 是 $S$ 中的专化。若 $x'$ 专化到 $X_s$
中的一个点，则 $x' \leadsto x$，其中
$x \in \text{Ass}_{X_s}(\mathcal{F}_s)$。
\end{lemma}

\begin{proof}
设 $x' \leadsto t$，其中 $t \in X_s$。则可以找到专化
$x' \leadsto x \leadsto t$，使得 $x$ 对应于
$\overline{\{x'\}} \cap f^{-1}(\{s\})$ 的某个不可约分支的泛点。按假设，
$\mathcal{F}$ 在 $x$ 处在 $S$ 上平坦。由《态射进阶》，引理
\ref{more-morphisms-lemma-associated-point-specializes}，可知
$x \in \text{Ass}_{X/S}(\mathcal{F})$，即得所需结论。
\end{proof}

\begin{lemma}
\label{flat-lemma-criterion}
设 $f : X \to S$ 是有限型概形态射。设 $\mathcal{F}$ 是有限型拟凝聚
$\mathcal{O}_X$-模。设 $s \in S$。设 $(S', s') \to (S, s)$ 是一个
初等 \'etale 邻域，并设
$$
\xymatrix{
X \ar[d] & X' \ar[l]^g \ar[d] \\
S & S' \ar[l]
}
$$
是概形态射的交换图表。假设
\begin{enumerate}
\item $\mathcal{F}$ 在 $X_s$ 的所有点处都在 $S$ 上平坦；
\item $X' \to S'$ 是有限型的；
\item $g^*\mathcal{F}$ 沿 $X'_{s'}$ 纯；
\item $g : X' \to X$ 是 \'etale；
\item $g(X')$ 包含 $\text{Ass}_{X_s}(\mathcal{F}_s)$。
\end{enumerate}
在此情形，$\mathcal{F}$ 沿 $X_s$ 纯，当且仅当 $X' \to X \times_S S'$ 的像
包含 $\text{Ass}_{X \times_S S'/S'}(\mathcal{F} \times_S S')$ 中位于
$S'$ 内专化到 $s'$ 的点之上的那些点。
\end{lemma}

\begin{proof}
由于态射 $S' \to S$ 是 \'etale，若 $\mathcal{F}$ 沿 $X_s$ 纯，则
$\mathcal{F} \times_S S'$ 沿 $X_s$ 纯，见引理
\ref{flat-lemma-quasi-finite-base-change}。由于纯性满足平坦下降，见引理
\ref{flat-lemma-flat-descend-pure}，若 $\mathcal{F} \times_S S'$ 沿 $X_{s'}$ 纯，
则 $\mathcal{F}$ 沿 $X_s$ 纯。因此可以用 $S'$ 替换 $S$，并假设 $S = S'$，
于是 $g : X' \to X$ 是 $S$ 上有限型概形之间的 \'etale 态射。此外，可以用
$\Spec(\mathcal{O}_{S, s})$ 替换 $S$，并假设 $S$ 是局部的。

\medskip\noindent
首先假设 $\mathcal{F}$ 沿 $X_s$ 纯。在此情形，由纯性
$\text{Ass}_{X/S}(\mathcal{F})$ 的每个点都专化到 $X_s$ 中的某个点。因此由
引理 \ref{flat-lemma-associated-point-specializes}，可知
$\text{Ass}_{X/S}(\mathcal{F})$ 的每个点都专化到
$\text{Ass}_{X_s}(\mathcal{F}_s)$ 中的某个点。于是
$\text{Ass}_{X/S}(\mathcal{F})$ 的每个点都在 $g$ 的像中（因为该像是开集，
并包含 $\text{Ass}_{X_s}(\mathcal{F}_s)$）。

\medskip\noindent
反之，假设 $g(X')$ 包含 $\text{Ass}_{X/S}(\mathcal{F})$。令
$S^h = \Spec(\mathcal{O}_{S, s}^h)$ 为 $S$ 在 $s$ 处的 Hensel 化。记
$g^h : (X')^h \to X^h$ 为沿 $S^h \to S$ 对 $g$ 作的基变换，并记
$\mathcal{F}^h$ 为 $\mathcal{F}$ 到 $X^h$ 的拉回。由《除子》，引理
\ref{divisors-lemma-base-change-relative-assassin} 和注记
\ref{divisors-remark-base-change-relative-assassin}，相对伴随点集
$\text{Ass}_{X^h/S^h}(\mathcal{F}^h)$ 是经投影 $X^h \to X$ 的
$\text{Ass}_{X/S}(\mathcal{F})$ 的逆像。由于假设 $g(X')$ 包含
$\text{Ass}_{X/S}(\mathcal{F})$，可知基变换
$g^h((X')^h) = g(X') \times_S S^h$ 包含
$\text{Ass}_{X^h/S^h}(\mathcal{F}^h)$。因此可以归约到 $S$ 是 Hensel 局部环的谱的情形。
令 $x \in \text{Ass}_{X/S}(\mathcal{F})$。要完成引理的证明，需要证明 $x$ 专化到
$X_s$ 中的一个点；这正是 \ref{flat-definition-pure} 之后关于纯性的讨论中的
判据 (4)。按假设，存在 $x' \in X'$ 使得 $g(x') = x$。由于
$g : X' \to X$ 是 \'etale，可知 $x' \in \text{Ass}_{X'/S}(g^*\mathcal{F})$，见
引理 \ref{flat-lemma-etale-weak-assassin-up-down}（应用于纤维态射
$X'_w \to X_w$，其中 $w \in S$ 是 $x'$ 的像）。由于 $g^*\mathcal{F}$ 沿
$X'_s$ 纯，可知对某个 $y \in X'_s$ 有 $x' \leadsto y$。于是
$x = g(x') \leadsto g(y)$，且 $g(y) \in X_s$，即得所需结论。
\end{proof}

\begin{lemma}
\label{flat-lemma-finite-type-flat-pure-along-fibre-is-universal}
设 $f : X \to S$ 是概形态射。设 $\mathcal{F}$ 是拟凝聚
$\mathcal{O}_X$-模。设 $s \in S$。假设
\begin{enumerate}
\item $f$ 是有限型的；
\item $\mathcal{F}$ 是有限型的；
\item $\mathcal{F}$ 在 $X_s$ 的所有点处都在 $S$ 上平坦；
\item $\mathcal{F}$ 沿 $X_s$ 纯。
\end{enumerate}
则 $\mathcal{F}$ 沿 $X_s$ 普遍纯。
\end{lemma}

\begin{proof}
先作一个预备说明。设 $(S', s') \to (S, s)$ 是初等 \'etale 邻域。记
$\mathcal{F}'$ 为 $\mathcal{F}$ 到 $X' = X \times_S S'$ 的拉回。
由 \ref{flat-definition-pure} 之后的讨论可知，$\mathcal{F}'$ 沿 $X'_{s'}$ 纯。
此外，$\mathcal{F}'$ 沿 $X'_{s'}$ 在 $S'$ 上平坦。因此只需证明
$\mathcal{F}'$ 沿 $X'_{s'}$ 普遍纯。事实上，给定带点概形的任意态射
$(T, t) \to (S, s)$，纤维积 $(T', t') = (T \times_S S', (t, s'))$ 在
$(T, t)$ 上平坦；因此若 $\mathcal{F}_{T'}$ 沿 $X_{t'}$ 纯，则由引理
\ref{flat-lemma-flat-descend-pure}，$\mathcal{F}_T$ 沿 $X_t$ 纯。
所以在证明过程中总可以用初等 \'etale 邻域替换 $(s, S)$。
由于问题具有局部性，也可以用 $\Spec(\mathcal{O}_{S, s})$ 替换 $S$。

\medskip\noindent
选取初等 \'etale 邻域 $(S', s') \to (S, s)$ 以及交换图表
$$
\xymatrix{
X \ar[d] & X' \ar[l]^g \ar[d] \\
S & \Spec(\mathcal{O}_{S', s'}) \ar[l]
}
$$
使得 $X' \to X \times_S \Spec(\mathcal{O}_{S', s'})$ 是 \'etale，
$X_s = g((X')_{s'})$，概形 $X'$ 仿射，并且
$\Gamma(X', g^*\mathcal{F})$ 是自由的 $\mathcal{O}_{S', s'}$-模；见引理
\ref{flat-lemma-finite-type-flat-along-fibre-free-variant}。注意，
$X' \to \Spec(\mathcal{O}_{S', s'})$ 是有限型的（因为它是一个拟紧态射，
由一个 \'etale 态射与一个有限型态射的基变换复合而成）。根据证明第一段的
预备说明，可以用 $\Spec(\mathcal{O}_{S', s'})$ 替换 $S$。因此可以假设存在
交换图表
$$
\xymatrix{
X \ar[dr] & & X' \ar[ll]^g \ar[ld] \\
& S &
}
$$
其中的概形都是 $S$ 上有限型的，$g$ 是 \'etale，$X_s \subset g(X')$，
$S$ 是闭点为 $s$ 的局部概形，$X'$ 仿射，并且
$\Gamma(X', g^*\mathcal{F})$ 是自由的 $\Gamma(S, \mathcal{O}_S)$-模。
注意此时 $g^*\mathcal{F}$ 在 $S$ 上普遍纯，见引理
\ref{flat-lemma-affine-locally-projective-pure}。

\medskip\noindent
在此情形，应用引理 \ref{flat-lemma-criterion}，由 $\mathcal{F}$ 沿 $X_s$ 纯且
沿 $X_s$ 在 $S$ 上平坦的假设，可推出
$\text{Ass}_{X/S}(\mathcal{F}) \subset g(X')$。由《除子》，引理
\ref{divisors-lemma-base-change-relative-assassin} 及注记
\ref{divisors-remark-base-change-relative-assassin}，对任意带点概形态射
$(T, t) \to (S, s)$，有
$$
\text{Ass}_{X_T/T}(\mathcal{F}_T) \subset
(X_T \to X)^{-1}(\text{Ass}_{X/S}(\mathcal{F})) \subset
g(X') \times_S T = g_T(X'_T).
$$
因此，将所显示图表基变换到 $(T, t)$ 后应用引理 \ref{flat-lemma-criterion}，
可得 $\mathcal{F}_T$ 沿 $X_t$ 纯，这就是所需结论。
\end{proof}

\begin{lemma}
\label{flat-lemma-finite-type-flat-pure-is-universal}
设 $f : X \to S$ 是有限型概形态射。设 $\mathcal{F}$ 是有限型拟凝聚
$\mathcal{O}_X$-模。假设 $\mathcal{F}$ 在 $S$ 上平坦。在此情形，
$\mathcal{F}$ 相对于 $S$ 纯，当且仅当 $\mathcal{F}$ 相对于 $S$ 普遍纯。
\end{lemma}

\begin{proof}
这是引理 \ref{flat-lemma-finite-type-flat-pure-along-fibre-is-universal} 与定义的
直接推论。
\end{proof}

\begin{lemma}
\label{flat-lemma-limit-purity}
设 $I$ 是有向集。设 $(S_i, g_{ii'})$ 是 $I$ 上的仿射概形逆系统。
置 $S = \lim_i S_i$，并令 $s \in S$。记 $g_i : S \to S_i$ 为投影，并置
$s_i = g_i(s)$。假设 $f : X \to S$ 是有限表示态射，$\mathcal{F}$ 是有限表示的
拟凝聚 $\mathcal{O}_X$-模，沿 $X_s$ 纯，并且在 $X_s$ 的所有点处都在 $S$ 上平坦。
则存在 $i \in I$、有限表示态射 $X_i \to S_i$ 以及有限表示的拟凝聚
$\mathcal{O}_{X_i}$-模 $\mathcal{F}_i$，使得沿 $(X_i)_{s_i}$ 纯，
并在 $(X_i)_{s_i}$ 的所有点处在 $S_i$ 上平坦，且
$X \cong X_i \times_{S_i} S$，并且 $\mathcal{F}_i$ 到 $X$ 的拉回同构于
$\mathcal{F}$。
\end{lemma}

\begin{proof}
令 $U \subset X$ 为 $\mathcal{F}$ 在 $S$ 上平坦的点集。由《态射进阶》，定理
\ref{more-morphisms-theorem-openness-flatness}，这是 $X$ 的开子概形。按假设
$X_s \subset U$。由于 $X_s$ 拟紧，可以找到一个拟紧开集 $U' \subset U$，满足
$X_s \subset U'$。由《极限》，引理 \ref{limits-lemma-descend-finite-presentation}，
可以找到 $i \in I$ 以及有限表示态射 $f_i : X_i \to S_i$，其到 $S$ 的基变换
同构于 $f_i$。固定这样的选择，并置 $X_{i'} = X_i \times_{S_i} S_{i'}$。
于是 $X = \lim_{i'} X_{i'}$，过渡态射是仿射的。由《极限》，引理
\ref{limits-lemma-descend-modules-finite-presentation}，增大 $i$ 后可以假设存在
有限表示的拟凝聚 $\mathcal{O}_{X_i}$-模 $\mathcal{F}_i$，其到 $S$ 的基变换同构于
$\mathcal{F}$。由《极限》，引理 \ref{limits-lemma-descend-opens}，再增大 $i$ 后可以
假设存在开集 $U'_i \subset X_i$，其在 $X$ 中的逆像为 $U'$。特别注意
$(X_i)_{s_i} \subset U'_i$。由《极限》，引理
\ref{limits-lemma-descend-module-flat-finite-presentation}（再一次增大 $i$ 后），
可以假设 $\mathcal{F}_i$ 在 $U'_i$ 上平坦。特别地，$\mathcal{F}_i$ 沿
$(X_i)_{s_i}$ 平坦。

\medskip\noindent
接下来，使用引理 \ref{flat-lemma-finite-presentation-flat-along-fibre}，选取初等
\'etale 邻域 $(S_i', s_i') \to (S_i, s_i)$ 以及概形的交换图表
$$
\xymatrix{
X_i \ar[d] & X_i' \ar[l]^{g_i} \ar[d] \\
S_i & S_i' \ar[l]
}
$$
使得 $g_i$ 是 \'etale，$(X_i)_{s_i} \subset g_i(X_i')$，概形 $X_i'$、$S_i'$ 都
仿射，并且 $\Gamma(X_i', g_i^*\mathcal{F}_i)$ 是
$\Gamma(S_i', \mathcal{O}_{S_i'})$-投射模。注意 $g_i^*\mathcal{F}_i$ 在 $S'_i$
上普遍纯，见引理 \ref{flat-lemma-affine-locally-projective-pure}。对任意 $i' \geq i$，
可以将上述图表基变换为 $S_{i'}$ 上带有态射
$(S'_{i'}, s'_{i'}) \to (S_{i'}, s_{i'})$ 和
$g_{i'} : X'_{i'} \to X_{i'}$ 的图表；也可以将图表基变换为 $S$ 上带有态射
$(S', s') \to (S, s)$ 和 $g : X' \to X$ 的图表。

\medskip\noindent
此时可以使用纯性判据。令 $W'_i \subset X_i \times_{S_i} S'_i$ 等于
\'etale 态射 $X'_i \to X_i \times_{S_i} S'_i$ 的像。对每个 $i' \geq i$，
类似地有像 $W'_{i'} \subset X_{i'} \times_{S_{i'}} S'_{i'}$，并有像
$W' \subset X \times_S S'$。取像与基变换交换，故
$W'_{i'} = W'_i \times_{S'_i} S'_{i'}$ 且 $W' = W_i \times_{S'_i} S'$。
由于 $\mathcal{F}$ 沿 $X_s$ 纯，引理 \ref{flat-lemma-criterion} 蕴含
\begin{equation}
\label{flat-equation-inclusion}
f^{-1}(\Spec(\mathcal{O}_{S', s'})) \cap
\text{Ass}_{X \times_S S'/S'}(\mathcal{F} \times_S S') \subset W'
\end{equation}
由《态射进阶》，引理
\ref{more-morphisms-lemma-relative-assassin-constructible}，可知
$$
E = \{t \in S' \mid \text{Ass}_{X_t}(\mathcal{F}_t) \subset W' \}
\quad\text{且}\quad
E_{i'} = \{t \in S'_{i'} \mid
\text{Ass}_{X_t}(\mathcal{F}_{i', t}) \subset W'_{i'} \}
$$
是 $S'$ 与 $S'_{i'}$ 的局部可构造子集。由《态射进阶》，引理
\ref{more-morphisms-lemma-base-change-assassin-in-U}，可知 $E_{i'}$ 是关于态射
$S'_{i'} \to S'_i$ 的 $E_i$ 逆像，而 $E$ 是关于态射 $S' \to S'_i$ 的 $E_i$
逆像。因此，等式（\ref{flat-equation-inclusion}）等价于
$\Spec(\mathcal{O}_{S', s'})$ 映入 $E_i$。由于
$\mathcal{O}_{S', s'} = \colim_{i' \geq i} \mathcal{O}_{S'_{i'}, s'_{i'}}$，
可知对某个 $i' \geq i$，$\Spec(\mathcal{O}_{S'_{i'}, s'_{i'}})$ 映入 $E_i$；见
《极限》，引理 \ref{limits-lemma-limit-contained-in-constructible}。然后将引理
\ref{flat-lemma-criterion} 应用于 $S_{i'}$ 上的情形，得到
$\mathcal{F}_{i'}$ 沿 $(X_{i'})_{s_{i'}}$ 纯。
\end{proof}

\begin{lemma}
\label{flat-lemma-flat-finite-presentation-purity-open}
设 $f : X \to S$ 是有限表示态射。设 $\mathcal{F}$ 是在 $S$ 上平坦的有限表示
拟凝聚 $\mathcal{O}_X$-模。则集合
$$
U = \{s \in S \mid \mathcal{F}\text{ 沿 }X_s\text{ 纯}\}
$$
在 $S$ 中是开的。
\end{lemma}

\begin{proof}
设 $s \in U$。使用引理 \ref{flat-lemma-finite-presentation-flat-along-fibre}，可以
找到初等 \'etale 邻域 $(S', s') \to (S, s)$ 以及交换图表
$$
\xymatrix{
X \ar[d] & X' \ar[l]^g \ar[d] \\
S & S' \ar[l]
}
$$
使得 $g$ 是 \'etale，$X_s \subset g(X')$，概形 $X'$、$S'$ 都仿射，且
$\Gamma(X', g^*\mathcal{F})$ 是 $\Gamma(S', \mathcal{O}_{S'})$-投射模。
注意 $g^*\mathcal{F}$ 在 $S'$ 上普遍纯，见引理
\ref{flat-lemma-affine-locally-projective-pure}。令 $W' \subset X \times_S S'$ 等于
\'etale 态射 $X' \to X \times_S S'$ 的像。% P05-EN-CH38-0046: Frozen English line 5517 names $W$ although the defined open is $W'$; the canonical target preserves the frozen symbol and records the source inconsistency.
注意 $W$ 在 $S'$ 上是开且拟紧的。
置
$$
E = \{t \in S' \mid \text{Ass}_{X_t}(\mathcal{F}_t) \subset W' \}.
$$
由《态射进阶》，引理
\ref{more-morphisms-lemma-relative-assassin-constructible}，$E$ 是 $S'$ 中的
可构造子集。由引理 \ref{flat-lemma-criterion}，可知
$\Spec(\mathcal{O}_{S', s'}) \subset E$。由《态射》，引理
\ref{morphisms-lemma-constructible-containing-open}，可知 $E$ 包含 $s'$ 的
一个开邻域 $V'$。再次应用引理 \ref{flat-lemma-criterion}，可知对 $V'$ 在 $S$ 中像内
的任意点 $s_1$，层 $\mathcal{F}$ 沿 $X_{s_1}$ 纯。由于 $S' \to S$ 是 \'etale，
$V'$ 在 $S$ 中的像是开的，结论得证。
\end{proof}


\section{纯性的应用}
\label{flat-section-applications-purity}

\noindent
下面给出一些使用纯性的例子。第一个引理实际上使用了稍弱形式的纯性。

\begin{lemma}
\label{flat-lemma-injectivity-map-source-flat-pure}
设 $f : X \to S$ 是有限型概形态射。设 $\mathcal{F}$ 是 $X$ 上有限型拟凝聚层。
假设 $S$ 是闭点为 $s$ 的局部概形。假设 $\mathcal{F}$ 沿 $X_s$ 纯，且
$\mathcal{F}$ 在 $S$ 上平坦。设 $\varphi : \mathcal{F} \to \mathcal{G}$ 是
拟凝聚 $\mathcal{O}_X$-模的态射。下列条件等价：
\begin{enumerate}
\item 对所有 $x \in \text{Ass}_{X_s}(\mathcal{F}_s)$，茎上的映射
$\varphi_x$ 都是单射；
\item $\varphi$ 是单射。
\end{enumerate}
\end{lemma}

\begin{proof}
令 $\mathcal{K} = \Ker(\varphi)$。我们的目标是证明 $\mathcal{K} = 0$。为此，
只需证明 $\text{WeakAss}_X(\mathcal{K}) = \emptyset$，见《除子》，引理
\ref{divisors-lemma-weakly-ass-zero}。有
$\text{WeakAss}_X(\mathcal{K}) \subset \text{WeakAss}_X(\mathcal{F})$，见
《除子》，引理 \ref{divisors-lemma-ses-weakly-ass}。由于 $\mathcal{F}$ 平坦，
由引理 \ref{flat-lemma-bourbaki-finite-type-general-base} 可知
$\text{WeakAss}_X(\mathcal{F}) \subset \text{Ass}_{X/S}(\mathcal{F})$。
由纯性，$\text{Ass}_{X/S}(\mathcal{F})$ 的任意点 $x'$ 都是 $X_s$ 中某点的
一般化，因此由引理 \ref{flat-lemma-associated-point-specializes}，它是某个
$x \in \text{Ass}_{X_s}(\mathcal{F}_s)$ 的专化。于是 $\varphi_x$ 的单射性蕴含
$\varphi_{x'}$ 的单射性，从而 $\mathcal{K}_{x'} = 0$。
\end{proof}

\begin{proposition}
\label{flat-proposition-finite-presentation-flat-pure-is-projective}
设 $f : X \to S$ 是概形的仿射有限表示态射。设 $\mathcal{F}$ 是有限表示、
在 $S$ 上平坦的拟凝聚 $\mathcal{O}_X$-模。下列条件等价：
\begin{enumerate}
\item $f_*\mathcal{F}$ 在 $S$ 上局部投射；
\item $\mathcal{F}$ 相对于 $S$ 纯。
\end{enumerate}
特别地，给定有限表示环映射 $A \to B$ 以及在 $A$ 上平坦的有限表示
$B$-模 $N$，有：$N$ 作为 $A$-模投射，当且仅当
$\widetilde{N}$ 在 $\Spec(B)$ 上相对于 $\Spec(A)$ 纯。
\end{proposition}

\begin{proof}
蕴含 (1) $\Rightarrow$ (2) 正是引理 \ref{flat-lemma-affine-locally-projective-pure}。
假设 $\mathcal{F}$ 相对于 $S$ 纯。注意，由引理
\ref{flat-lemma-finite-type-flat-pure-along-fibre-is-universal}，这蕴含 $\mathcal{F}$
在任意基变换后仍然纯。由《下降》，引理
\ref{descent-lemma-locally-projective-descends}，只需证明 $f_*\mathcal{F}$ 在
$S$ 上 fpqc 局部投射。取 $s \in S$。我们证明 $f_*\mathcal{F}$ 限制到 $s$ 的
某个 \'etale 邻域后局部投射。具体地，由引理
\ref{flat-lemma-finite-presentation-flat-along-fibre}，用 $s$ 的一个仿射初等 \'etale
邻域替换 $S$ 后，可以假设存在图表
$$
\xymatrix{
X \ar[dr] & & X' \ar[ll]^g \ar[ld] \\
& S &
}
$$
其中概形都是 $S$ 上仿射且有限表示的，$g$ 是 \'etale，$X_s \subset g(X')$，且
$\Gamma(X', g^*\mathcal{F})$ 是 $\Gamma(S, \mathcal{O}_S)$-投射模。
注意此时 $g^*\mathcal{F}$ 在 $S$ 上普遍纯，见引理
\ref{flat-lemma-affine-locally-projective-pure}。因此由引理 \ref{flat-lemma-criterion}，
可知开集 $g(X')$ 包含 $\text{Ass}_{X/S}(\mathcal{F})$ 中位于
$\Spec(\mathcal{O}_{S, s})$ 之上的点。置
$$
E = \{t \in S \mid \text{Ass}_{X_t}(\mathcal{F}_t) \subset g(X') \}.
$$
由《态射进阶》，引理
\ref{more-morphisms-lemma-relative-assassin-constructible}，$E$ 是 $S$ 中的
可构造子集。我们已经看到 $\Spec(\mathcal{O}_{S, s}) \subset E$。由《态射》，
引理 \ref{morphisms-lemma-constructible-containing-open}，可知 $E$ 包含 $s$ 的
一个开邻域。因此用 $s$ 的仿射邻域替换 $S$ 后，可以假设
$\text{Ass}_{X/S}(\mathcal{F}) \subset g(X')$。由引理
\ref{flat-lemma-base-change-universally-flat}，这意味着
$$
\Gamma(X, \mathcal{F}) \longrightarrow \Gamma(X', g^*\mathcal{F})
$$
是 $\Gamma(S, \mathcal{O}_S)$-普遍单射。由《代数》，引理
\ref{algebra-lemma-pure-submodule-ML}，可知 $\Gamma(X, \mathcal{F})$ 作为
$\Gamma(S, \mathcal{O}_S)$-模是 Mittag-Leffler 的。由于
$\Gamma(X, \mathcal{F})$ 是可数生成且作为 $\Gamma(S, \mathcal{O}_S)$-模平坦的，
由《代数》，引理 \ref{algebra-lemma-countgen-projective}，可知它是投射模。
\end{proof}

\noindent
我们可以利用该命题改进之前的一些结果。下面的引理改进了命题
\ref{flat-proposition-finite-presentation-flat-at-point}。

\begin{lemma}
\label{flat-lemma-flat-finite-presentation-affine-neighbourhood-projective}
设 $f : X \to S$ 是局部有限表示态射。设 $\mathcal{F}$ 是有限表示的
拟凝聚 $\mathcal{O}_X$-模。设 $x \in X$，且 $s = f(x) \in S$。
若 $\mathcal{F}$ 在 $x$ 处在 $S$ 上平坦，则存在仿射初等 \'etale 邻域
$(S', s') \to (S, s)$ 以及包含 $x' = (x, s')$ 的仿射开集
$U' \subset X \times_S S'$，使得 $\Gamma(U', \mathcal{F}|_{U'})$ 是
$\Gamma(S', \mathcal{O}_{S'})$-投射模。
\end{lemma}

\begin{proof}
在证明过程中，可以用 $x$ 的一个开邻域替换 $X$，并用 $s$ 的初等 \'etale
邻域替换 $S$。因此，由平坦性的开性（见《态射进阶》，定理
\ref{more-morphisms-theorem-openness-flatness}），可以假设 $\mathcal{F}$ 在
$S$ 上平坦。还可以假设 $S$ 和 $X$ 都仿射。再缩小 $X$ 后，可以假设
$\text{Ass}_{X_s}(\mathcal{F}_s)$ 的每个点都是 $x$ 的一般化。将 $(S, s)$
替换为初等 \'etale 邻域时这一性质保持。因此可以应用引理
\ref{flat-lemma-finite-presentation-flat-along-fibre}，归约到存在图表的情形
$$
\xymatrix{
X \ar[dr] & & X' \ar[ll]^g \ar[ld] \\
& S &
}
$$
其中概形都是 $S$ 上仿射且有限表示的，$g$ 是 \'etale，$X_s \subset g(X')$，且
$\Gamma(X', g^*\mathcal{F})$ 是 $\Gamma(S, \mathcal{O}_S)$-投射模。
注意此时 $g^*\mathcal{F}$ 在 $S$ 上普遍纯，见引理
\ref{flat-lemma-affine-locally-projective-pure}。

\medskip\noindent
令 $U \subset g(X')$ 为包含 $x$ 的仿射开邻域。我们断言 $\mathcal{F}|_U$ 沿
$U_s$ 纯。若证得此点，则由引理
\ref{flat-lemma-flat-finite-presentation-purity-open}，缩小 $S$ 后
$\mathcal{F}|_U$ 相对于 $S$ 纯，从而由命题
\ref{flat-proposition-finite-presentation-flat-pure-is-projective} 得到投射性。
为证明断言，需要证明：将 $(S, s)$ 替换为任意初等 \'etale 邻域后，
$\text{Ass}_{U/S}(\mathcal{F}|_U)$ 中位于某个 $s' \in S$ 之上的任意点
$\xi$（其中 $s' \leadsto s$）都专化到 $U_s$ 中的点。由于 $U \subset g(X')$，
可找到 $\xi' \in X'$ 使 $g(\xi') = \xi$。由于 $g^*\mathcal{F}$ 在 $S$ 上纯，
应用引理 \ref{flat-lemma-associated-point-specializes}，可知存在专化
$\xi' \leadsto x'$，其中 $x' \in \text{Ass}_{X'_s}(g^*\mathcal{F}_s)$。
于是 $g(x') \in \text{Ass}_{X_s}(\mathcal{F}_s)$（例如见引理
\ref{flat-lemma-etale-weak-assassin-up-down}，应用于 Noether 概形的 \'etale 态射
$X'_s \to X_s$），因而按上面对 $X$ 的选择有 $g(x') \leadsto x$！由于 $x \in U$，
可知 $g(x') \in U$。于是 $\xi = g(\xi') \leadsto g(x') \in U_s$，即得所需结论。
\end{proof}

\noindent
下面的引理改进了引理 \ref{flat-lemma-finite-type-flat-at-point-free-variant}。

\begin{lemma}
\label{flat-lemma-flat-finite-type-affine-neighbourhood-projective}
设 $f : X \to S$ 是局部有限型态射。设 $\mathcal{F}$ 是有限型拟凝聚
$\mathcal{O}_X$-模。设 $x \in X$，且 $s = f(x) \in S$。若 $\mathcal{F}$ 在 $x$
处在 $S$ 上平坦，则存在仿射初等 \'etale 邻域 $(S', s') \to (S, s)$ 以及
包含 $x' = (x, s')$ 的仿射开集
$U' \subset X \times_S \Spec(\mathcal{O}_{S', s'})$，使得
$\Gamma(U', \mathcal{F}|_{U'})$ 是自由的 $\mathcal{O}_{S', s'}$-模。
\end{lemma}

\begin{proof}
问题在 $X$ 和 $S$ 上都是 Zariski 局部的。因此可以假设 $X$ 和 $S$ 仿射。
于是可找到 $S$ 上的闭浸入 $i : X \to \mathbf{A}^n_S$。显然，只需对
$\mathbf{A}^n_S$ 上的层 $i_*\mathcal{F}$ 和点 $i(x)$ 证明引理即可。这样问题
归约到 $X \to S$ 有限表示的情形。用 $\Spec(\mathcal{O}_{S', s'})$ 替换 $S$，
并用 $X \times_S \Spec(\mathcal{O}_{S', s'})$ 的一个开子集替换 $X$ 后，可以假设
$\mathcal{F}$ 有限表示，见命题 \ref{flat-proposition-finite-type-flat-at-point}。
此时应用引理 \ref{flat-lemma-flat-finite-presentation-affine-neighbourhood-projective}
和《代数》，定理 \ref{algebra-theorem-projective-free-over-local-ring}，即可得出结论。
\end{proof}

\begin{lemma}
\label{flat-lemma-flat-finite-type-local-colimit-free}
设 $A \to B$ 是本质有限型局部环之间的局部环映射。设 $N$ 是有限 $B$-模，
作为 $A$-模平坦。若 $A$ 是 Hensel 环，则 $N$ 是自由 $A$-模 $F_i$ 的滤过余极限
$$
N = \colim_i F_i
$$
其中系统的所有过渡映射 $u_i : F_i \to F_{i'}$ 都诱导单射
$\overline{u}_i : F_i/\mathfrak m_AF_i \to F_{i'}/\mathfrak m_AF_{i'}$。
此外，$N$ 是 Mittag-Leffler $A$-模。
\end{lemma}

\begin{proof}
可以找到有限型态射 $X \to S = \Spec(A)$、位于 $S$ 闭点 $s$ 之上的点
$x \in X$，以及有限型拟凝聚 $\mathcal{O}_X$-模 $\mathcal{F}$，使得作为 $A$-模
$\mathcal{F}_x \cong N$。缩小 $X$ 后，可以假设
$\text{Ass}_{X_s}(\mathcal{F}_s)$ 的每个点都专化到 $x$。由引理
\ref{flat-lemma-flat-finite-type-affine-neighbourhood-projective}，存在 $x$ 的仿射开邻域
基本系 $U_i \subset X$，使得 $\Gamma(U_i, \mathcal{F})$ 是自由 $A$-模 $F_i$。
注意，若 $U_{i'} \subset U_i$，则
$$
F_i/\mathfrak m_AF_i = \Gamma(U_{i, s}, \mathcal{F}_s)
\longrightarrow
\Gamma(U_{i', s}, \mathcal{F}_s) = F_{i'}/\mathfrak m_AF_{i'}
$$
是单射，因为核中的一个截面会支撑在 $X_s$ 的某个不含 $x$ 的闭子集上，
这与上面对 $X$ 的选择矛盾。由于映射 $F_i \to F_{i'}$ 是 $A$-普遍单射
（引理 \ref{flat-lemma-universally-injective-local}），由《代数》，引理
\ref{algebra-lemma-colimit-universally-injective-ML}，可知 $N$ 是
Mittag-Leffler 模。
\end{proof}

\noindent
如果是第一次通读，可以跳过下面的引理。

\begin{lemma}
\label{flat-lemma-flat-finite-type-local-valuation-ring-has-content}
设 $A \to B$ 是本质有限型局部环之间的局部环映射。设 $N$ 是有限 $B$-模，
作为 $A$-模平坦。若 $A$ 是赋值环，则 $N$ 的任意元素都有内容理想
$I \subset A$（见《代数进阶》，定义
\ref{more-algebra-definition-content-ideal}）。此外，$I$ 是主理想。
\end{lemma}

\begin{proof}
最后一个断言由 $I$ 是有限生成理想这一事实得出；见《代数进阶》，引理
\ref{more-algebra-lemma-content-finitely-generated} 与《代数》，引理
\ref{algebra-lemma-characterize-valuation-ring}。

\medskip\noindent
$I$ 存在性的证明。令 $A \subset A^h$ 为 Hensel 化。令 $B'$ 为
$B \otimes_A A^h$ 在极大理想
$\mathfrak m_B \otimes A^h + B \otimes \mathfrak m_{A^h}$ 处的局部化。
则 $B \to B'$ 平坦，因而忠实平坦。令 $N' = N \otimes_B B'$。令
$x \in N$，并令 $x' \in N'$ 为其像。我们断言，对理想 $I \subset A$，有
$x \in IN \Leftrightarrow x' \in IN'$。事实上，$N/IN \to N'/IN'$ 是将
$B \to B'$ 与 $N/IN$ 张量所得的映射，而由《代数》，引理
\ref{algebra-lemma-faithfully-flat-universally-injective}，$B \to B'$ 普遍单射。
由《代数进阶》，引理 \ref{more-algebra-lemma-henselization-valuation-ring}，
以及《代数》，引理 \ref{algebra-lemma-ideals-valuation-ring}，映射 $A \to A^h$
在理想集上给出保持包含关系的双射 $I \mapsto IA^h$。因此，$x$ 在 $A$ 中有内容
理想，当且仅当 $x'$ 在 $A^h$ 中有内容理想。对于 $x' \in N'$ 的断言由引理
\ref{flat-lemma-flat-finite-type-local-colimit-free} 和《代数》，引理
\ref{algebra-lemma-minimal-contains} 得出。
\end{proof}

\noindent
应用如下。

\begin{lemma}
\label{flat-lemma-proper-flat-over-dvr-reduced-fibre}
设 $X \to \Spec(R)$ 是适当平坦态射，其中 $R$ 是赋值环。若特殊纤维约化，
则 $X$ 以及 $X \to \Spec(R)$ 的每个纤维都约化。
\end{lemma}

\begin{proof}
假设特殊纤维 $X_s$ 约化。令 $x \in X$ 为任意点，我们证明
$\mathcal{O}_{X, x}$ 约化；这将证明 $X$ 约化。令 $x \leadsto x'$ 是一个专化，
其中 $x'$ 位于特殊纤维中；这样的专化存在，因为适当态射是闭的。考虑局部环
$A = \mathcal{O}_{X, x'}$。则 $\mathcal{O}_{X, x}$ 是 $A$ 的局部化，所以只需证明
$A$ 约化。令 $a \in A$ 非零，并令 $I = (\pi) \subset R$ 为其内容理想，见引理
\ref{flat-lemma-flat-finite-type-local-valuation-ring-has-content}。于是 $a = \pi a'$，
且 $a'$ 映到 $A/\mathfrak mA$ 中的非零元素，其中 $\mathfrak m \subset R$ 是极大理想。
若 $a$ 幂零，则 $a'$ 也幂零，因为由 $A$ 在 $R$ 上平坦，$\pi$ 不是零因子。
但是 $a'$ 映到约化环 $A/\mathfrak m A = \mathcal{O}_{X_s, x'}$ 中的非零元素。
除非 $A$ 约化，否则这矛盾；而这正是要证明的结论。

\medskip\noindent
当然，若 $X$ 约化，则 $X$ 在 $R$ 上的泛纤维也约化。若 $\mathfrak p \subset R$
是素理想，则由《代数》，引理 \ref{algebra-lemma-make-valuation-rings}，
$R/\mathfrak p$ 是赋值环。因此将 $X$ 基变换到 $R/\mathfrak p$ 后重复上述论证，
即可证明位于 $\mathfrak p$ 之上的纤维约化。
\end{proof}






\section{扁平化函子}
\label{flat-section-flattening-functors}

\noindent
设 $S$ 是概形。回顾：若对每个以仿射概形组成、极限为 $T$ 的有向逆系统
$\{T_i\}_{i \in I}$，都有 $F(T) = \colim_i F(T_i)$，则称函子
$F : (\Sch/S)^{opp} \to \textit{Sets}$ 保持极限。

\begin{situation}
\label{flat-situation-iso}
设 $f : X \to S$ 是概形态射。设 $u : \mathcal{F} \to \mathcal{G}$ 是拟凝聚
$\mathcal{O}_X$-模的同态。对于任意在 $S$ 上的概形 $T$，记
$u_T : \mathcal{F}_T \to \mathcal{G}_T$ 为 $u$ 到 $T$ 的基变换；换言之，
$u_T$ 是沿投影态射 $X_T = X \times_S T \to X$ 对 $u$ 的拉回。
在此情形可以考虑函子
\begin{equation}
\label{flat-equation-iso}
F_{iso} : (\Sch/S)^{opp} \longrightarrow \textit{Sets}, \quad
T \longrightarrow \left\{
\begin{matrix}
\{*\} & \text{若} & u_T \text{ 是同构}, \\
\emptyset & \text{否则。} &
\end{matrix}
\right.
\end{equation}
还有变体 $F_{inj}$、$F_{surj}$、$F_{zero}$，分别要求 $u_T$ 是单射、满射或零映射。
\end{situation}

\begin{lemma}
\label{flat-lemma-iso-sheaf}
在情形 \ref{flat-situation-iso} 中。
\begin{enumerate}
\item 函子 $F_{iso}$、$F_{inj}$、$F_{surj}$、$F_{zero}$ 都满足 fpqc 拓扑的层性质。
\item 若 $f$ 拟紧且 $\mathcal{G}$ 有限型，则 $F_{surj}$ 保持极限。
\item 若 $f$ 拟紧且 $\mathcal{F}$ 有限型，则 $F_{zero}$ 保持极限。
\item 若 $f$ 拟紧、$\mathcal{F}$ 有限型且 $\mathcal{G}$ 有限表示，则
$F_{iso}$ 保持极限。
\end{enumerate}
\end{lemma}

\begin{proof}
设 $\{T_i \to T\}_{i \in I}$ 是 $S$ 上概形的 fpqc 覆盖。置
$X_i = X_{T_i} = X \times_S T_i$，并置 $u_i = u_{T_i}$。注意
$\{X_i \to X_T\}_{i \in I}$ 是 $X_T$ 的 fpqc 覆盖，见《拓扑》，引理
\ref{topologies-lemma-fpqc}。特别地，对每个 $x \in X_T$，存在 $i \in I$ 以及
映到 $x$ 的 $x_i \in X_i$。由于
$\mathcal{O}_{X_T, x} \to \mathcal{O}_{X_i, x_i}$ 平坦，因而忠实平坦（见《代数》，
引理 \ref{algebra-lemma-local-flat-ff}），可知 $(u_i)_{x_i}$ 是单射、满射、双射
或零映射，当且仅当 $(u_T)_x$ 具有相应性质。这就证明了引理的第 (1) 部分。

\medskip\noindent
(2) 的证明。假设 $f$ 拟紧且 $\mathcal{G}$ 有限型。令
$T = \lim_{i \in I} T_i$ 为仿射 $S$-概形的有向极限，并假设 $u_T$ 是满射。
置 $X_i = X_{T_i} = X \times_S T_i$，以及
$u_i = u_{T_i} : \mathcal{F}_i = \mathcal{F}_{T_i}
\to \mathcal{G}_i = \mathcal{G}_{T_i}$。为证明 (2)，需证明对某个 $i$，$u_i$ 是满射。
取 $i_0 \in I$，并用 $\{i \mid i \geq i_0\}$ 替换 $I$。由于 $f$ 拟紧，概形
$X_{i_0}$ 拟紧。因此可以选择仿射开集 $W_1, \ldots, W_m \subset X$ 以及仿射开覆盖
$X_{i_0} = U_{1, i_0} \cup \ldots \cup U_{m, i_0}$，使得 $U_{j, i_0}$
在投影态射 $X_{i_0} \to X$ 下映入 $W_j$。对任意 $i \in I$，令 $U_{j, i}$
为 $U_{j, i_0}$ 的逆像。置 $U_j = \lim_i U_{j, i}$，则
$X_T = U_1 \cup \ldots \cup U_m$ 是 $X_T$ 的仿射开覆盖。现在只需证明：对给定
$j \in \{1, \ldots, m\}$，存在 $i = i(j) \in I$ 使 $u_i|_{U_{j, i}}$
为满射。由《性质》，引理 \ref{properties-lemma-finite-type-module}，这转化为下列代数问题：
设 $A$ 是环，$u : M \to N$ 是 $A$-模映射。假设
$R = \colim_{i \in I} R_i$ 是 $A$-代数的有向余极限。若 $N$ 是有限 $A$-模，且
$u \otimes 1 : M \otimes_A R \to N \otimes_A R$ 是满射，则对某个 $i$，映射
$u \otimes 1 : M \otimes_A R_i \to N \otimes_A R_i$ 是满射。这就是《代数》，
引理 \ref{algebra-lemma-module-map-property-in-colimit} 的第 (2) 部分。

\medskip\noindent
(3) 的证明。与 (2) 的证明完全相同的论证将其归约为下列代数问题：
设 $A$ 是环，$u : M \to N$ 是 $A$-模映射。假设
$R = \colim_{i \in I} R_i$ 是 $A$-代数的有向余极限。若 $M$ 是有限 $A$-模，且
$u \otimes 1 : M \otimes_A R \to N \otimes_A R$ 为零，则对某个 $i$，映射
$u \otimes 1 : M \otimes_A R_i \to N \otimes_A R_i$ 为零。这就是《代数》，
引理 \ref{algebra-lemma-module-map-property-in-colimit} 的第 (1) 部分。

\medskip\noindent
(4) 的证明。假设 $f$ 拟紧且 $\mathcal{F}, \mathcal{G}$ 有限表示。以与上一段
完全相同的方式论证（另外使用《性质》，引理
\ref{properties-lemma-finite-presentation-module}），(3) 转化为下列代数断言：
设 $A$ 是环，$u : M \to N$ 是 $A$-模映射。假设
$R = \colim_{i \in I} R_i$ 是 $A$-代数的有向余极限。假设 $M$ 是有限 $A$-模，
$N$ 是有限表示 $A$-模，且
$u \otimes 1 : M \otimes_A R \to N \otimes_A R$ 是同构。则对某个 $i$，映射
$u \otimes 1 : M \otimes_A R_i \to N \otimes_A R_i$ 是同构。这就是《代数》，
引理 \ref{algebra-lemma-module-map-property-in-colimit} 的第 (3) 部分。
\end{proof}

\begin{situation}
\label{flat-situation-flat-at-point}
设 $(A, \mathfrak m_A)$ 为局部环。
记 $\mathcal{C}$ 为如下范畴：其对象是 $A$-代数 $A'$，
其中这些环是局部环，且代数结构
$A \to A'$ 是局部环之间的局部同态。
$\mathcal{C}$ 中对象 $A', A''$ 之间的态射是 $A$-代数的
局部同态 $A' \to A''$。
设 $A \to B$ 是局部环之间的局部环映射，且 $M$ 是 $B$-模。
若 $A'$ 是 $\mathcal{C}$ 的对象，则置 $B' = B \otimes_A A'$，
并将 $M' = M \otimes_A A'$ 视为 $B'$-模。
给定 $A' \in \Ob(\mathcal{C})$，考虑条件
\begin{equation}
\label{flat-equation-flat-at-primes}
\forall \mathfrak q \in V(\mathfrak m_{A'}B' + \mathfrak m_B B')
\subset \Spec(B') :
M'_{\mathfrak q}\text{ 在 }A'\text{ 上平坦}.
\end{equation}
注意这与《代数进阶》，方程
(\ref{more-algebra-equation-flat-at-primes}) 相似。
特别地，若 $A' \to A''$ 是 $\mathcal{C}$ 中的态射，且
(\ref{flat-equation-flat-at-primes}) 对 $A'$ 成立，则它对 $A''$ 也成立，见《代数进阶》，
引理 \ref{more-algebra-lemma-base-change-flat-at-primes}。
因此得到函子
\begin{equation}
\label{flat-equation-flat-at-point}
F_{lf} : \mathcal{C} \longrightarrow \textit{Sets}, \quad
A' \longrightarrow \left\{
\begin{matrix}
\{*\} & \text{若 }(\ref{flat-equation-flat-at-primes})\text{ 成立}, \\
\emptyset & \text{否则。} &
\end{matrix}
\right.
\end{equation}
\end{situation}

\begin{lemma}
\label{flat-lemma-flat-at-point}
在情形 \ref{flat-situation-flat-at-point} 中。
\begin{enumerate}
\item 若 $A' \to A''$ 是 $\mathcal{C}$ 中的平坦态射，
则 $F_{lf}(A') = F_{lf}(A'')$。
\item 若 $A \to B$ 本质有限表示，且
$M$ 是有限表示的 $B$-模，则 $F_{lf}$ 保持极限：
若 $\{A_i\}_{i \in I}$ 是 $\mathcal{C}$ 中对象的有向系统，则
$F_{lf}(\colim_i A_i) = \colim_i F_{lf}(A_i)$.
\end{enumerate}
\end{lemma}

\begin{proof}
第 (1) 部分是《代数进阶》引理
\ref{more-algebra-lemma-flat-descent-flat-at-primes} 的特例。
第 (2) 部分是《代数进阶》引理
\ref{more-algebra-lemma-limit-preserving-flat-at-primes} 的特例。
\end{proof}

\begin{lemma}
\label{flat-lemma-flat-at-point-finite}
在情形 \ref{flat-situation-flat-at-point} 中。设 $B \to C$ 是局部
$A$-代数之间的局部映射，且 $N$ 是 $C$-模。记
$F'_{lf} : \mathcal{C} \to \textit{Sets}$ 为与对 $(C, N)$ 关联的函子。
若 $M \cong N$ 作为 $B$-模且 $B \to C$ 有限，则
$F_{lf} = F'_{lf}$.
\end{lemma}

\begin{proof}
设 $A'$ 是 $\mathcal{C}$ 的对象。置 $C' = C \otimes_A A'$，
并仿照情形 \ref{flat-situation-flat-at-point} 中 $B'$、$M'$ 的定义
置 $N' = N \otimes_A A'$。
注意 $M' \cong N'$ 作为 $B'$-模。
假设 $B \to C$ 有限，则有两个结论：
(a) $\mathfrak m_C = \sqrt{\mathfrak m_B C}$，且 (b)
$B' \to C'$ 有限。结论 (a) 蕴含
$$
V(\mathfrak m_{A'}C' + \mathfrak m_C C')
=
\left(
\Spec(C') \to \Spec(B')
\right)^{-1}V(\mathfrak m_{A'}B' + \mathfrak m_B B').
$$
设 $\mathfrak q \subset V(\mathfrak m_{A'}B' + \mathfrak m_B B')$。
则 $M'_{\mathfrak q}$ 在 $A'$ 上平坦，当且仅当
$C'_{\mathfrak q}$-模 $N'_{\mathfrak q}$ 在 $A'$ 上平坦
（因为它们作为 $A'$-模同构），当且仅当
$C'_{\mathfrak q}$ 的每个极大理想 $\mathfrak r$ 都满足
$N'_{\mathfrak r}$ 在 $A'$ 上平坦（见《代数》，引理
\ref{algebra-lemma-flat-localization}）。
由 (b)，$B'_{\mathfrak q} \to C'_{\mathfrak q}$ 有限，
$C'_{\mathfrak q}$ 的极大理想恰好对应于
位于 $\mathfrak q$ 之上的 $C'$ 的素理想（见《代数》，引理
\ref{algebra-lemma-integral-going-up}），
而根据上面的显示方程，这些素理想全都包含在
$V(\mathfrak m_{A'}C' + \mathfrak m_C C')$ 中。因此引理得证。
\end{proof}

\begin{lemma}
\label{flat-lemma-flat-at-point-go-up}
在情形
\ref{flat-situation-flat-at-point} 中，设 $B \to C$ 是局部
环之间的平坦局部同态。置 $N = M \otimes_B C$。记
$F'_{lf} : \mathcal{C} \to \textit{Sets}$ 为与对 $(C, N)$ 关联
的函子。则 $F_{lf} = F'_{lf}$。
\end{lemma}

\begin{proof}
设 $A'$ 是 $\mathcal{C}$ 的对象。置 $C' = C \otimes_A A'$，
并仿照情形 \ref{flat-situation-flat-at-point} 中 $B'$、$M'$ 的定义
置 $N' = N \otimes_A A' = M' \otimes_{B'} C'$。注意
$$
V(\mathfrak m_{A'}B' + \mathfrak m_B B')
=
\Spec( \kappa(\mathfrak m_B) \otimes_A \kappa(\mathfrak m_{A'}) )
$$
并且对 $V(\mathfrak m_{A'}C' + \mathfrak m_C C')$ 也有类似等式。
环映射
$$
\kappa(\mathfrak m_B) \otimes_A \kappa(\mathfrak m_{A'})
\longrightarrow
\kappa(\mathfrak m_C) \otimes_A \kappa(\mathfrak m_{A'})
$$
是忠实平坦的，因此
$V(\mathfrak m_{A'}C' + \mathfrak m_C C') \to
V(\mathfrak m_{A'}B' + \mathfrak m_B B')$ 是满射。
最后，若 $\mathfrak r \in V(\mathfrak m_{A'}C' + \mathfrak m_C C')$
映到 $\mathfrak q \in V(\mathfrak m_{A'}B' + \mathfrak m_B B')$，则
$M'_{\mathfrak q}$ 在 $A'$ 上平坦，当且仅当
$N'_{\mathfrak r}$ 在 $A'$ 上平坦，因为 $B' \to C'$ 平坦，见
《代数》，引理 \ref{algebra-lemma-flatness-descends-more-general}。
引理由这些观察形式地推出。
\end{proof}

\begin{situation}
\label{flat-situation-free-at-generic-points}
设 $f : X \to S$ 是具有几何不可约纤维的光滑态射。
设 $\mathcal{F}$ 是有限型拟凝聚 $\mathcal{O}_X$-模。对任意
$S$ 上的概形 $T$，记 $\mathcal{F}_T$ 为 $\mathcal{F}$ 到 $T$ 的基变换，
换言之，$\mathcal{F}_T$ 是沿投影态射
$X_T = X \times_S T \to X$ 的 $\mathcal{F}$ 的拉回。注意 $X_T \to T$
是具有几何不可约纤维的光滑态射，见《态射》引理
\ref{morphisms-lemma-base-change-smooth} 及《态射进阶》引理
\ref{more-morphisms-lemma-base-change-fibres-geometrically-irreducible}。
设 $p \geq 0$ 为整数。给定点 $t \in T$，考虑条件
\begin{equation}
\label{flat-equation-free-at-generic-point-fibre}
\mathcal{F}_T\text{ 在 }\xi_t\text{ 的一个邻域内自由且秩为 }p
\end{equation}
其中 $\xi_t$ 是纤维 $X_t$ 的泛点。对所有 $t \in T$，此条件在基变换下稳定，
因此得到函子
\begin{equation}
\label{flat-equation-free-at-generic-points}
H_p : (\Sch/S)^{opp} \longrightarrow \textit{Sets}, \quad
T \longrightarrow \left\{
\begin{matrix}
\{*\} & \text{若 }\mathcal{F}_T\text{ 满足
(\ref{flat-equation-free-at-generic-point-fibre}) }\forall t\in T, \\
\emptyset & \text{否则。}
\end{matrix}
\right.
\end{equation}
\end{situation}

\begin{lemma}
\label{flat-lemma-free-at-generic-points}
在情形 \ref{flat-situation-free-at-generic-points} 中。
\begin{enumerate}
\item 函子 $H_p$ 对 fpqc 拓扑满足层性质。
\item 若 $\mathcal{F}$ 有限表示，则函子 $H_p$ 保持极限。
\end{enumerate}
\end{lemma}

\begin{proof}
设 $\{T_i \to T\}_{i \in I}$ 是概形在 $S$ 上的 fpqc\footnote{利用有限表示的平坦态射是开映射，很容易证明 $H_p$ 是 fppf 拓扑下的层。这正是我们稍后真正需要的。不过，直接证明它也满足 fpqc 拓扑的层条件也颇有趣。} 覆盖。
置 $X_i = X_{T_i} = X \times_S T_i$，记 $\mathcal{F}_i$ 为 $\mathcal{F}$ 到 $X_i$ 的拉回。
假设对所有 $i$，$\mathcal{F}_i$ 都满足
(\ref{flat-equation-free-at-generic-point-fibre})。
取 $t \in T$，令 $\xi_t \in X_T$ 表示 $X_t$ 的泛点。
我们需要证明 $\mathcal{F}$ 在 $\xi_t$ 的某个邻域内自由。
对某个 $i \in I$，可找到映到 $t$ 的 $t_i \in T_i$。
令 $\xi_i \in X_i$ 表示 $X_{t_i}$ 的泛点，于是 $\xi_i$ 映到 $\xi_t$。
由于 $\mathcal{F}_i$ 在 $\xi_i$ 的某个邻域内秩为 $p$ 的自由性，得到
$(\mathcal{F}_i)_{x_i} \cong \mathcal{O}_{X_i, x_i}^{\oplus p}$
这蕴含
$\mathcal{F}_{T, \xi_t} \cong \mathcal{O}_{X_T, \xi_t}^{\oplus p}$
因为 $\mathcal{O}_{X_T, \xi_t} \to \mathcal{O}_{X_i, x_i}$ 平坦，见例如
《代数》，引理 \ref{algebra-lemma-finite-projective-descends}。
因此在 $X_T$ 中存在 $\xi_t$ 的仿射邻域 $U$ 以及满射
$\mathcal{O}_U^{\oplus p} \to \mathcal{F}_U = \mathcal{F}_T|_U$，见
《模》，引理 \ref{modules-lemma-finite-type-surjective-on-stalk}。
缩小 $T$ 后可假设 $U \to T$ 是满射。因此 $U \to T$ 是具有几何不可约
纤维的仿射光滑态射。此外，对每个 $t' \in T$，诱导映射
$$
\alpha :
\mathcal{O}_{U, \xi_{t'}}^{\oplus p}
\longrightarrow
\mathcal{F}_{U, \xi_{t'}}
$$
是同构（因为由同样的论证，右侧的模是秩为 $p$ 的自由模）。由引理
\ref{flat-lemma-induction-step} 得
$$
\Gamma(U, \mathcal{O}_U^{\oplus p})
\otimes_{\Gamma(T, \mathcal{O}_T)} \mathcal{O}_{T, t'}
\longrightarrow
\Gamma(U, \mathcal{F}_U)
\otimes_{\Gamma(T, \mathcal{O}_T)} \mathcal{O}_{T, t'}
$$
对每个 $t' \in T$ 都是单射。因此满射 $\alpha$ 是同构。这就完成了 (1) 的证明。

\medskip\noindent
假设 $\mathcal{F}$ 有限表示。
设 $T = \lim_{i \in I} T_i$ 是仿射 $S$-概形的有向极限，
并假设 $\mathcal{F}_T$ 满足
(\ref{flat-equation-free-at-generic-point-fibre}).
置 $X_i = X_{T_i} = X \times_S T_i$，记 $\mathcal{F}_i$ 为
$\mathcal{F}$ 到 $X_i$ 的拉回。
记 $U \subset X_T$ 为 $\mathcal{F}_T$ 在 $T$ 上平坦的点所成的开子概形，见
《态射进阶》，定理 \ref{more-morphisms-theorem-openness-flatness}。
由假设，每个纤维的泛点都是 $U$ 中的点，即
$U \to T$ 是具有几何不可约纤维的光滑满射。可稍微缩小 $U$，并假设 $U$ 是拟紧的。
利用《极限》引理 \ref{limits-lemma-descend-opens}，
可找到 $i \in I$ 及拟紧开集 $U_i \subset X_i$，
使其在 $X_T$ 中的逆像为 $U$。增大 $i$ 后可假设
$\mathcal{F}_i|_{U_i}$ 在 $T_i$ 上平坦，见《极限》引理
\ref{limits-lemma-descend-module-flat-finite-presentation}。
特别地，$\mathcal{F}_i|_{U_i}$ 有限局部自由，
因此定义局部常值的
秩函数 $\rho : U_i \to \{0, 1, 2, \ldots \}$。
令 $(U_i)_p \subset U_i$ 表示 $\rho$ 取值为 $p$ 的开闭子集。令
$V_i \subset T_i$ 为 $(U_i)_p$ 的像；注意 $V_i$ 是开且拟紧的。
由假设，$T \to T_i$ 的像包含于 $V_i$ 中。
因此由《极限》引理 \ref{limits-lemma-descend-opens}，
存在 $i' \geq i$ 使 $T_{i'} \to T_i$ 经由 $V_i$ 分解。
于是 $\mathcal{F}_{i'}$ 满足
(\ref{flat-equation-free-at-generic-point-fibre})，如所需。细节略去。
\end{proof}

\begin{lemma}
\label{flat-lemma-pre-flat-dimension-n}
设 $f : X \to S$ 是局部有限型的概形态射。
设 $\mathcal{F}$ 是有限型拟凝聚 $\mathcal{O}_X$-模。设 $n \geq 0$。
以下条件等价：
\begin{enumerate}
\item 对于 $s \in S$，令 $Z \subset X_s$ 为 $\mathcal{F}$ 在 $S$ 上不平坦的点所成的闭子集（见引理
\ref{flat-lemma-open-in-fibre-where-flat}），则 $\dim(Z) < n$，并且
\item 对于 $x \in X$，若 $\mathcal{F}$ 在 $x$ 处不是在 $S$ 上平坦的，则
$\text{trdeg}_{\kappa(f(x))}(\kappa(x)) < n$。
\end{enumerate}
若这些条件成立，则在任意基变换后仍成立。
\end{lemma}

\begin{proof}
设 $x \in X$ 是位于 $s \in S$ 之上的点。
由《簇》引理 \ref{varieties-lemma-dimension-locally-algebraic}，
$\{x\}$ 在 $X_s$ 中的闭包的维数为
$\text{trdeg}_{\kappa(s)}(\kappa(x))$。
反之，若 $Z \subset X_s$ 是维数为 $d$ 的闭子集，则存在 $x \in Z$，
使 $\text{trdeg}_{\kappa(s)}(\kappa(x)) = d$（同一引理）。
因此 (1) 与 (2) 等价（甚至逐纤维成立）。
基变换的断言由《态射》引理
\ref{morphisms-lemma-base-change-module-flat} 和
\ref{morphisms-lemma-dimension-fibre-after-base-change} 得出。
\end{proof}

\begin{definition}
\label{flat-definition-flat-dimension-n}
设 $f : X \to S$ 是局部有限型的概形态射。
设 $\mathcal{F}$ 是有限型拟凝聚 $\mathcal{O}_X$-模。设 $n \geq 0$。
若引理 \ref{flat-lemma-pre-flat-dimension-n} 的等价条件成立，
则称 {\it $\mathcal{F}$ 在维数 $\geq n$ 上关于 $S$ 平坦}。
\end{definition}

\begin{situation}
\label{flat-situation-flat-dimension-n}
设 $f : X \to S$ 是局部有限型的概形态射。
设 $\mathcal{F}$ 是有限型拟凝聚 $\mathcal{O}_X$-模。对任意
$S$ 上的概形 $T$，记 $\mathcal{F}_T$ 为 $\mathcal{F}$ 到 $T$ 的基变换，
换言之，$\mathcal{F}_T$ 是沿投影态射
$X_T = X \times_S T \to X$ 的 $\mathcal{F}$ 的拉回。注意 $X_T \to T$ 是有限型的，
且 $\mathcal{F}_T$ 是有限型 $\mathcal{O}_{X_T}$-模（《态射》引理
\ref{morphisms-lemma-base-change-finite-type} 及《模》引理
\ref{modules-lemma-pullback-finite-type}）。
设 $n \geq 0$。由定义 \ref{flat-definition-flat-dimension-n} 与引理
\ref{flat-lemma-pre-flat-dimension-n} 得到函子
\begin{equation}
\label{flat-equation-flat-dimension-n}
F_n : (\Sch/S)^{opp} \longrightarrow \textit{Sets}, \quad
T \longrightarrow \left\{
\begin{matrix}
\{*\} & \text{若 }\mathcal{F}_T\text{ 在 }T\text{ 上于 }\dim \geq n\text{ 的部分平坦}, \\
\emptyset & \text{否则。}
\end{matrix}
\right.
\end{equation}
\end{situation}

\begin{lemma}
\label{flat-lemma-flat-dimension-n}
在情形 \ref{flat-situation-flat-dimension-n} 中。
\begin{enumerate}
\item 函子 $F_n$ 对 fpqc 拓扑满足层性质。
\item 若 $f$ 拟紧且局部有限表示，并且 $\mathcal{F}$ 有限表示，
则函子 $F_n$ 保持极限。
\end{enumerate}
\end{lemma}

\begin{proof}
设 $\{T_i \to T\}_{i \in I}$ 是 $S$ 上概形的 fpqc 覆盖。
置 $X_i = X_{T_i} = X \times_S T_i$，记 $\mathcal{F}_i$ 为
$\mathcal{F}$ 到 $X_i$ 的拉回。假设对所有 $i$，$\mathcal{F}_i$
在维数 $\geq n$ 上关于 $T_i$ 平坦。取 $t \in T$。
选取指标 $i$ 及映到 $t$ 的点 $t_i \in T_i$。
考虑笛卡尔图
$$
\xymatrix{
X_{\Spec(\mathcal{O}_{T, t})} \ar[d] &
X_{\Spec(\mathcal{O}_{T_i, t_i})} \ar[d] \ar[l] \\
\Spec(\mathcal{O}_{T, t}) &
\Spec(\mathcal{O}_{T_i, t_i}) \ar[l]
}
$$
由于下方水平态射平坦，由《态射进阶》引理
\ref{more-morphisms-lemma-flat-locus-base-change}，$\mathcal{F}_i$ 在
$T_i$ 上不平坦的集合 $Z_i \subset X_{t_i}$ 与 $\mathcal{F}_T$ 在 $T$ 上不平坦的
集合 $Z \subset X_t$ 满足关系 $Z_i = Z_{\kappa(t_i)}$。
于是由《态射》引理 \ref{morphisms-lemma-dimension-fibre-after-base-change}，
得到 $\mathcal{F}_T$ 在维数 $\geq n$ 上关于 $T$ 平坦。

\medskip\noindent
假设 $f$ 拟紧且局部有限表示，并且 $\mathcal{F}$ 有限表示。
本段先将 (2) 的证明约化到 $f$ 有限表示的情形。
设 $T = \lim_{i \in I} T_i$ 是仿射 $S$-概形的有向极限，
并假设 $\mathcal{F}_T$ 在维数 $\geq n$ 上平坦。置
$X_i = X_{T_i} = X \times_S T_i$，记 $\mathcal{F}_i$ 为 $\mathcal{F}$ 到 $X_i$ 的拉回。
我们需要证明对某个 $i$，$\mathcal{F}_i$ 在维数 $\geq n$ 上平坦。
取 $i_0 \in I$，并将 $I$ 换为 $\{i \mid i \geq i_0\}$。
由于 $T_{i_0}$ 仿射（因而拟紧），存在有限个仿射开集
$W_j \subset S$（$j = 1, \ldots, m$）以及一个仿射开覆盖
$T_{i_0} = \bigcup_{j = 1, \ldots, m} V_{j, i_0}$，
使得 $T_{i_0} \to S$ 将 $V_{j, i_0}$ 映入 $W_j$。
对 $i \geq i_0$，记 $V_{j, i}$ 为 $V_{j, i_0}$ 在 $T_i$ 中的逆像。
若能对每个 $j$ 证明存在某个 $i$ 使 $\mathcal{F}_{V_{j, i_0}}$
在维数 $\geq n$ 上平坦，则结论成立。这样便可约化到 $S$ 仿射的情形。
此时 $X$ 拟紧，可取有限仿射开覆盖 $X = W_1 \cup \ldots \cup W_m$。
此时 $(X \to S, \mathcal{F})$ 的结论等价于
$(\coprod W_j, \coprod \mathcal{F}|_{W_j})$ 的结论。因此可假设
$f$ 有限表示。

\medskip\noindent
假设 $f$ 有限表示且 $\mathcal{F}$ 有限表示。记 $U \subset X_T$ 为
$\mathcal{F}_T$ 在 $T$ 上平坦的点所成的开子概形，见《态射进阶》定理
\ref{more-morphisms-theorem-openness-flatness}。
由假设，$Z = X_T \setminus U$ 在 $T$ 上每个纤维的维数都为 $< n$。
由《极限》引理 \ref{limits-lemma-approximate-given-relative-dimension}，
可找到闭子概形 $Z \subset Z' \subset X_T$，使对所有 $t \in T$ 都有
$\dim(Z'_t) < n$，且 $Z' \to X_T$ 有限表示。由《极限》引理
\ref{limits-lemma-descend-finite-presentation} 和
\ref{limits-lemma-descend-closed-immersion-finite-presentation}，
存在 $i \in I$ 及有限表示闭子概形 $Z'_i \subset X_i$，
其到 $T$ 的基变换为 $Z'$。由《极限》引理
\ref{limits-lemma-limit-dimension}，可假设 $Z'_i \to T_i$ 的所有纤维维数都为 $< n$。
由《极限》引理 \ref{limits-lemma-descend-module-flat-finite-presentation}，
可假设 $\mathcal{F}_i|_{X_i \setminus T'_i}$ 在 $T_i$ 上平坦。
这蕴含 $\mathcal{F}_i$ 在维数 $\geq n$ 上平坦；这里使用
$Z' \to X_T$ 有限表示，因而补集 $X_T \setminus Z'$ 拟紧！
于是 (2) 得证，引理的证明完成。
\end{proof}

\begin{situation}
\label{flat-situation-flat}
设 $f : X \to S$ 是概形态射。
设 $\mathcal{F}$ 是拟凝聚 $\mathcal{O}_X$-模。
对任意 $S$ 上的概形 $T$，记 $\mathcal{F}_T$ 为 $\mathcal{F}$ 到 $T$ 的
基变换，换言之，$\mathcal{F}_T$ 是沿投影态射
$X_T = X \times_S T \to X$ 的 $\mathcal{F}$ 的拉回。由于平坦模的基变换
仍平坦，得到函子
\begin{equation}
\label{flat-equation-flat}
F_{flat} : (\Sch/S)^{opp} \longrightarrow \textit{Sets}, \quad
T \longrightarrow \left\{
\begin{matrix}
\{*\} & \text{若 } \mathcal{F}_T\text{ 在 }T\text{ 上平坦}, \\
\emptyset & \text{否则。}
\end{matrix}
\right.
\end{equation}
\end{situation}

\begin{lemma}
\label{flat-lemma-flat}
在情形 \ref{flat-situation-flat} 中。
\begin{enumerate}
\item 函子 $F_{flat}$ 对 fpqc 拓扑满足层性质。
\item 若 $f$ 拟紧且局部有限表示，并且 $\mathcal{F}$ 有限表示，
则函子 $F_{flat}$ 保持极限。
\end{enumerate}
\end{lemma}

\begin{proof}
第 (1) 部分由以下断言推出：若 $T' \to T$ 是 $S$ 上概形的满射平坦态射，
则 $\mathcal{F}_{T'}$ 在 $T'$ 上平坦，当且仅当 $\mathcal{F}_T$ 在 $T$ 上平坦，见
《态射进阶》引理 \ref{more-morphisms-lemma-flat-locus-base-change}。
第 (2) 部分在将 $X$、$S$ 约化到仿射情形后，由《极限》引理
\ref{limits-lemma-descend-module-flat-finite-presentation} 得出（比较引理
\ref{flat-lemma-flat-dimension-n} 的证明）。
\end{proof}




\section{扁平化分层}
\label{flat-section-flattening}

\noindent
这里只给出定义。寻找“泛平坦性分层”的读者应参阅
《态射进阶》第
\ref{more-morphisms-section-generic-flatness-stratification} 节。

\begin{definition}
\label{flat-definition-flattening}
设 $X \to S$ 是概形态射。
设 $\mathcal{F}$ 是拟凝聚 $\mathcal{O}_X$-模。
若情形 \ref{flat-situation-flat} 中定义的函子 $F_{flat}$
可由 $S$ 上的概形 $S'$ 表示，则称 {\it $\mathcal{F}$ 的泛扁平化存在}。
若 $\mathcal{O}_X$ 的泛扁平化存在，则称 {\it $X$ 的泛扁平化存在}。
\end{definition}

\noindent
注意，若 $\mathcal{F}$ 的泛扁平化 $S'$\footnote{概形 $S'$ 有时称为 {\it universal flatificator}；在 \cite{GruRay} 中称为 {\it platificateur universel}。泛扁平化的存在性不应与《代数进阶》第 \ref{more-algebra-section-blowup-flat} 节讨论的结果类型混淆。} 存在，
则态射 $S' \to S$ 是概形的满射单射，$\mathcal{F}_{S'}$ 在 $S'$ 上平坦，
并且态射 $T \to S$ 穿过 $S'$ 分解，当且仅当 $\mathcal{F}_T$ 在 $T$ 上平坦。

\begin{example}
\label{flat-example-no-universal-flattening}
设 $X = S = \Spec(k[x, y])$，其中 $k$ 是域。令
$\mathcal{F} = \widetilde{M}$，其中 $M = k[x, x^{-1}, y]/(y)$。
对 $k[x, y]$-代数 $A$，置 $F_{flat}(A) = F_{flat}(\Spec(A))$。
则对所有 $n$ 有 $F_{flat}(k[x, y]/(x, y)^n) = \{*\}$，但
$F_{flat}(k[[x, y]]) = \emptyset$。这意味着 $F_{flat}$ 不可表示
（甚至不能由代数空间表示，见《形式概形》引理
\ref{formal-spaces-lemma-map-into-algebraic-space}）。
因此在此情形泛扁平化不存在。
\end{example}

\noindent
我们定义（比较《拓扑》注记
\ref{topology-remark-locally-finite-stratification}）：概形 $S$ 的（局部有限的、
概形论的）{\it 分层} 由偏序集 $I$ 所指标的一族闭子概形 $Z_i \subset S$ 给出，
满足 $S = \bigcup Z_i$（集合论意义下），$S$ 的每个点都有一个邻域只与有限多个
$Z_i$ 相交，并且满足
$$
Z_i \cap Z_j = \bigcup\nolimits_{k \leq i, j} Z_k.
$$
置 $S_i = Z_i \setminus \bigcup_{j < i} Z_j$，实际的分层是将
该概形分解为局部闭子概形的 $S = \coprod S_i$。我们常只标出分层
$S_i$，把闭子概形 $Z_i$ 的构造留给读者。给定一个分层，可得到单射
$$
S' = \coprod\nolimits_{i \in I} S_i \longrightarrow S.
$$
称之为与该分层关联的 {\it 单射}。用此术语可定义扁平化分层的含义。

\begin{definition}
\label{flat-definition-flattening-stratification}
设 $X \to S$ 是概形态射。
设 $\mathcal{F}$ 是拟凝聚 $\mathcal{O}_X$-模。
若情形 \ref{flat-situation-flat} 中定义的函子 $F_{flat}$
可由单射 $S' \to S$ 表示，且该单射与由局部闭子概形给出的
$S$ 的一个分层关联，则称 $\mathcal{F}$ 有 {\it 扁平化分层}。
若 $\mathcal{O}_X$ 有扁平化分层，则称 $X$ 有 {\it 扁平化分层}。
\end{definition}

\noindent
当扁平化分层存在时，理解给分层编号的指标集及其偏序通常很重要。
这往往与模的秩有关。例如，当 $X = S$ 且 $\mathcal{F}$ 是有限表示的
$\mathcal{O}_S$-模时，扁平化分层存在，并由 $\mathcal{F}$ 的 Fitting 理想给出，见
《因子》引理 \ref{divisors-lemma-finite-presentation-module}。




\section{Artin 环上的扁平化分层}
\label{flat-section-flattening-artinian}

\noindent
当底概形是 Artin 环的谱时，扁平化分层存在。

\begin{lemma}
\label{flat-lemma-flattening-stratification-artinian}
设 $S$ 是 Artin 环的谱。
对 $S$ 上任意概形 $X$ 及任意拟凝聚 $\mathcal{O}_X$-模，
泛扁平化存在。事实上，泛扁平化由闭浸入 $S' \to S$ 给出，
因而也是 $\mathcal{F}$ 的扁平化分层。
\end{lemma}

\begin{proof}
取仿射开覆盖 $X = \bigcup U_i$。
则 $F_{flat}$ 是与各对 $(U_i, \mathcal{F}|_{U_i})$ 关联的函子的乘积。
因此只需对每个 $(U_i, \mathcal{F}|_{U_i})$ 证明结论。
在仿射情形，引理由《代数进阶》引理
\ref{more-algebra-lemma-flattening-artinian-universal-property} 立即推出。
\end{proof}






\section{态射的扁平化}
\label{flat-section-flattening-map}

\noindent
定理 \ref{flat-theorem-flattening-map}
是后续扁平化论断的关键。

\begin{lemma}
\label{flat-lemma-universally-separating}
设 $S$ 是概形。
设 $g : X' \to X$ 是 $S$ 上概形的平坦态射，且 $X$ 在 $S$ 上局部有限型。
设 $\mathcal{F}$ 是有限型拟凝聚 $\mathcal{O}_X$-模，并且在 $S$ 上平坦。
若 $\text{Ass}_{X/S}(\mathcal{F}) \subset g(X')$，则典范映射
$$
\mathcal{F} \longrightarrow g_*g^*\mathcal{F}
$$
是单射，并且在任意基变换后仍为单射。
\end{lemma}

\begin{proof}
最后的断言意味着，对任意态射 $T \to S$，映射
$\mathcal{F}_T \to (g_T)_*g_T^*\mathcal{F}_T$ 是单射。条件
$\text{Ass}_{X/S}(\mathcal{F}) \subset g(X')$ 在基变换下保持，见《因子》引理
\ref{divisors-lemma-base-change-relative-assassin} 及注记
\ref{divisors-remark-base-change-relative-assassin}。
平坦性与有限型的假设也具有同样性质。
因此只需证明所显示箭头的单射性。
令 $\mathcal{K} = \Ker(\mathcal{F} \to g_*g^*\mathcal{F})$。
目标是证明 $\mathcal{K} = 0$。
为此只需证明 $\text{WeakAss}_X(\mathcal{K}) = \emptyset$，见《因子》引理
\ref{divisors-lemma-weakly-ass-zero}。
我们有
$\text{WeakAss}_X(\mathcal{K}) \subset \text{WeakAss}_X(\mathcal{F})$，见《因子》引理
\ref{divisors-lemma-ses-weakly-ass}。
由于 $\mathcal{F}$ 平坦，由引理 \ref{flat-lemma-bourbaki-finite-type-general-base} 得
$\text{WeakAss}_X(\mathcal{F}) \subset \text{Ass}_{X/S}(\mathcal{F})$。
由假设，$\text{Ass}_{X/S}(\mathcal{F})$ 的任意点 $x$ 都是某个 $x' \in X'$ 的像。
由于 $g$ 平坦，局部环映射 $\mathcal{O}_{X, x} \to \mathcal{O}_{X', x'}$
忠实平坦，故映射
$$
\mathcal{F}_x
\longrightarrow
g^*\mathcal{F}_{x'} =
\mathcal{F}_x \otimes_{\mathcal{O}_{X, x}} \mathcal{O}_{X', x'}
$$
是单射（见《代数》引理
\ref{algebra-lemma-faithfully-flat-universally-injective}）。
这就蕴含 $\mathcal{K}_x = 0$，即得所需结论。
\end{proof}

\begin{lemma}
\label{flat-lemma-flattening-module-map}
设 $A$ 是环，$u : M \to N$ 是 $A$-模的满射。
若 $M$ 是投射 $A$-模，则存在理想 $I \subset A$，使对任意环映射
$\varphi : A \to B$，以下条件等价：
\begin{enumerate}
\item $u \otimes 1 : M \otimes_A B \to N \otimes_A B$ 是同构；
\item $\varphi(I) = 0$.
\end{enumerate}
\end{lemma}

\begin{proof}
由于 $M$ 投射，可取投射 $A$-模 $C$，使 $F = M \oplus C$ 是自由 $A$-模。
将 $u$ 替换为 $u \oplus 1 : F = M \oplus C \to N \oplus C$ 后，
可假设 $M$ 自由。在此情形，令 $I$ 为相对于 $M$ 的某个（固定）基，
$\Ker(u)$ 中所有元素的系数在 $A$ 中生成的理想。
这是因为 $u$ 满射且 $\otimes_A B$ 右正合，所以
$\Ker(u \otimes 1)$ 是 $\Ker(u) \otimes_A B$ 在 $M \otimes_A B$ 中的像。
\end{proof}

\begin{theorem}
\label{flat-theorem-flattening-map}
在情形
\ref{flat-situation-iso} 中，假设
\begin{enumerate}
\item $f$ 有限表示；
\item $\mathcal{F}$ 有限表示、在 $S$ 上平坦且相对于 $S$ 纯；
\item $u$ 是满射。
\end{enumerate}
则 $F_{iso}$ 可由闭浸入 $Z \to S$ 表示。
此外，若 $\mathcal{G}$ 有限表示，则 $Z \to S$ 有限表示。
\end{theorem}

\begin{proof}
下文不再特别说明，我们使用 $\mathcal{F}$ 在 $S$ 上普遍纯这一事实，见引理
\ref{flat-lemma-finite-type-flat-pure-along-fibre-is-universal}。
由引理 \ref{flat-lemma-iso-sheaf} 以及《下降》引理
\ref{descent-lemma-closed-immersion} 和
\ref{descent-lemma-descent-data-sheaves}，该问题对 $S$ 上的 \'etale 拓扑是局部的。
因此只需证明：给定 $s \in S$，存在 $(S, s)$ 的一个 \'etale 邻域，使定理成立。

\medskip\noindent
利用引理 \ref{flat-lemma-finite-presentation-flat-along-fibre}，
并将 $S$ 替换为 $s$ 的一个初等 \'etale 邻域后，可假设存在交换图
$$
\xymatrix{
X \ar[dr] & & X' \ar[ll]^g \ar[ld] \\
& S &
}
$$
其中概形均在 $S$ 上有限表示，且 $g$ 是 \'etale，$X_s \subset g(X')$，
概形 $X'$ 与 $S$ 仿射，并且
$\Gamma(X', g^*\mathcal{F})$ 是 $\Gamma(S, \mathcal{O}_S)$-投射模。
注意 $g^*\mathcal{F}$ 在 $S$ 上普遍纯，见引理
\ref{flat-lemma-affine-locally-projective-pure}。
因此由引理 \ref{flat-lemma-criterion}，开集 $g(X')$ 包含
位于 $\Spec(\mathcal{O}_{S, s})$ 之上的
$\text{Ass}_{X/S}(\mathcal{F})$ 中的点。置
$$
E = \{t \in S \mid \text{Ass}_{X_t}(\mathcal{F}_t) \subset g(X') \}.
$$
由《态射进阶》引理
\ref{more-morphisms-lemma-relative-assassin-constructible}，$E$ 是 $S$ 的可构造子集。
我们已知 $\Spec(\mathcal{O}_{S, s}) \subset E$。由《态射》引理
\ref{morphisms-lemma-constructible-containing-open}，$E$ 包含 $s$ 的一个开邻域。
因此将 $S$ 替换为 $s$ 的较小仿射邻域后，可假设
$\text{Ass}_{X/S}(\mathcal{F}) \subset g(X')$。

\medskip\noindent
由于假设 $u$ 是满射，有 $F_{iso} = F_{inj}$。由引理
\ref{flat-lemma-universally-separating}，$u : \mathcal{F} \to \mathcal{G}$ 是单射，
当且仅当 $g^*u : g^*\mathcal{F} \to g^*\mathcal{G}$ 是单射；在任意基变换后
同样成立。因此可约化到如下情形：除定理假设外，$X \to S$ 是仿射概形的态射，
且 $\Gamma(X, \mathcal{F})$ 是 $\Gamma(S, \mathcal{O}_S)$-投射模。
此情形立即由引理 \ref{flat-lemma-flattening-module-map} 得出。

\medskip\noindent
若 $\mathcal{G}$ 有限表示，则将引理
\ref{flat-lemma-iso-sheaf} 第 (4) 部分与《极限》注记
\ref{limits-remark-limit-preserving} 合并，即得 $Z$ 有限表示。
\end{proof}

\begin{lemma}
\label{flat-lemma-Weil-restriction-closed-subschemes}
设 $f:X\to S$ 是有限表示、平坦且纯的概形态射。设 $Y$ 是 $X$ 的闭子概形。
令 $F=f_*Y$ 为沿 $f$ 的 $Y$ 的 Weil 限制函子，定义为
$$
F : (\Sch/S)^{opp} \to \textit{Sets}, \quad
T \mapsto
\left\{
\begin{matrix}
\{*\} & \text{若} & Y_T\to X_T \text{ 是同构，}\\
\emptyset & \text{否则。} &
\end{matrix}
\right.
$$
则 $F$ 可由闭浸入 $Z\to S$ 表示。此外，若 $Y\to S$ 有限表示，
则 $Z\to S$ 有限表示。
\end{lemma}

\begin{proof}
令 $\mathcal{I}$ 是 $X$ 中定义 $Y$ 的理想层，并令
$u:\mathcal{O}_X\to\mathcal{O}_X/\mathcal{I}$ 为满射。
对 $S$-概形 $T$，闭浸入 $Y_T\to X_T$ 是同构，当且仅当
$u_T$ 是同构。因此结论由定理 \ref{flat-theorem-flattening-map} 得出。
\end{proof}



\section{局部情形下的扁平化}
\label{flat-section-flattening-local}

\noindent
本节开始应用前面的材料，得到扁平化分层的一个影子。

\begin{theorem}
\label{flat-theorem-flattening-local}
在情形 \ref{flat-situation-flat-at-point} 中，假设 $A$ 是 Hensel 环，
$B$ 是 $A$ 上本质有限型的，且 $M$ 是有限 $B$-模。
则存在理想 $I \subset A$，使 $A/I$ 在范畴 $\mathcal{C}$ 上共表示函子 $F_{lf}$。
换言之，给定局部环的局部同态
$\varphi : A \to A'$，其中 $B' = B \otimes_A A'$ 且 $M' = M \otimes_A A'$，
以下条件等价：
\begin{enumerate}
\item $\forall \mathfrak q \in V(\mathfrak m_{A'}B' + \mathfrak m_B B')
\subset \Spec(B') :
M'_{\mathfrak q}\text{ 在 }A'\text{ 上平坦}$，且
\item $\varphi(I) = 0$。
\end{enumerate}
若 $B$ 在 $A$ 上本质有限表示且 $M$ 在 $B$ 上有限表示，
则 $I$ 是有限生成理想。
\end{theorem}

\begin{proof}
取有限型环映射 $A \to C$、有限 $C$-模 $N$ 及 $C$ 的素理想
$\mathfrak q$，使 $B = C_{\mathfrak q}$ 且 $M = N_{\mathfrak q}$。
下文中，当我们说
“对 $(N/C/A, \mathfrak q)$ 定理成立”是指
是指它对 $(A \to B, M)$ 成立，其中 $B = C_{\mathfrak q}$ 且
$M = N_{\mathfrak q}$。由引理 \ref{flat-lemma-flat-at-point-go-up}，
若将 $B$ 替换为在 $B$ 上平坦的局部环，函子 $F_{lf}$ 不变。
因此由于 $A$ 是 Hensel 环，可应用引理 \ref{flat-lemma-existence-algebra}，
并假设 $N/C/A$ 在 $\mathfrak q$ 处存在完全 d\'evissage。

\medskip\noindent
令 $(A_i, B_i, M_i, \alpha_i, \mathfrak q_i)_{i = 1, \ldots, n}$
为 $N/C/A$ 在 $\mathfrak q$ 处的一个完全 d\'evissage。令
$\mathfrak q'_i \subset A_i$ 为位于 $\mathfrak q_i \subset B_i$ 之上的唯一素理想，
如定义 \ref{flat-definition-complete-devissage-at-x-algebra} 所述。
由于 $C \to A_1$ 满射且 $N \cong M_1$ 作为 $C$-模，由引理
\ref{flat-lemma-flat-at-point-finite}，只需证明定理对
$(M_1/A_1/A, \mathfrak q'_1)$ 成立。
由于 $B_1 \to A_1$ 有限，且 $\mathfrak q_1$ 是位于 $\mathfrak q'_1$ 之上的
$B_1$ 的唯一素理想，故 $(A_1)_{\mathfrak q'_1} \to (B_1)_{\mathfrak q_1}$ 有限
（见《代数》引理 \ref{algebra-lemma-unique-prime-over-localize-below} 或《态射进阶》引理
\ref{more-morphisms-lemma-finite-morphism-single-point-in-fibre}）。
因此由引理 \ref{flat-lemma-flat-at-point-finite}，只需证明定理对
$(M_1/B_1/A, \mathfrak q_1)$ 成立。

\medskip\noindent
此时可对 d\'evissage 的长度 $n$ 作归纳，假设定理对
$(M_2/B_2/A, \mathfrak q_2)$ 成立。（若 $n = 1$，则 $M_2 = 0$，
它在 $A$ 上平坦。）反向执行上一段的最后几个步骤，并使用
$M_2 \cong \Coker(\alpha_2)$ 作为 $B_1$-模这一事实，可知定理对
$(\Coker(\alpha_1)/B_1/A, \mathfrak q_1)$ 成立。

\medskip\noindent
设 $A'$ 是 $\mathcal{C}$ 的对象。此时应用引理
\ref{flat-lemma-induction-step}：若对 $B_1 \otimes_A A'$ 中位于
$\mathfrak m_{A'}$ 之上的素理想 $\mathfrak q'$，有
$(M_1 \otimes_A A')_{\mathfrak q'}$ 在 $A'$ 上平坦，则
$(\Coker(\alpha_1) \otimes_A A')_{\mathfrak q'}$ 在 $A'$ 上平坦。
因此 $F_{lf}$ 是与 $(B_1)_{\mathfrak q_1}$ 上的模
$\Coker(\alpha_1)_{\mathfrak q_1}$ 关联的函子 $F'_{lf}$ 的子函子。
由上一段，$F'_{lf}$ 由某个理想 $J \subset A$ 的商环 $A/J$ 共表示。
故可将 $A$ 替换为 $A/J$，并假设
$\Coker(\alpha_1)_{\mathfrak q_1}$ 在 $A$ 上平坦。

\medskip\noindent
由于 $\Coker(\alpha_1)$ 是 $B_1$-模，且 $N_1/B_1/A$ 在 $\mathfrak q_1$ 处
存在完全 d\'evissage，又由引理
\ref{flat-lemma-complete-devissage-flat-finite-type-module}，
$\Coker(\alpha_1)_{\mathfrak q_1}$ 在 $A$ 上平坦，
可知 $\Coker(\alpha_1)$ 是自由 $A$-模，特别地在 $A$ 上平坦。因此由引理
\ref{flat-lemma-induction-step}，$F_{lf}(A')$ 非空，当且仅当
$\alpha \otimes 1_{A'}$ 是单射。令 $N_1 = \Im(\alpha_1) \subset M_1$，于是有正合列
$$
0 \to N_1 \to M_1 \to \Coker(\alpha_1) \to 0
\quad\text{且}\quad
B_1^{\oplus r_1} \to N_1 \to 0
$$
由于 $\Coker(\alpha_1)$ 平坦，第一列在普遍意义下正合（见《代数》引理
\ref{algebra-lemma-flat-universally-injective}）。
因此 $\alpha \otimes 1_{A'}$ 是单射，当且仅当
$B_1^{\oplus r_1} \otimes_A A' \to N_1 \otimes_A A'$ 是同构。
最后，应用定理 \ref{flat-theorem-flattening-map}，可知该函子由某个理想
$I$ 的商环 $A/I$ 共表示，从而 $F_{lf}$ 也由 $A/I$ 共表示。

\medskip\noindent
为证明最后的断言，假设 $A \to B$ 本质有限表示，且 $M$ 在 $B$ 上有限表示。
令 $I \subset A$ 为使 $F_{lf}$ 由 $A/I$ 共表示的理想。
写成 $I = \bigcup I_\lambda$，其中 $I_\lambda$ 遍历包含于 $I$ 中的有限生成理想。
由于 $F_{lf}(A/I) = \{*\}$，由引理 \ref{flat-lemma-flat-at-point} 第 (2) 部分可知，
对某个 $\lambda$ 有 $F_{lf}(A/I_\lambda) = \{*\}$。
显然这蕴含 $I = I_\lambda$。
\end{proof}

\begin{remark}
\label{flat-remark-flattening-local-scheme-theoretic}
这是定理 \ref{flat-theorem-flattening-local} 的一个概形论重述。
设 $(X, x) \to (S, s)$ 是带点概形的态射，且该态射局部有限型。
设 $\mathcal{F}$ 是有限型拟凝聚 $\mathcal{O}_X$-模。
假设 $S$ 是闭点为 $s$ 的 Hensel 局部概形。
存在闭子概形 $Z \subset S$，具有如下性质：对任意带点概形态射
$(T, t) \to (S, s)$，以下两个条件等价：
\begin{enumerate}
\item 在纤维 $X_t$ 中映到 $x \in X_s$ 的所有点处，$\mathcal{F}_T$ 都在 $T$ 上平坦；
\item $\Spec(\mathcal{O}_{T, t}) \to S$ 经由 $Z$ 分解。
\end{enumerate}
此外，如果 $X \to S$ 在 $x$ 处有限表示，且 $\mathcal{F}_x$
作为 $\mathcal{O}_{X, x}$-模有限表示，那么 $Z \to S$ 有限表示。
\end{remark}

\noindent
现在可以完全免费地从上述结果得到一些非常一般的结果。注意，
最有趣的情形可能是 $E = X_s$！

\begin{lemma}
\label{flat-lemma-freebie}
设 $S$ 是闭点为 $s$ 的 Hensel 局部环的谱。
设 $X \to S$ 是局部有限型的概形态射。
设 $\mathcal{F}$ 是有限型拟凝聚 $\mathcal{O}_X$-模。
设 $E \subset X_s$ 是一个子集。存在闭子概形 $Z \subset S$，具有如下性质：
对任意带点概形态射 $(T, t) \to (S, s)$，以下两个条件等价：
\begin{enumerate}
\item 在纤维 $X_t$ 中映到 $E \subset X_s$ 中某点的所有点处，
$\mathcal{F}_T$ 都在 $T$ 上平坦；
\item $\Spec(\mathcal{O}_{T, t}) \to S$ 经由 $Z$ 分解。
\end{enumerate}
此外，如果 $X \to S$ 局部有限表示，$\mathcal{F}$ 有限表示，且
$E \subset X_s$ 是闭且拟紧的，那么 $Z \to S$ 有限表示。
\end{lemma}

\begin{proof}
对 $x \in X_s$，记在注记
\ref{flat-remark-flattening-local-scheme-theoretic} 中得到的闭子概形为 $Z_x \subset S$。
显然，$Z = \bigcap_{x \in E} Z_x$ 即可！

\medskip\noindent
为证明最后一个断言，假设 $X$ 局部有限表示，$\mathcal{F}$ 有限表示，
且 $Z$ 闭且拟紧。首先选取有限多个仿射开集 $W_j \subset X$，使得
$E \subset \bigcup W_j$。显然，只需对每个态射 $W_j \to S$、层
$\mathcal{F}|_{X_j}$ 以及闭子集 $E \cap W_j$ 证明该结果即可。因此不妨设 $X$ 仿射。
此时，《代数进阶》引理
\ref{more-algebra-lemma-limit-preserving-flat-at-primes} 表明，由 (1) 定义的函子
“保持极限”。于是，完全按照定理 \ref{flat-theorem-flattening-local} 证明的最后一部分，
可以证明 $Z \to S$ 有限表示。
\end{proof}

\begin{remark}
\label{flat-remark-flattening-complete-noetherian}
追溯引理 \ref{flat-lemma-freebie} 的证明直至其根源，会发现这是一条漫长而曲折的道路。
但如果假设
\begin{enumerate}
\item $f$ 有限型；
\item $\mathcal{F}$ 是有限型 $\mathcal{O}_X$-模；
\item $E = X_s$；以及
\item $S$ 是 Noether 完备局部环的谱。
\end{enumerate}
那么可以完全依靠更初等的代数给出如下证明：首先取有限仿射开覆盖，
将情形约化到 $X$ 仿射。在此情形，$Z$ 的存在性由《代数进阶》引理
\ref{more-algebra-lemma-flattening-complete-local-universal-property} 给出。
证明的关键步骤是在截断
$\Spec(\mathcal{O}_{S, s}/\mathfrak m_s^n)$ 中逐步构造闭子概形 $Z$。
这依赖于底为阿廷概形时扁平化分层总是存在这一事实，以及
$\mathcal{O}_{S, s} = \lim \mathcal{O}_{S, s}/\mathfrak m_s^n$ 这一事实。
\end{remark}




\section{一个引理的变体}
\label{flat-section-variants-mod-injective}

\noindent
本节讨论《代数》引理 \ref{algebra-lemma-mod-injective-general} 与
\ref{algebra-lemma-mod-injective} 的变体。
最一般的版本是命题
\ref{flat-proposition-finite-type-injective-into-flat-mod-m}；该命题曾作为
\cite[引理 4.2.2]{GruRay} 陈述，但其中的证明只给出了
引理 \ref{flat-lemma-upstairs-finite-type-injective-into-flat-mod-m} 所陈述的较弱结果。
命题 \ref{flat-proposition-finite-type-injective-into-flat-mod-m} 的精巧证明归功于 Ofer Gabber。
由于目前没有命题的应用，我们建议读者读完引理
\ref{flat-lemma-upstairs-finite-type-injective-into-flat-mod-m} 的证明后跳到下一节；
下一节将用此引理证明定理
\ref{flat-theorem-check-flatness-at-associated-points}。

\begin{situation}
\label{flat-situation-mod-injective}
设 $\varphi : A \to B$ 是局部环之间的局部环同态，且本质有限型。
设 $M$ 是平坦 $A$-模，$N$ 是有限 $B$-模，且
$u : N \to M$ 是一个 $A$-模映射，满足
$\overline{u} : N/\mathfrak m_AN \to M/\mathfrak m_AM$ 是单射。
\end{situation}

\noindent
在此情形，我们的目标是证明 $u$ 是 $A$-普遍单射，
$N$ 作为 $B$-模有限表示，并且 $N$ 作为 $A$-模平坦。
若这些结论成立，就称该情形中{\it 引理成立}。

\begin{lemma}
\label{flat-lemma-Noetherian-finite-type-injective-into-flat-mod-m}
若在情形 \ref{flat-situation-mod-injective} 中环 $A$ 是 Noether 环，则引理成立。
\end{lemma}

\begin{proof}
应用《代数》引理 \ref{algebra-lemma-mod-injective}，可知 $u$ 是单射且
$N/u(M)$ 在 $A$ 上平坦。于是 $u$ 是 $A$-普遍单射（《代数》引理
\ref{algebra-lemma-flat-tor-zero}），并且 $N$ 在 $A$ 上平坦（《代数》引理
\ref{algebra-lemma-flat-ses}）。此时 $B$ 是 Noether 环，故 $N$ 有限表示。
\end{proof}

\begin{lemma}
\label{flat-lemma-reduce-finite-type-injective-into-flat-mod-m}
设 $A_0$ 是局部环。若对每个满足下列条件的情形
\ref{flat-situation-mod-injective}，引理均成立：$A = A_0$，$B$ 是 $A$ 上多项式代数的
一个局部化，且 $N$ 作为 $B$-模有限表示；那么对每个满足 $A = A_0$ 的情形
\ref{flat-situation-mod-injective}，引理都成立。
\end{lemma}

\begin{proof}
设 $A \to B$、$u : N \to M$ 如情形 \ref{flat-situation-mod-injective}。
写成 $B = C/I$，其中 $C$ 是在某个素理想处对 $A$ 上多项式代数作的局部化。
如果能证明 $N$ 作为 $C$-模有限表示，那么当然也能证明 $N$ 作为 $B$-模有限表示（见
《代数》引理 \ref{algebra-lemma-finitely-presented-over-subring}）。
因此不妨设 $B$ 是多项式代数的局部化。
接着写 $N = B^{\oplus n}/K$，其中 $K \subset B^{\oplus n}$ 是某个子模。
由于 $B/\mathfrak m_AB$ 是 Noether 环（它本质有限型于一个域），存在有限多个元素
$k_1, \ldots, k_s \in K$，使得令 $K' = \sum Bk_i$、$N' = B^{\oplus n}/K'$ 时，
典范满射 $N' \to N$ 在商模上诱导同构
$N'/\mathfrak m_AN' \cong N/\mathfrak m_AN$。
现在，若引理对复合映射 $u' : N' \to M$ 成立，则 $u'$ 是单射，
从而 $N' = N$ 且 $u' = u$。因此原情形中的引理也成立。
\end{proof}

\begin{lemma}
\label{flat-lemma-henselian-finite-type-injective-into-flat-mod-m}
若在情形 \ref{flat-situation-mod-injective} 中环 $A$ 是 Hensel 的，
则引理成立。
\end{lemma}

\begin{proof}
只需在 $B$ 本质有限表示于 $A$ 且 $N$ 作为 $B$-模有限表示时证明即可，见
引理 \ref{flat-lemma-reduce-finite-type-injective-into-flat-mod-m}。
暂时再假设 $N$ 在 $A$ 上平坦。于是 $N$ 是自由 $A$-模 $F_i$ 的滤余极限
$N = \colim_i F_i$，其中转移映射
$u_{ii'} : F_i \to F_{i'}$ 模去 $\mathfrak m_A$ 后为单射，见
引理 \ref{flat-lemma-flat-finite-type-local-colimit-free}。
由引理 \ref{flat-lemma-universally-injective-local}，每个复合映射
$u_i : F_i \to M$ 都是 $A$-普遍单射，
因此 $u = \colim u_i$ 如所需是 $A$-普遍单射。

\medskip\noindent
假设 $A$ 是 Hensel 局部环，$B$ 本质有限表示于 $A$，且 $N$ 作为 $B$-模有限表示。
由定理 \ref{flat-theorem-flattening-local}，存在有限生成理想 $I \subset A$，使得
$N/IN$ 在 $A/I$ 上平坦，并且除非 $I = 0$，否则 $N/I^2N$ 不在 $A/I^2$ 上平坦。
上一段的结果表明，引理对 $A/I$ 上的
$u \bmod I : N/IN \to M/IM$ 成立。
考虑交换图
$$
\xymatrix{
0 \ar[r] &
M \otimes_A I/I^2 \ar[r] &
M/I^2M \ar[r] &
M/IM \ar[r] & 0 \\
&
N \otimes_A I/I^2 \ar[r] \ar[u]^u &
N/I^2N \ar[r] \ar[u]^u &
N/IN \ar[r] \ar[u]^u & 0
}
$$
由于 $\otimes$ 的右正合性以及 $M$ 在 $A$ 上平坦，图中的两行正合。
注意左侧竖直箭头是映射
$N/IN \otimes_{A/I} I/I^2 \to M/IM \otimes_{A/I} I/I^2$，因而是单射。
作图追逐可知左下方箭头是单射，即
$\text{Tor}^1_{A/I^2}(I/I^2, M/I^2) = 0$；见
《代数》注记 \ref{algebra-remark-Tor-ring-mod-ideal}。
于是由《代数》引理 \ref{algebra-lemma-what-does-it-mean}，$N/I^2N$ 在 $A/I^2$ 上平坦，
除非 $I = 0$，这便矛盾。
\end{proof}

\noindent
下面的引理讨论情形 \ref{flat-situation-mod-injective} 的特殊情形：$M$ 带有
$B$-模结构且 $u$ 是 $B$-线性的。这是实践中最常用的情形，
证明也比一般情形容易得多。

\begin{lemma}
\label{flat-lemma-upstairs-finite-type-injective-into-flat-mod-m}
设 $A \to B$ 是局部环之间本质有限型的局部环同态，且 $u : N \to M$ 是
$B$-模映射。若 $N$ 是有限 $B$-模，$M$ 在 $A$ 上平坦，并且
$\overline{u} : N/\mathfrak m_A N \to M/\mathfrak m_A M$ 是单射，
那么 $u$ 是 $A$-普遍单射，$N$ 作为 $B$-模有限表示，且 $N$ 在 $A$ 上平坦。
\end{lemma}

\begin{proof}
设 $A \to A^h$ 是 $A$ 的 Hensel 化。令 $B'$ 为
$B \otimes_A A^h$ 在极大理想
$\mathfrak m_B \otimes A^h + B \otimes \mathfrak m_{A^h}$ 处的局部化。
由于 $B \to B'$ 平坦（因而忠实平坦，见《代数》引理
\ref{algebra-lemma-local-flat-ff}），我们可以将 $A \to B$ 替换为
$A^h \to B'$，将模 $M$ 替换为 $M \otimes_B B'$，将模 $N$ 替换为
$N \otimes_B B'$，并将 $u$ 替换为 $u \otimes \text{id}_{B'}$；见《代数》引理
\ref{algebra-lemma-descend-properties-modules} 与
\ref{algebra-lemma-flatness-descends-more-general}。
因此不妨设 $A$ 是 Hensel 局部环。在此情形，我们的引理由更一般的
引理 \ref{flat-lemma-henselian-finite-type-injective-into-flat-mod-m} 推出。
\end{proof}

\begin{lemma}
\label{flat-lemma-valuation-ring-finite-type-injective-into-flat-mod-m}
若在情形 \ref{flat-situation-mod-injective} 中环 $A$ 是赋值环，则引理成立。
\end{lemma}

\begin{proof}
回忆一下，$A$-模平坦当且仅当无挠，见《代数进阶》引理
\ref{more-algebra-lemma-valuation-ring-torsion-free-flat}。
令 $T \subset N$ 为 $A$-挠。则 $u(T) = 0$，且 $N/T$ 在 $A$ 上平坦。
因此 $N/T$ 作为 $B$-模有限表示，见《代数进阶》引理
\ref{more-algebra-lemma-flat-finite-type-valuation-ring-finite-presentation}。
于是 $T$ 是有限 $B$-模，见《代数》引理 \ref{algebra-lemma-extension}。
由于 $N/T$ 在 $A$ 上平坦，我们有
$T/\mathfrak m_A T \subset N/\mathfrak m_A N$，见《代数》引理
\ref{algebra-lemma-flat-tor-zero}。
因为 $\overline{u}$ 是单射而 $u(T) = 0$，可得 $T/\mathfrak m_A T = 0$。
由 Nakayama 引理，$T = 0$，见《代数》引理 \ref{algebra-lemma-NAK}。
至此三个断言中的两个已证得（$N$ 在 $A$ 上平坦且作为 $B$-模有限表示），
剩下的是证明 $u$ 普遍单射。

\medskip\noindent
由《代数》定理 \ref{algebra-theorem-universally-exact-criteria}，只需证明对每个有限表示
$A$-模 $Q$，映射 $N \otimes_A Q \to M \otimes_A Q$ 是单射。由《代数进阶》引理
\ref{more-algebra-lemma-generalized-valuation-ring-modules}，不妨设
$Q = A/(a)$，其中 $a \in \mathfrak m_A$ 非零。因此只需证明
$N/aN \to M/aM$ 是单射。
设 $x \in N$ 且 $u(x) \in aM$。由引理
\ref{flat-lemma-flat-finite-type-local-valuation-ring-has-content}，$x$ 有一个内容理想
$I \subset A$。由于 $I$ 有限生成（《代数进阶》引理
\ref{more-algebra-lemma-content-finitely-generated}），且 $A$ 是赋值环，
有 $I = (b)$，其中某个 $b$（由《代数》引理
\ref{algebra-lemma-characterize-valuation-ring}）。由《代数进阶》引理
\ref{more-algebra-lemma-equal-content}，元素 $u(x)$ 也有内容理想 $I$。
由于 $u(x) \in aM$，由《代数进阶》定义 \ref{more-algebra-definition-content-ideal} 可知
$(b) \subset (a)$。又因 $x \in bN$，遂有 $x \in aN$，即得所需。
\end{proof}

\noindent
考虑以下情形
\begin{equation}
\label{flat-equation-star}
\begin{matrix}
A \to B\text{ 为有限表示，}S \subset B
\text{ 是乘法闭子集，并且}\\
N\text{ 是有限表示的 }S^{-1}B\text{-模}
\end{matrix}
\end{equation}
在此情形，{\it 纯展开}是指一个仿射开集
$U \subset \Spec(B)$，满足 $\Spec(S^{-1}B) \subset U$，以及一个延拓 $N$ 的
有限表示 $\mathcal{O}(U)$-模 $N'$，使得 $N'$ 是 $A$-投射的，且
$N' \to N = S^{-1}N'$ 是 $A$-普遍单射。

\medskip\noindent
在 (
\ref{flat-equation-star}) 中，若 $A \to A_1$ 是环映射，则可以作基变换：取
$B_1 = B \otimes_A A_1$，令 $S_1 \subset B_1$ 为 $S$ 的像，并令
$N_1 = N \otimes_A A_1$。这是可行的，因为
$S_1^{-1}B_1 = S^{-1}B \otimes_A A_1$。下文将不再特别说明而使用这一点。

\begin{lemma}
\label{flat-lemma-properties-pure-spreadout}
在 (\ref{flat-equation-star}) 中，若存在纯展开，则
\begin{enumerate}
\item $N$ 的元素在 $A$ 中都有内容理想；以及
\item 若 $u : N \to M$ 是到平坦 $A$-模 $M$ 的态射，且对 $A$ 的每个极大理想
$\mathfrak m$，映射 $N/\mathfrak m N \to M/\mathfrak m M$ 都是单射，
则 $u$ 是 $A$-普遍单射。
\end{enumerate}
\end{lemma}

\begin{proof}
取纯展开定义中的 $U$、$N'$。任意 $x' \in N'$ 在 $A$ 中都有内容理想，
因为 $N'$ 是 $A$-投射的（这可以直接看出，也可由《代数进阶》引理
\ref{more-algebra-lemma-content-exists-flat-Mittag-Leffler} 与《代数》例
\ref{algebra-example-ML} 推出）。由于 $N' \to N$ 是 $A$-普遍单射，
任意 $x' \in N'$ 的像 $x \in N$ 也在 $A$ 中有内容理想
（它与 $x'$ 的内容理想相同）。对一般的 $x \in N$，选取 $s \in S$ 使得
$s x$ 属于 $N' \to N$ 的像，并利用 $x$ 与 $sx$ 具有相同内容理想这一点。

\medskip\noindent
令 $u : N \to M$ 如 (2) 中所述。为证明 $u$ 是 $A$-普遍单射，
可将 $A$ 替换为某个极大理想处的局部化（略去一个小细节）。
假设 $A$ 是极大理想为 $\mathfrak m$ 的局部环。取 $s \in S$，考虑复合映射
$$
N' \to N \xrightarrow{1/s} N \xrightarrow{u} M
$$
这些映射中的每一个模去 $\mathfrak m$ 后都是单射，因此由引理
\ref{flat-lemma-universally-injective-local}，其复合是 $A$-普遍单射。
由于 $N = \colim_{s \in S} (1/s)N'$，作为普遍单射映射的余极限，$u$
是 $A$-“逆向单射”的。
% P05-EN-CH38-0047: frozen source line 7266 says “inversally injective” after universally-injective component maps; candidate only, do not silently repair.
\end{proof}

\begin{lemma}
\label{flat-lemma-find-pure-spreadout}
在 (\ref{flat-equation-star}) 中，对每个 $\mathfrak p \in \Spec(A)$，
存在有限生成理想 $I \subset \mathfrak pA_\mathfrak p$，使得在 $A_\mathfrak p/I$ 上有纯展开。
\end{lemma}

\begin{proof}
可将 $A$ 替换为 $A_\mathfrak p$。因此不妨设 $A$ 是局部环，且
$\mathfrak p$ 是 $A$ 的极大理想 $\mathfrak m$。通过清除 $S^{-1}B$ 上 $N$ 的表示中的分母，
可写成 $N = S^{-1}N'$，其中 $N'$ 是有限表示 $B$-模。
由于 $B/\mathfrak m B$ 是 Noether 环，映射
$N'/\mathfrak m N' \to N/\mathfrak m N$ 的核 $K$ 有限生成。
于是可取 $s \in S$ 使 $K$ 被 $s$ 消灭。将 $B$ 替换为 $B_s$ 后（这只相当于
经过 $\Spec(B)$ 的一个仿射开子概形，因而是允许的），可知 $S$ 的元素在
$N'/\mathfrak m N'$ 上作用为单射。此时选取局部子环 $A_0 \subset A$，使其本质有限型于
$\mathbf{Z}$；选取有限型环映射 $A_0 \to B_0$，使得
$B = A \otimes_{A_0} B_0$；再选取有限 $B_0$-模 $N'_0$，使得
$N' = B \otimes_{B_0} N'_0 = A \otimes_{A_0} N'_0$。
断言 $I = \mathfrak m_{A_0} A$ 可行。确实有
$$
N'/IN' = N'_0/\mathfrak m_{A_0} N'_0 \otimes_{\kappa_{A_0}} A/I
$$
该模在 $A/I$ 上自由。$S$ 的元素模去极大理想后作用为单射，
因此例如由引理 \ref{flat-lemma-invert-universally-injective}，
$N'/IN' \to N/IN$ 是普遍单射。
\end{proof}

\begin{lemma}
\label{flat-lemma-universally-injective-if-flat}
在 (\ref{flat-equation-star}) 中，假设 $N$ 在 $A$ 上平坦，$M$ 是平坦 $A$-模，
且 $u : N \to M$ 是 $A$-模映射，满足对所有
$\mathfrak p \in \Spec(A)$，$u \otimes \text{id}_{\kappa(\mathfrak p)}$ 都是单射。
则 $u$ 是 $A$-普遍单射。
\end{lemma}

\begin{proof}
由《代数》引理 \ref{algebra-lemma-check-universally-injective-into-flat}，
只需检查对每个理想 $I \subset A$，映射 $N/IN \to M/IM$ 是单射。
将 $A$ 替换为 $A/I$ 后，可见只需证明 $u$ 是单射。

\medskip\noindent
证明 $u$ 为单射。设 $x \in N$ 是 $u$ 核中的非零元素。则模 $Ax$ 有一个弱伴随素理想
$\mathfrak p$，见《代数》引理 \ref{algebra-lemma-weakly-ass-zero}。
将 $A$ 替换为 $A_\mathfrak p$，不妨设 $A$ 局部，并可找到非零元素
$y \in Ax$，使其湮灭子的根基等于 $\mathfrak m_A$，见
《代数》引理 \ref{algebra-lemma-weakly-ass-local}。
因此 $\text{Supp}(y) \subset \Spec(S^{-1}B)$ 非空，且包含于
$\Spec(S^{-1}B) \to \Spec(A)$ 的闭纤维中。
取 $I \subset \mathfrak m_A$ 为有限生成理想，使得在 $A/I$ 上有纯展开，见
引理 \ref{flat-lemma-find-pure-spreadout}。于是对某个 $n$ 有 $I^n y = 0$。由平坦性，
$y \in \text{Ann}_M(I^n) = \text{Ann}_A(I^n) \otimes_R N$。
因此，为得到所需矛盾，只需证明
$$
\text{Ann}_A(I^n) \otimes_R N
\longrightarrow
\text{Ann}_A(I^n) \otimes_R M
$$
是单射。由于 $N$、$M$ 都平坦，且 $\text{Ann}_A(I^n)$ 被 $I^n$ 消灭，
只需证明对每个被 $I$ 消灭的 $A$-模 $Q$，映射
$Q \otimes_A N \to Q \otimes_A M$ 是单射。这由 $I$ 的选取以及
引理 \ref{flat-lemma-properties-pure-spreadout} 的 (2) 得到。
\end{proof}

\begin{lemma}
\label{flat-lemma-big-intersection-is-zero}
设 $A$ 是非域的局部整环。设 $S$ 是 $A$ 的有限生成理想组成的集合。
假设 $S$ 对乘积封闭，且
$\bigcup_{I \in S} V(I)$ 是 $\Spec(A)$ 的泛点的补集。
则 $\bigcap_{I \in S} I = (0)$。
\end{lemma}

\begin{proof}
由于 $\mathfrak m_A \subset A$ 不是 $\Spec(A)$ 的泛点，可知至少有一个 $I \in S$
满足 $I \subset \mathfrak m_A$。故 $\bigcap_{I \in S} I \subset \mathfrak m_A$。
设 $f \in \mathfrak m_A$ 非零。则
$V(f) \subset \bigcup_{I \in S} V(I)$。
由于 $V(f)$ 上的可构造拓扑拟紧（拓扑学引理
\ref{topology-lemma-constructible-hausdorff-quasi-compact} 以及
《代数》引理 \ref{algebra-lemma-spec-spectral}），可得
$V(f) \subset V(I_1) \cup \ldots \cup V(I_n)$，其中某些 $I_j \in S$。
因为 $I_1 \ldots I_n \in S$，可知对某个 $I$ 有 $V(f) \subset V(I)$。
由于 $I$ 有限生成，这意味着对某个 $m$ 有 $I^m \subset (f)$；又因 $S$ 对乘积封闭，
可知对某个 $I \in S$ 有 $I \subset (f^2)$。于是不能有 $f \in I$。
\end{proof}

\begin{lemma}
\label{flat-lemma-closed-points-complement}
设 $A$ 是局部环，设 $I, J \subset A$ 是理想。
若 $J$ 有限生成且对所有 $n \geq 1$ 都有 $I \subset J^n$，
则 $V(I)$ 包含 $\Spec(A) \setminus V(J)$ 的闭点。
\end{lemma}

\begin{proof}
设 $\mathfrak p \subset A$ 是 $\Spec(A) \setminus V(J)$ 的闭点。
我们要证明 $I \subset \mathfrak p$。否则，某个 $f \in I$ 在 $A/\mathfrak p$ 中的像非零。
注意 $V(J) \cap \Spec(A/\mathfrak p)$ 是所有非泛点组成的集合。
因此将引理 \ref{flat-lemma-big-intersection-is-zero} 应用于理想族
$J^nA/\mathfrak p$，可知 $f$ 在 $A/\mathfrak p$ 中的像为零。
\end{proof}

\begin{lemma}
\label{flat-lemma-make-smaller-flatness-ideal}
设 $A$ 是局部环，设 $I \subset A$ 是理想。
设 $U \subset \Spec(A)$ 是拟紧开集。
设 $M$ 是 $A$-模。假设
\begin{enumerate}
\item $M/IM$ 在 $A/I$ 上平坦；
\item $M$ 在 $U$ 上平坦；
\end{enumerate}
则 $M/I_2M$ 在 $A/I_2$ 上平坦，其中
$I_2 = \Ker(I \to \Gamma(U, I/I^2))$。
\end{lemma}

\begin{proof}
只需证明
$M \otimes_A I/I_2 \to IM/I_2M$ 是单射，见《代数》引理
\ref{algebra-lemma-what-does-it-mean-again}。
由假设 (2)，该映射在 $U$ 上是单射。因此只需证明
$M \otimes_A I/I_2$ 到其在 $U$ 上的截面的映射是单射。
我们有 $M \otimes_A I/I_2 = M/IM \otimes_A I/I_2$，且
$M/IM$ 是有限自由 $A/I$-模的滤余极限
（《代数》定理 \ref{algebra-theorem-lazard}）。
所以只需证明 $I/I_2$ 到其在 $U$ 上截面的映射是单射，
这由 $I_2$ 的构造立即得到。
\end{proof}

\begin{proposition}
\label{flat-proposition-finite-type-injective-into-flat-mod-m}
设 $A \to B$ 是局部环之间本质有限型的局部环同态。设 $M$ 是平坦 $A$-模，
$N$ 是有限 $B$-模，且 $u : N \to M$ 是 $A$-模映射，满足
$\overline{u} : N/\mathfrak m_AN \to M/\mathfrak m_AM$ 是单射。
则 $u$ 是 $A$-普遍单射，$N$ 作为 $B$-模有限表示，且 $N$ 在 $A$ 上平坦。
\end{proposition}

\begin{proof}
不妨设 $B$ 是有限表示 $A$-代数 $B_0$ 的局部化，且 $N$ 是有限表示
$B_0$-模 $M_0$ 的局部化，见引理
\ref{flat-lemma-reduce-finite-type-injective-into-flat-mod-m}。
由《态射进阶》引理
\ref{more-morphisms-lemma-generic-flatness-stratification}，在
$\Spec(B_0)$ 上相对于 $\Spec(A)$ 存在 $\widetilde{M_0}$ 的“泛平坦性分层”。
翻译回 $N$，得到一列闭子概形
$$
S = \Spec(A) \supset S_0 \supset S_1 \supset \ldots \supset S_t = \emptyset
$$
其中 $S_i \subset S$ 由 $A$ 的有限生成理想截出，并且
$\widetilde{N}$ 拉回到
$\Spec(B) \times_S (S_i \setminus S_{i + 1})$ 后，在
$S_i \setminus S_{i + 1}$ 上平坦。我们对 $t$ 作归纳来证明命题
（基例 $t = 1$ 将与其他步骤并行证明）。令 $\Spec(A/J_i)$ 为
$S_i \setminus S_{i + 1}$ 的概形论闭包。

\medskip\noindent
{\bf 断言 1。} $N/J_iN$ 在 $A/J_i$ 上平坦。对 $i = t - 1$ 这是直接的，
而对 $i > 0$ 则由归纳假设得到。因此不妨设 $t > 1$、$S_{t - 1} \not = \emptyset$，
且 $J_0 = 0$；我们需要证明 $N$ 平坦。令 $J \subset A$ 为定义 $S_1$ 的理想。
再次对 $t$ 归纳，我们也知道模去 $J$ 的幂后平坦。令 $A^h$ 为 $A$ 的 Hensel 化，
令 $B'$ 为 $B \otimes_A A^h$ 在极大理想
$\mathfrak m_B \otimes A^h + B \otimes \mathfrak m_{A^h}$ 处的局部化。
则 $B \to B'$ 忠实平坦。置 $N' = N \otimes_B B'$。注意，$N'$ 在 $A^h$ 上平坦
当且仅当 $N$ 在 $A$ 上平坦。由定理 \ref{flat-theorem-flattening-local}，存在最小理想
$I \subset A^h$，使得 $N'/IN'$ 在 $A^h/I$ 上平坦，且 $I$ 有限生成。
由上面的结论，对所有 $n \geq 1$ 都有 $I \subset J^nA^h$。令
$S_i^h \subset \Spec(A^h)$ 为 $S_i \subset \Spec(A)$ 的逆像。由引理
\ref{flat-lemma-closed-points-complement}，$V(I)$ 包含
$U = \Spec(A^h) - S_1^h$ 的闭点。按构造，$N'$ 在 $U$ 上对 $A^h$ 平坦。
由引理 \ref{flat-lemma-make-smaller-flatness-ideal}，$N'/I_2N'$ 在 $A/I_2$ 上平坦，其中
$I_2 = \Ker(I \to \Gamma(U, I/I^2))$。由 $I$ 的极小性，$I = I_2$。
这意味着 $I = I^2$ 在 $U$ 上局部成立，即对所有 $u \in U$，有
$I\mathcal{O}_{U, u} = (0)$ 或 $I\mathcal{O}_{U, u} = (1)$。
由于 $V(I)$ 包含 $U$ 的闭点，可知 $I = 0$ 在 $U$ 上成立。
又因 $U \subset \Spec(A^h)$ 概形论稠密（因为本段开头已将 $A$ 替换为 $A/J_0$），
可知 $I = 0$。于是 $N'$ 在 $A^h$ 上平坦，断言 1 得证。

\medskip\noindent
回到断言 1 之前所述的情形。令 $A^h$ 为 $A$ 的 Hensel 化，令 $B'$ 为
$B \otimes_A A^h$ 在极大理想
$\mathfrak m_B \otimes A^h + B \otimes \mathfrak m_{A^h}$ 处的局部化，并令
$N' = N \otimes_B B'$。现在可知，定理 \ref{flat-theorem-flattening-local} 的扁平化理想
$I \subset A^h$ 是幂零的。若 $nil(A^h)$ 表示幂零元素理想，则
$nil(A^h) = nil(A) A^h$
（《代数进阶》引理 \ref{more-algebra-lemma-henselization-nil}）。因此存在有限生成的
幂零理想 $I_0 \subset A$，使得 $N/I_0N$ 在 $A/I_0$ 上平坦。

\medskip\noindent
{\bf 断言 2。} 对每个素理想 $\mathfrak p \subset A$，映射
$\kappa(\mathfrak p) \otimes_A N \to \kappa(\mathfrak p) \otimes_A M$
是单射。若该映射不是单射，则称 $\mathfrak p$ 为坏素理想。
% P05-EN-CH38-0048: frozen source line 7490 says “We say p is bad it this is false”; candidate only, do not silently repair.
设 $C$ 是非空坏素理想链，并令
$\mathfrak p^* = \bigcup_{\mathfrak p \in C} \mathfrak p$.
由引理 \ref{flat-lemma-find-pure-spreadout}，存在有限生成理想
$\mathfrak a \subset \mathfrak p^*A_{\mathfrak p^*}$
使得在 $V(\mathfrak a)$ 上存在纯展开。
若 $\mathfrak p^*$ 是好的，则由
引理 \ref{flat-lemma-properties-pure-spreadout}
可知 $V(\mathfrak a)$ 的点都是好的。
然而，由于 $\mathfrak a$ 有限生成，且
$\mathfrak p^*A_{\mathfrak p^*} = \bigcup_{\mathfrak p \in C}A_{\mathfrak p^*}$
可知 $V(\mathfrak a)$ 包含某个 $\mathfrak p \in C$，矛盾。
所以 $\mathfrak p^*$ 是坏的。由 Zorn 引理，若存在坏素理想，则存在一个极大的坏素理想，
记为 $\mathfrak p$。换言之，不妨设每个
$\mathfrak p' \supset \mathfrak p$ 且 $\mathfrak p' \not = \mathfrak p$ 都是好的。
在此情形，对每个 $f \in A$ 且 $f \not \in \mathfrak p$，映射
$u \otimes \text{id}_{A/(\mathfrak p + f)}$ 都是普遍单射，见引理
\ref{flat-lemma-universally-injective-if-flat}。
因此只需证明 $N/\mathfrak p N$ 关于由子模 $f(N/\mathfrak pN)$ 定义的拓扑是分离的。
由于 $B \to B'$ 忠实平坦，只需对模 $N'/\mathfrak p N'$ 证明同样结论。
由引理 \ref{flat-lemma-flat-finite-type-local-colimit-free} 以及《代数进阶》引理
\ref{more-algebra-lemma-content-exists-flat-Mittag-Leffler}，$N'/\mathfrak pN'$ 的元素
在 $A^h/\mathfrak pA^h$ 中都有内容理想。因此只需证明
$\bigcap_{f \in A, f \not \in \mathfrak p} f(A^h/\mathfrak p A^h) = 0$.
然后对 $A^h/\mathfrak q A^h$ 也证明同样结论，其中每个素理想
$\mathfrak q \subset A^h$ 都在 $\mathfrak p A^h$ 上极小。由于 $A \to A^h$ 是 Hensel 化，
每个 $\mathfrak q$ 都收缩到 $\mathfrak p$，而每个
$\mathfrak q' \supset \mathfrak q$ 且 $\mathfrak q' \not = \mathfrak q$ 都收缩到严格包含
$\mathfrak p$ 的素理想 $\mathfrak p'$。因此由引理
\ref{flat-lemma-big-intersection-is-zero} 得到交的消失。

\medskip\noindent
现在可以把一切合并起来。具体地，利用断言 1 和断言 2，可知
$N/I_0 N \to M/I_0M$ 由引理
\ref{flat-lemma-universally-injective-if-flat} 是 $A/I_0$-普遍单射。
于是下图
$$
\xymatrix{
N \otimes_A (I_0^n/I_0^{n + 1}) \ar[r] \ar[d] &
M \otimes_A (I_0^n/I_0^{n + 1}) \ar@{=}[d] \\
I_0^n N /I_0^{n + 1} N \ar[r] &
I_0^n M /I_0^{n + 1} M
}
$$
表明左侧竖直箭头是单射。因此由《代数》引理
\ref{algebra-lemma-what-does-it-mean-again}，可知 $N$ 平坦。类似地，
$u$ 的普遍单射性可以约化到模去 $I_0$ 的情形（甚至无需先证明 $N$ 平坦）。证明完毕。
\end{proof}


\section{有限型平坦模，第三部分}
\label{flat-section-finite-type-flat-III}

\noindent
下面的结果是本章的主要结果之一。

\begin{theorem}
\label{flat-theorem-check-flatness-at-associated-points}
设 $f : X \to S$ 局部有限型。
设 $\mathcal{F}$ 是有限型拟凝聚 $\mathcal{O}_X$-模。
设 $x \in X$ 的像为 $s \in S$。以下条件等价：
\begin{enumerate}
\item $\mathcal{F}$ 在 $x$ 处关于 $S$ 平坦；
\item 对每个专化到 $x$ 的 $x' \in \text{Ass}_{X_s}(\mathcal{F}_s)$，
$\mathcal{F}$ 在 $x'$ 处关于 $S$ 平坦。
\end{enumerate}
\end{theorem}

\begin{proof}
显然 (1) 蕴含 (2)，因为对每个专化到 $x$ 的点，$\mathcal{F}_{x'}$ 都是
$\mathcal{F}_x$ 的局部化。置 $A = \mathcal{O}_{S, s}$、$B = \mathcal{O}_{X, x}$，以及
$N = \mathcal{F}_x$。令 $\Sigma \subset B$ 为 $B$ 中在 $N/\mathfrak m_AN$ 上
作为非零因子作用的元素组成的乘法子集。由
引理 \ref{flat-lemma-homothety-spectrum} 对 $\Spec(\Sigma^{-1}N)$ 的描述，
假设 (2) 蕴含 $\Sigma^{-1}N$ 在 $A$ 上平坦。
另一方面，按构造映射 $N \to \Sigma^{-1}N$ 模去 $\mathfrak m_A$ 后是单射。
因此应用引理 \ref{flat-lemma-upstairs-finite-type-injective-into-flat-mod-m} 即得结论。
\end{proof}

\noindent
现在直接应用这一点，得到以下有用结果。

\begin{lemma}
\label{flat-lemma-check-along-closed-fibre}
设 $S$ 是闭点为 $s$ 的局部概形。
设 $f : X \to S$ 局部有限型。
设 $\mathcal{F}$ 是有限型拟凝聚 $\mathcal{O}_X$-模。
假设
\begin{enumerate}
\item $\text{Ass}_{X/S}(\mathcal{F})$ 的每个点都专化到闭纤维 $X_s$ 中的某点
\footnote{例如，当 $f$ 有限型且 $\mathcal{F}$ 沿 $X_s$ 纯时，或当 $f$ 适当时，这成立。}；
\item $\mathcal{F}$ 在 $X_s$ 的每个点处关于 $S$ 平坦。
\end{enumerate}
则 $\mathcal{F}$ 在 $S$ 上平坦。
\end{lemma}

\begin{proof}
这是直接的，因为由定理
\ref{flat-theorem-check-flatness-at-associated-points}，只需在 $\mathcal{F}$ 关于 $S$ 的相对伴随点处
检查平坦性即可。
\end{proof}











\section{泛扁平化}
\label{flat-section-flattening-final}

\noindent
若 $f : X \to S$ 是概形的适当有限表示态射，则可以找到 $f$ 的泛扁平化。
本节讨论这一点及其若干变体。

\begin{lemma}
\label{flat-lemma-free-at-generic-points-representable}
在情形 \ref{flat-situation-free-at-generic-points} 中，对每个 $p \geq 0$，函子 $H_p$
（\ref{flat-equation-free-at-generic-points}）由局部闭浸入 $S_p \to S$ 表示。
若 $\mathcal{F}$ 有限表示，则 $S_p \to S$ 有限表示。
\end{lemma}

\begin{proof}
对每个 $S$，我们将同时对所有 $p \geq 0$ 证明该断言。
由引理 \ref{flat-lemma-free-at-generic-points}，函子 $H_p$ 是 fppf 拓扑下的层。
因此结合
《下降》，引理 \ref{descent-lemma-descent-data-sheaves},
《态射进阶》，引理
\ref{more-morphisms-lemma-separated-locally-quasi-finite-morphisms-fppf-descend}
以及
下降理论引理 \ref{descent-lemma-descending-fppf-property-immersion}，可知该问题
关于 $S$ 上的 \'etale 拓扑是局部的。特别地，该问题关于 $S$ 的 Zariski 拓扑是局部的。

\medskip\noindent
对 $s \in S$，记 $\xi_s$ 为纤维 $X_s$ 的唯一泛点。
注意，对每个 $s \in S$，$\mathcal{F}$ 的限制 $\mathcal{F}_s$
在 $\xi_s$ 的某个邻域内局部自由，秩为某个 $p(s) \geq 0$。
（由于 $X_s$ 既约且光滑，这由 $\mathcal{F}_s$ 关于 $X_s$ 的泛平坦性得到，见
《代数》引理 \ref{algebra-lemma-generic-flatness-Noetherian}，尽管这是过强的条件。）
为今后引用，注意
$$
p(s) =
\dim_{\kappa(\xi_s)}(
\mathcal{F}_{\xi_s} \otimes_{\mathcal{O}_{X, \xi_s}} \kappa(\xi_s)
).
$$
特别地，$H_{p(s)}(s)$ 非空，而当 $q \not = p(s)$ 时 $H_q(s)$ 为空。

\medskip\noindent
设 $U \subset X$ 是开子概形。由于 $f : X \to S$ 光滑，故为开态射。
由 (\ref{flat-equation-free-at-generic-points}) 立即可知，对于成对对象
$(f|_U : U \to f(U), \mathcal{F}|_U)$ 的函子 $H_p$ 与对于成对对象
$(f|_{f^{-1}(f(U))}, \mathcal{F}|_{f^{-1}(f(U))})$ 的函子 $H_p$ 相同。
因此，为证明 $S_p$ 在 $f(U)$ 上的存在性，总可以将 $X$ 替换为 $U$。

\medskip\noindent
取 $s \in S$。存在 $\xi_s$ 的仿射开邻域 $U$，使得 $\mathcal{F}|_U$
至多由 $p(s)$ 个元素生成。由上述论证，为了在 $s$ 的一个邻域内证明
$H_{p(s)}$ 的断言，不妨设 $\mathcal{F}$ 由 $p(s)$ 个元素生成，即存在满射
$$
u : \mathcal{O}_X^{\oplus p(s)} \longrightarrow \mathcal{F}
$$
此时显然 $H_{p(s)}$ 等于映射 $u$ 的 $F_{iso}$
（\ref{flat-equation-iso}）（这立即由引理
\ref{flat-lemma-injectivity-map-source-flat-pure} 得到；也可在再缩小一些、使 $S$ 与 $X$ 都仿射后，
由引理 \ref{flat-lemma-induction-step-fp} 得到）。
因此可应用定理 \ref{flat-theorem-flattening-map}，得到 $H_{p(s)}$
在 $s$ 的某个邻域内由闭浸入表示。

\medskip\noindent
结果形式上由上述内容推出。具体地，上述论证表明在 $S$ 的局部，函数
$s \mapsto p(s)$ 有界。因此可对 $p = \max_{s \in S} p(s)$ 作归纳。
由上文，函子 $H_p$ 由闭浸入 $S_p \to S$ 表示。将 $S$ 替换为
$S \setminus S_p$，最大值至少降低一，遂由归纳假设得证。

\medskip\noindent
假设 $\mathcal{F}$ 有限表示。由引理
\ref{flat-lemma-free-at-generic-points} 的 (2) 与极限论注记
\ref{limits-remark-limit-preserving}，$S_p \to S$ 局部有限表示。
然后重做本段中的归纳论证，可知当 $S$ 仿射时每个 $S_p$ 都拟紧：
首先，若 $p = \max_{s \in S} p(s)$，则 $S_p \subset S$ 闭（见上文），因而拟紧。
由于 $S_p \to S$ 是有限表示的闭浸入，$U = S \setminus S_p$ 是 $S$ 中拟紧开集
（例如见态射论第 \ref{morphisms-section-constructible} 节的讨论）。
接着 $S_{p - 1} \to U$ 是有限表示的闭浸入，故 $S_{p - 1}$ 拟紧，且
$U' = S \setminus (S_p \cup S_{p - 1})$ 拟紧。依此类推。
\end{proof}

\begin{lemma}
\label{flat-lemma-localize-flat-dimension-n}
在情形 \ref{flat-situation-flat-dimension-n} 中，设 $h : X' \to X$ 是 \'etale 态射。
置 $\mathcal{F}' = h^*\mathcal{F}$ 及 $f' = f \circ h$。
令 $F_n'$ 为与 $(f' : X' \to S, \mathcal{F}')$ 相关的
(\ref{flat-equation-flat-dimension-n})。则 $F_n$ 是 $F_n'$ 的子函子；若
$h(X') \supset \text{Ass}_{X/S}(\mathcal{F})$，则 $F_n = F'_n$。
\end{lemma}

\begin{proof}
设 $T \to S$ 是任意态射。则 $h_T : X'_T \to X_T$ 是 \'etale 的，因为它是 \'etale 态射 $g$ 的基变换。对 $t \in T$，记
% P05-EN-CH38-0049: frozen source line 7739 names g although this lemma defines h; candidate only, do not silently repair.
$Z \subset X_t$ 为 $\mathcal{F}_T$ 在 $T$ 上不平坦的点集；类似地，记
$Z' \subset X'_t$ 为 $\mathcal{F}'_T$ 在 $T$ 上不平坦的点集。由于
$\mathcal{F}'_T = h_T^*\mathcal{F}_T$，有
$Z' = h_t^{-1}(Z)$，见态射论引理
\ref{morphisms-lemma-flat-permanence}。
因此 $Z' \to Z$ 是 \'etale 态射，从而 $\dim(Z') \leq \dim(Z)$
（例如由下降理论引理
\ref{descent-lemma-dimension-at-point-local}，或直接因为 \'etale 态射是相对维数为 $0$ 的光滑态射）。
这就得到 $F_n \subset F_n'$。


\medskip\noindent
最后，假设 $h(X') \supset \text{Ass}_{X/S}(\mathcal{F})$，且 $T \to S$ 是满足
$F_n'(T)$ 非空的态射，即 $\mathcal{F}'_T$ 在 $T$ 上维数 $\geq n$ 处平坦。
取 $t \in T$，令 $Z \subset X_t$ 和 $Z' \subset X'_t$ 如上。为得到矛盾，假设
$\dim(Z) \geq n$。取泛点 $\xi \in Z$，它对应于一个维数 $\geq n$ 的不可约分支。
令 $x \in \text{Ass}_{X_t}(\mathcal{F}_t)$ 是 $\xi$ 的一般化点。由
除子论引理 \ref{divisors-lemma-base-change-relative-assassin} 及
注记 \ref{divisors-remark-base-change-relative-assassin}，$x$ 映到
$\text{Ass}_{X/S}(\mathcal{F})$ 中的一个点。因此 $x$ 在 $h_T$ 的像中，
即存在 $x' \in X'_T$ 使 $x = h_T(x')$。但因 $x \leadsto \xi$ 且 $\dim(Z') < n$，
$x' \not \in Z'$。于是 $\mathcal{F}'_T$ 在 $x'$ 处关于 $T$ 平坦，
从而 $\mathcal{F}_T$ 在 $x$ 处关于 $T$ 平坦（由态射论引理
\ref{morphisms-lemma-flat-permanence}）。由于这对每个这样的 $x$ 都成立，
由定理 \ref{flat-theorem-check-flatness-at-associated-points}，$\mathcal{F}_T$ 在 $\xi$ 处
关于 $T$ 平坦，得到所需矛盾。
\end{proof}

\begin{lemma}
\label{flat-lemma-compare-H-F}
假设 $X \to S$ 是仿射概形之间的光滑态射，纤维几何不可约且维数为 $d$，并且
$\mathcal{F}$ 是有限表示的拟凝聚 $\mathcal{O}_X$-模。则存在某个 $c \geq 0$，使得
$F_d = \coprod_{p = 0, \ldots, c} H_p$，其中 $F_d$ 如
(\ref{flat-equation-flat-dimension-n})，而 $H_p$ 如
(\ref{flat-equation-free-at-generic-points})。
\end{lemma}

\begin{proof}
由于 $X$ 仿射且 $\mathcal{F}$ 是有限表示的拟凝聚模，$\mathcal{F}$ 可由
$c \geq 0$ 个元素生成。因此对任意 $x \in X$，都有
$\dim_{\kappa(x)}(\mathcal{F}_x \otimes \kappa(x))$ 不超过 $c$。
特别地，当 $p > c$ 时 $H_p = \emptyset$。此外，显然有包含
$\coprod H_p \to F_d$。引理的内容是：若基变换 $\mathcal{F}_T$ 在 $T$ 上维数
$\geq d$ 处平坦，且 $t \in T$，那么 $\mathcal{F}_T$ 在 $X_t$ 的唯一泛点
$\xi$ 的某个开邻域 $U \subset X_T$ 上自由，秩为某个 $r$。也就是说，$H_r$
包含 $U$ 的像，而该像是 $t$ 的一个开邻域。$U$ 的存在性由《态射进阶》引理
\ref{more-morphisms-lemma-flat-and-free-at-point-fibre} 得到。
\end{proof}

\begin{lemma}
\label{flat-lemma-flat-dimension-n-representable}
在情形 \ref{flat-situation-flat-dimension-n} 中，设 $s \in S$，$d \geq 0$。假设
\begin{enumerate}
\item 在某个点 $s \in S$ 上存在 $\mathcal{F}/X/S$ 的完全 d\'evissage；
\item $X$ 在 $S$ 上有限表示；
\item $\mathcal{F}$ 是有限表示的 $\mathcal{O}_X$-模；以及
\item $\mathcal{F}$ 在 $S$ 上维数 $\geq d + 1$ 处平坦。
\end{enumerate}
那么经过可能将 $S$ 替换为 $s$ 的一个开邻域后，函子
$F_d$（\ref{flat-equation-flat-dimension-n}）由有限表示的单态射
$Z_d \to S$ 表示。
\end{lemma}

\begin{proof}
先作一个预备说明：$X$、$S$ 是仿射概形，而且只需在 $S$ 上仿射概形范畴中证明
$F_d$ 由有限表示的单态射 $Z_d \to S$ 表示。（当然不要求 $Z_d$ 仿射。）
因此在本引理的整个证明中，我们都在 $S$ 上的仿射概形范畴中工作。

\medskip\noindent
令 $(Z_k, Y_k, i_k, \pi_k, \mathcal{G}_k, \alpha_k)_{k = 1, \ldots, n}$
为 $s$ 上 $\mathcal{F}/X/S$ 的一个完全 d\'evissage，见定义
\ref{flat-definition-complete-devissage}。我们对 d\'evissage 的长度 $n$ 作归纳。
回忆 $Y_k \to S$ 光滑且纤维几何不可约，见定义
\ref{flat-definition-one-step-devissage}。令 $d_k$ 为 $Y_k$ 关于 $S$ 的相对维数。
回忆 $i_{k, *}\mathcal{G}_k = \Coker(\alpha_k)$，且 $i_k$ 是闭浸入。
由上述定义有
$d_1 = \dim(\text{Supp}(\mathcal{F}_s))$，且
$$
d_k = \dim(\text{Supp}(\Coker(\alpha_{k - 1})_s))
= \dim(\text{Supp}(\mathcal{G}_{k, s}))
$$
对 $k = 2, \ldots, n$ 成立。由于 $\alpha_k$ 在 $(Y_k)_s$ 的泛点处是同构，故
$d_1 > d_2 > \ldots > d_n \geq 0$。

\medskip\noindent
注意 $i_1$ 是闭浸入且 $\mathcal{F} = i_{1, *}\mathcal{G}_1$。
因此对任意 $T \to S$（其中 $T$ 仿射），有
$\mathcal{F}_T = i_{1, T, *}\mathcal{G}_{1, T}$，且 $i_{1, T}$ 仍是
在 $T$ 上的概形之间的闭浸入。于是 $\mathcal{F}_T$ 在 $T$ 上维数 $\geq d$ 处平坦
当且仅当 $\mathcal{G}_{1, T}$ 在 $T$ 上维数 $\geq d$ 处平坦。
由于 $\pi_1 : Z_1 \to Y_1$ 有限，同理可知 $\mathcal{G}_{1, T}$ 在 $T$ 上维数
$\geq d$ 处平坦，当且仅当 $\pi_{1, T, *}\mathcal{G}_{1, T}$ 在 $T$ 上维数
$\geq d$ 处平坦。对
“在维数 $\geq d + 1$ 处平坦”也有同样论证，
由 $\mathcal{F}$ 的假设，特别得到 $\pi_{1, *}\mathcal{G}_1$ 在 $S$ 上维数
$\geq d + 1$ 处平坦。

\medskip\noindent
假设 $d_1 > d$。由上面的讨论，特别可知 $\pi_{1, *}\mathcal{G}_1$ 在 $S$ 上
于 $(Y_1)_s$ 的泛点处平坦。由引理
\ref{flat-lemma-induction-step-fp}，可将 $S$ 替换为 $s$ 的仿射邻域，并假设
$\alpha_1$ 是 $S$-普遍单射。由于 $\alpha_1$ 是 $S$-普遍单射，对任意仿射 $T$ 上的
态射 $T \to S$，有短正合列
$$
0 \to \mathcal{O}_{Y_{1, T}}^{\oplus r_1}
\to \pi_{1, T, *}\mathcal{G}_{1, T} \to \Coker(\alpha_1)_T \to 0
$$
且第一个箭头仍是 $T$-普遍单射。因此 $(Y_1)_T$ 中
$\pi_{1, T, *}\mathcal{G}_{1, T}$ 在 $T$ 上平坦的点集，与 $(Y_1)_T$ 中
$\Coker(\alpha_1)_T$ 在 $S$ 上平坦的点集相同。这样问题约化为层
$\Coker(\alpha_1)$；它有长度为 $n - 1$ 的完全 d\'evissage，故由归纳得证。

\medskip\noindent
若 $d_1 < d$，则 $F_d$ 由 $S$ 表示，结论成立。

\medskip\noindent
最后一种情形是 $d_1 = d$。此情形由引理
\ref{flat-lemma-compare-H-F} 与引理
\ref{flat-lemma-free-at-generic-points-representable} 的组合得到。
\end{proof}

\begin{theorem}
\label{flat-theorem-flat-dimension-n-representable}
在情形 \ref{flat-situation-flat-dimension-n} 中，进一步假设 $f$ 有限表示，
$\mathcal{F}$ 是有限表示的 $\mathcal{O}_X$-模，且 $\mathcal{F}$ 关于 $S$ 纯。
则 $F_n$ 由有限表示的单态射 $Z_n \to S$ 表示。
\end{theorem}

\begin{proof}
由引理 \ref{flat-lemma-flat-dimension-n}，函子 $F_n$ 是 fppf 拓扑下的层。
注意，有限表示的单态射是分离且拟有限的（态射论引理
\ref{morphisms-lemma-monomorphism-loc-finite-type-loc-quasi-finite}）。
因此结合
《下降》，引理 \ref{descent-lemma-descent-data-sheaves},
《态射进阶》，引理
\ref{more-morphisms-lemma-separated-locally-quasi-finite-morphisms-fppf-descend}
以及下降理论引理 \ref{descent-lemma-descending-property-monomorphism} 和
\ref{descent-lemma-descending-property-finite-presentation}，可知该问题关于 $S$ 的
\'etale 拓扑是局部的。

\medskip\noindent
特别地，该情形关于 $S$ 上的 Zariski 拓扑是局部的，因此可设 $S$ 为仿射。
在此情形，$f$ 的纤维维数有上界，所以当 $n$ 足够大时 $F_n$ 可表示。
因此我们可以对 $n$ 使用递降归纳。
设我们已知 $F_{n + 1}$ 由有限表示的单态射
$Z_{n + 1} \to S$ 表示。考虑基变换
$X_{n + 1} = Z_{n + 1} \times_S X$，以及 $\mathcal{F}$ 到
$X_{n + 1}$ 的拉回 $\mathcal{F}_{n + 1}$。
由于 $Z_{n + 1} \to S$ 是有限表示的单态射，故它拟有限；于是由引理
\ref{flat-lemma-quasi-finite-base-change} 可知，$\mathcal{F}_{n + 1}$ 关于
$Z_{n + 1}$ 是纯的。
由于 $F_n$ 是 $F_{n + 1}$ 的子函子，我们得出：要证明 $F_n$ 的结论，
只需证明情形 $\mathcal{F}_{n + 1}/X_{n + 1}/Z_{n + 1}$ 所对应的函子
具有该结论。
这样就把问题化为在 $S_{n + 1} = S$ 时证明 $F_n$ 的结论；也就是说，
我们可以假设 $\mathcal{F}$ 在 $S$ 上的维数 $\geq n + 1$ 的部分平坦。

\medskip\noindent
固定 $n$，并假设 $\mathcal{F}$ 在 $S$ 上的维数 $\geq n + 1$ 的部分平坦。
为了完成证明，我们必须说明 $F_n$ 由有限表示的单态射
$Z_n \to S$ 表示。
由于问题关于 $S$ 的 \'etale 拓扑是局部的，只需证明：对每个 $s \in S$，
存在一个初等 \'etale 邻域 $(S', s') \to (S, s)$，使得基变换到 $S'$ 后
结论成立。
因此由引理 \ref{flat-lemma-existence-complete}，我们可以假设存在
\'etale 态射 $h_j : Y_j \to X$，其中 $j = 1, \ldots, m$，并且对每个
$j$，在 $s$ 上存在 $\mathcal{F}_j/Y_j/S$ 的完全 d\'evissage；
这里 $\mathcal{F}_j$ 是 $\mathcal{F}$ 到 $Y_j$ 的拉回，并且
$X_s \subset \bigcup h_j(Y_j)$。注意，由引理
\ref{flat-lemma-localize-flat-dimension-n}，层 $\mathcal{F}_j$ 在 $S$ 上的
维数 $\geq n + 1$ 的部分仍然平坦。
置 $W = \bigcup h_j(Y_j)$，这是 $X$ 的拟紧开子集。
由于 $\mathcal{F}$ 沿着 $X_s$ 是纯的，我们得到
$$
E = \{t \in S \mid \text{Ass}_{X_t}(\mathcal{F}_t) \subset W \}.
$$
它包含 $s$ 的所有泛化点。由态射论进一步的引理
\ref{more-morphisms-lemma-relative-assassin-constructible}，
$E$ 是 $S$ 的可构造子集。我们已经看到
$\Spec(\mathcal{O}_{S, s}) \subset E$。由态射论引理
\ref{morphisms-lemma-constructible-containing-open}，可知 $E$ 包含 $s$ 的
一个开邻域。因此缩小 $S$ 后可设 $E = S$。
由引理 \ref{flat-lemma-localize-flat-dimension-n} 可知，只需对
$X = \coprod Y_j$ 且 $\mathcal{F} = \coprod \mathcal{F}_j$ 所对应的函子
$F_n$ 证明该引理。
% P05-EN-CH38-0050: source line 7974 uses \mathcal{F}_i although the indexed sheaf is \mathcal{F}_j; preserved literally.
若 $F_{j, n}$ 表示对应于 $Y_j \to S$ 及层 $\mathcal{F}_i$ 的函子，
则 $F_n = \prod F_{j, n}$。因此，只需证明每个 $F_{j, n}$ 都由某个有限表示的
单态射 $Z_{j, n} \to S$ 表示，因为此时
$$
Z_n = Z_{1, n} \times_S \ldots \times_S Z_{m, n}
$$
于是定理化归为引理 \ref{flat-lemma-flat-dimension-n-representable} 所处理的
特殊情形。
\end{proof}

\noindent
下面的引理用泛扁平化明确说明该定理的含义。

\begin{lemma}
\label{flat-lemma-when-universal-flattening}
设 $f : X \to S$ 是概形态射。
设 $\mathcal{F}$ 是拟凝聚的 $\mathcal{O}_X$-模。
\begin{enumerate}
\item 若 $f$ 是有限表示的，$\mathcal{F}$ 是有限表示的
$\mathcal{O}_X$-模，且 $\mathcal{F}$ 关于 $S$ 纯，则存在 $\mathcal{F}$ 的泛扁平化
$S' \to S$。此外，$S' \to S$ 是有限表示的单态射。
\item 若 $f$ 是有限表示的且 $X$ 关于 $S$ 纯，则存在 $X$ 的泛扁平化
$S' \to S$。此外，$S' \to S$ 是有限表示的单态射。
\item 若 $f$ 适当且有限表示，且 $\mathcal{F}$ 是有限表示的
$\mathcal{O}_X$-模，则存在 $\mathcal{F}$ 的泛扁平化 $S' \to S$。
此外，$S' \to S$ 是有限表示的单态射。
\item 若 $f$ 适当且有限表示，则存在 $X$ 的泛扁平化 $S' \to S$。
\end{enumerate}
\end{lemma}

\begin{proof}
这些陈述立即由定理
\ref{flat-theorem-flat-dimension-n-representable} 应用于 $F_0 = F_{flat}$ 得出；
另外，若 $f$ 适当，则 $\mathcal{F}$ 自动相对于基纯，见引理
\ref{flat-lemma-proper-pure}。
\end{proof}







\section{Grothendieck 存在定理 IV}
\label{flat-section-existence}

\noindent
本节继续讨论概形上同调中第
\ref{coherent-section-existence}、\ref{coherent-section-existence-proper} 和
\ref{coherent-section-existence-proper-support} 节的内容。
我们将在以下情形中工作。

\begin{situation}
\label{flat-situation-existence}
这里给定环的逆系统 $(A_n)$，其过渡映射是满射且核局部幂零。
置 $A = \lim A_n$。给定一个在 $A$ 上分离且有限表示的概形 $X$。
置 $X_n = X \times_{\Spec(A)} \Spec(A_n)$，并将其视为 $X$ 的闭子概形。
进一步给定一个系统 $(\mathcal{F}_n, \varphi_n)$，其中 $\mathcal{F}_n$ 是
有限表示的 $\mathcal{O}_{X_n}$-模，在 $A_n$ 上平坦，且其支集在 $A_n$ 上适当，
并且
$$
\varphi_n :
\mathcal{F}_n \otimes_{\mathcal{O}_{X_n}} \mathcal{O}_{X_{n - 1}}
\longrightarrow
\mathcal{F}_{n - 1}
$$
是同构（记号使用态射论引理
\ref{morphisms-lemma-i-star-equivalence} 中的等价）。
\end{situation}

 \noindent
我们的目标是考察能否在 $X$ 上找到拟凝聚层 $\mathcal{F}$，使得
$\mathcal{F}_n = \mathcal{F} \otimes_{\mathcal{O}_X} \mathcal{O}_{X_n}$
对所有 $n$ 都有。

\begin{lemma}
\label{flat-lemma-compute-what-it-should-be}
在情形 \ref{flat-situation-existence} 中，考虑
$$
K = R\lim_{D_\QCoh(\mathcal{O}_X)}(\mathcal{F}_n) =
DQ_X(R\lim_{D(\mathcal{O}_X)}\mathcal{F}_n)
$$
则 $K$ 属于 $D^b_{\QCoh}(\mathcal{O}_X)$；事实上，$K$ 的非零上同调层
只出现在次数 $\geq 0$ 中。
\end{lemma}

\begin{proof}
这是概形导出范畴中例
\ref{perfect-example-inverse-limit-quasi-coherent}.
\end{proof}

\begin{lemma}
\label{flat-lemma-compute-against-perfect}
在情形 \ref{flat-situation-existence} 中，令 $K$ 如引理
\ref{flat-lemma-compute-what-it-should-be} 所述。对 $D(\mathcal{O}_X)$ 的任意完全对象
$E$，有
\begin{enumerate}
\item $M = R\Gamma(X, K \otimes^\mathbf{L} E)$ 是 $D(A)$ 的完美对象，
并且存在规范同构
$R\Gamma(X_n, \mathcal{F}_n \otimes^\mathbf{L} E|_{X_n}) =
M \otimes_A^\mathbf{L} A_n$
在 $D(A_n)$ 中成立，
\item $N = R\Hom_X(E, K)$ 是 $D(A)$ 的完美对象，
并且存在规范同构
$R\Hom_{X_n}(E|_{X_n}, \mathcal{F}_n) = N \otimes_A^\mathbf{L} A_n$
在 $D(A_n)$ 中成立。
\end{enumerate}
在两个陈述中，$E|_{X_n}$ 表示 $E$ 到 $X_n$ 的导出拉回。
\end{lemma}

\begin{proof}
证明 (2)。记 $E_n = E|_{X_n}$，并令
$N_n = R\Hom_{X_n}(E_n, \mathcal{F}_n)$.
回忆 $R\Hom_{X_n}(-, -)$ 等于
$R\Gamma(X_n, R\SheafHom(-, -))$，见概形上同调第
\ref{cohomology-section-global-RHom} 节。因此由概形导出范畴引理
\ref{perfect-lemma-base-change-RHom-perfect}
可知 $N_n$ 是 $D(A_n)$ 的完全对象，且其构造与基变换相容。
因此，由 $\varphi_n$ 给出的映射
$N_n \otimes_{A_n}^\mathbf{L} A_{n - 1} \to N_{n - 1}$
是同构。由《代数进阶》引理
\ref{more-algebra-lemma-Rlim-perfect-gives-perfect}，可知 $R\lim N_n$ 是完全的，
且其基变换回 $A_n$ 恢复 $N_n$。
另一方面，正合函子
$R\Hom_X(E, -) : D_\QCoh(\mathcal{O}_X) \to D(A)$
作为三角范畴之间的函子与乘积交换，因而也与导出极限交换，于是
$$
R\Hom_X(E, K) =
R\lim R\Hom_X(E, \mathcal{F}_n) =
R\lim R\Hom_X(E_n, \mathcal{F}_n) =
R\lim N_n
$$
这证明了 (2)。要看 (1) 成立，使用概形上同调引理
\ref{cohomology-lemma-dual-perfect-complex} 将其转化为 (2)。
\end{proof}

\begin{lemma}
\label{flat-lemma-relative-pseudo-coherence}
在情形 \ref{flat-situation-existence} 中，令 $K$ 如引理
\ref{flat-lemma-compute-what-it-should-be} 所述。则 $K$ 关于 $A$ 拟相干。
\end{lemma}

\begin{proof}
% P05-EN-CH38-0051: source line 8135 spells Combining as “Combinging”; frozen source wording is recorded without a rendered correction.
结合引理 \ref{flat-lemma-compute-against-perfect} 和概形导出范畴引理
\ref{perfect-lemma-perfect-enough}，可知对 $D(\mathcal{O}_X)$ 中的所有拟相干
$E$，$R\Gamma(X, K \otimes^\mathbf{L} E)$ 在 $D(A)$ 中拟相干。
因此由态射论进一步的引理
\ref{more-morphisms-lemma-characterize-pseudo-coherent}.
\end{proof}

\begin{lemma}
\label{flat-lemma-compute-over-affine}
在情形 \ref{flat-situation-existence} 中，令 $K$ 如引理
\ref{flat-lemma-compute-what-it-should-be} 所述。对任意拟紧开集 $U \subset X$，有
$$
R\Gamma(U, K) \otimes_A^\mathbf{L} A_n =
R\Gamma(U_n, \mathcal{F}_n)
$$
在 $D(A_n)$ 中成立，其中 $U_n = U \cap X_n$。
\end{lemma}

\begin{proof}
固定 $n$。由概形导出范畴引理
\ref{perfect-lemma-computing-sections-as-colim}
存在 $X$ 上完全复形的系统 $E_m$，使得
$R\Gamma(U, K) = \text{hocolim} R\Gamma(X, K \otimes^\mathbf{L} E_m)$.
事实上，该公式不仅对 $K$ 成立，而且对
$D_\QCoh(\mathcal{O}_X)$ 的每个对象都成立。将其应用于 $\mathcal{F}_n$，得到
\begin{align*}
R\Gamma(U_n, \mathcal{F}_n)
& =
R\Gamma(U, \mathcal{F}_n) \\
& =
\text{hocolim}_m R\Gamma(X, \mathcal{F}_n \otimes^\mathbf{L} E_m) \\
& =
\text{hocolim}_m R\Gamma(X_n, \mathcal{F}_n \otimes^\mathbf{L} E_m|_{X_n})
\end{align*}
使用引理 \ref{flat-lemma-compute-against-perfect} 以及
$- \otimes_A^\mathbf{L} A_n$ 与同伦余极限交换这一事实，即得结论。
\end{proof}

\begin{lemma}
\label{flat-lemma-finitely-presented}
在情形 \ref{flat-situation-existence} 中，令 $K$ 如引理
\ref{flat-lemma-compute-what-it-should-be} 所述。
记 $X_0 \subset X$ 为由位于 $\Spec(A)$ 的闭子集
$\Spec(A_1) = \Spec(A_2) = \ldots$ 之上的点组成的闭子集。
存在包含 $X_0$ 的开子集 $W \subset X$，使得
\begin{enumerate}
\item 除非 $i = 0$，否则 $H^i(K)|_W$ 为零；
\item $\mathcal{F} = H^0(K)|_W$ 是有限表示的；
\item $\mathcal{F}_n = \mathcal{F} \otimes_{\mathcal{O}_X} \mathcal{O}_{X_n}$。
\end{enumerate}
\end{lemma}

\begin{proof}
固定 $n \geq 1$。按构造，在 $D_\QCoh(\mathcal{O}_X)$ 中存在规范映射
$K \to \mathcal{F}_n$，从而存在拟凝聚层之间的规范映射
$H^0(K) \to \mathcal{F}_n$。这解释了 (3) 的含义。

\medskip\noindent
设 $x \in X_0$ 为一点。我们将找到 $x$ 的一个开邻域 $W$，
使得 (1)、(2) 和 (3) 成立。由于 $X_0$ 拟紧，这将证明引理。
取 $x$ 的一个仿射开邻域 $U \subset X$，记 $U = \Spec(B)$。
取一个满射 $P \to B$，其中 $P$ 在 $A$ 上光滑。
由引理 \ref{flat-lemma-relative-pseudo-coherence} 及相对拟相干的定义，
存在一个由有限自由 $P$-模组成的有界上复形 $F^\bullet$ 表示 $Ri_*K$，
其中 $i : U \to \Spec(P)$ 是该表示诱导的闭浸入。
令 $M_n$ 为与 $\mathcal{F}_n|_U$ 对应的 $B$-模。
由引理 \ref{flat-lemma-compute-over-affine}
$$
H^i(F^\bullet \otimes_A A_n) =
\left\{
\begin{matrix}
0 & \text{若} & i \not = 0 \\
M_n & \text{若} & i = 0
\end{matrix}
\right.
$$
令 $i$ 为使 $F^i$ 非零的最大指标。若 $i \leq 0$，则 (1)、(2) 和 (3) 成立。
否则 $i > 0$，且我们看到映射
$$
F^{i - 1} \to F^i
$$
在点 $x$ 处的秩最大。因此在 $\Spec(P)$ 中 $x$ 的某个开邻域上秩最大。
于是将 $P$ 替换为主元局部化后，可设上面显示的映射是满射。
由于 $F^i$ 有限自由，我们可以选择分裂 $F^{i - 1} = F' \oplus F^i$。
然后可以将 $F^\bullet$ 替换为复形
$$
\ldots \to F^{i - 2} \to F' \to 0 \to \ldots
$$
并对 $i$ 作归纳即可。
\end{proof}

\begin{lemma}
\label{flat-lemma-proper-support}
在情形 \ref{flat-situation-existence} 中，令 $K$ 如引理
\ref{flat-lemma-compute-what-it-should-be} 所述。令 $W \subset X$
如引理 \ref{flat-lemma-finitely-presented} 所述。置
$\mathcal{F} = H^0(K)|_W$。则在必要时缩小开集 $W$ 后，
$\mathcal{F}$ 的支集在 $A$ 上适当。
\end{lemma}

\begin{proof}
固定 $n \geq 1$。令 $I_n = \Ker(A \to A_n)$。
由《代数进阶》引理 \ref{more-algebra-lemma-limit-henselian}，
$(A, I_n)$ 是 Hensel 对。
令 $Z \subset W$ 为 $\mathcal{F}$ 的支集。由于 $\mathcal{F}$ 有限表示，
这是一个闭子集。由引理 \ref{flat-lemma-finitely-presented} 的 (3)，可知
$Z \times_{\Spec(A)} \Spec(A_n)$ 等于 $\mathcal{F}_n$ 的支集，因而在
% P05-EN-CH38-0052: source line 8262 uses A/I although the construction is indexed by I_n; frozen token is retained literally.
$\Spec(A/I)$ 上适当。
由态射论进一步的引理
\ref{more-morphisms-lemma-split-off-proper-part-henselian}
可写成 $Z = Z_1 \amalg Z_2$，其中 $Z_1, Z_2$ 在 $Z$ 中既开又闭，
$Z_1$ 在 $A$ 上适当，并且
$Z_1 \times_{\Spec(A)} \Spec(A/I_n)$ 等于 $\mathcal{F}_n$ 的支集。
换言之，$Z_2$ 不与 $X_0$ 相交。因此将 $W$ 替换为
$W \setminus Z_2$ 后即得引理。
\end{proof}

\begin{lemma}
\label{flat-lemma-monomorphism-isomorphism}
令 $A = \lim A_n$ 为环系统的极限，其过渡映射为满射且核局部幂零。
令 $S = \Spec(A)$。令 $T \to S$ 为局部有限型的单态射。
若对所有 $n$，$\Spec(A_n) \to S$ 都经由 $T$ 分解，则 $T = S$。
\end{lemma}

\begin{proof}
置 $S_n = \Spec(A_n)$。令 $T_0 \subset T$ 为分解 $S_n \to T$ 的公共像。
则 $T_0$ 拟紧。取包含 $T_0$ 的拟紧开集 $T' \subset T$。
于是 $S_n \to T$ 经由 $T'$ 分解。若能证明 $T' = S$，则
$T' = T = S$。因此可设 $T$ 拟紧。

\medskip\noindent
假设 $T$ 拟紧。在此情形，$T \to S$ 分离且拟有限（态射论引理
\ref{morphisms-lemma-monomorphism-loc-finite-type-loc-quasi-finite}).
利用 Zariski 主定理（采用态射论进一步引理
\ref{more-morphisms-lemma-quasi-finite-separated-pass-through-finite})
可选取分解 $T \to W \to S$，其中 $W \to S$ 有限且
$T \to W$ 为开浸入。写作 $W = \Spec(B)$。
分解 $S_n \to T$ 是唯一的，也可视为到 $W$ 的态射，从而得到
$$
A \longrightarrow B \longrightarrow \lim A_n = A
$$
考虑由右侧箭头得到的态射
$h : S = \Spec(A) \to \Spec(B) = W$。则
$$
T \times_{W, h} S
$$
是 $S$ 的开子概形，并包含所有 $n$ 的 $S_n \to S$ 的像。
完成证明只需说明：若开集 $U \subset S$ 包含某个 $n \geq 1$ 的
$S_n \to S$ 的像，则该开集等于 $S$。这是因为
$(A, \Ker(A \to A_n))$ 是 Hensel 对（《代数进阶》引理
\ref{more-algebra-lemma-limit-henselian}），从而 $S$ 的每个闭点都包含在
$S_n \to S$ 的像中。
\end{proof}

\begin{theorem}[Grothendieck 存在定理]
\label{flat-theorem-existence}
在情形 \ref{flat-situation-existence} 中，存在有限表示的
$\mathcal{O}_X$-模 $\mathcal{F}$，其在 $A$ 上平坦且支集在 $A$ 上适当，并满足
$\mathcal{F}_n = \mathcal{F} \otimes_{\mathcal{O}_X} \mathcal{O}_{X_n}$
对所有 $n$ 均成立，且与映射 $\varphi_n$ 相容。
\end{theorem}

\begin{proof}
应用引理 \ref{flat-lemma-compute-what-it-should-be}、
\ref{flat-lemma-compute-against-perfect},
\ref{flat-lemma-relative-pseudo-coherence},
\ref{flat-lemma-compute-over-affine},
\ref{flat-lemma-finitely-presented} 和
\ref{flat-lemma-proper-support}
，得到一个开子概形 $W \subset X$，包含所有位于 $\Spec(A_n)$
之上的点，以及一个有限表示的 $\mathcal{O}_W$-模 $\mathcal{F}$，
其支集在 $A$ 上适当，并且
$\mathcal{F}_n = \mathcal{F} \otimes_{\mathcal{O}_W} \mathcal{O}_{X_n}$
对所有 $n \geq 1$ 均成立。（这是有意义的，因为 $X_n \subset W$。）
由引理 \ref{flat-lemma-proper-pure}，$\mathcal{F}$ 关于 $\Spec(A)$ 普遍纯。
由定理 \ref{flat-theorem-flat-dimension-n-representable}（解释见引理
\ref{flat-lemma-when-universal-flattening}），存在 $\mathcal{F}$ 的泛扁平化
$S' \to \Spec(A)$，且态射 $S' \to \Spec(A)$ 是有限表示的单态射。
由于 $\mathcal{F}$ 到 $\Spec(A_n)$ 的基变换是 $\mathcal{F}_n$，可知对每个
$n$，$\Spec(A_n) \to \Spec(A)$ 唯一地经由 $S'$ 分解。
由引理 \ref{flat-lemma-monomorphism-isomorphism}，可知 $S' = \Spec(A)$。
这意味着 $\mathcal{F}$ 在 $A$ 上平坦。
最后，由于 $\mathcal{F}$ 的概形论支集 $Z$ 在 $\Spec(A)$ 上适当，
态射 $Z \to X$ 是闭的。因此推前 $(W \to X)_*\mathcal{F}$ 支持在 $W$ 上，
并具有所需的全部性质。
\end{proof}






\section{Grothendieck 存在定理 V}
\label{flat-section-existence-derived}

\noindent
本节在导出范畴中证明 Grothendieck 存在定理的一个类比，沿用第
\ref{flat-section-existence} 节处理拟凝聚模的方法。
经典情形在概形上同调的
\ref{coherent-section-existence}、\ref{coherent-section-existence-proper} 和
\ref{coherent-section-existence-proper-support} 节中讨论。
我们将在以下情形中工作。

\begin{situation}
\label{flat-situation-existence-derived}
这里给定环的逆系统 $(A_n)$，其过渡映射为满射且核局部幂零。
置 $A = \lim A_n$。给定一个在 $A$ 上适当、平坦且有限表示的概形 $X$。
置 $X_n = X \times_{\Spec(A)} \Spec(A_n)$，并将其视为 $X$ 的闭子概形。
进一步给定系统 $(K_n, \varphi_n)$，其中 $K_n$ 是
$D(\mathcal{O}_{X_n})$ 的拟相干对象，并且
$$
\varphi_n : K_n \longrightarrow K_{n - 1}
$$
是 $D(\mathcal{O}_{X_n})$ 中的映射，并诱导 $D(\mathcal{O}_{X_{n - 1}})$ 中的同构
$K_n \otimes_{\mathcal{O}_{X_n}}^\mathbf{L} \mathcal{O}_{X_{n - 1}}
\to K_{n - 1}$。
\end{situation}

 \noindent
更准确地说，应写成
$\varphi_n : K_n \to Ri_{n - 1, *}K_{n - 1}$
其中 $i_{n - 1} : X_{n - 1} \to X_n$ 是包含态射；在此记号下，条件是伴随映射
$Li_{n - 1}^*K_n \to K_{n - 1}$ 为同构。
我们的目标是找到拟相干对象 $K \in D(\mathcal{O}_X)$，使得
$K_n = K \otimes_{\mathcal{O}_X}^\mathbf{L} \mathcal{O}_{X_n}$
对所有 $n$ 成立（仍以同样方式滥用记号）。

\begin{lemma}
\label{flat-lemma-compute-what-it-should-be-derived}
在情形 \ref{flat-situation-existence-derived} 中，考虑
$$
K = R\lim_{D_\QCoh(\mathcal{O}_X)}(K_n) =
DQ_X(R\lim_{D(\mathcal{O}_X)} K_n)
$$
则 $K$ 属于 $D^-_{\QCoh}(\mathcal{O}_X)$。
\end{lemma}

\begin{proof}
函子 $DQ_X$ 存在，因为 $X$ 拟紧且拟分离，见概形导出范畴引理
\ref{perfect-lemma-better-coherator}.
由于 $DQ_X$ 是右伴随，它与乘积交换，因而也与导出极限交换。
这就得到引理陈述中的等式。

\medskip\noindent
由概形导出范畴引理 \ref{perfect-lemma-boundedness-better-coherator}，
函子 $DQ_X$ 的上同调维数有界。因此只需证明
$R\lim K_n \in D^-(\mathcal{O}_X)$。
为此，取 $U \subset X$ 为仿射开集。则存在规范正合序列
$$
0 \to
R^1\lim H^{m - 1}(U, K_n) \to H^m(U, R\lim K_n) \to
\lim H^m(U, K_n) \to 0
$$
由概形上同调引理 \ref{cohomology-lemma-RGamma-commutes-with-Rlim} 得到。
由于 $U$ 仿射且 $K_n$ 拟相干（因而由概形导出范畴引理
\ref{perfect-lemma-pseudo-coherent} 具有拟凝聚上同调层），由概形导出范畴引理
\ref{perfect-lemma-affine-compare-bounded} 可知
$H^m(U, K_n) = H^m(K_n)(U)$。
因此只需证明 $K_n$ 的上界与 $n$ 无关。

\medskip\noindent
由于 $K_n$ 拟相干，有 $K_n \in D^-(\mathcal{O}_{X_n})$。
设 $a_n$ 是使 $H^{a_n}(K_n)$ 非零的最大指标。当然
$a_1 \leq a_2 \leq a_3 \leq \ldots$。
注意，$H^{a_n}(K_n)$ 是一个
$\mathcal{O}_{X_n}$-有限表示模
(概形上同调引理 \ref{cohomology-lemma-finite-cohomology})。
我们有 $H^{a_n}(K_{n - 1}) =
H^{a_n}(K_n) \otimes_{\mathcal{O}_{X_n}} \mathcal{O}_{X_{n - 1}}$.
由于 $X_{n - 1} \to X_n$ 是加厚，由 Nakayama 引理（代数引理
\ref{algebra-lemma-NAK}）可知，若
$H^{a_n}(K_n) \otimes_{\mathcal{O}_{X_n}} \mathcal{O}_{X_{n - 1}}$
为零，则 $H^{a_n}(K_n)$ 也为零。因此对所有 $n$ 有
$a_n = a_{n - 1}$，结论得证。
\end{proof}

\begin{lemma}
\label{flat-lemma-compute-against-perfect-derived}
在情形 \ref{flat-situation-existence-derived} 中，令 $K$ 如引理
\ref{flat-lemma-compute-what-it-should-be-derived} 所述。对
$D(\mathcal{O}_X)$ 的任意完全对象 $E$，上同调
$$
M = R\Gamma(X, K \otimes^\mathbf{L} E)
$$
是 $D(A)$ 的拟相干对象，且存在规范同构
$$
R\Gamma(X_n, K_n \otimes^\mathbf{L} E|_{X_n}) = M \otimes_A^\mathbf{L} A_n
$$
在 $D(A_n)$ 中成立。这里 $E|_{X_n}$ 表示 $E$ 到 $X_n$ 的导出拉回。
\end{lemma}

\begin{proof}
记 $E_n = E|_{X_n}$，并令
$M_n = R\Gamma(X_n, K_n \otimes^\mathbf{L} E|_{X_n})$.
由概形导出范畴引理
\ref{perfect-lemma-flat-proper-pseudo-coherent-direct-image-general}
可知 $M_n$ 是 $D(A_n)$ 的拟相干对象，且其构造与基变换相容。
因此，由 $\varphi_n$ 给出的映射
$M_n \otimes_{A_n}^\mathbf{L} A_{n - 1} \to M_{n - 1}$
是同构。由《代数进阶》引理
\ref{more-algebra-lemma-Rlim-pseudo-coherent-gives-pseudo-coherent}
可知 $R\lim M_n$ 拟相干，且其基变换回 $A_n$ 恢复 $M_n$。
另一方面，正合函子
$R\Gamma(X, -) : D_\QCoh(\mathcal{O}_X) \to D(A)$
作为三角范畴之间的函子与乘积交换，因而也与导出极限交换，于是
$$
R\Gamma(X, E \otimes^\mathbf{L} K) =
R\lim R\Gamma(X,  E \otimes^\mathbf{L} K_n) =
R\lim R\Gamma(X_n, E_n \otimes^\mathbf{L} K_n) =
R\lim M_n
$$
这正是所需结论。
\end{proof}

\begin{lemma}
\label{flat-lemma-relative-pseudo-coherence-derived}
在情形 \ref{flat-situation-existence-derived} 中，令 $K$ 如引理
\ref{flat-lemma-compute-what-it-should-be-derived} 所述。则 $K$ 在 $X$ 上拟相干。
\end{lemma}

\begin{proof}
% P05-EN-CH38-0053: source line 8510 spells Combining as “Combinging”; translated by intended meaning.
结合引理 \ref{flat-lemma-compute-against-perfect-derived} 和概形导出范畴引理
\ref{perfect-lemma-perfect-enough}
可知，对 $D(\mathcal{O}_X)$ 中所有拟相干的 $E$，
$R\Gamma(X, K \otimes^\mathbf{L} E)$ 在 $D(A)$ 中拟相干。
因此由态射论进一步的引理
\ref{more-morphisms-lemma-characterize-pseudo-coherent}
可知 $K$ 关于 $A$ 拟相干。
% P05-EN-CH38-0054: source line 8519 has “is of flat”; translated by intended “is flat”.
由于 $X$ 在 $A$ 上平坦且有限表示，这等价于在 $X$ 上拟相干，见
态射论进一步的引理
\ref{more-morphisms-lemma-check-relative-pseudo-coherence-on-charts}.
\end{proof}

\begin{lemma}
\label{flat-lemma-compute-over-affine-derived}
在情形 \ref{flat-situation-existence-derived} 中，令 $K$ 如引理
\ref{flat-lemma-compute-what-it-should-be-derived} 所述。对任意拟紧开集
$U \subset X$，有
$$
R\Gamma(U, K) \otimes_A^\mathbf{L} A_n =
R\Gamma(U_n, K_n)
$$
在 $D(A_n)$ 中成立，其中 $U_n = U \cap X_n$。
\end{lemma}

\begin{proof}
固定 $n$。由概形导出范畴引理
\ref{perfect-lemma-computing-sections-as-colim}
存在 $X$ 上完全复形的系统 $E_m$，使得
$R\Gamma(U, K) = \text{hocolim} R\Gamma(X, K \otimes^\mathbf{L} E_m)$.
事实上，该公式不仅对 $K$ 成立，而且对
$D_\QCoh(\mathcal{O}_X)$ 的每个对象都成立。将其应用于 $K_n$，得到
\begin{align*}
R\Gamma(U_n, K_n)
& =
R\Gamma(U, K_n) \\
& =
\text{hocolim}_m R\Gamma(X, K_n \otimes^\mathbf{L} E_m) \\
& =
\text{hocolim}_m R\Gamma(X_n, K_n \otimes^\mathbf{L} E_m|_{X_n})
\end{align*}
利用引理 \ref{flat-lemma-compute-against-perfect-derived} 以及
$- \otimes_A^\mathbf{L} A_n$ 与同伦余极限交换这一事实，即得结论。
\end{proof}

\begin{theorem}[导出 Grothendieck 存在定理]
\label{flat-theorem-existence-derived}
在情形 \ref{flat-situation-existence-derived} 中，存在
$D(\mathcal{O}_X)$ 中的拟相干对象 $K$，使得
$K_n = K \otimes_{\mathcal{O}_X}^\mathbf{L} \mathcal{O}_{X_n}$
对所有 $n$ 均成立，且与映射 $\varphi_n$ 相容。
\end{theorem}

\begin{proof}
应用引理 \ref{flat-lemma-compute-what-it-should-be-derived}、
\ref{flat-lemma-compute-against-perfect-derived},
\ref{flat-lemma-relative-pseudo-coherence-derived}
，得到 $D(\mathcal{O}_X)$ 的拟相干对象 $K$。
在引理 \ref{flat-lemma-compute-over-affine-derived} 中取仿射开集，
立即可知 $K$ 在 $X_n$ 上的限制为 $K_n$。
\end{proof}

\begin{remark}
\label{flat-remark-correct-generality}
本节的结果可以推广。若只假设 $X \to \Spec(A)$ 分离且有限表示，
$K_n$ 关于 $A_n$ 拟相干，并支持在 $X_n$ 的一个在 $A_n$ 上适当的闭子集上，
结论很可能仍然成立。所得 $K$ 将关于 $A$ 拟相干，并支持在一个
在 $A$ 上适当的闭子集上。若将来需要这一结果，我们会在此给出精确陈述并证明。
\end{remark}






\section{爆破与平坦性}
\label{flat-section-blowup-flat}

\noindent
本节继续讨论如下形式的结果：“爆破后严格变换变得平坦”，见《代数进阶》第
\ref{more-algebra-section-blowup-flat} 节和除子论第
\ref{divisors-section-blowup-flat} 节。
本节采用以下（多少算是标准的）记号。若 $X \to S$ 是概形态射，
$\mathcal{F}$ 是 $X$ 上的拟凝聚模，且 $T \to S$ 是概形态射，
则用 $\mathcal{F}_T$ 表示 $\mathcal{F}$ 到基变换
$X_T = X \times_S T$ 的拉回。

\begin{remark}
\label{flat-remark-successive-blowups}
设 $S$ 为拟紧拟分离概形。设 $f : X \to S$ 为概形态射。
设 $\mathcal{F}$ 为 $X$ 上的拟凝聚模。设 $U \subset S$ 为拟紧开子概形。
给定一个 $U$-容许爆破 $S' \to S$，用 $X'$ 表示 $X$ 的严格变换，
用 $\mathcal{F}'$ 表示 $\mathcal{F}$ 的严格变换，并通过除子论引理
\ref{divisors-lemma-strict-transform} 将其视为 $X'$ 上的拟凝聚模。
设 $P$ 是 $\mathcal{F}/X/S$ 的一个性质，它在 $U$-容许爆破下作上述严格变换时
保持不变。本节的一般问题是：在 $\mathcal{F}/X/S$ 满足辅助条件时，证明存在
$U$-容许爆破 $S' \to S$，使严格变换 $\mathcal{F}'/X'/S'$ 具有性质 $P$。

\medskip\noindent
一般策略是使用如下事实：$U$-容许爆破的复合仍是 $U$-容许爆破，见除子论引理
\ref{divisors-lemma-composition-admissible-blowups}。
实际上，我们将使用更精确的除子论引理
\ref{divisors-lemma-composition-finite-type-blowups}，并结合除子论引理
\ref{divisors-lemma-strict-transform-composition-blowups}。
因此，只需找到一列 $U$-容许爆破
$$
S = S_0 \leftarrow S_1 \leftarrow \ldots \leftarrow S_n
$$
使得置 $\mathcal{F}_0 = \mathcal{F}$、$X_0 = X$，并令
$\mathcal{F}_i/X_i$ 等于 $\mathcal{F}_{i - 1}/X_{i - 1}$ 的严格变换后，
最终得到具有性质 $P$ 的 $\mathcal{F}_n/X_n/S_n$。

\medskip\noindent
特别地，选择有限型拟凝聚理想层
$\mathcal{I} \subset \mathcal{O}_S$，使得 $V(\mathcal{I}) = S \setminus U$，
见性质论引理 \ref{properties-lemma-quasi-coherent-finite-type-ideals}。
令 $S' \to S$ 为沿 $\mathcal{I}$ 的爆破，令 $E \subset S'$ 为例外除子
（除子论引理
\ref{divisors-lemma-blowing-up-gives-effective-Cartier-divisor}）。
于是问题化为存在有效 Cartier 除子 $D \subset S$ 且其支集为
% P05-EN-CH38-0055: source line 8651 uses X \setminus U although D and U lie in S; frozen token retained.
$X \setminus U$ 的情形。特别地，可设 $U$ 在 $S$ 中概形论稠密
（除子论引理 \ref{divisors-lemma-complement-effective-Cartier-divisor}）。

\medskip\noindent
设 $P$ 关于 $S$ 是局部的：若 $S = \bigcup S_i$ 是由拟紧开集组成的有限开覆盖，
且 $P$ 对 $\mathcal{F}_{S_i}/X_{S_i}/S_i$ 成立，则 $P$ 对
$\mathcal{F}/X/S$ 成立。
在此情形，上述一般问题也关于 $S$ 局部；也就是说，若给定 $s \in S$，
可以找到 $s$ 的拟紧开邻域 $W$，使 $\mathcal{F}_W/X_W/W$ 的问题可解，
则 $\mathcal{F}/X/S$ 的问题可解。这由除子论引理
\ref{divisors-lemma-extend-admissible-blowups} 和
\ref{divisors-lemma-dominate-admissible-blowups} 得出。
\end{remark}

\begin{lemma}
\label{flat-lemma-helper-blowup-affine-space}
设 $R$ 为环，$f \in R$。设 $r\geq 0$ 为整数。
设 $R \to S$ 为环同态，$M$ 为 $S$-模。假设
\begin{enumerate}
\item $R \to S$ 有限表示且平坦；
% P05-EN-CH38-0056: source line 8674 says a fibre over kappa(p) is geometrically integral over R; frozen base token retained.
\item 每个纤维环 $S \otimes_R \kappa(\mathfrak p)$ 在 $R$ 上几何整；
\item $M$ 是有限 $S$-模；
\item $M_f$ 是有限表示的 $S_f$-模；
\item 对所有 $\mathfrak p \in R$，若 $f \not \in \mathfrak p$ 且
$\mathfrak q = \mathfrak pS$，则模 $M_{\mathfrak q}$ 是秩为 $r$ 的
自由 $S_\mathfrak q$-模。
\end{enumerate}
则存在有限生成理想 $I \subset R$，满足 $V(f) = V(I)$，并且对所有
$a \in I$，若 $R' = R[\frac{I}{a}]$，则商
$$
M' = (M \otimes_R R')/a\text{-幂挠部分}
$$
作为 $S' = S \otimes_R R'$-模满足如下性质：对每个素理想
$\mathfrak p' \subset R'$，存在 $g \in S'$，使
$g \not \in \mathfrak p'S'$，且 $M'_g$ 是秩为 $r$ 的自由
$S'_g$-模。
\end{lemma}

\begin{proof}
本引理推广了《代数进阶》引理 \ref{more-algebra-lemma-blowup-module}；
我们建议读者先阅读该证明。
% P05-EN-CH38-0057: source line 8697 attributes finite generation of M to (1), while hypothesis (3) supplies it; frozen citation retained.
选择满射 $S^{\oplus n} \to M$，由 (1) 这是可行的。
选择有限子模 $K \subset \Ker(S^{\oplus n} \to M)$，使得
$S^{\oplus n}/K \to M$ 在逆转 $f$ 后成为同构。由 (4) 这是可行的。
置 $M_1 = S^{\oplus n}/K$，并假设可以对 $M_1$ 证明本引理。
设 $I \subset R$ 为相应理想。则对 $a \in I$，映射
$$
M_1' = (M_1 \otimes_R R')/a\text{-幂挠部分}
\longrightarrow
M' = (M \otimes_R R')/a\text{-幂挠部分}
$$
是满射。由于 $R'_a = R_f$，它在 $R'$ 中逆转 $a$ 后也是同构，见代数引理
\ref{algebra-lemma-blowup-in-principal}。
但 $a$ 在 $M'_1$ 上是非零因子，故上面显示的映射是同构。
因此只需在 $M$ 为有限表示的 $S$-模时证明本引理。

\medskip\noindent
假设 $M$ 是满足 (3) 的有限表示 $S$-模。
则 $J = \text{Fit}_r(M) \subset S$ 是有限生成理想。
由引理 \ref{flat-lemma-fibres-irreducible-flat-projective-nonnoetherian}，
可以将 $S$ 写成自由 $R$-模的直和项：
$\bigoplus_{\alpha \in A} R = S \oplus C$。
对任意元素 $h \in S$，在上述分解中写成 $h = \sum a_\alpha$，
称 $a_\alpha$ 为 $h$ 的系数。
令 $I' \subset R$ 为由 $J$ 中元素的系数组成的理想。
乘以 $S$ 的元素定义 $R$-线性映射 $S \to S$，故 $I'$ 由 $J$ 的生成元的
系数生成；也就是说，$I'$ 是有限生成理想。我们断言 $I = fI'$ 可行。

\medskip\noindent
先验证 $V(f) = V(I)$。包含关系 $V(f) \subset V(I)$ 是显然的。
反过来，若 $f \not \in \mathfrak p$，则由性质 (5) 和《代数进阶》引理
\ref{more-algebra-lemma-fitting-ideal-generate-locally}，
$\mathfrak q =  \mathfrak p S$ 不属于 $V(J)$。
因此，$J$ 中存在一个元素，其在 $S \otimes_R \kappa(\mathfrak p)$ 中的像不为零。
所以 $I'$ 中存在不包含于 $\mathfrak p$ 的元素，故
$\mathfrak p \not \in V(fI') = V(I)$。

\medskip\noindent
令 $a \in I$，并置 $R' = R[\frac{I}{a}]$。
对某个 $a' \in I'$，可写成 $a = fa'$。
由代数引理 \ref{algebra-lemma-affine-blowup} 和
\ref{algebra-lemma-blowup-add-principal}，可知 $I' R' = a'R'$，
且 $a'$ 在 $R'$ 中是非零因子。
% P05-EN-CH38-0058: source line 8743 bases S over itself rather than R; frozen tensor base is retained.
置 $S' = S \otimes_S R'$。
$JS' = \text{Fit}_r(M \otimes_S S')$ 的每个元素 $g$ 都可写成
$g = \sum_\alpha c_\alpha$，其中某些 $c_\alpha \in I'R'$。
由于 $I'R' = a'R'$，可对某些 $c'_\alpha \in R'$ 写成
$c_\alpha = a'c'_\alpha$，并且在 $S'$ 中有
$g = (\sum c'_\alpha)a' = g' a'$。
此外，存在 $g_0 \in J$，使某个 $\alpha$ 满足 $a' = c_\alpha$。
对该元素，在 $S'$ 中有 $g_0 = g'_0 a'$，其中 $g'_0$ 是 $S'$ 中的单位。
令 $\mathfrak p' \subset R'$ 为素理想，且
$\mathfrak q' = \mathfrak p'S'$。
由上可知，$JS'_{\mathfrak q'}$ 是由非零因子 $a'$ 生成的主理想。
由《代数进阶》引理
\ref{more-algebra-lemma-principal-fitting-ideal}
可知 $M'_{\mathfrak q'}$ 可由 $r$ 个元素生成。
由于 $M'$ 有限，存在 $m_1, \ldots, m_r \in M'$ 和
$g \in S'$，且 $g \not \in \mathfrak q'$，使得相应映射
$(S')^{\oplus r} \to M'$ 在逆转 $g$ 后成为满射。

\medskip\noindent
最后，考虑理想 $J' = \text{Fit}_{r - 1}(M')$。
注意 $J'S'_g$ 由 $m_1, \ldots, m_r$ 之间关系的系数生成
（Fitting 理想与基变换相容）。因此只需证明 $J' = 0$，见《代数进阶》引理
\ref{more-algebra-lemma-fitting-ideal-finite-locally-free}.
由于 $R'_a = R_f$（代数引理 \ref{algebra-lemma-blowup-in-principal}）
且 $M'_a = M_f$，由 (5) 可知，对任意形如
$\mathfrak q'' = \mathfrak p''S'$ 的素理想 $\mathfrak q'' \subset S'$，
其中 $\mathfrak p'' \subset R'_a$，$J'_a$ 在
$S_{\mathfrak q''}$ 中的像为零。由于
$S'_a \subset \prod_{\mathfrak q''\text{ 如上}} S'_{\mathfrak q''}$
（由引理 \ref{flat-lemma-base-change-universally-flat}，
$(S'_a)_{\mathfrak p''} \subset S'_{\mathfrak q''}$），
% P05-EN-CH38-0059: source line 8777 multiplies the S'-ideal J' by R'_a; frozen expression retained.
可知 $J'R'_a = 0$。由于 $a$ 在 $R'$ 中是非零因子，遂得
$J' = 0$，证明完成。
\end{proof}

\begin{lemma}
\label{flat-lemma-flatten-module-pre}
设 $S$ 为拟紧拟分离概形。设 $X \to S$ 为概形态射。
设 $\mathcal{F}$ 为 $X$ 上的拟凝聚模。设 $U \subset S$ 为拟紧开集。假设
\begin{enumerate}
\item $X \to S$ 仿射、有限表示、平坦，且纤维几何整；
\item $\mathcal{F}$ 是有限型模；
\item $\mathcal{F}_U$ 有限表示；
\item $\mathcal{F}$ 在位于 $U$ 中各点之上的纤维的所有一般点处在 $S$ 上平坦。
\end{enumerate}
则存在 $U$-容许爆破 $S' \to S$ 和开子概形 $V \subset X_{S'}$，使得：
(a) $\mathcal{F}$ 的严格变换 $\mathcal{F}'$ 限制为有限局部自由
$\mathcal{O}_V$-模；(b) $V \to S'$ 为满射。
\end{lemma}

\begin{proof}
给定满足引理假设的 $\mathcal{F}/X/S$ 和 $U \subset S$，以 $P$ 表示性质：
“$\mathcal{F}$ 在纤维的所有一般点处在 $S$ 上平坦”。
显然，$P$ 在严格变换下保持不变，见除子论引理
\ref{divisors-lemma-strict-transform-flat} 和态射论引理
\ref{morphisms-lemma-base-change-module-flat}。同样显然，$P$ 关于 $S$ 是局部的。
因此注 \ref{flat-remark-successive-blowups} 的全部观察均适用于本引理提出的问题。

\medskip\noindent
考虑函数 $r : U \to \mathbf{Z}_{\geq 0}$，它把 $u \in U$ 映为整数
$$
r(u) = \dim_{\kappa(\xi_u)}(\mathcal{F}_{\xi_u} \otimes \kappa(\xi_u))
$$
其中 $\xi_u$ 是纤维 $X_u$ 的一般点。
由态射论进一步的引理
\ref{more-morphisms-lemma-flat-and-free-at-point-fibre}
以及 $X_S$ 中开集在 $S$ 中的像为开集这一事实，可知 $r(u)$ 局部常值。
因此 $U = U_0 \amalg U_1 \amalg \ldots \amalg U_c$ 是开闭子概形的有限不交并，
且 $r$ 在 $U_i$ 上恒等于 $i$。
由除子论引理 \ref{divisors-lemma-separate-disjoint-opens-by-blowing-up}，
可以找到一个 $U$-容许爆破，将 $S$ 分解为两个概形的不交并，其中第一个包含
$U_0$，第二个包含 $U_1 \cup \ldots \cup U_c$。
再重复此过程 $c - 1$ 次，可设 $S$ 是不交并
$S = S_0 \amalg S_1 \amalg \ldots \amalg S_c$，且 $U_i \subset S_i$。
因此可设上述函数 $r$ 为常值，记其值为 $r$。

\medskip\noindent
由注 \ref{flat-remark-successive-blowups}，可设存在有效 Cartier 除子
$D \subset S$，其支集为 $S \setminus U$。
再次应用注 \ref{flat-remark-successive-blowups}，并结合除子论引理
\ref{divisors-lemma-characterize-effective-Cartier-divisor}，可设
$S = \Spec(R)$ 且 $D = \Spec(R/(f))$，其中 $f \in R$ 为非零因子。
此情形由引理 \ref{flat-lemma-helper-blowup-affine-space} 处理。
\end{proof}

\begin{lemma}
\label{flat-lemma-trick-fitting-ideal}
设 $A \to C$ 为秩为 $d$ 的有限局部自由环同态。
设 $h \in C$，使 $C_h$ 在 $A$ 上 \'etale。
设 $J \subset C$ 为理想。置 $I = \text{Fit}_0(C/J)$，其中将
$C/J$ 视为有限 $A$-模。则对某个理想 $J' \subset C_h$，有
$IC_h = JJ'$。若 $J$ 有限生成，则 $I$ 和 $J'$ 也有限生成。
\end{lemma}

\begin{proof}
我们将使用 Fitting 理想的基本性质，见《代数进阶》引理
\ref{more-algebra-lemma-fitting-ideal-basics}。
于是 $IC$ 是 $C/J \otimes_A C$ 的 Fitting 理想。
注意，$C \to C \otimes_A C$，$c \mapsto 1 \otimes c$ 有一个截面
（乘法映射）。按假设，$C \to C \otimes_A C$ 在 $\Spec(C_h)$ 经该截面所得像中
的每个素理想处都是 \'etale 的。
因此乘法映射 $C \otimes_A C_h \to C_h$ 是 \'etale 的，特别地是平坦的，
见代数引理 \ref{algebra-lemma-map-between-etale}。
因此存在一个这样的 $C_h$-代数，使得作为 $C_h$-代数有
$C \otimes_A C_h \cong C_h \oplus C'$，见代数引理
\ref{algebra-lemma-surjective-flat-finitely-presented}。
于是对某个理想 $I' \subset C'$，作为 $C_h$-模有
$(C/J) \otimes_A C_h \cong (C_h/J_h) \oplus C'/I'$。
故 $IC_h = JJ'$，其中 $J' = \text{Fit}_0(C'/I')$，并且
% P05-EN-CH38-0060: source line 8872 switches from C'/I' to C'/J'; frozen quotient retained.
将 $C'/J'$ 视为 $C_h$-模。
\end{proof}

\begin{lemma}
\label{flat-lemma-push-ideal}
设 $A \to B$ 为 \'etale 环同态。设 $a \in A$ 为非零因子。
设 $J \subset B$ 为有限型理想，满足 $V(J) \subset V(aB)$。
对每个 $\mathfrak q \subset B$，存在有限型理想 $I \subset A$，
满足 $V(I) \subset V(a)$，并存在 $g \in B$，$g \not \in \mathfrak q$，
使得对某个有限型理想 $J' \subset B_g$ 有 $IB_g = JJ'$。
\end{lemma}

\begin{proof}
可将 $B$ 替换为在某个元素 $g \in B$ 处的主元局部化，其中
$g \not \in \mathfrak q$。因此可设 $B$ 为标准 \'etale，见代数命题
\ref{algebra-proposition-etale-locally-standard}。
于是可设 $B$ 是 $C = A[x]/(f)$ 的一个局部化，其中
$f \in A[x]$ 为某个次数 $d$ 的首一多项式。
设对某个 $h \in C$ 有 $B = C_h$。选择在 $B$ 上生成 $J$ 的元素
$h_1, \ldots, h_n \in C$。
条件 $V(J) \subset V(aB)$ 意味着对某个充分大的 $m$，在 $B$ 中有
$a^m = \sum b_i h_i$。置 $h_{n + 1} = a^m$。
如引理 \ref{flat-lemma-trick-fitting-ideal} 中一样，取
% P05-EN-CH38-0061: source line 8897 uses r+1 after defining n generators; frozen indices retained.
$I = \text{Fit}_0(C/(h_1, \ldots, h_{r + 1}))$。
由于模 $C/(h_1, \ldots, h_{r + 1})$ 被 $a^m$ 零化，
可知 $a^{dm} \in I$，从而 $V(I) \subset V(a)$。
\end{proof}

\begin{lemma}
\label{flat-lemma-flatten-module-etale-localize}
设 $S$ 为拟紧拟分离概形。设 $X \to S$ 为概形态射。
设 $\mathcal{F}$ 为 $X$ 上的拟凝聚模。设 $U \subset S$ 为拟紧开集。
假设存在有限多个交换图
$$
\xymatrix{
& X_i \ar[r]_{j_i} \ar[d] & X \ar[d] \\
S_i^* \ar[r] & S_i \ar[r]^{e_i} & S
}
$$
其中
\begin{enumerate}
\item $e_i : S_i \to S$ 是拟紧 \'etale 态射，且
$S = \bigcup e_i(S_i)$；
\item $j_i : X_i \to X$ 是 \'etale 态射，且
$X = \bigcup j_i(X_i)$；
\item $S^*_i \to S_i$ 是 $e_i^{-1}(U)$-容许爆破，使
$j_i^*\mathcal{F}$ 的严格变换 $\mathcal{F}_i^*$ 在 $S^*_i$ 上平坦。
\end{enumerate}
则存在 $U$-容许爆破 $S' \to S$，使 $\mathcal{F}$ 的严格变换在 $S'$ 上平坦。
\end{lemma}

\begin{proof}
我们断言，本引理的假设在 $U$-容许爆破下保持不变。
具体地，设 $b : S' \to S$ 是沿拟凝聚理想层 $\mathcal{I}$ 的
$U$-容许爆破。此外，设 $S^*_i \to S_i$ 是沿拟凝聚理想层
$\mathcal{J}_i$ 的爆破。
则对 $\mathcal{F}$ 关于爆破 $S' \to S$ 的严格变换 $\mathcal{F}'$，
态射族 $e'_i : S'_i = S_i \times_S S' \to S'$ 和
$j'_i : X_i' = X_i \times_S S' \to X \times_S S'$ 满足条件 (1)、(2)、(3)。
首先，注意 $S_i'$ 是 $S_i$ 沿 $\mathcal{I}$ 的拉回的爆破，见除子论引理
\ref{divisors-lemma-flat-base-change-blowing-up}。
其次，考虑 $S_i'$ 沿理想 $\mathcal{J}_i$ 的拉回的爆破
$S_i^{\prime *} \to S_i'$。
由除子论引理 \ref{divisors-lemma-blowing-up-two-ideals}，得到交换图
$$
\xymatrix{
S_i^{\prime *} \ar[r] \ar[rd] \ar[d] & S'_i \ar[d] \\
S_i^* \ar[r] & S_i
}
$$
且上图中的所有态射都是爆破。因此由除子论引理
\ref{divisors-lemma-strict-transform-flat} 和
\ref{divisors-lemma-strict-transform-composition-blowups}，可知
% P05-EN-CH38-0062: source lines 8963-8967 use an undefined F_i' and mismatched blowup base; formula retained.
\begin{align*}
& (j'_i)^*\mathcal{F}'\text{ 关于 }
S_i^{\prime *} \to S_i'\text{ 的严格变换} \\
= &
j_i^*\mathcal{F}\text{ 关于 }
S_i^{\prime *} \to S_i\text{ 的严格变换} \\
= &
\mathcal{F}_i'\text{ 关于 }
S_i^{\prime *} \to S_i'\text{ 的严格变换} \\
= &
\mathcal{F}_i^*\text{ 经由 }
X_i \times_{S_i} S_i^{\prime *} \to X_i\text{ 的拉回}
\end{align*}
因此它在 $S_i^{\prime *}$ 上平坦
（态射论引理 \ref{morphisms-lemma-base-change-module-flat}）。
由此可见，注 \ref{flat-remark-successive-blowups} 的全部观察均适用于如下问题：
在本引理的假设下，寻找一个 $U$-容许爆破，使 $\mathcal{F}$ 的严格变换
在基上变得平坦。特别地，可设 $S \setminus U$ 是某个有效 Cartier 除子
$D \subset S$ 的支集。再次应用注 \ref{flat-remark-successive-blowups}，并结合
除子论引理 \ref{divisors-lemma-characterize-effective-Cartier-divisor}，可设
$S = \Spec(A)$ 且 $D = \Spec(A/(a))$，其中 $a \in A$ 为非零因子。

\medskip\noindent
选取一个 $i$ 和 $s \in S_i$。由引理 \ref{flat-lemma-push-ideal}，可以找到开邻域
$s \in W_i \subset S_i$ 和有限型拟凝聚理想
$\mathcal{I} \subset \mathcal{O}_S$，使得对某个有限型拟凝聚理想
$\mathcal{J}'_i \subset \mathcal{O}_{W_i}$ 有
$\mathcal{I} \cdot \mathcal{O}_{W_i} = \mathcal{J}_i \mathcal{J}'_i$，
并且 $V(\mathcal{I}) \subset V(a) = S \setminus U$。
由于 $S_i$ 拟紧，可将 $S_i$ 替换为由这类开集组成的有限族
$W_1, \ldots, W_n$，并假设对每个 $i$ 存在拟凝聚理想层
$\mathcal{I}_i \subset \mathcal{O}_S$，使得对某个有限型拟凝聚理想
$\mathcal{J}'_i \subset \mathcal{O}_{S_i}$ 有
$\mathcal{I}_i \cdot \mathcal{O}_{S_i} = \mathcal{J}_i \mathcal{J}'_i$。
如证明第一段的讨论，考虑 $S$ 沿乘积
$\mathcal{I}_1 \ldots \mathcal{I}_n$ 的爆破 $S'$
（按构造，该爆破是 $U$-容许的）。$S' \to S$ 到 $S_i$ 的基变换是沿
$$
\mathcal{J}_i \cdot
\mathcal{J}'_i \mathcal{I}_1 \ldots \hat{\mathcal{I}_i} \ldots \mathcal{I}_n
$$
的爆破；它经由给定爆破 $S_i^* \to S_i$ 分解
（除子论引理 \ref{divisors-lemma-blowing-up-two-ideals}）。
在上图的记号中，这意味着 $S_i^{\prime *} = S_i'$。
因此以 $S'$ 替换 $S$ 后，达到 $j_i^*\mathcal{F}$ 在 $S_i$ 上平坦的情形。
于是 $j_i^*\mathcal{F}$ 在 $S$ 上平坦，见引理
\ref{flat-lemma-etale-flat-up-down}。由态射论引理
\ref{morphisms-lemma-flat-permanence}，可知 $\mathcal{F}$ 在 $S$ 上平坦。
\end{proof}

\begin{theorem}
\label{flat-theorem-flatten-module}
设 $S$ 为拟紧拟分离概形。设 $X$ 为 $S$ 上的概形。
设 $\mathcal{F}$ 为 $X$ 上的拟凝聚模。设 $U \subset S$ 为拟紧开集。假设
\begin{enumerate}
\item $X$ 拟紧；
\item $X$ 在 $S$ 上局部有限表示；
\item $\mathcal{F}$ 是有限型模；
\item $\mathcal{F}_U$ 有限表示；
\item $\mathcal{F}_U$ 在 $U$ 上平坦。
\end{enumerate}
则存在 $U$-容许爆破 $S' \to S$，使 $\mathcal{F}$ 的严格变换
$\mathcal{F}'$ 是有限表示的 $\mathcal{O}_{X \times_S S'}$-模，
且在 $S'$ 上平坦。
\end{theorem}

\begin{proof}
先证明可以找到一个 $U$-容许爆破，使严格变换平坦。
该问题在源和目标上关于 \'etale 拓扑局部；精确陈述见引理
\ref{flat-lemma-flatten-module-etale-localize}。
特别地，可设 $S = \Spec(R)$ 和 $X = \Spec(A)$ 为仿射。
对 $s \in S$，记 $\mathcal{F}_s = \mathcal{F}|_{X_s}$
（即 $\mathcal{F}$ 到纤维的拉回）。由于 $X \to S$ 有限型，
$d = \max_{s \in S} \dim(\text{Supp}(\mathcal{F}_s))$ 是整数。
我们对 $d$ 作归纳。

\medskip\noindent
设 $x \in X$ 是位于 $s \in S$ 之上的 $X$ 中一点，且
$\dim_x(\text{Supp}(\mathcal{F}_s)) = d$。
应用引理 \ref{flat-lemma-elementary-devissage}，得到
$g : X' \to X$、$e : S' \to S$、$i : Z' \to X'$ 和 $\pi : Z' \to Y'$。
注意 $Y' \to S'$ 是仿射概形之间的光滑态射，其纤维几何不可约且维数为 $d$。
由于问题关于 \'etale 拓扑局部，只需对 $g^*\mathcal{F}/X'/S'$ 证明定理。
由于 $i : Z' \to X'$ 是有限表示的闭浸入（且严格变换与仿射推前交换，见除子论引理
\ref{divisors-lemma-strict-transform-affine}），只需对 $\mathcal{G}$ 证明扁平化结论。
由于 $\pi$ 有限（因而也仿射），只需对
$\pi_*\mathcal{G}/Y'/S'$ 证明扁平化结论。因此可设 $X \to S$
是仿射概形之间的光滑态射，其纤维几何不可约且维数为 $d$。

\medskip\noindent
接着，应用引理 \ref{flat-lemma-flatten-module-pre} 中的爆破。
这样达到如下情形：存在满射到 $S$ 的开集 $V \subset X$，使
$\mathcal{F}|_V$ 有限局部自由。设 $\xi \in X$ 为 $X_s$ 的一般点。令
$r = \dim_{\kappa(\xi)} \mathcal{F}_\xi \otimes \kappa(\xi)$。
选择映射 $\alpha : \mathcal{O}_X^{\oplus r} \to \mathcal{F}$，它诱导同构
$\kappa(\xi)^{\oplus r} \to \mathcal{F}_\xi \otimes \kappa(\xi)$。
由于 $\mathcal{F}$ 在 $V$ 上局部自由，可找到 $\xi$ 的一个开邻域 $W$，
使 $\alpha$ 在其上为同构。将 $S$ 缩小为 $s$ 的一个仿射开邻域，使
$W \to S$ 为满射。设 $\mathcal{F}$ 是由 $A$-模 $N$ 所伴随的拟凝聚模。
由于 $\mathcal{F}$ 在纤维的所有一般点处在 $S$ 上平坦
（事实上在 $W$ 的所有点处如此），可知
$$
\alpha_\mathfrak p : A_\mathfrak p^{\oplus r} \to N_\mathfrak p
$$
对 $R$ 的所有素理想 $\mathfrak p$ 都普遍单射，见引理
\ref{flat-lemma-induction-step}。因此 $\alpha$ 普遍单射，见代数引理
\ref{algebra-lemma-universally-injective-check-stalks}。
置 $\mathcal{H} = \Coker(\alpha)$。
由除子论引理 \ref{divisors-lemma-strict-transform-universally-injective}，
给定 $U$-容许爆破 $S' \to S$ 后，严格变换
$\mathcal{F}'$ 和 $\mathcal{H}'$ 嵌入正合序列
$$
0 \to \mathcal{O}_{X \times_S S'}^{\oplus r} \to \mathcal{F}'
\to \mathcal{H}' \to 0
$$
因此引理 \ref{flat-lemma-induction-step} 还说明：$\mathcal{F}'$ 在一点
$x'$ 处平坦，当且仅当 $\mathcal{H}'$ 在该点处平坦。
特别地，$\mathcal{H}_U$ 在 $U$ 上平坦，且 $\mathcal{H}_U$ 是有限表示的模。
可将归纳假设应用于 $\mathcal{H}$，从而存在一个 $U$-容许爆破，
使严格变换 $\mathcal{H}'$ 如所需般平坦。

\medskip\noindent
为完成定理证明，还必须说明 $\mathcal{F}'$ 是有限表示的模
（可能还要再作一个 $U$-容许爆破）。由于可设 $U \subset S$ 概形论稠密
（见注 \ref{flat-remark-successive-blowups} 第三段），这由引理
\ref{flat-lemma-flat-finite-type-finitely-presented-over-dense-open} 得出。
定理证毕。
\end{proof}





\section{应用}
\label{flat-section-applications}

\noindent
本节应用上面的一些结果。

\begin{lemma}
\label{flat-lemma-flat-after-blowing-up}
设 $S$ 为拟紧拟分离概形。设 $X$ 为 $S$ 上的概形。
设 $U \subset S$ 为拟紧开集。假设
\begin{enumerate}
\item $X \to S$ 有限型且拟分离；
\item $X_U \to U$ 平坦且局部有限表示。
\end{enumerate}
则存在 $U$-容许爆破 $S' \to S$，使 $X$ 的严格变换在 $S'$ 上平坦且有限表示。
\end{lemma}

\begin{proof}
按假设，$X \to S$ 拟紧且拟分离，所以 $X$ 关于爆破 $S' \to S$ 的严格变换
也拟紧且拟分离。因此，为证明引理，只需找到一个 $U$-容许爆破，
使严格变换平坦且局部有限表示。
令 $X = W_1 \cup \ldots \cup W_n$ 为有限仿射开覆盖。
若能找到 $U$-容许爆破 $S_i \to S$，使 $W_i$ 的严格变换平坦且局部有限表示，
则存在支配所有 $S_i \to S$ 的 $U$-容许爆破 $S' \to S$，并满足要求
（见除子论引理 \ref{divisors-lemma-dominate-admissible-blowups}；
另见注 \ref{flat-remark-successive-blowups}）。因此可设 $X$ 仿射。

\medskip\noindent
假设 $X$ 仿射。由态射论引理
\ref{morphisms-lemma-quasi-affine-finite-type-over-S}，可选取 $S$ 上的浸入
$j : X \to \mathbf{A}^n_S$。
取拟紧开子概形 $V \subset \mathbf{A}^n_S$，使 $j$ 诱导 $S$ 上的闭浸入
$i : X \to V$。将定理 \ref{flat-theorem-flatten-module} 应用于
$V \to S$ 和拟凝聚模 $i_*\mathcal{O}_X$，得到 $U$-容许爆破
$S' \to S$，使 $i_*\mathcal{O}_X$ 的严格变换在 $S'$ 上平坦，且作为
$\mathcal{O}_{V \times_S S'}$-模有限表示。
令 $X'$ 为 $X$ 关于 $S' \to S$ 的严格变换，令
$i' : X' \to V \times_S S'$ 为诱导态射。
由于严格变换与沿仿射态射的推前交换（除子论引理
\ref{divisors-lemma-strict-transform-affine}），可知
% P05-EN-CH38-0063: source line 9167 says flat over S after a base change to S'; frozen base retained.
$i'_*\mathcal{O}_{X'}$ 在 $S$ 上平坦，且作为
$\mathcal{O}_{V \times_S S'}$-模有限表示。这推出引理。
\end{proof}

\begin{lemma}
\label{flat-lemma-finite-after-blowing-up}
设 $S$ 为拟紧拟分离概形。设 $X$ 为 $S$ 上的概形。
设 $U \subset S$ 为拟紧开集。假设
\begin{enumerate}
\item $X \to S$ 适当；
\item $X_U \to U$ 有限局部自由。
\end{enumerate}
则存在 $U$-容许爆破 $S' \to S$，使 $X$ 的严格变换在 $S'$ 上有限局部自由。
\end{lemma}

\begin{proof}
由引理 \ref{flat-lemma-flat-after-blowing-up}，可设 $X \to S$ 平坦且有限表示。
必要时以一个 $U$-容许爆破替换 $S$ 后，可设 $U \subset S$ 概形论稠密。
% P05-EN-CH38-0064: source lines 9190-9193 use an undefined f for the morphism X -> S; symbol retained.
于是由引理 \ref{flat-lemma-proper-flat-finite-over-dense-open}，$f$ 有限。
进而由态射论引理 \ref{morphisms-lemma-finite-flat}，$f$ 有限局部自由。
\end{proof}

\begin{lemma}
\label{flat-lemma-zariski-after-blowup}
设 $\varphi : X \to S$ 为有限型分离态射，其中 $S$ 拟紧且拟分离。
设 $U \subset S$ 为拟紧开集，使 $\varphi^{-1}U \to U$ 为同构。
则存在 $U$-容许爆破 $S' \to S$，使 $X$ 的严格变换 $X'$
同构于 $S'$ 的一个开子概形。
\end{lemma}

\begin{proof}
注 \ref{flat-remark-successive-blowups} 的讨论适用。
因此可以先作一次 $U$-容许爆破，并假设补集 $S \setminus U$ 是某个
有效 Cartier 除子 $D$ 的支集。特别地，$U$ 在 $S$ 中概形论稠密。
接着再作一次 $U$-容许爆破，达到 $X \to S$ 平坦且有限表示的情形，
见引理 \ref{flat-lemma-flat-after-blowing-up}。此时结论由引理
\ref{flat-lemma-zariski} 得出。
\end{proof}

\noindent
下一个引理说明，适当修改可由一个爆破支配。

\begin{lemma}
\label{flat-lemma-dominate-modification-by-blowup}
设 $\varphi : X \to S$ 为适当态射，其中 $S$ 拟紧且拟分离。
设 $U \subset S$ 为拟紧开集，使 $\varphi^{-1}U \to U$ 为同构。
则存在支配 $X$ 的 $U$-容许爆破 $S' \to S$；也就是说，爆破态射存在分解
$S' \to X \to S$。
\end{lemma}

\begin{proof}
注 \ref{flat-remark-successive-blowups} 的讨论适用。
因此可以先作一次 $U$-容许爆破，并假设补集 $S \setminus U$ 是某个
有效 Cartier 除子 $D$ 的支集。特别地，$U$ 在 $S$ 中概形论稠密。
选择另一个 $U$-容许爆破 $S' \to S$，使 $X$ 的严格变换 $X'$
是 $S'$ 的开子概形，见引理 \ref{flat-lemma-zariski-after-blowup}。
由于 $X' \to S'$ 适当，且 $U \subset S'$ 稠密，可知 $X' = S'$。
略去一些细节。
\end{proof}

\begin{lemma}
\label{flat-lemma-get-section-after-blowup}
设 $S$ 为概形。设 $U \subset W \subset S$ 为开子概形。
设 $f : X \to W$ 和 $s : U \to X$ 为态射，且
$f \circ s = \text{id}_U$。假设
\begin{enumerate}
\item $f$ 适当；
\item $S$ 拟紧且拟分离；
\item $U$ 和 $W$ 拟紧。
\end{enumerate}
则存在 $U$-容许爆破 $b : S' \to S$，以及延拓 $s$ 的态射
$s' : b^{-1}(W) \to X$，满足 $f \circ s' = b|_{b^{-1}(W)}$。
\end{lemma}

\begin{proof}
我们可以并且确实用 $s$ 的概形论像替换 $X$。
于是 $X \to W$ 在 $U$ 上是同构，见态射论引理
\ref{morphisms-lemma-scheme-theoretic-image-of-partial-section}.
由引理 \ref{flat-lemma-dominate-modification-by-blowup}，存在 $U$-容许爆破
$W' \to W$ 和 $s$ 的延拓 $W' \to X$。
最后应用除子论引理 \ref{divisors-lemma-extend-admissible-blowups}，
将 $W' \to W$ 延拓为 $S$ 的 $U$-容许爆破，证明完成。
\end{proof}









\section{紧化}
\label{flat-section-compactify}

\noindent
设 $S$ 为拟紧拟分离概形。若存在到一个在 $S$ 上适当的概形
$\overline{X}$ 的拟紧开浸入 $X \to \overline{X}$，则称 $S$ 上的概形
$X$ {\it 在 $S$ 上有紧化}，或称 {\it 在 $S$ 上可紧化}。
若 $X$ 在 $S$ 上有紧化，则 $X \to S$ 分离且有限型。
Nagata 的一个定理说明逆命题也成立，见 \cite{Lutkebohmert}、
\cite{Conrad-Nagata}、\cite{Nagata-1}、\cite{Nagata-2}、
\cite{Nagata-3} 和 \cite{Nagata-4}。
我们将在下一节证明该定理，见定理 \ref{flat-theorem-nagata}。

\medskip\noindent
设 $S$ 为拟紧拟分离概形。设 $X \to S$ 为有限型分离态射。
{\it $X$ 在 $S$ 上的紧化范畴} 定义如下：
\begin{enumerate}
\item 对象是 $S$ 上的开浸入 $j : X \to \overline{X}$，其中
$\overline{X} \to S$ 适当。
\item 态射 $(j' : X \to \overline{X}') \to (j : X \to \overline{X})$
是 $S$ 上的概形态射 $f : \overline{X}' \to \overline{X}$，满足
$f \circ j' = j$。
\end{enumerate}
若 $j : X \to \overline{X}$ 是紧化，则 $j$ 是拟紧开浸入，见概形论注
\ref{schemes-remark-quasi-compact-and-quasi-separated}。

\medskip\noindent
{\bf 警告。} 我们{\it 不}假设紧化 $j : X \to \overline{X}$ 的像稠密。
因此，若 $f : \overline{X}' \to \overline{X}$ 是紧化之间的态射，
未必有 $f^{-1}(j(X)) = j'(X)$。

\begin{lemma}
\label{flat-lemma-compactifications-cofiltered}
设 $S$ 为拟紧拟分离概形。设 $X$ 为 $S$ 上可紧化的概形。
\begin{enumerate}
\item[(a)] $X$ 在 $S$ 上的紧化范畴是余滤的。
\item[(b)] 由满足下述条件的紧化 $j : X \to \overline{X}$ 组成的
全子范畴是初始的：$j(X)$ 在 $\overline{X}$ 中稠密且概形论稠密
（范畴论定义 \ref{categories-definition-initial}）。
\item[(c)] 若 $f : \overline{X}' \to \overline{X}$ 是 $X$ 的紧化之间的态射，
且 $j'(X)$ 在 $\overline{X}'$ 中稠密，则
$f^{-1}(j(X)) = j'(X)$。
\end{enumerate}
\end{lemma}

\begin{proof}
为证明 (a)，必须验证范畴论定义
\ref{categories-definition-codirected} 的条件 (1)、(2)、(3)。
条件 (1) 成立，正是因为我们假设 $X$ 可紧化。
设 $j_i : X \to \overline{X}_i$，$i = 1, 2$ 为两个紧化。
考虑 $(j_1, j_2) : X \to \overline{X}_1 \times_S \overline{X}_2$
的概形论像 $\overline{X}$。这确定了第三个紧化
$j : X \to \overline{X}$，它支配两个 $j_i$：
$$
\xymatrix{
(X, \overline{X}_1) & (X, \overline{X}) \ar[l] \ar[r] & (X, \overline{X}_2)
}
$$
因此 (2) 成立。设 $f_1, f_2 : \overline{X}_1 \to \overline{X}_2$
为紧化 $j_i : X \to \overline{X}_i$，$i = 1, 2$ 之间的两个态射。
令 $\overline{X} \subset \overline{X}_1$ 为 $f_1$ 和 $f_2$ 的等化子。
由于 $\overline{X}_2 \to S$ 分离，$\overline{X}$ 是
$\overline{X}_1$ 的闭子概形，因而在 $S$ 上适当。
此外，由于 $f_1|_X = f_2|_X = \text{id}_X$，得到开浸入
$X \to \overline{X}$。由闭浸入
$\overline{X} \to \overline{X}_1$ 给出的态射
$(X \to \overline{X}) \to (j_1 : X \to \overline{X}_1)$
等化 $f_1$ 和 $f_2$，这证明条件 (3)。

\medskip\noindent
证明 (b)。设 $j : X \to \overline{X}$ 为紧化。
若 $\overline{X}'$ 表示 $X$ 在 $\overline{X}$ 中的概形论闭包，
则由态射论引理 \ref{morphisms-lemma-quasi-compact-immersion}，
$X$ 在 $\overline{X}'$ 中稠密且概形论稠密。
这证明了范畴论定义 \ref{categories-definition-initial} 的第一个条件。
由于已证明 $X$ 的紧化范畴余滤，范畴论定义
\ref{categories-definition-initial} 的第二个条件由第一个条件推出
（我们略去这一范畴论练习的解答）。

\medskip\noindent
证明 (c)。用 $j'(X)$ 的概形论闭包替换 $\overline{X}'$
（这不改变底层拓扑空间）后，结论由态射论引理
\ref{morphisms-lemma-scheme-theoretic-image-of-partial-section} 得出。
\end{proof}

\noindent
我们还可以考虑所有紧化的范畴（其中 $X$ 可变）。
事实表明，将此范畴对在内部诱导同构的态射集合做局部化，
所得范畴等价于 $S$ 上可紧化概形的范畴。

\begin{lemma}
\label{flat-lemma-compactifyable}
设 $S$ 为拟紧拟分离概形。设 $f : X \to Y$ 为 $S$ 上的概形态射，
其中 $Y$ 在 $S$ 上分离且有限型，而 $X$ 在 $S$ 上可紧化。
则 $X$ 在 $Y$ 上有紧化。
\end{lemma}

\begin{proof}
设 $j : X \to \overline{X}$ 为 $X$ 在 $S$ 上的紧化。
令 $\overline{X}'$ 为
$(j, f) : X \to \overline{X} \times_S Y$ 的概形论像。
态射 $\overline{X}' \to Y$ 适当，因为
$\overline{X} \times_S Y \to Y$ 是 $\overline{X} \to S$ 的基变换。
另一方面，由于 $Y$ 在 $S$ 上分离，态射
$(1, f) : X \to X \times_S Y$ 是闭浸入
（概形论引理 \ref{schemes-lemma-semi-diagonal}）。
因此，将态射论引理
\ref{morphisms-lemma-scheme-theoretic-image-of-partial-section}
应用于投影 $\overline{X} \times_S Y \to \overline{X}$ 的“部分截面”
$s = (j, f)$，可知 $X \to \overline{X}'$ 是开浸入。
\end{proof}

\noindent
设 $S$ 为拟紧拟分离概形。定义{\it 紧化范畴}如下：
其对象是偶 $(X, \overline{X})$，其中 $\overline{X}$ 是在 $S$ 上适当的概形，
而 $X \subset \overline{X}$ 为拟紧开集；其态射是 $S$ 上概形态射的交换图
$$
\xymatrix{
X \ar[d] \ar[r]_f & Y \ar[d] \\
\overline{X} \ar[r]^{\overline{f}} & \overline{Y}
}
$$

\begin{lemma}
\label{flat-lemma-right-multiplicative-system}
设 $S$ 为拟紧拟分离概形。满足下述条件的态射族
$(u, \overline{u}) : (X', \overline{X}') \to (X, \overline{X})$
构成紧化范畴中箭头的一个右乘法系
（范畴论定义 \ref{categories-definition-multiplicative-system}）：
$u$ 是同构。
\end{lemma}

\begin{proof}
公理 RMS1 的验证是平凡的。下面验证 RMS2。设给定图
$$
\xymatrix{
& (X', \overline{X}') \ar[d]_{(u, \overline{u})} \\
(Y, \overline{Y}) \ar[r]^{(f, \overline{f})} & (X, \overline{X})
}
$$
其中 $u : X' \to X$ 是同构。令 $Y' = Y \times_X X'$，
其投影映射为 $v : Y' \to Y$（这是同构）。
再置 $\overline{Y}' = \overline{Y} \times_{\overline{X}} \overline{X}'$，
其投影映射为 $\overline{v} : \overline{Y}' \to \overline{Y}$。
显然，$Y' \to \overline{Y}'$ 是开浸入。图
$$
\xymatrix{
(Y', \overline{Y}') \ar[r]_{(g, \overline{g})} \ar[d]_{(v, \overline{v})} &
(X', \overline{X}') \ar[d]_{(u, \overline{u})} \\
(Y, \overline{Y}) \ar[r]^{(f, \overline{f})} & (X, \overline{X})
}
$$
说明公理 RMS2 成立。

\medskip\noindent
下面验证 RMS3。设给定紧化之间的一对态射
$(f, \overline{f}), (g, \overline{g}) :
(X, \overline{X}) \to (Y, \overline{Y})$
以及态射
$(v, \overline{v}) : (Y, \overline{Y}) \to (Y', \overline{Y}')$
使得 $v$ 是同构，并且
$(v, \overline{v}) \circ (f, \overline{f}) =
(v, \overline{v}) \circ (g, \overline{g})$。则 $f = g$。
因此，若令 $\overline{X}' \subset \overline{X}$ 为
$\overline{f}$ 和 $\overline{g}$ 的等化子，则
$(u, \overline{u}) : (X, \overline{X}') \to (X, \overline{X})$
是紧化范畴中的态射，并且如所需地满足
$(f, \overline{f}) \circ (u, \overline{u}) =
(g, \overline{g}) \circ (u, \overline{u})$。
\end{proof}

\begin{lemma}
\label{flat-lemma-invert-right-multiplicative-system}
设 $S$ 为拟紧拟分离概形。函子
$(X, \overline{X}) \mapsto X$ 定义了一个等价：从紧化范畴在引理
\ref{flat-lemma-right-multiplicative-system} 的右乘法系处的局部化
（范畴论引理 \ref{categories-lemma-right-localization}），
到 $S$ 上可紧化概形的范畴。
\end{lemma}

\begin{proof}
记 $\mathcal{C}$ 为紧化范畴，记
$Q : \mathcal{C} \to \mathcal{C}'$ 为范畴论引理
\ref{categories-lemma-properties-right-localization} 的局部化函子。
记 $\mathcal{D}$ 为 $S$ 上可紧化概形的范畴。
由刚才引用的引理及我们选择的乘法系，显然得到函子
$\mathcal{C}' \to \mathcal{D}$。该函子显然本质满。
若 $f : X \to Y$ 是可紧化概形之间的态射，则选择到一个在 $S$ 上适当的概形
中的开浸入 $Y \to \overline{Y}$，再选择到一个在 $\overline{Y}$ 上适当的概形
$\overline{X}$ 中的嵌入 $X \to \overline{X}$
（将引理 \ref{flat-lemma-compactifyable} 应用于 $X \to \overline{Y}$ 即可）。
这给出紧化之间的态射
$(X, \overline{X}) \to (Y, \overline{Y})$，它产生给定态射 $X \to Y$。
最后，设局部化范畴中给定一对源和目标相同的态射，记为
$$
a = ((f, \overline{f}) : (X', \overline{X}') \to (Y, \overline{Y}),
(u, \overline{u}) : (X', \overline{X}') \to (X, \overline{X}))
$$
以及
$$
b = ((g, \overline{g}) : (X'', \overline{X}'') \to (Y, \overline{Y}),
(v, \overline{v}) : (X'', \overline{X}'') \to (X, \overline{X}))
$$
它们在 $S$ 上产生同一个态射 $X \to Y$；换言之，
$f \circ u^{-1} = g \circ v^{-1}$。
由范畴论引理 \ref{categories-lemma-morphisms-right-localization}，可设
$(X', \overline{X}') = (X'', \overline{X}'')$ 且
$(u, \overline{u}) = (v, \overline{v})$。
此时可考虑 $\overline{f}$ 和 $\overline{g}$ 的等化子
$\overline{X}''' \subset \overline{X}'$。
态射 $(w, \overline{w}) : (X', \overline{X}''') \to (X', \overline{X}')$
属于该乘法子集；通过预复合 $(w, \overline{w})$，可知局部化范畴中
$a = b$。
\end{proof}






\section{Nagata 紧化}
\label{flat-section-compactifications}

\noindent
本节证明在第 \ref{flat-section-compactify} 节中宣布的定理。

\begin{lemma}
\label{flat-lemma-check-separated}
设 $X \to S$ 是概形态射。若 $X = U \cup V$ 是一个开覆盖，
其中 $U \to S$ 和 $V \to S$ 均为分离态射，且
$U \cap V \to U \times_S V$ 为闭态射，则
$X \to S$ 为分离态射。
\end{lemma}

\begin{proof}
证明从略。提示：利用 $X \times_S X$ 的开覆盖
$U \times_S U$、$U \times_S V$、$V \times_S U$ 和
$V \times_S V$，验证 $\Delta : X \to X \times_S X$ 是闭态射。
\end{proof}

\begin{lemma}
\label{flat-lemma-separate-disjoint-locally-closed-by-blowing-up}
设 $X$ 是拟紧且拟分离概形，$U \subset X$ 是拟紧开子集。
\begin{enumerate}
\item 若 $Z_1, Z_2 \subset X$ 是有限表示的闭子概形，且
$Z_1 \cap Z_2 \cap U = \emptyset$，则存在一个
$U$-容许爆破 $X' \to X$，使 $Z_1$ 与 $Z_2$ 的严格变换不相交。
\item 若 $T_1, T_2 \subset U$ 是互不相交的构造性闭子集，
则存在一个 $U$-容许爆破 $X' \to X$，使 $T_1$ 与 $T_2$
的闭包不相交。
\end{enumerate}
\end{lemma}

\begin{proof}
(1) 的证明。假设 $Z_i \to X$ 为有限表示，意味着 $Z_i$ 的
拟相干理想层 $\mathcal{I}_i$ 为有限型；参见
《态射》，引理 \ref{morphisms-lemma-closed-immersion-finite-presentation}。
令 $Z \subset X$ 为由乘积 $\mathcal{I}_1 \mathcal{I}_2$
切出的闭子概形。注意，$Z \cap U$ 是 $Z_1 \cap U$ 与
$Z_2 \cap U$ 的不交并。由《因子》，引理
\ref{divisors-lemma-separate-disjoint-opens-by-blowing-up}
，存在一个 $U \cap Z$-容许爆破 $Z' \to Z$，使 $Z_1$ 与
$Z_2$ 的严格变换不相交。记 $Y \subset Z$ 为此爆破的中心。
作为 $Y \to Z$ 与 $Z \to X$ 的复合，$Y \to X$
是有限表示的闭浸入（《因子》，定义
\ref{divisors-definition-admissible-blowup}
和 《态射》，引理
\ref{morphisms-lemma-composition-finite-presentation}）。
因此，以 $Y$ 为中心的爆破 $X' \to X$ 是 $U$-容许爆破。
由严格变换的一般性质，$Z_1, Z_2$ 关于 $X' \to X$
的严格变换与 $Z_1, Z_2$ 关于 $Z' \to Z$ 的严格变换相同；参见
《因子》，引理 \ref{divisors-lemma-strict-transform}。
这就证明了 (1)。

\medskip\noindent
(2) 的证明。由《性质》，引理
\ref{properties-lemma-quasi-coherent-finite-type-ideals}
，存在有限型拟相干理想层
$\mathcal{J}_i \subset \mathcal{O}_U$，使
$T_i = V(\mathcal{J}_i)$（作为集合）。由《性质》，引理
\ref{properties-lemma-extend}，存在有限型拟相干理想层
$\mathcal{I}_i \subset \mathcal{O}_X$，其在 $U$ 上的限制为
$\mathcal{J}_i$。将 (1) 的结论应用于闭子概形
$Z_i = V(\mathcal{I}_i)$ 即得结论。
\end{proof}

\begin{lemma}
\label{flat-lemma-blowup-iso-along}
设 $f : X \to Y$ 是拟紧且拟分离概形之间的适当态射。
设 $V \subset Y$ 是拟紧开子集，并令 $U = f^{-1}(V)$。
设 $T \subset V$ 是闭子集，并且 $f|_U : U \to V$
在 $V$ 中 $T$ 的某个开邻域上是同构。则存在一个
$V$-容许爆破 $Y' \to Y$，使 $f$ 的严格变换
$f' : X' \to Y'$ 在 $T$ 于 $Y'$ 中的闭包的某个开邻域上是同构。
\end{lemma}

\begin{proof}
令 $T' \subset V$ 为使 $f|_U$ 成为同构的最大开集之补。
于是 $T', T$ 在 $V$ 中均为闭集，且 $T \cap T' = \emptyset$。
由于 $V$ 是谱拓扑空间，可以找到构造性闭子集 $T_c, T'_c$，
使 $T \subset T_c$、$T' \subset T'_c$ 且
$T_c \cap T'_c = \emptyset$（先取 $V$ 中一个包含 $T'$ 而不与
$T$ 相交的拟紧开集 $W$，并令 $T_c = V \setminus W$；再取
$V$ 中一个包含 $T_c$ 而不与 $T'$ 相交的拟紧开集 $W'$，
并令 $T'_c = V \setminus W'$）。
由引理 \ref{flat-lemma-separate-disjoint-locally-closed-by-blowing-up}，
在用一个 $V$-容许爆破替换 $Y$ 后，可以假设 $T_c$ 与 $T'_c$
在 $Y$ 中的闭包不相交。令
$Y_0 = Y \setminus \overline{T}'_c$、$V_0 = V \setminus T'_c$、
$U_0 = U \times_V V_0$ 以及 $X_0 = X \times_Y Y_0$。
由于 $U_0 \to V_0$ 是同构，由引理
\ref{flat-lemma-zariski-after-blowup}，可以找到一个
$V_0$-容许爆破 $Y'_0 \to Y_0$，使 $X_0$ 的严格变换
$X'_0$ 同构地映到 $Y'_0$。由《因子》，引理
\ref{divisors-lemma-extend-admissible-blowups}，存在一个
$V$-容许爆破 $Y' \to Y$，其在 $Y_0$ 上的限制是
$Y'_0 \to Y_0$。若以 $f' : X' \to Y'$ 表示 $f$ 的严格变换，
则所需结论成立，因为 $f'$ 在 $Y'_0$ 上的限制是同构。
\end{proof}

\begin{lemma}
\label{flat-lemma-find-common-blowups}
设 $S$ 是拟紧且拟分离概形。设 $U \to X_1$ 和
$U \to X_2$ 是 $S$ 上概形的开浸入，并假设
$U$、$X_1$、$X_2$ 在 $S$ 上均为有限型且分离。
则存在 $S$ 上概形的交换图
$$
\xymatrix{
X_1' \ar[d] \ar[r] & X & X_2' \ar[l] \ar[d] \\
X_1 & U \ar[l] \ar[lu] \ar[u] \ar[ru] \ar[r] & X_2
}
$$
其中 $X_i' \to X_i$ 是 $U$-容许爆破，$X_i' \to X$
是开浸入，且 $X$ 在 $S$ 上分离并为有限型。
\end{lemma}

\begin{proof}
在整个证明中，所有概形都在 $S$ 上分离且为有限型。
特别地，这些概形拟紧且拟分离，而它们之间的态射拟紧且分离。
参见《概形》，第 \ref{schemes-section-quasi-compact} 节和第 \ref{schemes-section-separation-axioms} 节。
我们将使用如下事实：若 $U \to W$ 是这类 $S$ 上概形的浸入，
则 $U$ 在 $W$ 中的概形论像 $Z$ 是 $W$ 的闭子概形，
$U \to Z$ 是开浸入，$U \subset Z$ 在概形论意义下稠密，
并且 $U \subset Z$ 在拓扑意义下稠密。参见《态射》，引理
\ref{morphisms-lemma-quasi-compact-immersion}。

\medskip\noindent
令 $X_{12} \subset X_1 \times_S X_2$ 为
$U \to X_1 \times_S X_2$ 的概形论像。由《态射》，引理
\ref{morphisms-lemma-scheme-theoretic-image-of-partial-section}，
投影 $p_i : X_{12} \to X_i$ 诱导同构
$p_i^{-1}(U) \to U$。选取 $U$-容许爆破
$X_i^i \to X_i$，使 $X_{12}$ 的严格变换 $X_{12}^i$
同构于 $X_i^i$ 的一个开子概形；参见引理
\ref{flat-lemma-zariski-after-blowup}。令
$\mathcal{I}_i \subset \mathcal{O}_{X_i}$ 为相应的有限型拟相干理想层。
回顾 $X_{12}^i \to X_{12}$ 是以
$p_i^{-1}\mathcal{I}_i \mathcal{O}_{X_{12}}$ 为中心的爆破；
参见《因子》，引理 \ref{divisors-lemma-strict-transform}。
令 $X_{12}'$ 为 $X_{12}$ 以
$p_1^{-1}\mathcal{I}_1 p_2^{-1}\mathcal{I}_2 \mathcal{O}_{X_{12}}$
为中心的爆破；其含义参见《因子》，引理
\ref{divisors-lemma-blowing-up-two-ideals}。特别地，我们得到交换图
$$
\xymatrix{
X_{12}' \ar[d] \ar[r] & X_{12}^2 \ar[d] \\
X_{12}^1 \ar[r] & X_{12}
}
$$
其中所有态射均为 $U$-容许爆破。由于
$X_{12}^i \subset X_i^i$ 是开子概形，由《因子》，引理
\ref{divisors-lemma-extend-admissible-blowups}，可以选取
$U$-容许爆破 $X_i' \to X_i^i$，其限制为
$X_{12}' \to X_{12}^i$。于是 $X_{12}' \subset X_i'$
是开子概形，且交换图
$$
\xymatrix{
X_{12}' \ar[d] \ar[r] & X_i' \ar[d] \\
X_{12}^i \ar[r] & X_i^i
}
$$
的竖直箭头是爆破，水平箭头是开浸入。注意，
$X'_{12} \to X_1' \times_S X_2'$ 是浸入且适当
（使用 $X'_{12} \to X_{12}$ 为适当态射、
$X_{12} \to X_1 \times_S X_2$ 为闭态射，并且
$X_1' \times_S X_2' \to X_1 \times_S X_2$ 为分离态射，
再应用 《态射》，引理
\ref{morphisms-lemma-image-proper-scheme-closed}）。
因此 $X'_{12} \to  X_1' \times_S X_2'$ 是闭浸入。
沿公共开子概形 $X_{12}'$ 黏合 $X_1'$ 和 $X_2'$ 来定义 $X$，
则 $X \to S$ 为有限型且分离（引理 \ref{flat-lemma-check-separated}）。
由于 $U$-容许爆破的复合仍是 $U$-容许爆破
（《因子》，引理
\ref{divisors-lemma-composition-admissible-blowups}），引理得证。
\end{proof}

\begin{lemma}
\label{flat-lemma-replaced-by-strict-transform}
设 $X \to S$ 和 $Y \to S$ 是概形态射。设 $U \subset X$
是开子概形。设 $V \to X \times_S Y$ 是拟紧态射，且它与第一投影
的复合映入 $U$。设 $Z \subset X \times_S Y$ 是
$V \to X \times_S Y$ 的概形论像，并设
$X' \to X$ 是 $U$-容许爆破。则 $V \to X' \times_S Y$
的概形论像是 $Z$ 关于该爆破的严格变换。
\end{lemma}

\begin{proof}
以 $Z' \to Z$ 表示严格变换。态射 $Z' \to X'$ 诱导态射
$Z' \to X' \times_S Y$，后者是闭浸入
（因为依定义，$Z'$ 是 $X' \times_X Z$ 的闭子概形）。
因此只需证明 $V \to Z'$ 的概形论像 $Z''$ 就是 $Z'$。
注意，$Z'' \subset Z'$ 是闭子概形，且 $V \to Z'$
经 $Z''$ 分解。态射 $V \to X \times_S Y$ 与
$V \to X' \times_S Y$ 均拟紧；对后者，这由《概形》，引理
\ref{schemes-lemma-quasi-compact-permanence}
以及 $X' \times_S Y \to X \times_S Y$ 作为适当态射的基变换为分离态射而得。
故由《态射》，引理
\ref{morphisms-lemma-quasi-compact-scheme-theoretic-image}
可知 $Z \cap (U \times_S Y) = Z'' \cap (U \times_S Y)$。
因此，浸入态射 $Z'' \to Z'$ 在 $Z' \to Z$ 的例外除子
$E$ 之外是同构。然而，$Z'$ 的结构层没有支集在 $E$ 上的非零截面
（这是严格变换的定义），故满射
$\mathcal{O}_{Z'} \to \mathcal{O}_{Z''}$ 必为同构。
\end{proof}

\begin{lemma}
\label{flat-lemma-compactification-dominates}
设 $S$ 是拟紧且拟分离概形。设 $U$ 是在 $S$ 上有限型且分离的概形，
$V \subset U$ 是拟紧开子集。若 $V$ 在 $S$ 上有紧化
$V \subset Y$，则存在一个 $V$-容许爆破 $Y' \to Y$
和开子集 $V \subset V' \subset Y'$，使 $V \to U$
延拓为适当态射 $V' \to U$。
\end{lemma}

\begin{proof}
考虑“对角”态射 $V \to Y \times_S U$ 的概形论像
$Z \subset Y \times_S U$。若用一个 $V$-容许爆破替换 $Y$，
则 $Z$ 被替换为它关于该爆破的严格变换；参见引理
\ref{flat-lemma-replaced-by-strict-transform}。因此由引理
\ref{flat-lemma-zariski-after-blowup}，可以假设 $Z \to Y$ 是开浸入。
若 $V' \subset Y$ 表示其像，则诱导态射 $V' \to U$ 是适当态射：
投影 $Y \times_S U \to U$ 是适当态射，而
$V' \cong Z$ 是 $Y \times_S U$ 的闭子概形。
\end{proof}

\noindent
下述引理仅在 Noether 情形下表述。拟紧且拟分离概形的版本也成立，
但它将直接由本节的主定理推出。

\begin{lemma}
\label{flat-lemma-two-compactifications}
设 $S$ 是 Noether 概形，$U$ 是在 $S$ 上有限型且分离的概形。
设 $U = U_1 \cup U_2$ 是开覆盖，$U_1$ 与 $U_2$ 在 $S$ 上均有紧化，
并且 $U_1 \cap U_2$ 在 $U$ 中稠密。则 $U$ 在 $S$ 上有紧化。
\end{lemma}

\begin{proof}
对 $i = 1, 2$，选取紧化 $U_i \subset X_i$。可以假设
$U_i$ 在 $X_i$ 中概形论稠密。由引理
\ref{flat-lemma-compactification-dominates}，还可以假设存在开子集
$V_i \subset X_i$ 和适当态射 $\psi_i : V_i \to U$，
它延拓 $\text{id} : U_i \to U_i$。图示为
$$
\xymatrix{
U_i \ar[r] \ar[d] & V_i \ar[r] \ar[dl]^{\psi_i} & X_i \\
U
}
$$
若 $\{i, j\} = \{1, 2\}$，记
$Z_i = U \setminus U_j = U_i \setminus (U_1 \cap U_2)$
以及
$Z_j = U \setminus U_i = U_j \setminus (U_1 \cap U_2)$.
于是
$$
U = U_1 \amalg Z_2 = Z_1 \amalg U_2 = Z_1 \amalg (U_1 \cap U_2) \amalg Z_2
$$
记 $Z_{i, i} \subset V_i$ 为 $Z_i$ 在 $\psi_i$ 下的逆像。
注意，$\psi_i$ 在 $Z_i$ 的某个开邻域上是同构。
记 $Z_{i, j} \subset V_i$ 为 $Z_j$ 在 $\psi_i$ 下的逆像。
注意，$\psi_i : Z_{i, j} \to Z_j$ 是适当态射。
由于 $Z_i$ 与 $Z_j$ 是 $U$ 中互不相交的闭子集，
$Z_{i, i}$ 与 $Z_{i, j}$ 是 $V_i$ 中互不相交的闭子集。

\medskip\noindent
以 $\overline{Z}_{i, i}$ 和 $\overline{Z}_{i, j}$ 表示
$Z_{i, i}$ 和 $Z_{i, j}$ 在 $X_i$ 中的闭包。
用一个 $V_i$-容许爆破替换 $X_i$ 后，可以假设
$\overline{Z}_{i, i}$ 与 $\overline{Z}_{i, j}$ 不相交；
参见引理 \ref{flat-lemma-separate-disjoint-locally-closed-by-blowing-up}。
我们假设这对 $X_1$ 和 $X_2$ 都成立。注意，用进一步的
$V_i$-容许爆破替换 $X_i$ 时，此性质仍得保持。

\medskip\noindent
令 $V_{12} = V_1 \times_U V_2$。有一个浸入
$V_{12} \to X_1 \times_S X_2$，它是闭浸入
$V_{12} = V_1 \times_U V_2 \to V_1 \times_S V_2$
（《概形》，引理 \ref{schemes-lemma-fibre-product-after-map}）
与开浸入 $V_1 \times_S V_2 \to X_1 \times_S X_2$ 的复合。
令 $X_{12} \subset X_1 \times_S X_2$ 为
$V_{12} \to X_1 \times_S X_2$ 的概形论像。投影态射
$$
p_1 : X_{12} \to X_1
\quad\text{且}\quad
p_2 : X_{12} \to X_2
$$
是适当态射，因为 $X_1$ 和 $X_2$ 在 $S$ 上适当。
若用一个 $V_1$-容许爆破替换 $X_1$，则 $X_{12}$
被替换为它关于该爆破的严格变换；参见引理
\ref{flat-lemma-replaced-by-strict-transform}。

\medskip\noindent
以 $\psi : V_{12} \to U$ 表示复合
$\psi = \psi_1 \circ p_1|_{V_{12}} = \psi_2 \circ p_2|_{V_{12}}$.
考虑闭子概形
$$
Z_{12, 2} =
(p_1|_{V_{12}})^{-1}(Z_{1, 2}) =
(p_2|_{V_{12}})^{-1}(Z_{2, 2}) =
\psi^{-1}(Z_2) \subset V_{12}
$$
由于 $\psi_2 : V_2 \to U$ 在 $Z_2$ 的某个开邻域上是同构，
且 $V_{12} = V_1 \times_U V_2$，态射
$p_1|_{V_{12}} : V_{12} \to V_1$ 在 $Z_{1, 2}$
的某个开邻域上是同构。由引理 \ref{flat-lemma-blowup-iso-along}，
存在一个 $V_1$-容许爆破 $X_1' \to X_1$，使 $p_1$
的严格变换 $p'_1 : X'_{12} \to X'_1$ 在
$Z_{1, 2}$ 于 $X'_1$ 中的闭包的某个开邻域上是同构。
用 $X'_1$ 替换 $X_1$，并用 $X'_{12}$ 替换 $X_{12}$ 后，
可以假设 $p_1$ 在 $\overline{Z}_{1, 2}$ 的某个开邻域上是同构。

\medskip\noindent
上一段的化简告诉我们
$$
X_{12} \cap (\overline{Z}_{1, 2} \times_S \overline{Z}_{2, 1}) = \emptyset
$$
% P05-EN-CH38-0065: frozen source omits “is” in “where the intersection taken”; candidate “where the intersection is taken”.
，其中交集取于 $X_1 \times_S X_2$ 中。事实上，
$X_{12}$ 中的逆像 $p_1^{-1}(\overline{Z}_{1, 2})$
同构地映到 $\overline{Z}_{1, 2}$。特别地，$Z_{12, 2}$
在 $p_1^{-1}(\overline{Z}_{1, 2})$ 中稠密。因此 $p_2$
将 $p_1^{-1}(\overline{Z}_{1, 2})$ 映入
$\overline{Z}_{2, 2}$。由于
$\overline{Z}_{2, 2} \cap \overline{Z}_{2, 1} = \emptyset$，结论成立。

\medskip\noindent
考虑由黏合得到的概形
$$
W_i = U \coprod\nolimits_{U_i} (X_i \setminus \overline{Z}_{i, j}),
\quad i = 1, 2
$$
应用引理 \ref{flat-lemma-check-separated} 可见 $W_i \to S$ 是分离态射。
首先，$U \to S$ 与 $X_i \to S$ 均为分离态射。浸入
$U_i \to U \times_S (X_i \setminus \overline{Z}_{i, j})$
是闭浸入：任何满足 $u_i \in U_i$ 且
$u \in U \setminus U_i$ 的特化 $u_i \leadsto u$，都能沿适当态射
$\psi_i : V_i \to U$ 唯一提升为 $V_i$ 中的特化
$u_i \leadsto v_i$，而 $v_i$ 必属于 $Z_{i, j}$。
因此该浸入的像为闭集，从而该浸入是闭浸入。

\medskip\noindent
另一方面，设 $A$ 是 $S$ 上的任意赋值环，分式域为 $K$，
并设 $\gamma : \Spec(K) \to (U_1 \cap U_2)$ 是 $S$ 上态射。
则存在某个 $i$，使 $\gamma$ 可延拓为态射
$h_i : \Spec(A) \to W_i$。事实上，对 $i = 1, 2$，
由 $X_i$ 在 $S$ 上适当的赋值判别法，存在延拓 $\gamma$ 的态射
$g_i : \Spec(A) \to X_i$；参见《态射》，引理
\ref{morphisms-lemma-characterize-proper}。
因此，唯一可能出现的问题是对 $i = 1, 2$ 均有
$g_i(\mathfrak m_A) \in \overline{Z}_{i, j}$。
% P05-EN-CH38-0066: frozen source misspells “emptiness” as “emptyness”.
但这与上面已证明的 $X_{12}$ 和
$\overline{Z}_{1, 2} \times_S \overline{Z}_{2, 1}$ 交为空相矛盾。

\medskip\noindent
考虑如引理 \ref{flat-lemma-find-common-blowups} 中的图
$$
\xymatrix{
W_1' \ar[d] \ar[r] & W & W_2' \ar[l] \ar[d] \\
W_1 & U \ar[l] \ar[lu] \ar[u] \ar[ru] \ar[r] & W_2
}
$$
由上一段，对每个实线图
$$
\xymatrix{
\Spec(K) \ar[r]_\gamma  \ar[d] & W \ar[d] \\
\Spec(A) \ar@{..>}[ru] \ar[r] & S
}
$$
若 $\Im(\gamma) \subset U_1 \cap U_2$，则存在某个 $i$
以及 $\gamma$ 的延拓 $h_i : \Spec(A) \to W_i$。
利用 $W'_i \to W_i$ 为适当态射的赋值判别法，可以进一步将
$h_i$ 提升为 $h'_i : \Spec(A) \to W'_i$。故图中的虚线箭头存在。
由于 $W$ 在 $S$ 上分离，该箭头也唯一。由《态射》，引理
\ref{morphisms-lemma-refined-valuative-criterion-universally-closed}，
这说明 $W \to S$ 是普遍闭态射。又因 $W \to S$
已经有限型且分离，故它是适当态射，结论得证。
\end{proof}

\begin{theorem}
\label{flat-theorem-nagata}
\begin{reference}
参见 \cite{Lutkebohmert}、\cite{Conrad-Nagata}、\cite{Nagata-1}、
\cite{Nagata-2}、\cite{Nagata-3} 以及 \cite{Nagata-4}
\end{reference}
设 $S$ 是拟紧且拟分离概形，$X \to S$ 是分离的有限型态射。
则 $X$ 在 $S$ 上有紧化。
\end{theorem}

\begin{proof}
首先化归到 Noether 情形。我们强烈建议读者跳过本段。
存在闭浸入 $X \to X'$，其中 $X' \to S$ 分离且为有限表示；
参见《极限》，命题
\ref{limits-proposition-separated-closed-in-finite-presentation}。
若能找到 $X'$ 在 $S$ 上的一个紧化，则取 $X$ 在其中的概形论像，
便得到 $X$ 在 $S$ 上的紧化。因此可以假设
$X \to S$ 分离且为有限表示。可以将
$S = \lim S_i$ 写成一族 Noether 概形的有向极限，
其转移态射均为仿射态射；参见《极限》，命题
\ref{limits-proposition-approximate}。可以选取某个 $i$
和有限表示态射 $X_i \to S_i$，使它到 $S$ 的基变换为
$X \to S$；参见《极限》，引理
\ref{limits-lemma-descend-finite-presentation}。增大 $i$ 后，
可以假设 $X_i \to S_i$ 分离；参见《极限》，引理
\ref{limits-lemma-descend-separated-finite-presentation}。
若能找到 $X_i$ 在 $S_i$ 上的紧化，则它到 $S$ 的基变换
就是 $X$ 在 $S$ 上的紧化。这就化归到下一段讨论的情形。

\medskip\noindent
现设 $S$ 为 Noether 概形。可以选取有限仿射开覆盖
$X = \bigcup_{i = 1, \ldots, n} U_i$，使
$U_1 \cap \ldots \cap U_n$ 在 $X$ 中稠密。这由
《性质》，引理 \ref{properties-lemma-point-and-maximal-points-affine}
以及 $X$ 拟紧且只有有限多个不可约分支这一事实推出。
对每个 $i$，由《态射》，引理
\ref{morphisms-lemma-quasi-affine-finite-type-over-S}，
可以选取 $n_i \geq 0$ 和浸入 $U_i \to \mathbf{A}^{n_i}_S$。
在 $\mathbf{P}^{n_i}_S$ 中取概形论像，可知对
$i = 1, \ldots, n$，$U_i$ 在 $S$ 上有紧化。
应用引理 \ref{flat-lemma-two-compactifications} 共
$(n - 1)$ 次，即得定理。
\end{proof}






\section{h 拓扑}
\label{flat-section-h}

\noindent
粗略地说，在本文中，h 层是 Zariski 拓扑下的层，并且对满射的
有限表示适当态射满足层条件；参见引理
\ref{flat-lemma-characterize-sheaf-h}。不过值得指出的是，
概形范畴上 h 拓扑的定义随参考文献而异。

\medskip\noindent
Voevodsky 最初将 h 覆盖定义为有限族有限型态射，并要求它们联合地
普遍为商态射（《态射》，定义
\ref{morphisms-definition-submersive}）；参见
\cite[定义 3.1.2]{Voevodsky}。当所用概形范畴限制为某个固定
Noether 基概形上的全体有限型概形时，这一定义最为适用。
在此情形中，Voevodsky 将 h 覆盖与 ph 覆盖联系起来。
ph 拓扑由 Zariski 覆盖和适当满射生成。更多内容参见
《拓扑》，第 \ref{topologies-section-ph} 节。

\medskip\noindent
我们在《拓扑》，第 \ref{topologies-section-V} 节中研究 V 拓扑。
若 $Y$ 中点的任意特化都是 $X$ 中点的某个特化之像，
且任意基变换后仍如此，则拟紧态射 $X \to Y$ 定义一个 V 覆盖
（《拓扑》，引理 \ref{topologies-lemma-refine-qcqs-V}）。
此时 $X \to Y$ 是普遍商态射
（《拓扑》，引理
\ref{topologies-lemma-V-covering-universally-submersive}）。
处理非 Noether 概形时，V 覆盖这一概念可以很好地代替
“具有固定靶且联合地普遍为商态射的态射族”。

\medskip\noindent
我们的方法是先对局部有限表示的态射族（允许无限族）证明
ph 覆盖与 V 覆盖等价。随后在 Stacks Project 中以这些态射族作为
h 覆盖的定义。对 Noether 概形和有限族而言，这些覆盖与
Voevodsky 的定义一致；参见引理 \ref{flat-lemma-Noetherian-h-covering}。
在固定拟紧且拟分离概形 $S$ 上有限表示概形的范畴中，
这些覆盖所确定的拓扑与 \cite[定义 2.7]{Witt-Grass}
中的拓扑相同。

\begin{lemma}
\label{flat-lemma-equivalence-h-v-locally-finite-presentation}
设 $\{f_i : X_i \to X\}_{i \in I}$ 是具有固定靶的概形态射族，
并且对每个 $i$，$f_i$ 都局部为有限表示。下列条件等价：
\begin{enumerate}
\item $\{X_i \to X\}$ 是 ph 覆盖；
\item $\{X_i \to X\}$ 是 V 覆盖。
\end{enumerate}
\end{lemma}

\begin{proof}
设 $U \subset X$ 为仿射开子集。由《拓扑》，定义
\ref{topologies-definition-ph-covering} 和
\ref{topologies-definition-V-covering}，只需证明基变换
$\{X_i \times_X U \to U\}$ 可由标准 ph 覆盖加细，当且仅当
它可由标准 V 覆盖加细。因此可以假设 $X$ 仿射，而需证明
$\{X_i \to X\}$ 可由标准 ph 覆盖加细，当且仅当它可由
标准 V 覆盖加细。标准 ph 覆盖是标准 V 覆盖；参见
《拓扑》，引理 \ref{topologies-lemma-standard-ph-standard-V}。
故只需证明反向蕴含。

\medskip\noindent
假设 $X$ 仿射，且 $\{f_i : X_i \to X\}_{i \in I}$
可由标准 V 覆盖
$\{g_j : Y_j \to X\}_{j = 1, \ldots, m}$ 加细。
对每个 $j$，选取 $i_j$ 和态射 $h_j : Y_j \to X_{i_j}$，
使 $g_j = f_{i_j} \circ h_j$。由于 $Y_j$ 仿射，因而拟紧，
对每个 $j$ 可以找到有限多个仿射开子集
$U_{j, k} \subset X_{i_j}$，使
$\Im(h_j) \subset \bigcup U_{j, k}$。于是
$\{U_{j, k} \to X\}_{j, k}$ 加细 $\{X_i \to X\}$，
并且是标准 V 覆盖（它是仿射概形之间态射的有限族，
而且从 $\{Y_j \to X\}$ 的相应性质继承赋值环提升性质）。
这就化归到下一段讨论的情形。

\medskip\noindent
假设 $\{f_i : X_i \to X\}_{i = 1, \ldots, n}$
是标准 V 覆盖，且 $f_i$ 为有限表示。需证明
$\{X_i \to X\}$ 可由标准 ph 覆盖加细。选取一般平坦性分层
$$
X = S \supset S_0 \supset S_1 \supset \ldots \supset S_t = \emptyset
$$
它是 《态射进阶》，引理
\ref{more-morphisms-lemma-generic-flatness-stratification-scheme}，
对仿射概形之间的有限表示态射
$$
\coprod\nolimits_{i = 1, \ldots, n} f_i :
\coprod\nolimits_{i = 1, \ldots, n} X_i
\longrightarrow
X
$$
所给出的分层。下文将不再说明地使用该分层的全部性质。
依构造，每个 $f_i$ 到 $U_k = S_k \setminus S_{k + 1}$
的基变换都是平坦态射。以 $Y_k$ 表示 $U_k$ 在 $S_k$
中的概形论闭包。由于 $U_k \to S_k$ 是拟紧开浸入
（参见《性质》，引理
\ref{properties-lemma-quasi-coherent-finite-type-ideals}），
$U_k \subset Y_k$ 是拟紧稠密（且概形论稠密）的开浸入；
参见《态射》，引理
\ref{morphisms-lemma-quasi-compact-scheme-theoretic-image}。
态射 $\coprod_{k = 0, \ldots, t - 1} Y_k \to X$ 有限且满，
故 $\{Y_k \to X\}$ 是标准 ph 覆盖，因而也是标准 V 覆盖
（见上文）。由标准 V 覆盖的传递性质
（《拓扑》，引理 \ref{topologies-lemma-composition-standard-V}），
只需证明覆盖 $\{X_i \to X\}$ 到每个 $Y_k$ 的拉回
可由标准 V 覆盖加细。这就化归到下一段的情形。

\medskip\noindent
假设 $\{f_i : X_i \to X\}_{i = 1, \ldots, n}$ 是标准 V 覆盖，
$f_i$ 为有限表示，并且存在稠密拟紧开子集 $U \subset X$，
使 $X_i \times_X U \to U$ 平坦。由定理
\ref{flat-theorem-flatten-module}，存在 $U$-容许爆破
$X' \to X$，使 $f_i$ 的严格变换
$f'_i : X'_i \to X'$ 平坦。注意，投射（因而闭）态射
$X' \to X$ 是满射，因为 $U \subset X$ 稠密且 $U$
被视为 $X'$ 的开子集。必要时用进一步的 $U$-容许爆破替换
$X'$，还可假设 $U \subset X'$ 概形论稠密
（参见注 \ref{flat-remark-successive-blowups}）。
因此，对每个点 $x \in X'$，存在赋值环 $V$ 和态射
$g : \Spec(V) \to X'$，使 $\Spec(V)$ 的一般点映入 $U$，
而 $\Spec(V)$ 的闭点映到 $x$；参见《态射》，引理
\ref{morphisms-lemma-reach-points-scheme-theoretic-image}。
由于 $\{X_i \to X\}$ 是标准 V 覆盖，可以选取赋值环扩张
$V \subset W$、指标 $i$ 和态射 $\Spec(W) \to X_i$，使图
$$
\xymatrix{
\Spec(W) \ar[d] \ar[rr] & & X_i \ar[d] \\
\Spec(V) \ar[r] & X' \ar[r] & X
}
$$
交换。由于 $X'_i \subset X' \times_X X_i$ 是包含开子概形
$U \times_X X_i$ 的闭子概形，$\Spec(W)$ 是整概形，
而诱导态射 $h : \Spec(W) \to X' \times_X X_i$
将 $\Spec(W)$ 的一般点映入 $U \times_X X_i$，可知 $h$
经闭子概形 $X'_i \subset X' \times_X X_i$ 分解。
因此 $\{f'_i : X'_i \to X'\}$ 是 V 覆盖。特别地，
$\coprod f'_i$ 是满射，故 $\{X'_i \to X'\}$ 是 fppf 覆盖。
由于 fppf 覆盖是 ph 覆盖
（《态射进阶》，引理 \ref{more-morphisms-lemma-fppf-ph}），
可以找到加细 $\{X'_i \to X\}$ 的标准 ph 覆盖
$\{Y_j \to X'\}$。设此覆盖由适当满射 $Y \to X'$
和有限仿射开覆盖 $Y = \bigcup Y_j$ 给出。于是复合
$Y \to X$ 是适当满射，从而 $\{Y_j \to X\}$ 是标准 ph 覆盖。
证明完毕。
\end{proof}

\noindent
现给出我们的定义。

\begin{definition}
\label{flat-definition-h-covering}
设 $T$ 是概形。$T$ 的一个{\it h 覆盖}是态射族
$\{f_i : T_i \to T\}_{i \in I}$，其中每个 $f_i$
都局部为有限表示，并满足引理
\ref{flat-lemma-equivalence-h-v-locally-finite-presentation}
中的一个（从而全部）等价条件。
\end{definition}

\noindent
对 Noether 概形，这与 ph 覆盖相同
（我们将在下述引理 \ref{flat-lemma-Noetherian-h-ph} 中记录此事），
因而恢复了 Voevodsky 的概念。

\begin{lemma}
\label{flat-lemma-Noetherian-h-covering}
设 $X$ 是 Noether 概形，并设 $\{X_i \to X\}_{i \in I}$
是有限族有限型态射。下列条件等价：
\begin{enumerate}
\item $\coprod_{i \in I} X_i \to X$ 是普遍商态射
（《态射》，定义 \ref{morphisms-definition-submersive}）；
\item $\{X_i \to X\}_{i \in I}$ 是 h 覆盖。
\end{enumerate}
\end{lemma}

\begin{proof}
由更一般的 《拓扑》，引理
\ref{topologies-lemma-V-covering-universally-submersive}
和 h 覆盖的定义，(2) $\Rightarrow$ (1) 立即成立。
假设 $\coprod X_i \to X$ 是普遍商态射。我们将证明
$\{X_i \to X\}$ 可由 ph 覆盖加细；由《拓扑》，引理
\ref{topologies-lemma-refine-by-ph} 和 h 覆盖的定义，这已经足够。
论证与引理 \ref{flat-lemma-equivalence-h-v-locally-finite-presentation}
的证明中所用者相同。

\medskip\noindent
选取一般平坦性分层
$$
X = S \supset S_0 \supset S_1 \supset \ldots \supset S_t = \emptyset
$$
它是 《态射进阶》，引理
\ref{more-morphisms-lemma-generic-flatness-stratification-scheme}
中对有限表示态射
$$
\coprod\nolimits_{i = 1, \ldots, n} f_i :
\coprod\nolimits_{i = 1, \ldots, n} X_i
\longrightarrow
X
$$
所给出的分层。下文将不再说明地使用该分层的全部性质。
依构造，每个 $f_i$ 到 $U_k = S_k \setminus S_{k + 1}$
的基变换都是平坦态射。以 $Y_k$ 表示 $U_k$ 在 $S_k$
中的概形论闭包。由于 $U_k \to S_k$ 是拟紧开浸入
（本段所有概形均为 Noether 概形），$U_k \subset Y_k$
是拟紧稠密（且概形论稠密）的开浸入；参见《态射》，引理
\ref{morphisms-lemma-quasi-compact-scheme-theoretic-image}。
态射 $\coprod_{k = 0, \ldots, t - 1} Y_k \to X$ 有限且满，
故 $\{Y_k \to X\}$ 是 ph 覆盖。由 ph 覆盖的传递性质
（《拓扑》，引理 \ref{topologies-lemma-ph}），
只需证明覆盖 $\{X_i \to X\}$ 到每个 $Y_k$ 的拉回
可由 ph 覆盖加细。这就化归到下一段的情形。

\medskip\noindent
假设 $\coprod X_i \to X$ 是普遍商态射，并且存在稠密开子集
$U \subset X$，使对每个 $i$，$X_i \times_X U \to U$
均为平坦态射。由定理 \ref{flat-theorem-flatten-module}，
存在 $U$-容许爆破 $X' \to X$，使对每个 $i$，
$f_i$ 的严格变换 $f'_i : X'_i \to X'$ 均为平坦态射。
注意，投射（因而闭）态射 $X' \to X$ 是满射，因为
$U \subset X$ 稠密且 $U$ 被视为 $X'$ 的开子集。必要时用进一步的
$U$-容许爆破替换 $X'$，还可假设 $U \subset X'$ 稠密
（参见注 \ref{flat-remark-successive-blowups}）。
因此，对每个点 $x \in X'$，存在离散赋值环 $A$ 和态射
$g : \Spec(A) \to X'$，使 $\Spec(A)$ 的一般点映入 $U$，
而 $\Spec(A)$ 的闭点映到 $x$；参见《极限》，引理
\ref{limits-lemma-reach-point-closure-Noetherian}。令
$$
W = \Spec(A) \times_X \coprod X_i = \coprod \Spec(A) \times_X X_i
$$
由于 $\coprod X_i \to X$ 是普遍商态射，$W$ 中存在特化
$w' \leadsto w$，使 $w'$ 映到 $\Spec(A)$ 的一般点，
而 $w$ 映到 $\Spec(A)$ 的闭点。
（否则，$W \to \Spec(A)$ 的闭纤维在一般化下稳定，因而为开集，
这与 $W \to \Spec(A)$ 是商态射相矛盾。）
设 $w' \in \Spec(A) \times_X X_i$，则当然也有
$w \in \Spec(A) \times_X X_i$。令
$x'_i \leadsto x_i$ 为 $w' \leadsto w$ 在
$X' \times_X X_i$ 中的像。由于 $x'_i \in X'_i$，且
$X'_i \subset X' \times_X X_i$ 是闭子概形，可知
$x_i \in X'_i$。又因为 $x_i$ 映到 $x \in X'$，所以
$\coprod X'_i \to X'$ 是满射！特别地，
$\{X'_i \to X'\}$ 是 fppf 覆盖，而 fppf 覆盖是 ph 覆盖
（《态射进阶》，引理 \ref{more-morphisms-lemma-fppf-ph}）。
由于 $X' \to X$ 是适当满射，$\{X'_i \to X\}$ 是 ph 覆盖，
证明完毕。
\end{proof}

\begin{lemma}
\label{flat-lemma-Noetherian-h-ph}
设 $X$ 是局部 Noether 概形。以 $X$ 为靶的态射族
$\{f_i : X_i \to X\}_{i \in I}$ 是 h 覆盖，当且仅当它是 ph 覆盖。
\end{lemma}

\begin{proof}
由定义 \ref{flat-definition-h-covering}，h 覆盖是 ph 覆盖。
反之，若 $\{f_i : X_i \to X\}$ 是 ph 覆盖，则态射 $f_i$
局部为有限型（《拓扑》，定义
\ref{topologies-definition-ph-covering}）。由于 $X$ 局部 Noether，
每个 $f_i$ 都局部为有限表示，故依定义得到 h 覆盖。
\end{proof}

\noindent
% P05-EN-CH38-0067: frozen source has compound subject “lemma and ... theorem” with singular “shows”; candidate “show”.
下述引理与 \cite[定理 8.4]{rydh_descent} 表明，
我们的定义与 David Rydh 的论文 \cite{rydh_descent} 中的定义一致
（或至少密切相关）。为简便起见，我们限制在仿射基情形。

\begin{lemma}
\label{flat-lemma-approximate-h-cover}
设 $X$ 是仿射概形，$\{X_i \to X\}_{i \in I}$ 是 h 覆盖。
则存在满的适当态射
$$
Y \longrightarrow X
$$
它还是有限表示的（！），并且存在有限仿射开覆盖
$Y = \bigcup_{j = 1, \ldots, m} Y_j$，使
$\{Y_j \to X\}_{j = 1, \ldots, m}$ 加细
$\{X_i \to X\}_{i \in I}$。
\end{lemma}

\begin{proof}
依假设，存在适当满射 $Y \to X$ 和有限仿射开覆盖
$Y = \bigcup_{j = 1, \ldots, m} Y_j$，使
$\{Y_j \to X\}_{j = 1, \ldots, m}$ 加细
$\{X_i \to X\}_{i \in I}$。这意味着对每个 $j$，
存在指标 $i_j \in I$ 及 $X$ 上态射
$h_j : Y_j \to X_{i_j}$。参见定义 \ref{flat-definition-h-covering}
和 《拓扑》，定义 \ref{topologies-definition-ph-covering}。
问题在于尚不知道 $Y \to X$ 是否为有限表示。
由《极限》，引理
\ref{limits-lemma-proper-limit-of-proper-finite-presentation}，可将
$$
Y = \lim Y_\lambda
$$
写成 $X$ 上概形 $Y_\lambda$ 的有向极限，其中这些概形
适当且为有限表示，态射 $Y \to Y_\lambda$ 及转移态射均为闭浸入。
% P05-EN-CH38-0068: frozen source duplicates “the” before “transition morphisms”.
注意，每个 $Y_\lambda \to X$ 都是满射。由《极限》，引理
\ref{limits-lemma-descend-opens}，可以找到某个 $\lambda$
和覆盖 $Y_\lambda$ 的拟紧开子集
$Y_{\lambda, j} \subset Y_\lambda$，$j = 1, \ldots, m$，
它们在 $Y$ 上限制为 $Y_j$。于是
$Y_j = \lim Y_{\lambda, j}$。增大 $\lambda$ 后，可假设对所有
$j$，$Y_{\lambda, j}$ 均仿射；参见《极限》，引理
\ref{limits-lemma-limit-affine}。最后，由于 $X_i \to X$
局部为有限表示，可以使用局部有限表示态射的函子性刻画
（《极限》，命题
\ref{limits-proposition-characterize-locally-finite-presentation}）
找到某个 $\lambda$，使对每个 $j$ 都存在态射
$h_{\lambda, j} : Y_{\lambda, j} \to X_{i_j}$，
它在 $Y_j$ 上的限制就是上面选取的态射 $h_j$。
因此 $\{Y_{\lambda, j} \to X\}$ 加细
$\{X_i \to X\}$，证明完毕。
\end{proof}

\noindent
现在回到 h 覆盖一般理论的建立。

\begin{lemma}
\label{flat-lemma-zariski-h}
fppf 覆盖是 h 覆盖。因此 syntomic、光滑、\'etale 以及
Zariski 覆盖也都是 h 覆盖。
\end{lemma}

\begin{proof}
这是因为 fppf 覆盖中的态射按要求局部为有限表示，
且 fppf 覆盖是 ph 覆盖；参见《态射进阶》，引理
\ref{more-morphisms-lemma-fppf-ph}。第二个断言由第一个断言和
《拓扑》，引理
\ref{topologies-lemma-zariski-etale-smooth-syntomic-fppf} 推出。
\end{proof}

\begin{lemma}
\label{flat-lemma-surjective-proper-finite-presentation-h}
设 $f : Y \to X$ 是有限表示的适当满概形态射。
则 $\{Y \to X\}$ 是 h 覆盖。
\end{lemma}

\begin{proof}
合并应用 《拓扑》，引理
\ref{topologies-lemma-zariski-etale-smooth-syntomic-fppf-fpqc-ph-V} 和
\ref{topologies-lemma-surjective-proper-ph}.
\end{proof}

\begin{lemma}
\label{flat-lemma-refine-by-h}
设 $T$ 是概形，并设 $\{f_i : T_i \to T\}_{i \in I}$ 是态射族，
使对所有 $i$，$f_i$ 均局部为有限表示。下列条件等价：
\begin{enumerate}
\item $\{T_i \to T\}_{i \in I}$ 是 h 覆盖；
\item 存在加细 $\{T_i \to T\}_{i \in I}$ 的 h 覆盖；
\item $\{\coprod_{i \in I} T_i \to T\}$ 是 h 覆盖。
\end{enumerate}
\end{lemma}

\begin{proof}
这由 ph 覆盖的相应断言
（《拓扑》，引理 \ref{topologies-lemma-refine-by-ph}）
或 V 覆盖的相应断言
（《拓扑》，引理 \ref{topologies-lemma-refine-by-V}）推出。
\end{proof}

\noindent
下面证明我们的 h 覆盖概念满足 《站点》，定义
\ref{sites-definition-site} 的条件。

\begin{lemma}
\label{flat-lemma-h}
设 $T$ 是概形。
\begin{enumerate}
\item 若 $T' \to T$ 是同构，则 $\{T' \to T\}$ 是 $T$ 的 h 覆盖。
\item 若 $\{T_i \to T\}_{i\in I}$ 是 h 覆盖，且对每个 $i$
都有 h 覆盖 $\{T_{ij} \to T_i\}_{j\in J_i}$，则
$\{T_{ij} \to T\}_{i \in I, j\in J_i}$ 是 h 覆盖。
\item 若 $\{T_i \to T\}_{i\in I}$ 是 h 覆盖，而
$T' \to T$ 是概形态射，则
$\{T' \times_T T_i \to T'\}_{i\in I}$ 是 h 覆盖。
\end{enumerate}
\end{lemma}

\begin{proof}
这立即由 ph 覆盖或 V 覆盖的相应断言
（《拓扑》，引理 \ref{topologies-lemma-ph} 或
\ref{topologies-lemma-V}），以及局部有限表示态射类
在基变换和复合下保持这一事实推出。
\end{proof}

\noindent
下面定义我们将在 Stacks Project 中使用的大 h 站点。
继续之前，宜先阅读 《拓扑》，第 \ref{topologies-section-procedure} 节中的一般讨论。

\begin{definition}
\label{flat-definition-big-h-site}
一个{\it 大 h 站点}是按如下方式构造的、《站点》，定义
\ref{sites-definition-site} 意义下的任意站点 $\Sch_h$：
\begin{enumerate}
\item 任取一组概形 $S_0$，并任取这些概形之间的一组 h 覆盖
$\text{Cov}_0$。
\item 以任意范畴 $\Sch_\alpha$ 作为底层范畴；该范畴从集合
$S_0$ 出发，按《集合》，引理
\ref{sets-lemma-construct-category} 构造。
\item 从范畴 $\Sch_\alpha$、h 覆盖类和上面选取的集合
$\text{Cov}_0$ 出发，按《集合》，引理
\ref{sets-lemma-coverings-site} 任取一组覆盖。
\end{enumerate}
\end{definition}

\noindent
关于大站点定义的动机与说明，参见《拓扑》，定义
\ref{topologies-definition-big-zariski-site} 之后的注释。

\begin{definition}
\label{flat-definition-standard-h}
设 $T$ 是仿射概形。$T$ 的一个{\it 标准 h 覆盖}是态射族
$\{f_i : T_i \to T\}_{i = 1, \ldots, n}$，其中每个 $T_i$
都仿射、每个 $f_i$ 都为有限表示，并满足下列等价条件之一：
(1) $\{T_i \to T\}$ 可由标准 ph 覆盖加细；或
(2) $\{T_i \to T\}$ 是 V 覆盖。
\end{definition}

\noindent
这些条件的等价性由引理
\ref{flat-lemma-equivalence-h-v-locally-finite-presentation}、
《拓扑》，定义 \ref{topologies-definition-ph-covering}
和引理 \ref{topologies-lemma-refine-by-ph} 推出。

\medskip\noindent
在继续引入概形 $S$ 的大 h 站点之前，先指出：
大 h 站点 $\Sch_h$ 上的拓扑在某种意义下由全体概形范畴上的
h 拓扑诱导而来。

\begin{lemma}
\label{flat-lemma-h-induced}
设 $\Sch_h$ 是定义 \ref{flat-definition-big-h-site} 中的大 h 站点，
$T \in \Ob(\Sch_h)$，并设
$\{T_i \to T\}_{i \in I}$ 是 $T$ 的任意 h 覆盖。
\begin{enumerate}
\item 站点 $\Sch_h$ 中存在 $T$ 的覆盖
$\{U_j \to T\}_{j \in J}$，它加细 $\{T_i \to T\}_{i \in I}$。
\item 若 $\{T_i \to T\}_{i \in I}$ 是标准 h 覆盖，
则它按定义自然等价于 $\Sch_h$ 的一个覆盖。
\item 若 $\{T_i \to T\}_{i \in I}$ 是 Zariski 覆盖，
则它按定义自然等价于 $\Sch_h$ 的一个覆盖。
\end{enumerate}
\end{lemma}

\begin{proof}
证明从略。提示：这与 《拓扑》，引理
\ref{topologies-lemma-ph-induced} 的证明完全相同。
\end{proof}

\begin{definition}
\label{flat-definition-big-small-h}
设 $S$ 是概形，$\Sch_h$ 是包含 $S$ 的大 h 站点。
\begin{enumerate}
\item $S$ 的{\it 大 h 站点}记为 $(\Sch/S)_h$，
即 《站点》，第 \ref{sites-section-localize} 节中引入的站点
$\Sch_h/S$。
\item $S$ 的{\it 大仿射 h 站点}记为
$(\textit{Aff}/S)_h$，它是 $(\Sch/S)_h$ 的满子范畴，
其对象为仿射的 $U/S$。$(\textit{Aff}/S)_h$ 的覆盖是
$(\Sch/S)_h$ 的任意覆盖 $\{U_i \to U\}$，并要求它是标准 h 覆盖。
\end{enumerate}
\end{definition}

\noindent
我们明确陈述大仿射 h 站点确实是一个站点。

\begin{lemma}
\label{flat-lemma-verify-site-h}
设 $S$ 是概形，$\Sch_h$ 是包含 $S$ 的大 h 站点。
则 $(\textit{Aff}/S)_h$ 是一个站点。
\end{lemma}

\begin{proof}
仿照 《拓扑》，引理
\ref{topologies-lemma-verify-site-etale} 的证明，只需说明
标准 h 覆盖族满足 《站点》，定义
\ref{sites-definition-site} 的性质 (1)、(2)、(3)。
这很清楚。例如，给定标准 h 覆盖
$\{T_i \to T\}_{i\in I}$，并对每个 $i$ 给定标准 h 覆盖
$\{T_{ij} \to T_i\}_{j \in J_i}$，则
$\{T_{ij} \to T\}_{i \in I, j\in J_i}$ 是 h 覆盖
（引理 \ref{flat-lemma-h}），$\bigcup_{i\in I} J_i$ 是有限集，
而每个 $T_{ij}$ 都仿射。因此
$\{T_{ij} \to T\}_{i \in I, j\in J_i}$ 是标准 h 覆盖。
\end{proof}

\begin{lemma}
\label{flat-lemma-fibre-products-h}
设 $S$ 是概形，$\Sch_h$ 是包含 $S$ 的大 h 站点。
站点 $\Sch_h$、$(\Sch/S)_h$ 和 $(\textit{Aff}/S)_h$
的底层范畴都有纤维积。在每种情形中，到全体概形范畴
$\Sch$ 的显然函子均与取纤维积交换。范畴
$(\Sch/S)_h$ 有终对象，即 $S/S$。
\end{lemma}

\begin{proof}
对 $\Sch_h$，结论依构造成立；参见《集合》，引理
\ref{sets-lemma-what-is-in-it}。设有概形态射
$U \to S$、$V \to U$、$W \to U$，其中
$U, V, W \in \Ob(\Sch_h)$。$\Sch_h$ 中的纤维积
$V \times_U W$ 也是 $\Sch$ 中的纤维积；在全体 $S$ 上概形的范畴中，
它是 $V/S$ 与 $W/S$ 在 $U/S$ 上的纤维积，因而也是
$(\Sch/S)_h$ 中的纤维积。这证明了 $(\Sch/S)_h$ 的结论。
若 $U, V, W$ 均仿射，则 $V \times_U W$ 也仿射，
故 $(\textit{Aff}/S)_h$ 的结论也成立。
\end{proof}

\noindent
下面验证大仿射站点与大站点定义同一个拓扑斯。

\begin{lemma}
\label{flat-lemma-affine-big-site-h}
设 $S$ 是概形，$\Sch_h$ 是包含 $S$ 的大 h 站点。
函子 $(\textit{Aff}/S)_h \to (\Sch/S)_h$ 是余连续的，
并诱导从 $\Sh((\textit{Aff}/S)_h)$ 到
$\Sh((\Sch/S)_h)$ 的拓扑斯等价。
\end{lemma}

\begin{proof}
特殊余连续函子的概念见《站点》，定义
\ref{sites-definition-special-cocontinuous-functor}。
因此需验证 《站点》，引理 \ref{sites-lemma-equivalence}
的假设 (1) -- (5)。记包含函子为
$u : (\textit{Aff}/S)_h \to (\Sch/S)_h$。
若 $T$ 仿射，则 $T/S$ 的任意 h 覆盖都可由标准 h 覆盖加细，
例如可用引理 \ref{flat-lemma-approximate-h-cover}；故余连续性成立，
从而 (1) 成立。由于标准 h 覆盖本身就是 h 覆盖，$u$ 显然连续，
故 (2) 成立。(3) 和 (4) 立即由 $u$ 全忠实推出。
最后，每个概形都有仿射开覆盖（这是 h 覆盖），故 (5) 成立。
\end{proof}

\begin{lemma}
\label{flat-lemma-characterize-sheaf-h}
设 $\mathcal{F}$ 是 $(\Sch/S)_h$ 上的预层。
则 $\mathcal{F}$ 是层，当且仅当
\begin{enumerate}
\item $\mathcal{F}$ 对 Zariski 覆盖满足层条件；
\item 若 $f : V \to U$ 适当、满且为有限表示，则
$\mathcal{F}(U)$ 双射地映到两个映射
$\mathcal{F}(V) \to \mathcal{F}(V \times_U V)$ 的等化子。
\end{enumerate}
此外，在 (1) 成立时，性质 (2) 等价于性质
\begin{enumerate}
\item[(2')] 对 (2) 中 $U$ 仿射的 $\{V \to U\}$ 满足层性质。
\end{enumerate}
\end{lemma}

\begin{proof}
先证明若 (1) 和 (2) 成立，则 $\mathcal{F}$ 是层。
设 $\{T_i \to T\}$ 是 $(\Sch/S)_h$ 中的覆盖。
我们验证该覆盖的层条件。设 $s_i \in \mathcal{F}(T_i)$ 是一族截面，
它们在 $T_i \times_T T_{i'}$ 上限制为同一截面。
我们将证明存在唯一截面 $s \in \mathcal{F}(T)$，
它在 $T_i$ 上限制为 $s_i$。设
$T = \bigcup U_j$ 是仿射开覆盖。由性质 (1)，为构造 $s$，
只需构造在 $U_j \cap U_{j'}$ 上相符的截面
$s_j \in \mathcal{F}(U_j)$。考虑覆盖
$\{T_i \times_T U_j \to U_j\}$。则截面
$s_{ji} = s_i|_{T_i \times_T U_j}$ 在
$(T_i \times_T U_j) \times_{U_j} (T_{i'} \times_T U_j)$ 上相符。
选取有限表示的适当满射 $V_j \to U_j$ 和有限仿射开覆盖
$V_j = \bigcup V_{jk}$，使 $\{V_{jk} \to U_j\}$
加细 $\{T_i \times_T U_j \to U_j\}$；参见引理
\ref{flat-lemma-approximate-h-cover}。若
$s_{jk} \in \mathcal{F}(V_{jk})$ 表示通过隐含态射将 $s_{ji}$
拉回到 $V_{jk}$ 所得的截面，则 $s_{jk}$ 黏合为截面
$s'_j \in \mathcal{F}(V_j)$。再次使用交叠处的相符性，可见
$s'_j$ 属于两个映射
$\mathcal{F}(V_j) \to \mathcal{F}(V_j \times_{U_j} V_j)$ 的等化子。
故由 (2)，$s'_j$ 来自唯一截面
$s_j \in \mathcal{F}(U_j)$。这些截面 $s_j$
具有全部所需性质的验证从略。

\medskip\noindent
现证明在 (1) 成立时 (2) 与 (2') 等价。设
$V \to U$ 是 $(\Sch/S)_h$ 中适当、满且为有限表示的态射。
选取仿射开覆盖 $U = \bigcup U_i$，并令
$V_i = V \times_U U_i$。由 (2')，映射
$\mathcal{F}(U_i) \to \mathcal{F}(V_i)$ 为单射；
由 (1)，映射 $\mathcal{F}(U) \to \prod \mathcal{F}(U_i)$
为单射。因此 $\mathcal{F}(U) \to \mathcal{F}(V)$ 为单射。
最后，设给定 $t \in \mathcal{F}(V)$，它属于两个映射
$\mathcal{F}(V) \to \mathcal{F}(V \times_U V)$ 的等化子。
则对所有 $i$，$t|_{V_i}$ 属于两个映射
$\mathcal{F}(V_i) \to \mathcal{F}(V_i \times_{U_i} V_i)$
的等化子。故由 (2')，对所有 $i$，得到唯一截面
$s_i \in \mathcal{F}(U_i)$，它映到 $t|_{V_i}$。
对所有 $i, j$ 验证
$s_i|_{U_i \cap U_j} = s_j|_{U_i \cap U_j}$ 从略；
这里使用刚证明的唯一性。由覆盖 $U = \bigcup U_i$ 的层性质，
得到截面 $s \in \mathcal{F}(U)$。$s$ 在
$\mathcal{F}(V)$ 中映到 $t$ 的证明从略。
\end{proof}

\noindent
下面建立与这些站点相关的拓扑斯之间的一些关系。

\begin{lemma}
\label{flat-lemma-morphism-big-h}
设 $\Sch_h$ 是大 h 站点，$f : T \to S$ 是 $\Sch_h$ 中的态射。
函子
$$
u : (\Sch/T)_h \longrightarrow (\Sch/S)_h,
\quad
V/T \longmapsto V/S
$$
是余连续的，并有连续右伴随
$$
v : (\Sch/S)_h \longrightarrow (\Sch/T)_h,
\quad
(U \to S) \longmapsto (U \times_S T \to T).
$$
它们诱导同一个拓扑斯态射
$$
f_{big} :
\Sh((\Sch/T)_h)
\longrightarrow
\Sh((\Sch/S)_h)
$$
此外，
$f_{big}^{-1}(\mathcal{G})(U/T) = \mathcal{G}(U/S)$，
$f_{big, *}(\mathcal{F})(U/S) = \mathcal{F}(U \times_S T/T)$。
此外，$f_{big}^{-1}$ 有左伴随 $f_{big!}$，后者与纤维积和等化子交换。
\end{lemma}

\begin{proof}
函子 $u$ 余连续、连续，并与纤维积和等化子交换。因此可应用
《站点》，引理 \ref{sites-lemma-when-shriek} 和
\ref{sites-lemma-preserve-equalizers}，从而得到
$f_{big}^{-1}$ 的公式及 $f_{big!}$ 的存在性。此外，函子 $v$
是右伴随，因为给定 $U/T$ 与 $V/S$，有预期的等式
$\Mor_S(u(U), V) = \Mor_T(U, V \times_S T)$。
故可应用 《站点》，引理
\ref{sites-lemma-have-functor-other-way} 和
\ref{sites-lemma-have-functor-other-way-morphism}，
得到 $f_{big, *}$ 的公式。
\end{proof}

\begin{lemma}
\label{flat-lemma-composition-h}
% P05-EN-CH38-0069: frozen source lists X, Y, Y although the following morphism uses Z; candidate third scheme is Z.
给定 $(\Sch/S)_h$ 中的概形 $X$、$Y$、$Y$ 以及态射
$f : X \to Y$、$g : Y \to Z$，有
$g_{big} \circ f_{big} = (g \circ f)_{big}$.
\end{lemma}

\begin{proof}
这由引理 \ref{flat-lemma-morphism-big-h} 中大站点上函子的
推前与拉回的简单描述推出。
\end{proof}










\section{h 拓扑的进一步结果}
\label{flat-section-h-more}

\noindent
本节证明关于 h 拓扑的若干进一步结果。先看一些反例。

\begin{example}
\label{flat-example-structure-sheaf-not-h-sheaf}
“结构层”$\mathcal{O}$ 不是 h 拓扑中的层。例如，考虑有限表示的
满闭浸入 $X \to Y$。由引理
\ref{flat-lemma-surjective-proper-finite-presentation-h}，
$\{X \to Y\}$ 是 h 覆盖。还要注意 $X \times_Y X = X$。
% P05-EN-CH38-0070: frozen source has “where a sheaf”; candidate “were a sheaf”.
因此，若 $\mathcal{O}$ 是 h 拓扑中的层，则
$\mathcal{O}_Y(Y) \to \mathcal{O}_X(X)$ 应为双射。
一旦 $X$, $Y$ 仿射而态射 $X \to Y$ 不是同构，事实便非如此。
\end{example}

\begin{example}
\label{flat-example-representable-sheaf-not-h-sheaf}
在任意站点 $(\Sch/S)_h$ 上，该拓扑都不是次典范的；
换言之，可表预层未必是层。事实上，“结构层”$\mathcal{O}$
是可表的，因为在 $(\Sch/S)_h$ 中
$\mathcal{O}(X) = \Mor_S(X, \mathbf{A}^1_S)$；
而例 \ref{flat-example-structure-sheaf-not-h-sheaf} 已说明
$\mathcal{O}$ 不是层。
\end{example}

\begin{lemma}
\label{flat-lemma-limit-h-topology}
设 $T$ 是仿射概形，并被写成一族仿射概形所成有向逆系统的极限
$T = \lim_{i \in I} T_i$。
\begin{enumerate}
\item 设 $\mathcal{V} = \{V_j \to T\}_{j = 1, \ldots, m}$
是 $T$ 的标准 h 覆盖；参见定义 \ref{flat-definition-standard-h}。
则存在指标 $i$ 和标准 h 覆盖
$\mathcal{V}_i = \{V_{i, j} \to T_i\}_{j = 1, \ldots, m}$，
使它到 $T$ 的基变换 $T \times_{T_i} \mathcal{V}_i$
同构于 $\mathcal{V}$。
\item 设 $\mathcal{V}_i$, $\mathcal{V}'_i$ 是 $T_i$
的两个标准 h 覆盖。若
$f : T \times_{T_i} \mathcal{V}_i \to T \times_{T_i} \mathcal{V}'_i$
是 $T$ 的覆盖之间的态射，则存在指标 $i' \geq i$ 和态射
% P05-EN-CH38-0071: frozen source uses unindexed \mathcal{V} in the domain below although this item starts from \mathcal{V}_i.
$f_{i'} : T_{i'} \times_{T_i} \mathcal{V} \to
T_{i'} \times_{T_i} \mathcal{V}'_i$
，其到 $T$ 的基变换为 $f$。
\item 若
% P05-EN-CH38-0072: frozen source again uses \mathcal{V} rather than \mathcal{V}_i for a covering of T_i.
$f, g : \mathcal{V} \to \mathcal{V}'_i$
是 $T_i$ 的标准 h 覆盖之间的态射，并且它们到 $T$
的基变换 $f_T, g_T$ 相等，则存在指标 $i' \geq i$，
使 $f_{T_{i'}} = g_{T_{i'}}$。
\end{enumerate}
换言之，$T$ 的标准 h 覆盖范畴是 $T_i$ 的标准 h 覆盖范畴
在 $I$ 上的余极限。
\end{lemma}

\begin{proof}
由《极限》，引理 \ref{limits-lemma-descend-finite-presentation}，
$T$ 上有限表示概形的范畴是 $T_i$ 上有限表示概形范畴
在 $I$ 上的余极限。
% P05-EN-CH38-0073: frozen source says “categories of finite presentation over T_i”; candidate inserts “schemes”.
由《极限》，引理
\ref{limits-lemma-descend-affine-finite-presentation}，
对 $T$ 上仿射且有限表示的概形范畴也有相同结论。
为完成引理的证明，只需证明如下断言：若
$\{V_{j, i} \to T_i\}_{j = 1, \ldots, m}$ 是有限表示态射的有限族，
每个 $V_{j, i}$ 仿射，且基变换族
$\{T \times_{T_i} V_{j, i} \to T\}$ 是 h 覆盖，
则对某个 $i' \geq i$，态射族
$\{T_{i'} \times_{T_i} V_{j, i} \to T_{i'}\}$ 是 h 覆盖。
为此，利用引理 \ref{flat-lemma-approximate-h-cover} 选取有限表示的
适当满射 $Y \to T$ 和有限仿射开覆盖
$Y = \bigcup_{k = 1, \ldots, n} Y_k$，使
$\{Y_k \to T\}_{k = 1, \ldots, n}$ 加细
$\{T \times_{T_i} V_{j, i} \to T\}$。
利用上面的论证及 《极限》，引理
\ref{limits-lemma-eventually-proper}、
\ref{limits-lemma-descend-surjective} 和
\ref{limits-lemma-descend-opens}，可找到 $i' \geq i$、
有限表示的适当满射 $Y_{i'} \to T_{i'}$ 以及仿射开覆盖
$Y_{i'} = \bigcup_{k = 1, \ldots, n} Y_{i', k}$，
并使 $\{Y_{i', k} \to Y_{i'}\}$ 加细
$\{T_{i'} \times_{T_i} V_{j, i} \to T_{i'}\}$。
因此最后这个态射族是 h 覆盖，证明完毕。
\end{proof}

\begin{lemma}
\label{flat-lemma-extend-sheaf-h}
设 $S$ 是包含在大站点 $\Sch_h$ 中的概形。设
$F : (\Sch/S)_h^{opp} \to \textit{Sets}$ 是 h 层，
并在 $\mathcal{C} = (\Sch/S)_h$ 时满足 《拓扑》，引理
\ref{topologies-lemma-extend} 的性质 (b)。
则 $F$ 到全体 $S$ 上概形范畴的延拓 $F'$ 对所有 h 覆盖
满足层条件，并保持极限（《极限》，注记
\ref{limits-remark-limit-preserving}）。
\end{lemma}

\begin{proof}
使用引理 \ref{flat-lemma-limit-h-topology} 和 \ref{flat-lemma-h-induced}，
并采用 《拓扑》，引理
\ref{topologies-lemma-extend-sheaf-general} 和
\ref{topologies-lemma-extend-sheaf} 证明中的论证即可。细节从略。
\end{proof}








\section{爆破方块与 ph 拓扑}
\label{flat-section-blow-up-ph}

\noindent
设 $X$ 是概形，$Z \subset X$ 是闭子概形，且包含态射为有限表示；
换言之，与 $Z$ 对应的拟相干理想层为有限型。
设 $b : X' \to X$ 是 $X$ 沿 $Z$ 的爆破，并令
$E = b^{-1}(Z)$ 为例外除子。参见《因子》，第 \ref{divisors-section-blowing-up} 节。在此情形下，并在本节中，
称
\begin{equation}
\label{flat-equation-blow-up-square}
\vcenter{
\xymatrix{
E \ar[d] \ar[r] & X' \ar[d]^b \\
Z \ar[r] & X
}
}
\end{equation}
为一个{\it 爆破方块}。

\begin{lemma}
\label{flat-lemma-blow-up-square-ph}
设 $\mathcal{F}$ 是站点 $(\Sch/S)_{ph}$ 上的层；参见
《拓扑》，定义 \ref{topologies-definition-big-small-ph}。
则对范畴 $(\Sch/S)_{ph}$ 中的任意爆破方块
(\ref{flat-equation-blow-up-square})，图
$$
\xymatrix{
\mathcal{F}(E) & \mathcal{F}(X') \ar[l] \\
\mathcal{F}(Z) \ar[u] & \mathcal{F}(X) \ar[u] \ar[l]
}
$$
在集合范畴中为笛卡儿图。
\end{lemma}

\begin{proof}
由于 $Z \amalg X' \to X$ 是适当满射，
$\{Z \amalg X' \to X\}$ 是 ph 覆盖
（《拓扑》，引理 \ref{topologies-lemma-surjective-proper-ph}）。
并且
$$
(Z \amalg X') \times_X (Z \amalg X') =
Z \amalg E \amalg E \amalg X' \times_X X'
$$
由于 $\mathcal{F}$ 是 Zariski 层，$\mathcal{F}$ 将不交并映到乘积。
因此，覆盖 $\{Z \amalg X' \to X\}$ 的层条件说明
$\mathcal{F}(X) \to \mathcal{F}(Z) \times \mathcal{F}(X')$
为单射，其像是满足下列条件的对 $(t, s')$ 所成集合：
(a) $t|_E = s'|_E$；(b) $s'$ 属于两个映射
$\mathcal{F}(X') \to \mathcal{F}(X' \times_X X')$ 的等化子。
接着注意，显然态射
$$
E \times_Z E \amalg X' \longrightarrow X' \times_X X'
$$
是适当满射，因为 $b$ 诱导同构
$X' \setminus E \to X \setminus Z$。故
$\mathcal{F}(X' \times_X X') \to
\mathcal{F}(E \times_Z E) \times \mathcal{F}(X')$ 为单射。
于是 (a) $\Rightarrow$ (b)，引理得证。
\end{proof}

\begin{lemma}
\label{flat-lemma-thickening-ph}
设 $\mathcal{F}$ 是 《拓扑》，定义
\ref{topologies-definition-big-small-ph} 中站点
$(\Sch/S)_{ph}$ 上的层。设 $X \to X'$ 是
$(\Sch/S)_{ph}$ 中的态射，并且是增厚。则
$\mathcal{F}(X') \to \mathcal{F}(X)$ 为双射。
\end{lemma}

\begin{proof}
% P05-EN-CH38-0074: frozen source has stray “of” in “a proper surjective morphism of and”.
注意，$X \to X'$ 是适当满射，且 $X \times_{X'} X = X$。
对 ph 覆盖 $\{X \to X'\}$ 应用层性质
（《拓扑》，引理
\ref{topologies-lemma-surjective-proper-ph}）即得结论。
\end{proof}








\section{几乎爆破方块与 h 拓扑}
\label{flat-section-blow-up-h}

\noindent
考虑爆破方块 (\ref{flat-equation-blow-up-square})。虽然态射
$b : X' \to X$ 是投射态射
（《因子》，引理 \ref{divisors-lemma-blowing-up-projective}），
一般却没有简单方法保证 $b$ 为有限表示。由于 h 覆盖使用有限表示态射
构造，我们需要一个变体。具体地，称概形的交换图
\begin{equation}
\label{flat-equation-almost-blow-up-square}
\vcenter{
\xymatrix{
E \ar[d] \ar[r] & X' \ar[d]^b \\
Z \ar[r] & X
}
}
\end{equation}
为一个{\it 几乎爆破方块}，如果满足下列条件：
\begin{enumerate}
\item $Z \to X$ 是有限表示的闭浸入；
\item $E = b^{-1}(Z)$ 是 $X'$ 的局部主闭子概形；
\item $b$ 适当且为有限表示；
\item 由 $\mathcal{O}_{X'}$ 中支集在 $E$ 上的截面所成拟相干理想
切出的闭子概形 $X'' \subset X'$
（《性质》，引理
\ref{properties-lemma-sections-supported-on-closed-subset}）
是 $X$ 沿 $Z$ 的爆破。
\end{enumerate}
由此可知，态射 $b$ 诱导同构
$X' \setminus E \to X \setminus Z$。
几乎爆破方块的一些很简单的例子参见例
\ref{flat-example-one-generator} 和 \ref{flat-example-two-generators}。

\medskip\noindent
爆破的基变换通常不是爆破，但几乎爆破与基变换相容。

\begin{lemma}
\label{flat-lemma-base-change-almost-blow-up}
考虑几乎爆破方块 (\ref{flat-equation-almost-blow-up-square})，
并设 $Y \to X$ 为任意态射。则基变换
$$
\xymatrix{
Y \times_X E \ar[d] \ar[r] & Y \times_X X' \ar[d] \\
Y \times_X Z \ar[r] & Y
}
$$
也是几乎爆破方块。
\end{lemma}

\begin{proof}
由《态射》，引理 \ref{morphisms-lemma-base-change-proper} 和
\ref{morphisms-lemma-base-change-finite-presentation}，
态射 $Y \times_X X' \to Y$ 适当且为有限表示。态射
$Y \times_X Z \to Y$ 是有限表示的闭浸入
（《态射》，引理
\ref{morphisms-lemma-base-change-closed-immersion}）。
$Y \times_X Z$ 在 $Y \times_X X'$ 中的逆像等于 $E$
在 $Y \times_X X'$ 中的逆像，因而局部为主
（《因子》，引理
\ref{divisors-lemma-pullback-locally-principal}）。
分别令 $X'' \subset X'$ 和 $Y'' \subset Y \times_X X'$
为与 $\mathcal{O}_{X'}$ 和
% P05-EN-CH38-0075: frozen source has Y \times_Y X' where the ambient space is Y \times_X X'.
$\mathcal{O}_{Y \times_Y X'}$ 中分别支集在 $E$ 和
$Y \times_X E$ 上的截面所成拟相干理想对应的闭子概形。
显然，$Y'' \subset Y \times_X X''$ 是与
% P05-EN-CH38-0076: frozen source has Y \times_Y X'' where the ambient base change is Y \times_X X''.
$\mathcal{O}_{Y \times_Y X''}$ 中支集在
$Y \times_X (E \cap X'')$ 上的截面所成拟相干理想对应的闭子概形。
因此 $Y''$ 是 $Y$ 关于爆破 $X'' \to X$ 的严格变换；
参见《因子》，定义
\ref{divisors-definition-strict-transform}。故由《因子》，引理
\ref{divisors-lemma-strict-transform}，$Y''$
是 $Y$ 沿 $Y \times_X Z$ 的爆破。
\end{proof}

\noindent
可以缩小几乎爆破方块。

\begin{lemma}
\label{flat-lemma-shrink-almost-blow-up}
考虑几乎爆破方块 (\ref{flat-equation-almost-blow-up-square})，
并设 $W \to X'$ 是有限表示的闭浸入。下列条件等价：
\begin{enumerate}
\item $X' \setminus E$ 在概形论意义下包含于 $W$；
\item $X$ 沿 $Z$ 的爆破 $X''$ 在概形论意义下包含于 $W$；
\item 图
$$
\xymatrix{
E \cap W \ar[d] \ar[r] & W \ar[d] \\
Z \ar[r] & X
}
$$
是几乎爆破方块。这里 $E \cap W$ 是概形论交。
\end{enumerate}
\end{lemma}

\begin{proof}
假设 (1)。则满射 $\mathcal{O}_{X'} \to \mathcal{O}_W$
在开集
% P05-EN-CH38-0077: frozen source writes X' \subset E; condition (1) and the argument require X' \setminus E.
$X' \subset E$ 上是同构。由于 $X'' \subset X'$ 的理想层
由 $\mathcal{O}_{X'}$ 中支集在 $E$ 上的截面组成
（依我们对几乎爆破方块的定义），可知 (2) 成立。
若 (2) 成立，则 (3) 成立。若 (3) 成立，则 (1) 成立，
因为 $X'' \cap (X' \setminus E)$ 同构于 $X \setminus Z$，
而后者又同构于 $X' \setminus E$。
\end{proof}

\noindent
实际爆破是任意给定几乎爆破的各个缩小之极限。

\begin{lemma}
\label{flat-lemma-blow-up-limit-almost-blow-up}
考虑几乎爆破方块 (\ref{flat-equation-almost-blow-up-square})，
其中 $X$ 拟紧且拟分离。则 $X$ 沿 $Z$ 的爆破 $X''$ 可写成
$$
X'' = \lim X'_i
$$
这里的极限取遍由有限表示闭子概形 $X'_i \subset X'$
组成的有向系统；这些闭子概形满足引理
\ref{flat-lemma-shrink-almost-blow-up} 的等价条件。
\end{lemma}

\begin{proof}
令 $\mathcal{I} \subset \mathcal{O}_{X'}$ 为与 $X''$
对应的拟相干理想层。由《性质》，引理
\ref{properties-lemma-quasi-coherent-colimit-finite-type}，
可将 $\mathcal{I}$ 写成其有限型拟相干子模的滤过余极限
$\mathcal{I} = \colim \mathcal{I}_i$。这些模与闭子概形
$X'_i$ 按 $1$-对-$1$ 对应，故证明完毕。
\end{proof}

\noindent
几乎爆破方块确实存在。

\begin{lemma}
\label{flat-lemma-almost-blow-up-square}
设 $X$ 是拟紧且拟分离概形，$Z \subset X$ 是由有限型拟相干理想层
切出的闭子概形。则存在如
(\ref{flat-equation-almost-blow-up-square}) 所示的几乎爆破方块。
\end{lemma}

\begin{proof}
可以将 $X = \lim X_i$ 写成一族 Noether 概形所成逆系统的有向极限，
其转移态射均为仿射态射；参见《极限》，命题
\ref{limits-proposition-approximate}。可以找到指标 $i$
和闭浸入 $Z_i \to X_i$，其到 $X$ 的基变换是闭浸入
$Z \to X$。参见《极限》，引理
\ref{limits-lemma-descend-finite-presentation} 和
\ref{limits-lemma-descend-closed-immersion-finite-presentation}。
令 $b_i : X'_i \to X_i$ 为以 $Z_i$ 为中心的爆破。
这给出爆破方块
$$
\xymatrix{
E_i \ar[r] \ar[d] & X'_i \ar[d]^{b_i} \\
Z_i \ar[r] & X_i
}
$$
其中所有态射都是 Noether 概形之间的有限型态射，因而为有限表示。
所以这是几乎爆破方块。由引理
\ref{flat-lemma-base-change-almost-blow-up}，将此图基变换到 $X$
即可得到所需的几乎爆破方块。
\end{proof}

\noindent
几乎爆破方块在引理 \ref{flat-lemma-shrink-almost-blow-up}
所述的缩小意义下是唯一的。

\begin{lemma}
\label{flat-lemma-almost-blow-up-unique}
设 $X$ 是拟紧且拟分离概形，$Z \subset X$
是由有限型拟相干理想层切出的闭子概形。设给定几乎爆破方块
(\ref{flat-equation-almost-blow-up-square})
$$
\xymatrix{
E_k \ar[r] \ar[d] & X_k' \ar[d] \\
Z \ar[r] & X
}
$$
其中 $k = 1, 2$。则存在几乎爆破方块
$$
\xymatrix{
E \ar[r] \ar[d] & X' \ar[d] \\
Z \ar[r] & X
}
$$
以及 $X$ 上闭浸入 $i_k : X' \to X'_k$，使
$E = i_k^{-1}(E_k)$。
\end{lemma}

\begin{proof}
以 $X'' \to X$ 表示 $X$ 沿 $Z$ 的爆破。将 $X''$
同时视为 $X'_1$ 和 $X'_2$ 的闭子概形。如引理
\ref{flat-lemma-blow-up-limit-almost-blow-up} 中那样写
$X'' = \lim X'_{1, i}$。由《极限》，命题
\ref{limits-proposition-characterize-locally-finite-presentation}，
存在某个 $i$ 和态射 $h : X'_{1, i} \to X'_2$，
它与包含 $X'' \subset X'_{1, i}$ 及
$X'' \subset X'_2$ 相容。由《极限》，引理
\ref{limits-lemma-eventually-closed-immersion}，对某个
$i' \geq i$，$h$ 在 $X'_{1, i'}$ 上的限制是闭浸入。
证明完毕。
\end{proof}

\noindent
在适当情形下，我们关于爆破的扁平化技巧也适用于几乎爆破。

\begin{lemma}
\label{flat-lemma-flat-after-almost-blowing-up}
设 $Y$ 是拟紧且拟分离概形，$X$ 是 $Y$ 上有限表示概形。
设 $V \subset Y$ 是拟紧开子集，且 $X_V \to V$ 平坦。
则存在交换图
$$
\xymatrix{
E \ar[ddd] \ar[rd] & & & D \ar[lll] \ar[ddd] \ar[ld] \\
& Y' \ar[d] & X' \ar[l] \ar[d] \\
& Y & X \ar[l] \\
Z \ar[ru] & & & T \ar[lll] \ar[lu]
}
$$
其左右两个方块为几乎爆破方块，上下两个方块为笛卡儿方块，
并且 $Z \cap V = \emptyset$，而 $X' \to Y'$ 平坦
（且为有限表示）。
\end{lemma}

\begin{proof}
若 $Y$ 是 Noether 概形，则引理立即由引理
\ref{flat-lemma-flat-after-blowing-up} 推出，因为在此情形下，
爆破方块就是几乎爆破方块（还要用到严格变换是爆破）。
一般情形通过绝对 Noether 逼近化归到 Noether 情形。

\medskip\noindent
可以将 $Y = \lim Y_i$ 写成一族 Noether 概形所成逆系统的有向极限，
其转移态射均仿射；参见《极限》，命题
\ref{limits-proposition-approximate}。可以找到指标 $i$
和有限表示态射 $X_i \to Y_i$，其到 $Y$ 的基变换是
$X \to Y$；参见《极限》，引理
\ref{limits-lemma-descend-finite-presentation}。增大 $i$ 后，
可以假设 $V$ 是某个开子概形 $V_i \subset Y_i$ 的逆像；
参见《极限》，引理 \ref{limits-lemma-descend-opens}。
最后，再增大 $i$，可以假设 $X_{i, V_i} \to V_i$ 平坦；
参见《极限》，引理
\ref{limits-lemma-descend-flat-finite-presentation}。
由 Noether 情形，可对 $X_i \to Y_i \supset V_i$
构造引理中的图。沿 $Y \to Y_i$ 对此图作基变换即得所求。
这里使用了基变换保持相关态射性质；参见《态射》，引理
\ref{morphisms-lemma-base-change-proper},
\ref{morphisms-lemma-base-change-finite-presentation},
\ref{morphisms-lemma-base-change-closed-immersion} 和
\ref{morphisms-lemma-base-change-flat}
；还使用了几乎爆破方块的基变换仍为几乎爆破方块，
参见引理 \ref{flat-lemma-base-change-almost-blow-up}。
\end{proof}

\begin{lemma}
\label{flat-lemma-blow-up-square-h}
设 $\mathcal{F}$ 是定义 \ref{flat-definition-big-small-h}
中构造的某个站点 $(\Sch/S)_h$ 上的层。
则对范畴 $(\Sch/S)_h$ 中的任意几乎爆破方块
(\ref{flat-equation-almost-blow-up-square})，图
$$
\xymatrix{
\mathcal{F}(E) & \mathcal{F}(X') \ar[l] \\
\mathcal{F}(Z) \ar[u] & \mathcal{F}(X) \ar[u] \ar[l]
}
$$
在集合范畴中为笛卡儿图。
\end{lemma}

\begin{proof}
由于 $Z \amalg X' \to X$ 是有限表示的适当满射，
$\{Z \amalg X' \to X\}$ 是 h 覆盖
（引理 \ref{flat-lemma-surjective-proper-finite-presentation-h}）。
并且
$$
(Z \amalg X') \times_X (Z \amalg X') =
Z \amalg E \amalg E \amalg X' \times_X X'
$$
由于 $\mathcal{F}$ 是 Zariski 层，$\mathcal{F}$ 将不交并映到乘积。
因此覆盖 $\{Z \amalg X' \to X\}$ 的层条件说明
$\mathcal{F}(X) \to \mathcal{F}(Z) \times \mathcal{F}(X')$
为单射，其像是满足下列条件的对 $(t, s')$ 所成集合：
(a) $t|_E = s'|_E$；(b) $s'$ 属于两个映射
$\mathcal{F}(X') \to \mathcal{F}(X' \times_X X')$ 的等化子。
接着注意，显然态射
$$
E \times_Z E \amalg X' \longrightarrow X' \times_X X'
$$
是有限表示的适当满射，因为 $b$ 诱导同构
$X' \setminus E \to X \setminus Z$。故
$\mathcal{F}(X' \times_X X') \to
\mathcal{F}(E \times_Z E) \times \mathcal{F}(X')$ 为单射。
于是 (a) $\Rightarrow$ (b)，引理得证。
\end{proof}

\begin{lemma}
\label{flat-lemma-thickening-h}
设 $\mathcal{F}$ 是定义 \ref{flat-definition-big-small-h}
中构造的某个站点 $(\Sch/S)_h$ 上的层。设
$X \to X'$ 是 $(\Sch/S)_h$ 中有限表示的增厚。则
$\mathcal{F}(X') \to \mathcal{F}(X)$ 为双射。
\end{lemma}

\begin{proof}
第一种证明。注意，$X \to X'$ 是有限表示的适当满射，且
$X \times_{X'} X = X$。对 h 覆盖 $\{X \to X'\}$
应用层性质（引理
\ref{flat-lemma-surjective-proper-finite-presentation-h}）即得结论。

\medskip\noindent
第二种证明（有些可笑）。$X'$ 沿 $X$ 的爆破是空概形。
理由是，若 $a$ 是 $A$ 的幂零元素，则仿射爆破代数
$A[\frac{I}{a}]$（《代数》，第 \ref{algebra-section-blow-up} 节）为零。细节从略。
因此得到如下形式的几乎爆破方块：
$$
\xymatrix{
\emptyset \ar[r] \ar[d] & \emptyset \ar[d] \\
X \ar[r] & X'
}
$$
由于 $\mathcal{F}$ 是层，$\mathcal{F}(\emptyset)$ 是单点集。
应用引理 \ref{flat-lemma-blow-up-square-h} 即得结论。
\end{proof}

\begin{proposition}
\label{flat-proposition-check-h}
设 $\mathcal{F}$ 是定义 \ref{flat-definition-big-small-h}
中构造的某个站点 $(\Sch/S)_h$ 上的预层。
则 $\mathcal{F}$ 是层，当且仅当满足下列条件：
\begin{enumerate}
\item $\mathcal{F}$ 是 Zariski 拓扑下的层；
\item 给定 $(\Sch/S)_h$ 中的态射 $f : X \to Y$，
其中 $Y$ 仿射，而 $f$ 满、平坦、适当且为有限表示，
则 $\mathcal{F}(Y)$ 是两个映射
$\mathcal{F}(X) \to \mathcal{F}(X \times_Y X)$ 的等化子；
\item 给定范畴 $(\Sch/S)_h$ 中的几乎爆破方块
(\ref{flat-equation-almost-blow-up-square})，其中 $X$ 仿射，则图
$$
\xymatrix{
\mathcal{F}(E) & \mathcal{F}(X') \ar[l] \\
\mathcal{F}(Z) \ar[u] & \mathcal{F}(X) \ar[u] \ar[l]
}
$$
在集合范畴中为笛卡儿图。
\end{enumerate}
\end{proposition}

\begin{proof}
假设 $\mathcal{F}$ 是层。由于 Zariski 覆盖是 h 覆盖，
条件 (1) 成立；参见引理 \ref{flat-lemma-zariski-h}。
对 (2) 中的 $f$，$\{X \to Y\}$ 显然是 fppf 覆盖，
因而是 h 覆盖；参见引理 \ref{flat-lemma-zariski-h}。
故条件 (2) 成立。条件 (3) 由引理
\ref{flat-lemma-blow-up-square-h} 成立。

\medskip\noindent
反之，假设 $\mathcal{F}$ 满足 (1)、(2)、(3)。
我们将应用引理 \ref{flat-lemma-characterize-sheaf-h}
证明 $\mathcal{F}$ 是层。考虑 $(\Sch/S)_h$ 中满、有限表示且适当的态射
$f : X \to Y$，其中 $Y$ 仿射。只需证明
$\mathcal{F}(Y)$ 是两个映射
$\mathcal{F}(X) \to \mathcal{F}(X \times_Y X)$ 的等化子。

\medskip\noindent
先进一步假设 $f : X \to Y$ 是闭浸入
（换言之，$f$ 是增厚）。则 $Y$ 沿 $X$ 的爆破为空概形，
从而产生一个顶点为 $\emptyset, \emptyset, X, Y$ 的几乎爆破方块
% P05-EN-CH38-0078: frozen source says “consisting with”; candidate “consisting of”.
（比较引理 \ref{flat-lemma-thickening-h} 的第二种证明）。
因此条件 (3) 告诉我们图
$$
\xymatrix{
\mathcal{F}(\emptyset) & \mathcal{F}(\emptyset) \ar[l] \\
\mathcal{F}(X) \ar[u] & \mathcal{F}(Y) \ar[u] \ar[l]
}
$$
在集合范畴中为笛卡儿图。由于 $\mathcal{F}$ 是 Zariski 拓扑下的层，
$\mathcal{F}(\emptyset)$ 是单点集。因此
$\mathcal{F}(X) = \mathcal{F}(Y)$。

\medskip\noindent
插曲 A：设 $T \to T'$ 是 $(\Sch/S)_h$ 中有限表示的增厚。
则 $\mathcal{F}(T') \to \mathcal{F}(T)$ 为双射。
事实上，选取仿射开覆盖 $T' = \bigcup T'_i$，并令
$T_i = T \times_{T'} T'_i$ 为 $T$ 中相应的仿射开集。
由上一段的结论，对所有 $i$，映射
$\mathcal{F}(T'_i) \to \mathcal{F}(T_i)$ 都为双射。
利用 Zariski 层性质，可见
$\mathcal{F}(T') \to \mathcal{F}(T)$ 为单射。
重复该论证可知它是双射。略去少量细节。

\medskip\noindent
插曲 B：考虑范畴 $(\Sch/S)_h$ 中的几乎爆破方块
(\ref{flat-equation-almost-blow-up-square})。我们断言图
$$
\xymatrix{
\mathcal{F}(E) & \mathcal{F}(X') \ar[l] \\
\mathcal{F}(Z) \ar[u] & \mathcal{F}(X) \ar[u] \ar[l]
}
$$
在集合范畴中为笛卡儿图。选取 $X$ 的仿射开覆盖，
并仿照插曲 A 论证，即可由条件 (3) 推出此事。细节从略。

\medskip\noindent
接下来，设 $f : X \to Y$ 是 $(\Sch/S)_h$ 中满、有限表示且适当的态射，
其中 $Y$ 仿射。选取一般平坦性分层
$$
Y \supset Y_0 \supset Y_1 \supset \ldots \supset Y_t = \emptyset
$$
它是 《态射进阶》，引理
\ref{more-morphisms-lemma-generic-flatness-stratification-scheme}
中为 $f : X \to Y$ 给出的分层。下文将不再说明地使用该分层的全部性质。
令 $X_0 = X \times_Y Y_0$。由插曲 B，有
$\mathcal{F}(Y_0) = \mathcal{F}(Y)$、
$\mathcal{F}(X_0) = \mathcal{F}(X)$，且
$\mathcal{F}(X_0 \times_{Y_0} X_0) = \mathcal{F}(X \times_Y X)$.

\medskip\noindent
我们对 $t$ 归纳证明结论。若 $t = 1$，则
$X_0 \to Y_0$ 满、适当、平坦且为有限表示，故由性质 (2) 得到结论。
若 $t > 1$，则可用 $Y_0$ 替换 $Y$，并用 $X_0$ 替换 $X$
（见上文），从而假设 $Y = Y_0$。

\medskip\noindent
考虑拟紧开子概形
$V = Y \setminus Y_1 = Y_0 \setminus Y_1$.
为 $f : X \to Y \supset V$ 选取如引理
\ref{flat-lemma-flat-after-almost-blowing-up} 中的图
$$
\xymatrix{
E \ar[ddd] \ar[rd] & & & D \ar[lll] \ar[ddd] \ar[ld] \\
& Y' \ar[d] & X' \ar[l] \ar[d] \\
& Y & X \ar[l] \\
Z \ar[ru] & & & T \ar[lll] \ar[lu]
}
$$
于是 $f' : X' \to Y'$ 平坦且为有限表示。$f'$ 也适当
（使用《态射》，引理
\ref{morphisms-lemma-composition-proper} 和
\ref{morphisms-lemma-image-proper-scheme-closed}）。
因此，像 $W = f'(X') \subset Y'$ 是 $Y'$ 的开闭子概形
（对开性使用《态射》，引理
\ref{morphisms-lemma-fppf-open}）。注意，
$Y' \setminus E$ 包含于 $W$。由引理
\ref{flat-lemma-shrink-almost-blow-up}，这意味着可在上图中用
$W$ 替换 $Y'$。换言之，可以并且确实假设 $f'$ 为满射。
此时由插曲 B 可知
$$
\vcenter{
\xymatrix{
\mathcal{F}(E) & \mathcal{F}(Y') \ar[l] \\
\mathcal{F}(Z) \ar[u] & \mathcal{F}(Y) \ar[u] \ar[l]
}
}
\quad\text{且}\quad
\vcenter{
\xymatrix{
\mathcal{F}(D) & \mathcal{F}(X') \ar[l] \\
\mathcal{F}(T) \ar[u] & \mathcal{F}(X) \ar[u] \ar[l]
}
}
$$
中的两个方块均为笛卡儿图。注意，
$Z \cap Y_1 \to Z$ 是有限表示的增厚
（因为 $Z$ 作为 $Y$ 中与 $V$ 不相交的闭子概形，
在集合论意义下包含于 $Y_1$）。因此得到滤过
$$
% P05-EN-CH38-0079: frozen source reverses the last inclusions; a decreasing stratification requires \supset throughout.
Z \supset Z \cap Y_1 \supset Z \cap Y_2 \subset \ldots \subset
Z \cap Y_t = \emptyset
$$
该滤过用于 $f$ 在 $T$ 上的限制
$T = Z \times_Y X \to Z$。故由归纳假设，
$\mathcal{F}(Z) \to \mathcal{F}(T)$
是集合的单射，其像为两个映射
$\mathcal{F}(T) \to \mathcal{F}(T \times_Z T)$ 的等化子。

\medskip\noindent
设 $s \in \mathcal{F}(X)$ 属于两个映射
$\mathcal{F}(X) \to \mathcal{F}(X \times_Y X)$ 的等化子。
由上文，限制 $s|_T$ 来自唯一元素
$t \in \mathcal{F}(Z)$；同理，限制 $s|_{X'}$ 来自唯一元素
$t' \in \mathcal{F}(Y')$。沿上面巨大交换图中箭头所对应的
$\mathcal{F}$ 限制映射追踪截面，可知 $t$ 与 $t'$
在 $\mathcal{F}(E)$ 中限制为同一元素：它们在
$\mathcal{F}(D)$ 中限制为同一元素，并且有 (2)；
这里使用了 $D \to E$ 作为 $X' \to Y'$ 的限制，是满、平坦、适当
且有限表示的态射。因此，由上面两个笛卡儿方块中的第一个，
得到唯一截面
$u \in \mathcal{F}(Y)$
，它在 $Z$ 和 $Y'$ 上分别限制为 $t$ 和 $t'$。
% P05-EN-CH38-0080: frozen source has “u restrict”; candidate “u restricts”.
要验证 $u$ 在 $X$ 上限制为 $s$，使用第二个图即可。
\end{proof}

\begin{example}
\label{flat-example-one-generator}
设 $A$ 是环，$f \in A$ 是元素，$J \subset A$
是被 $f$ 的某个幂零化的有限生成理想。则
$$
\xymatrix{
E = \Spec(A/fA + J) \ar[r] \ar[d] & \Spec(A/J) = X' \ar[d] \\
Z = \Spec(A/fA) \ar[r] & \Spec(A) = X
}
$$
是几乎爆破方块。
\end{example}

\begin{example}
\label{flat-example-two-generators}
设 $A$ 是环，$f_1, f_2 \in A$ 是元素。
$$
\xymatrix{
E = \text{Proj}(A/(f_1, f_2)[T_0, T_1]) \ar[r] \ar[d] &
% P05-EN-CH38-0081: frozen source appears to omit the closing parenthesis of \text{Proj}(...) before = X'.
\text{Proj}(A[T_0, T_1]/(f_2 T_0 - f_1 T_1) = X' \ar[d] \\
Z = \Spec(A/(f_1, f_2)) \ar[r] & \Spec(A) = X
}
$$
是几乎爆破方块。
\end{example}

\begin{lemma}
\label{flat-lemma-refine-check-h}
设 $\mathcal{F}$ 是定义 \ref{flat-definition-big-small-h}
中构造的某个站点 $(\Sch/S)_h$ 上的预层。
则 $\mathcal{F}$ 是层，当且仅当满足下列条件：
\begin{enumerate}
\item $\mathcal{F}$ 是 Zariski 拓扑下的层；
\item 给定 $(\Sch/S)_h$ 中的态射 $f : X \to Y$，
其中 $Y$ 仿射，而 $f$ 满、平坦、适当且为有限表示，
则 $\mathcal{F}(Y)$ 是两个映射
$\mathcal{F}(X) \to \mathcal{F}(X \times_Y X)$ 的等化子；
\item $\mathcal{F}$ 将范畴 $(\Sch/S)_h$ 中如例
\ref{flat-example-one-generator} 的几乎爆破方块变为集合的笛卡儿图；
\item $\mathcal{F}$ 将范畴 $(\Sch/S)_h$ 中如例
\ref{flat-example-two-generators} 的几乎爆破方块变为集合的笛卡儿图。
\end{enumerate}
\end{lemma}

\begin{proof}
由命题 \ref{flat-proposition-check-h}，只需证明：给定范畴
$(\Sch/S)_h$ 中的几乎爆破方块
(\ref{flat-equation-almost-blow-up-square})，其中 $X$ 仿射，则图
$$
\xymatrix{
\mathcal{F}(E) & \mathcal{F}(X') \ar[l] \\
\mathcal{F}(Z) \ar[u] & \mathcal{F}(X) \ar[u] \ar[l]
}
$$
在集合范畴中为笛卡儿图。证明的大致思路是用其他几乎爆破方块
支配该态射，以便局部应用假设 (3) 和 (4)。

\medskip\noindent
假设有范畴 $(\Sch/S)_h$ 中的几乎爆破方块
(\ref{flat-equation-almost-blow-up-square})、开覆盖
$X = \bigcup U_i$，以及开覆盖
$U_i \cap U_j = \bigcup U_{ijk}$，使图
$$
\vcenter{
\xymatrix{
\mathcal{F}(E \cap b^{-1}(U_i)) & \mathcal{F}(b^{-1}(U_i)) \ar[l] \\
\mathcal{F}(Z \cap U_i) \ar[u] & \mathcal{F}(U_i) \ar[u] \ar[l]
}
}
\quad\text{且}\quad
\vcenter{
\xymatrix{
\mathcal{F}(E \cap b^{-1}(U_{ijk})) &
\mathcal{F}(b^{-1}(U_{ijk})) \ar[l] \\
\mathcal{F}(Z \cap U_{ijk}) \ar[u] &
\mathcal{F}(U_{ijk}) \ar[u] \ar[l]
}
}
$$
均为笛卡儿图。则图
$$
\xymatrix{
\mathcal{F}(E) & \mathcal{F}(X') \ar[l] \\
\mathcal{F}(Z) \ar[u] & \mathcal{F}(X) \ar[u] \ar[l]
}
$$
也为笛卡儿图。这是因为 $\mathcal{F}$ 是 Zariski 拓扑下的层。

\medskip\noindent
特别地，若有爆破方块
% P05-EN-CH38-0082: frozen source calls equation-almost-blow-up-square a “blow up square”; candidate “almost blow up square”.
(\ref{flat-equation-almost-blow-up-square})，其中 $b : X' \to X$
为闭浸入，而 $Z$ 是局部主闭子概形，则
$\mathcal{F}(X) = \mathcal{F}(X') \times_{\mathcal{F}(E)} \mathcal{F}(Z)$.
事实上，在 $X$ 上仿射局部地，可得到如 (3) 中的几乎爆破方块。

\medskip\noindent
令 $Z \subset X$、$E_k \subset X'_k \to X$、
$E \subset X' \to X$ 以及 $i_k : X' \to X'_k$
如引理 \ref{flat-lemma-almost-blow-up-unique} 的陈述中所述。则
$$
\xymatrix{
E \ar[d] \ar[r] & X' \ar[d] \\
E_k \ar[r] & X'_k
}
$$
是上一段所讨论类型的几乎爆破方块。因此
$$
\mathcal{F}(X'_k) = \mathcal{F}(X') \times_{\mathcal{F}(E)} \mathcal{F}(E_k)
$$
对 $k = 1, 2$ 成立。由此可知
$$
\mathcal{F}(X) \longrightarrow
\mathcal{F}(X'_k) \times_{\mathcal{F}(E_k)} \mathcal{F}(Z)
$$
对 $k = 1$ 为双射，当且仅当它对 $k = 2$ 为双射。
因此，给定有限表示的闭浸入 $Z \to X$，其中 $X$ 拟紧且拟分离，
等式
$\mathcal{F}(X) = \mathcal{F}(X') \times_{\mathcal{F}(E)} \mathcal{F}(Z)$
是否成立，与所选几乎爆破方块
(\ref{flat-equation-almost-blow-up-square}) 无关。
（而且由引理 \ref{flat-lemma-almost-blow-up-square}，
$Z \subset X$ 的几乎爆破方块确实存在。）

\medskip\noindent
最后，考虑 $(\Sch/S)_h$ 的仿射对象 $X$ 和有限表示闭浸入
$Z \to X$。对定义 $X$ 中 $Z$ 的理想的生成元数目 $r$ 归纳，
证明概形对 $(X, Z)$ 具有所需性质。若生成元数目为 $\leq 2$，
则可选取例 \ref{flat-example-two-generators} 中的几乎爆破方块，
并由假设 (4) 得到结论。

\medskip\noindent
归纳步骤。设 $X = \Spec(A)$ 且
$Z = \Spec(A/(f_1, \ldots, f_r))$，其中 $r > 2$。
为概形对 $(X, Z)$ 选取爆破方块
(\ref{flat-equation-almost-blow-up-square})。令
$Z_1 = \Spec(A/(f_1, f_2))$，并令
$$
\xymatrix{
E_1 \ar[d] \ar[r] & Y \ar[d] \\
Z_1 \ar[r] & X
}
$$
为例 \ref{flat-example-two-generators} 中构造的几乎爆破方块。
由引理 \ref{flat-lemma-base-change-almost-blow-up}，基变换
$$
(I)
\vcenter{
\xymatrix{
Y \times_X E \ar[r] \ar[d] & Y \times_X X' \ar[d] \\
Y \times_X Z \ar[r] & Y
}
}
\quad\text{且}\quad
(II)
\vcenter{
\xymatrix{
E \ar[r] \ar[d] & Z_1 \times_X X' \ar[d] \\
Z \ar[r] & Z_1
}
}
$$
都是几乎爆破方块。$Z_1$ 中 $Z$ 的理想由 $r - 2$ 个元素生成。
$Y \times_X Z$ 的理想由 $f_1, \ldots, f_r$ 到 $Y$ 的拉回生成。
在 $Y$ 上局部地，由 $f_1, f_2$ 生成的理想可由一个元素生成；
故 $Y \times_X Z$ 在 $Y$ 上仿射局部地由至多
$r - 1$ 个元素切出。由归纳假设和上面的讨论，
$$
\mathcal{F}(Y) =
\mathcal{F}(Y \times_X X') \times_{\mathcal{F}(Y \times_X E)}
\mathcal{F}(Y \times_X Z)
$$
并且
$$
\mathcal{F}(Z_1) =
\mathcal{F}(Z_1 \times_X X') \times_{\mathcal{F}(E)}
\mathcal{F}(Z)
$$
由假设 (4)，有
$$
\mathcal{F}(X) =
\mathcal{F}(Y) \times_{\mathcal{F}(E_1)} \mathcal{F}(Z_1)
$$
现设给定一对 $(s', t)$，其中 $s' \in \mathcal{F}(X')$、
$t \in \mathcal{F}(Z)$，且二者在 $\mathcal{F}(E)$ 中的限制相同。
则
% P05-EN-CH38-0083: frozen source writes s'|{...} rather than restriction syntax s'|_{...}.
$(s'|{Z_1 \times_X X'}, t)$ 是某个唯一元素
$t_1 \in \mathcal{F}(Z_1)$ 的像。同理，
$(s'|_{Y \times_X X'}, t|_{Y \times_X Z})$ 是某个唯一元素
$s_Y \in \mathcal{F}(Y)$ 的像。我们断言 $s_Y$ 与 $t_1$
在 $\mathcal{F}(E_1)$ 中限制为同一元素。这是因为几乎爆破方块
$$
\xymatrix{
E_1 \times_X E \ar[r] \ar[d] & E_1 \times_X X' \ar[d] \\
E_1 \times_X Z \ar[r] & E_1
}
$$
既是 (I) 沿 $E_1 \to Y$ 的基变换，又是 (II) 沿
$E_1 \to Z_1$ 的基变换；并且用于构造 $s_Y$ 与 $t_1$
的截面对彼此相符。因此，由第三个纤维积等式，存在唯一的
$s \in \mathcal{F}(X)$，它在 $\mathcal{F}(Y)$ 中映到 $s_Y$，
并在
% P05-EN-CH38-0084: frozen source says t_1 is reached in F(Z), but t_1 was constructed in F(Z_1) and the fibre product uses F(Z_1).
$\mathcal{F}(Z)$ 中映到 $t_1$。验证 $s$ 在
$\mathcal{F}(X')$ 中映到 $s'$，并在 $\mathcal{F}(Z)$ 中映到
$t$ 的过程从略；提示：使用刚构造的 $s$ 的唯一性，并在仿射局部工作。
\end{proof}

\begin{lemma}
\label{flat-lemma-refine-check-h-stack}
设 $p : \mathcal{S} \to (\Sch/S)_h$ 是群胚纤维化范畴。
则 $\mathcal{S}$ 是群胚叠，当且仅当满足下列条件：
\begin{enumerate}
\item $\mathcal{S}$ 是 Zariski 拓扑下的群胚叠；
\item 给定 $(\Sch/S)_h$ 中的态射 $f : X \to Y$，
其中 $Y$ 仿射，而 $f$ 满、平坦、适当且为有限表示，则
$$
\mathcal{S}_Y \longrightarrow
\mathcal{S}_X \times_{\mathcal{S}_{X \times_Y X}} \mathcal{S}_X
$$
是范畴等价；
\item 对范畴 $(\Sch/S)_h$ 中如例 \ref{flat-example-one-generator}
或 \ref{flat-example-two-generators} 的几乎爆破方块，函子
$$
\mathcal{S}_X \longrightarrow
\mathcal{S}_Z \times_{\mathcal{S}_E} \mathcal{S}_{X'}
$$
是范畴等价。
\end{enumerate}
\end{lemma}

\begin{proof}
本引理形式地由引理 \ref{flat-lemma-refine-check-h}
和我们对群胚叠的定义推出。例如，假设 (1)、(2)、(3)。
为证明 $\mathcal{S}$ 是叠，需证明态射和对象的下降；
参见《叠》，定义
\ref{stacks-definition-stack-in-groupoids}。

\medskip\noindent
若 $x, y$ 是 $\mathcal{S}$ 中位于 $(\Sch/S)_h$ 的对象 $U$
之上的对象，则我们的假设蕴含 $\mathit{Isom}(x, y)$
是 $(\Sch/U)_h$ 上满足引理 \ref{flat-lemma-refine-check-h}
中 (1)、(2)、(3)、(4) 的预层，因而是层。略去一些细节。

\medskip\noindent
设 $\{U_i \to U\}_{i \in I}$ 是 $(\Sch/S)_h$ 的覆盖，
$(x_i, \varphi_{ij})$ 是 $\mathcal{S}$ 中相对于态射族
$\{U_i \to U\}_{i \in I}$ 的下降资料；参见《叠》，定义
\ref{stacks-definition-descent-data}。考虑规则 $F$：
它把 $(\Sch/U)_h$ 中的 $V/U$ 映到由对 $(y, \psi_i)$ 组成的集合，
其中 $y$ 是 $\mathcal{S}_V$ 的对象，而
$\psi_i : y|_{U_i \times_U V} \to x_i|_{U_i \times_U V}$
是 $\mathcal{S}$ 中位于 $U_i \times_U V$ 之上的态射，并满足
$$
\varphi_{ij}|_{U_i \times_U U_j \times_U V}
\circ
\psi_i|_{U_i \times_U U_j \times_U V}
=
\psi_j|_{U_i \times_U U_j \times_U V}
$$
在同构意义下成立。由于已经有态射的下降，显然
$F(V/U)$ 或为空集，或为单点集。另一方面，
$F(U_{i_0}/U)$ 非空，因为它包含
$(x_{i_0}, \varphi_{i_0i})$。我们的目标是证明
$F(U/U)$ 非空，因此只需说明 $F$ 是 $(\Sch/U)_h$ 上的层。
可使用引理 \ref{flat-lemma-refine-check-h} 的判据。
而我们的假设 (1)、(2)、(3) 蕴含（画出一些略去的交换图即可）
引理 \ref{flat-lemma-refine-check-h} 的性质 (1)、(2)、(3)、(4)
对 $F$ 成立。

\medskip\noindent
若 $\mathcal{S}$ 是群胚叠，则 (1)、(2)、(3) 成立的验证从略。
\end{proof}





\section{绝对弱正规化与 h 覆盖}
\label{flat-section-h-and-weak-normalization}

\noindent
本节使用第 \ref{flat-section-blow-up-h} 节中的判据给出若干 h 层，
并将结构层的 h 层化与绝对弱正规化联系起来。
为此需要下述初等引理。

\begin{lemma}
\label{flat-lemma-funny-blow-up}
设 $Z, X, X', E$ 构成如例 \ref{flat-example-two-generators}
中的几乎爆破方块。则对 $p > 0$，
$H^p(X', \mathcal{O}_{X'}) = 0$，并且
$\Gamma(X, \mathcal{O}_X) \to \Gamma(X', \mathcal{O}_{X'})$
是环的满射，其核为平方零理想。
\end{lemma}

\begin{proof}
先假设 $A = \mathbf{Z}[f_1, f_2]$ 为多项式环。
此时，我们的几乎爆破方块就是 $X = \Spec(A)$ 沿闭子概形
$Z$ 的爆破；事实上，$X' \subset \mathbf{P}^1_X$
是由 $\mathcal{O}_{\mathbf{P}^1_X}(1)$ 的整体截面
$f_2T_0 - f_1 T_1$ 切出的有效 Cartier 除子。因此有消解
$$
0 \to
\mathcal{O}_{\mathbf{P}^1_X}(-1) \to
\mathcal{O}_{\mathbf{P}^1_X} \to
\mathcal{O}_{X'} \to
0
$$
使用《概形上同调》，第 \ref{coherent-section-cohomology-projective-space} 节
给出的上同调描述，可知在此情形下
$\Gamma(X, \mathcal{O}_X) \to \Gamma(X', \mathcal{O}_{X'})$
为同构，且 $H^1(X', \mathcal{O}_{X'}) = 0$。

\medskip\noindent
接着注意，如例 \ref{flat-example-two-generators} 中的任意图，
都是上一段之图沿环映射 $\mathbf{Z}[f_1, f_2] \to A$ 的基变换。
故由《态射进阶》，引理
\ref{more-morphisms-lemma-check-h1-fibre-zero},
\ref{more-morphisms-lemma-h1-fibre-zero} 和
\ref{more-morphisms-lemma-h1-fibre-zero-check-h0-kappa}
，可知一般地 $H^1(X', \mathcal{O}_{X'})$ 为零，并且映射
$H^0(X, \mathcal{O}_X) \to H^0(X', \mathcal{O}_{X'})$ 为满射。
% P05-EN-CH38-0085: frozen source ends this inference with the fragment “and the surjectivity ... in general”.

\medskip\noindent
接下来在一般情形中研究其核。若 $a \in A$ 映到零，
则在仿射坐标片上考察可得
$$
a = (f_1x - f_2)(a_0 + a_1x + \ldots + a_rx^r)\text{ 于 }A[x]\text{ 中}
$$
其中 $r \geq 0$ 且 $a_0, \ldots, a_r \in A$；同理
$$
a = (f_1 - f_2y)(b_0 + b_1y + \ldots + b_s y^s)\text{ 于 }A[y]\text{ 中}
$$
其中 $s \geq 0$ 且 $b_0, \ldots, b_s \in A$。这意味着
$$
a = f_2 a_0,\ f_1 a_0 = f_2 a_1,\ \ldots,\ f_1 a_r = 0,
\ a = f_1 b_0,\ f_2 b_0 = f_1 b_1,\ \ldots,\ f_2 b_s = 0
$$
若 $(a', r', a'_i, s', b'_j)$ 是另一组这样的数据，则
$$
aa' = f_1f_2a_0b'_0 = f_1f_2a_1b'_1 = f_1f_2a_2b'_2 = \ldots = 0
$$
这正合所需。
\end{proof}

\noindent
对 $\mathbf{F}_p$-代数 $A$，定义 $\colim_F A$ 为系统
$$
A \xrightarrow{F} A \xrightarrow{F} A \xrightarrow{F} \ldots
$$
的余极限，其中 $F : A \to A$、$a \mapsto a^p$
是 Frobenius 自同态。

\begin{lemma}
\label{flat-lemma-h-sheaf-colim-F}
设 $p$ 是素数，$S$ 是 $\mathbf{F}_p$ 上概形，
$(\Sch/S)_h$ 是定义 \ref{flat-definition-big-small-h} 中的站点。
则 $(\Sch/S)_h$ 上存在唯一的层 $\mathcal{F}$，使对
$(\Sch/S)_h$ 的任意拟紧且拟分离对象 $X$，
$$
\mathcal{F}(X) = \colim_F \Gamma(X, \mathcal{O}_X)
$$
\end{lemma}

\begin{proof}
将函子
$$
X \longrightarrow \colim_F \Gamma(X, \mathcal{O}_X)
$$
的 Zariski 层化记作 $\mathcal{F}$。
对于拟紧且拟分离概形 $X$，由《层》，引理
\ref{sheaves-lemma-directed-colimits-sections}
以及 $\mathcal{O}$ 是 Zariski 拓扑的层这一事实，有
$\mathcal{F}(X) = \colim_F \Gamma(X, \mathcal{O}_X)$。
% P05-EN-CH38-0086: source 11871--11874 has a full stop before lowercase ``by''.
因此，只需证明 $\mathcal{F}$ 是 h 层。为此，我们核验
引理 \ref{flat-lemma-refine-check-h} 的条件 (1)、(2)、(3) 和 (4)。
条件 (1) 成立，因为我们作了（几乎不必要的）Zariski 层化。
条件 (2) 成立，因为 $\mathcal{O}$ 是 fppf 层（《下降》，引理
\ref{descent-lemma-sheaf-condition-holds}），而且若 $A$ 是两个
$\mathbf{F}_p$-代数同态 $B \to C$ 的等化子，则 $\colim_F A$
是两个映射 $\colim_F B \to \colim_F C$ 的等化子。

\medskip\noindent
我们核验条件 (3)。令 $A, f, J$ 如
例 \ref{flat-example-one-generator} 中所述。我们必须证明
$$
\colim_F A = \colim_F A/J \times_{\colim_F A/fA + J} \colim_F A/fA
$$
这归结为如下代数问题：设 $a', a'' \in A$ 满足
$F^n(a' - a'') \in fA + J$。求出 $a \in A$ 和 $m \geq 0$，使得
$a - F^m(a') \in J$ 且 $a - F^m(a'') \in fA$，并证明序对
$(a, m)$ 在形如 $(a, m) \mapsto (F(a), m + 1)$ 的替换之下唯一确定。
为此，只须将 $F^n(a' - a'') = f h + g$ 写成 $h \in A$、$g \in J$
的形式，并令 $a = F^n(a') - g = F^n(a'') + fh$、$m = n$。
为证明唯一性，设 $(a_1, m_1)$ 是第二个解。作上述形式的替换后，
可以假设 $m = m_1$。于是 $a - a_1 \in J$ 且 $a - a_1 \in fA$。
由于 $J$ 被 $f$ 的某个幂零化，$a - a_1$ 是幂零元。因此对充分大的
$k$，$F^k(a - a_1)$ 为零。再作若干次替换后，便得到 $a = a_1$。

\medskip\noindent
我们核验条件 (4)。令 $X, X', Z, E$ 如
例 \ref{flat-example-two-generators} 中所述。由
引理 \ref{flat-lemma-funny-blow-up} 可知
$$
\mathcal{F}(X) = \colim_F \Gamma(X, \mathcal{O}_X)
\longrightarrow
\colim_F \Gamma(X', \mathcal{O}_{X'}) = \mathcal{F}(X')
$$
是双射。在此情形下 $E = \mathbf{P}^1_Z$，故
$\mathcal{F}(Z) \to \mathcal{F}(E)$ 也是双射。因此结论在此情形下同样成立。
\end{proof}

\noindent
设 $p$ 是素数。对于 $\mathbf{F}_p$-代数 $A$，令 $\lim_F A$
等于逆系统
$$
\ldots \xrightarrow{F} A \xrightarrow{F} A \xrightarrow{F} A
$$
的极限，其中 $F : A \to A$、$a \mapsto a^p$ 是 Frobenius 自同态。

\begin{lemma}
\label{flat-lemma-h-sheaf-lim-F}
设 $p$ 是素数，$S$ 是 $\mathbf{F}_p$ 上概形，
$(\Sch/S)_h$ 是定义 \ref{flat-definition-big-small-h} 中的站点。规则
$$
\mathcal{F}(X) = \lim_F \Gamma(X, \mathcal{O}_X)
$$
定义了 $(\Sch/S)_h$ 上的一个层。
\end{lemma}

\begin{proof}
为证明 $\mathcal{F}$ 是层，我们核验 引理 \ref{flat-lemma-refine-check-h}
的条件 (1)、(2)、(3) 和 (4)。条件 (1) 成立，因为层的极限仍是层，
且 $\mathcal{O}$ 是 Zariski 层。条件 (2) 成立，因为 $\mathcal{O}$
是 fppf 层（《下降》，引理
\ref{descent-lemma-sheaf-condition-holds}），而且若 $A$ 是两个
$\mathbf{F}_p$-代数同态 $B \to C$ 的等化子，则 $\lim_F A$
是两个映射 $\lim_F B \to \lim_F C$ 的等化子。

\medskip\noindent
我们核验条件 (3)。令 $A, f, J$ 如
例 \ref{flat-example-one-generator} 中所述。我们必须证明
\begin{align*}
\lim_F A
& \to
\lim_F A/J \times_{\lim_F A/fA + J} \lim_F A/fA \\
& =
\lim_F (A/J \times_{A/fA + J} A/fA) \\
& =
\lim_F A/(fA \cap J)
\end{align*}
是双射。由于 $J$ 被 $f$ 的某个幂零化，
$\mathfrak a = fA \cap J$ 是幂零理想，即存在 $n$ 使得
$\mathfrak a^n = 0$。不难验证，在此情形下
$\lim_F A \to \lim_F A/\mathfrak a$ 是双射。

\medskip\noindent
我们核验条件 (4)。令 $X, X', Z, E$ 如
例 \ref{flat-example-two-generators} 中所述。由
引理 \ref{flat-lemma-funny-blow-up} 和上述同一论证可知
$$
\mathcal{F}(X) = \lim_F \Gamma(X, \mathcal{O}_X)
\longrightarrow
\lim_F \Gamma(X', \mathcal{O}_{X'}) = \mathcal{F}(X')
$$
是双射。在此情形下 $E = \mathbf{P}^1_Z$，故
$\mathcal{F}(Z) \to \mathcal{F}(E)$ 也是双射。因此结论在此情形下同样成立。
\end{proof}

\noindent
在下一个引理中，我们使用概形 $X$ 的绝对弱正规化 $X^{awn}$；见
《态射》，第 \ref{morphisms-section-seminormalization} 节。

\begin{lemma}
\label{flat-lemma-weak-normalization-ph-sheaf}
设 $(\Sch/S)_{ph}$ 是 《拓扑》，定义
\ref{topologies-definition-big-small-ph} 中的站点。规则
$$
X \longmapsto \Gamma(X^{awn}, \mathcal{O}_{X^{awn}})
$$
是 $(\Sch/S)_{ph}$ 上的层。
\end{lemma}

\begin{proof}
为证明 $\mathcal{F}$ 是层，我们核验 《拓扑》，引理
\ref{topologies-lemma-characterize-sheaf} 的条件 (1) 和 (2)。
条件 (1) 成立，因为 $X^{awn}$ 的形成与开覆盖可交换；见《态射》，
引理 \ref{morphisms-lemma-seminormalization} 及其证明。

\medskip\noindent
设 $\pi : Y \to X$ 是满的固有态射。我们必须证明，两个映射
$$
\Gamma(Y^{awn}, \mathcal{O}_{Y^{awn}}) \to
\Gamma((Y \times_X Y)^{awn}, \mathcal{O}_{(Y \times_X Y)^{awn}})
$$
的等化子等于 $\Gamma(X^{awn}, \mathcal{O}_{X^{awn}})$。令 $f$
是该等化子的元素。考虑态射
$$
f : Y^{awn} \longrightarrow \mathbf{A}^1_X
$$
由于 $Y^{awn} \to X$ 是泛闭的，$f$ 的概形论像 $Z$ 是
$\mathbf{A}^1_X$ 的一个在 $X$ 上固有的闭子概形，并且
$f : Y^{awn} \to Z$ 是满态射；见《态射》，引理
\ref{morphisms-lemma-scheme-theoretic-image-is-proper}。因此 $Z \to X$
是有限满态射（《态射》，引理 \ref{morphisms-lemma-finite-proper}）。

\medskip\noindent
设 $k$ 是域，$z_1, z_2 : \Spec(k) \to Z$ 是两个经 $Z \to X$
等化的态射。我们断言 $z_1 = z_2$。只须证明像
$\lambda_i = z_i^*f \in k$ 相等（因为 $Z$ 的结构层在 $X$ 的结构层上
由 $f$ 生成）。为此，选取域扩张 $K/k$ 和态射
$y_1, y_2 : \Spec(K) \to Y^{awn}$，使得
$z_i \circ (\Spec(K) \to \Spec(k)) = f \circ y_i$。由于
$Y^{awn} \to Z$ 是满态射，这是可以做到的。选取一个基数非常大的
代数闭扩张 $\Omega/k$。对任意 $k$-代数同态
$\sigma_i : K \to \Omega$，得到
$$
\Spec(\Omega)
\xrightarrow{\sigma_1, \sigma_2}
\Spec(K \otimes_k K)
\xrightarrow{y_1, y_2}
Y^{awn} \times_X Y^{awn}
$$
由于典范态射
$(Y \times_X Y)^{awn} \to Y^{awn} \times_X Y^{awn}$ 是泛同胚，
且 $\Omega$ 代数闭，上述复合唯一提升为态射
$\Spec(\Omega) \to (Y \times_X Y)^{awn}$。由于 $f$ 属于上述等化子，
这证明 $\sigma_1(\lambda_1) = \sigma_2(\lambda_2)$。关于域扩张的一个
简单引理表明，这蕴含 $\lambda_1 = \lambda_2$；细节从略。

\medskip\noindent
由此可知 $Z \to X$ 是泛单射，即 $Z \to X$ 在点上是单射，并诱导
纯不可分的剩余域扩张（《态射》，引理
\ref{morphisms-lemma-universally-injective}）。
% P05-EN-CH38-0087: source 12060 says ``purely inseparated residue field extensions''.
综上，$Z \to X$ 是泛同胚；见《态射》，引理
\ref{morphisms-lemma-universal-homeomorphism}。

\medskip\noindent
令 $g : X^{awn} \to Z$ 是由 $X^{awn}$ 的泛性质得到的映射。于是
$Y^{awn} \to X^{awn} \to Z$ 与 $f : Y^{awn} \to Z$ 是两个 $X$-态射。
由 $Y^{awn} \to Y$ 的泛性质，相应的两个 $Y$-态射
$Y^{awn} \to Y \times_X Z$ 必须相等。这蕴含作为到
$\mathbf{A}^1_X$ 的态射有 $g \circ \pi^{wan} = f$，从而
$g \in \Gamma(X^{awn}, \mathcal{O}_{X^{awn}})$ 正是所求元素。
% P05-EN-CH38-0088: source 12071 uses \pi^{wan}, unlike the surrounding awn notation.
\end{proof}

\begin{lemma}
\label{flat-lemma-weak-normalization-h-sheaf}
设 $S$ 是概形。选取定义 \ref{flat-definition-big-small-h} 中的站点
$(\Sch/S)_h$。规则
$$
X \longmapsto \Gamma(X^{awn}, \mathcal{O}_{X^{awn}})
$$
是 $(\Sch/S)_h$ 上“结构层”$\mathcal{O}$ 的层化。对 ph 拓扑亦然。
\end{lemma}

\begin{proof}
在 引理 \ref{flat-lemma-weak-normalization-ph-sheaf} 中，我们已经看到
本引理的规则 $\mathcal{F}$ 定义了 ph 拓扑的层，因而更是 h 拓扑的层。
显然，存在典范的环预层态射 $\mathcal{O} \to \mathcal{F}$。
为完成证明，只须证明
\begin{enumerate}
\item 若 $f \in \mathcal{O}(X)$ 在 $\mathcal{F}(X)$ 中的像为零，
则存在 h 覆盖 $\{X_i \to X\}$，使得 $f|_{X_i} = 0$；
\item 给定 $f \in \mathcal{F}(X)$，存在 h 覆盖 $\{X_i \to X\}$，
使得 $f|_{X_i}$ 是某个 $f_i \in \mathcal{O}(X_i)$ 的像。
\end{enumerate}
令 $f$ 如 (1) 中所述。则 $f|_{X^{awn}} = 0$，这意味着 $f$ 局部幂零。
因此，若 $X' \subset X$ 是由 $f$ 截出的闭子概形，则 $X' \to X$
是有限表现的满闭浸入。故 $\{X' \to X\}$ 是所需的 h 覆盖。
令 $f$ 如 (2) 中所述。用一个仿射开覆盖的成员替换 $X$ 后，可以假设
$X = \Spec(A)$ 是仿射的。于是 $f \in A^{awn}$；见
《态射》，引理 \ref{morphisms-lemma-seminormalization-ring}。
由《态射》，引理 \ref{morphisms-lemma-universal-homeo-limit}，
可找到有限表现的环同态 $A \to B$，使得
$\Spec(B) \to \Spec(A)$ 是泛同胚，并且 $f$ 是某个 $b \in B$
在典范映射 $B \to A^{awn}$ 下的像。于是
$\{\Spec(B) \to \Spec(A)\}$ 是 h 覆盖，结论得证。关于 ph 拓扑的
断言可由相同方法推出（也可由 h 拓扑的断言推出）。
\end{proof}

\noindent
设 $p$ 是素数。若映射 $F : A \to A$、$x \mapsto x^p$ 是
$A$ 的自同构，则称 $\mathbf{F}_p$-代数 $A$ 是{\it 完美的}。

\begin{lemma}
\label{flat-lemma-perfect-weankly-normal}
设 $p$ 是素数。$\mathbf{F}_p$-代数 $A$ 绝对弱正规，当且仅当它是完美的。
\end{lemma}

\begin{proof}
由《态射》，定义 \ref{morphisms-definition-seminormal-ring}
中的条件 (2)(b) 立即可知，若 $A$ 绝对弱正规，则它是完美的。

\medskip\noindent
假设 $A$ 是完美的。设 $x, y \in A$ 且 $x^3 = y^2$。
若 $p > 3$，则可对某些 $n, m > 0$ 写成 $p = 2n + 3m$。
选取 $a, b \in A$，使得 $a^p = x$ 且 $b^p = y$。令
$c = a^n b^m$，则
$$
c^{2p} = x^{2n} y^{2m} = x^{2n + 3m} = x^p
$$
从而 $c^2 = x$。类似地，$c^3 = y$。若 $p = 2$，则将
$x = a^2$ 代入得到 $a^6 = y^2$，因而 $a^3 = y$。
若 $p = 3$，则将 $y = a^3$ 代入得到 $x^3 = a^6$，因而 $x = a^2$。

\medskip\noindent
设对某个素数 $\ell$，$x, y \in A$ 满足
$\ell^\ell x = y^\ell$。若 $\ell \not = p$，则 $a = y/\ell$
满足 $a^\ell = x$ 且 $\ell a = y$。若 $\ell = p$，则
$y = 0$，而对某个 $a$ 有 $x = a^p$。
\end{proof}

\begin{lemma}
\label{flat-lemma-char-p}
设 $p$ 是素数。
\begin{enumerate}
\item 若 $A$ 是 $\mathbf{F}_p$-代数，则 $\colim_F A = A^{awn}$。
\item 若 $S$ 是 $\mathbf{F}_p$ 上概形，则 $\mathcal{O}$ 的 h 层化
将拟紧且拟分离的 $X$ 映到 $\colim_F \Gamma(X, \mathcal{O}_X)$。
\end{enumerate}
\end{lemma}

\begin{proof}
证明 (1)。由《代数》，引理 \ref{algebra-lemma-p-ring-map}，
$A \to \colim_F A$ 在谱上诱导泛同胚。因此只须证明
$B = \colim_F A$ 绝对弱正规；见《态射》，引理
\ref{morphisms-lemma-seminormalization-ring}。注意，环同态
$F : B \to B$ 是自同构；换言之，$B$ 是完美环。因此可应用
引理 \ref{flat-lemma-perfect-weankly-normal}。

\medskip\noindent
证明 (2)。在仿射局部考察，由 (1) 以及 引理
\ref{flat-lemma-h-sheaf-colim-F} 和
\ref{flat-lemma-weak-normalization-h-sheaf} 即得。
\end{proof}






\section{正特征下向量丛的下降}
\label{flat-section-descent-vectorbundles-h}

\noindent
本节的参考文献是 \cite{Witt-Grass}。

\medskip\noindent
对概形 $S$，以 $\textit{Vect}(S)$ 表示有限局部自由
$\mathcal{O}_S$-模的范畴。设 $p$ 是素数，$S$ 是
$\mathbf{F}_p$ 上拟紧且拟分离的概形。本节将使用范畴
$$
\colim_F \textit{Vect}(S) =
\colim
\left(
\textit{Vect}(S) \xrightarrow{F^*}
\textit{Vect}(S) \xrightarrow{F^*}
\textit{Vect}(S) \xrightarrow{F^*} \ldots
\right)
$$
其中 $F : S \to S$ 是绝对 Frobenius 态射。具体而言，该范畴的对象是
序对 $(\mathcal{E}, n)$，其中 $\mathcal{E}$ 是有限局部自由
$\mathcal{O}_S$-模，$n \geq 0$ 是整数。其态射取为
% P05-EN-CH38-0089: source 12210--12211 uses unhyphenated “down to earth” attributively.
$$
\Hom_{\colim_F \textit{Vect}(S)}((\mathcal{E}, n), (\mathcal{G}, m)) =
\colim_N \Hom_S(F^{N - n, *}\mathcal{E}, F^{N - m, *}\mathcal{G})
$$
其中 $F : S \to S$ 是 $S$ 的绝对 Frobenius 态射。因此对象
$(\mathcal{E}, n)$ 同构于对象 $(F^*\mathcal{E}, n + 1)$。

\begin{lemma}
\label{flat-lemma-colim-F-Vect}
设 $p$ 是素数，$S$ 是 $\mathbf{F}_p$ 上拟紧且拟分离的概形。
范畴 $\colim_F \textit{Vect}(S)$ 等价于 $S$ 上环层
$\colim_F \mathcal{O}_S$ 的有限局部自由模范畴。
\end{lemma}

\begin{proof}
从略。
\end{proof}

\begin{lemma}
\label{flat-lemma-vector-bundle-I}
设 $p$ 是素数。考虑特征 $p$ 中如
例 \ref{flat-example-one-generator} 的几乎爆破方块 $X, X', Z, E$。
则函子
$$
\colim_F \textit{Vect}(X)
\longrightarrow
\colim_F \textit{Vect}(Z)
\times_{\colim_F \textit{Vect}(E)}
\colim_F \textit{Vect}(X')
$$
是等价。
\end{lemma}

\begin{proof}
令 $A, f, J$ 如 例 \ref{flat-example-one-generator} 中所述。
由于这里的所有概形都是仿射的，并且向量丛范畴中有内 Hom，
只要证明对有限投射 $A$-模 $P$ 有
% P05-EN-CH38-0090: source 12252 says “the fully faithfulness”; standard noun is “full faithfulness”.
$$
\colim P \otimes_{A, F^N} A =
\colim P \otimes_{A, F^N} A/J
\times_{\colim P \otimes_{A, F^N} A/fA + J}
\colim P \otimes_{A, F^N} A/fA
$$
由此即可推出该函子的全忠实性。将 $P$ 写成有限自由模的直和项后，
此式可由 $P$ 有限自由的情形推出，而该情形立即归结为 $P = A$。
$P = A$ 的情形由 引理 \ref{flat-lemma-h-sheaf-colim-F} 得到
（事实上，我们在该引理的证明中直接证明了这一情形）。

\medskip\noindent
本质满性。这里得到如下代数问题。设 $P_1$ 是有限投射 $A/J$-模，
$P_2$ 是有限投射 $A/fA$-模，并且
$$
\varphi :
P_1 \otimes_{A/J} A/fA + J
\longrightarrow
P_2 \otimes_{A/fA} A/fA + J
$$
是同构。目标是证明：存在 $N$、有限投射 $A$-模 $P$、同构
$\varphi_1 : P \otimes_A A/J \to P_1 \otimes_{A/J, F^N} A/J$,
以及同构
$\varphi_2 : P \otimes_A A/fA \to P_2 \otimes_{A/fA, F^N} A/fA$
，并且它们以显然的方式与 $\varphi$ 相容。证明如下。首先注意
$$
A/(J \cap fA) = A/J \times_{A/fA + J} A/fA
$$
因此由《代数进阶》，引理
\ref{more-algebra-lemma-finitely-presented-module-over-fibre-product}，
存在 $A/(J \cap fA)$ 上有限投射模 $P'$，连同同构
$\varphi'_1 : P' \otimes_A A/J \to P_1$ 和
$\varphi_2 : P' \otimes_A A/fA \to P_2$，它们与 $\varphi$ 相容。
% P05-EN-CH38-0091: source 12290--12291 primes only \varphi'_1 although both maps arise from P'.
由于 $J$ 是有限生成的 $f$-幂挠理想，$J \cap fA$ 是幂零理想。
因此对某个 $N$，$F^N$ 有分解
$$
A \xrightarrow{\alpha} A/(J \cap fA) \xrightarrow{\beta} A
$$
令 $P = P' \otimes_{A/(J \cap fA), \beta} A$，结论即得。
\end{proof}

\begin{lemma}
\label{flat-lemma-vector-bundle-II}
设 $p$ 是素数。考虑特征 $p$ 中如
例 \ref{flat-example-two-generators} 的几乎爆破方块 $X, X', Z, E$。
则函子
$$
G :
\colim_F \textit{Vect}(X)
\longrightarrow
\colim_F \textit{Vect}(Z)
\times_{\colim_F \textit{Vect}(E)}
\colim_F \textit{Vect}(X')
$$
是等价。
\end{lemma}

\begin{proof}
全忠实性。设 $(\mathcal{E}, n)$ 和 $(\mathcal{F}, m)$ 是
$\colim_F \textit{Vect}(X)$ 的对象。令
$(a, b) : G(\mathcal{E}, n) \to G(\mathcal{F}, m)$ 是右端中的态射。
可以选取 $N \gg 0$，并将 $a$ 视为映射
$a : F^{N - n, *}\mathcal{E}|_Z \to F^{N - m, *}\mathcal{F}|_Z$
，将 $b$ 视为映射
$b : F^{N - n, *}\mathcal{E}|_{X'} \to F^{N - m, *}\mathcal{F}|_{X'}$
，二者在 $E$ 上相同。选取有限仿射开覆盖
$X = X_1 \cup \ldots \cup X_n$，使得 $\mathcal{E}|_{X_i}$ 和
$\mathcal{F}|_{X_i}$ 都是有限自由 $\mathcal{O}_{X_i}$-模。
对每个 $i$，基变换
$$
\xymatrix{
E_i \ar[r] \ar[d] & X'_i \ar[d] \\
Z_i \ar[r] & X_i
}
$$
是另一个如 例 \ref{flat-example-two-generators} 的几乎爆破方块。
对于这些方块，由 引理 \ref{flat-lemma-h-sheaf-colim-F}
（见该引理的证明）可知
$$
\colim_F H^0(X_i, \mathcal{O}_{X_i}) =
\colim_F H^0(Z_i, \mathcal{O}_{Z_i})
\times_{\colim_F H^0(E_i, \mathcal{O}_{E_i})}
\colim_F H^0(X'_i, \mathcal{O}_{X'_i})
$$
因此增大 $N$ 后，可以假设映射 $a|_{Z_i}$ 和 $b|_{X'_i}$ 来自映射
$c_i : F^{N - n, *}\mathcal{E}|_{X_i} \to F^{N - m, *}\mathcal{F}|_{X_i}$。
必要时再增大 $N$，可以假设 $c_i$ 与 $c_j$ 在 $X_i \cap X_j$ 上相同。
故这些映射黏合成左端中所需的态射
$(\mathcal{E}, n) \to (\mathcal{F}, m)$。

\medskip\noindent
本质满性。令 $(\mathcal{F}, \mathcal{G}, \varphi)$ 是由有限局部自由
$\mathcal{O}_Z$-模 $\mathcal{F}$、有限局部自由
$\mathcal{O}_{X'}$-模 $\mathcal{G}$ 以及同构
$\varphi : \mathcal{F}|_E \to \mathcal{G}|_E$ 构成的三元组。
我们必须证明，用某个 Frobenius 幂拉回替换该三元组后，它来自一个
有限局部自由 $\mathcal{O}_X$-模。

\medskip\noindent
Noether 情形的归约；建议读者跳过本段。回顾
$X = \Spec(A)$、$Z = \Spec(A/(f_1, f_2))$、
$X' = \text{Proj}(A[T_0, T_1]/(f_2T_0 - f_1T_1))$、
$E = \mathbf{P}^1_Z$。由《极限》，引理
\ref{limits-lemma-descend-invertible-modules}，可找到包含 $f_1$ 和 $f_2$
的有限生成 $\mathbf{F}_p$-子代数 $A_0 \subset A$，使三元组
$(\mathcal{F}, \mathcal{G}, \varphi)$ 下降到
$X_0 = \Spec(A_0)$、$Z_0 = \Spec(A_0/(f_1, f_2))$、
$X_0' = \text{Proj}(A_0[T_0, T_1]/(f_2T_0 - f_1T_1))$、
$E_0 = \mathbf{P}^1_{Z_0}$。因此可以假设这些概形都是 Noether 的。

\medskip\noindent
假设 $X$ 是 Noether 的。可以选取有限仿射开覆盖
$X = X_1 \cup \ldots \cup X_n$，使 $\mathcal{F}|_{Z \cap X_i}$ 自由。
由于可以在 Zariski 拓扑中黏合 $\colim_F \textit{Vect}(X)$ 的对象
（引理 \ref{flat-lemma-colim-F-Vect}），并且我们已经知道在 $X_i$ 和
$X_i \cap X_j$ 上的全忠实性（见证明第一段），只须在每个 $X_i$
上证明存在性。这就归结到下一段讨论的情形。

\medskip\noindent
假设 $X$ 是 Noether 的，且 $\mathcal{F} = \mathcal{O}_Z^{\oplus r}$。
利用 $\varphi$ 得到同构 $\mathcal{O}_E^{\oplus r} \to \mathcal{G}|_E$。
令 $I = (f_1, f_2) \subset A$。令
$\mathcal{I} \subset \mathcal{O}_{X'}$ 是 $E$ 的理想层；它由
$f_1$ 和 $f_2$ 整体生成。对任意 $n$，存在满射
$$
(\mathcal{I}^n/\mathcal{I}^{n + 1})^{\oplus r} =
\mathcal{I}^n/\mathcal{I}^{n + 1} \otimes_{\mathcal{O}_E}
\mathcal{G}|_E \longrightarrow
\mathcal{I}^n\mathcal{G}/\mathcal{I}^{n + 1}\mathcal{G}
$$
因此该模的一阶上同调群为零。这里用到 $E = \mathbf{P}^1_Z$，故其结构层，
事实上任意整体生成的拟凝聚模，都有消失的 $H^1$。比较 《态射进阶》，
引理 \ref{more-morphisms-lemma-globally-generated-vanishing}。
再利用短正合列
$$
0 \to  \mathcal{I}^n\mathcal{G}/\mathcal{I}^{n + 1}\mathcal{G} \to
\mathcal{G}/\mathcal{I}^{n + 1}\mathcal{G} \to
\mathcal{G}/\mathcal{I}^n\mathcal{G} \to 0
$$
和归纳法，可知
$$
\lim H^0(X', \mathcal{G}/\mathcal{I}^n\mathcal{G})
\to
H^0(E, \mathcal{G}|_E) = H^0(E, \mathcal{O}_E^{\oplus r}) =
A/I^{\oplus r}
$$
是满射。由形式函数定理（《概形上同调》，定理
\ref{coherent-theorem-formal-functions}），这蕴含
$$
H^0(X', \mathcal{G}) \to
H^0(E, \mathcal{G}|_E) = H^0(E, \mathcal{O}_E^{\oplus r}) =
A/I^{\oplus r}
$$
是满射。因此可选取与给定的 $\mathcal{G}|_E$ 平凡化相容的映射
$\alpha : \mathcal{O}_{X'}^{\oplus r} \to \mathcal{G}$。
故 $\alpha$ 在 $X'$ 中 $E$ 的某个开邻域上是同构。于是 $Z$ 的每一点都有
一个可在其上解决问题的仿射开邻域。由于
$X' \setminus E \to X \setminus Z$ 是同构，对 $X$ 中不在 $Z$ 内的点亦然。
再作一次 Zariski 黏合即完成证明。
% P05-EN-CH38-0092: source 12418 and 12426 typesets A/I^{\oplus r}; the direct-sum exponent should apply to A/I.
\end{proof}

\begin{proposition}
\label{flat-proposition-h-descent-vector-bundles-p}
设 $p$ 是素数，$S$ 是特征 $p$ 的概形。则群胚纤维化范畴
$$
p : \mathcal{S} \longrightarrow (\Sch/S)_h
$$
是群胚叠，其中 $U$ 上的纤维范畴是 $U$ 上有限局部自由
$\colim_F \mathcal{O}_U$-模的范畴。此外，若 $U$ 拟紧且拟分离，
则 $\mathcal{S}_U$ 是 $\colim_F \textit{Vect}(U)$。
\end{proposition}

\begin{proof}
最后一个断言正是 引理 \ref{flat-lemma-colim-F-Vect} 的内容。
为证明该命题，我们核验 引理 \ref{flat-lemma-refine-check-h-stack}
的条件 (1)、(2) 和 (3)。

\medskip\noindent
条件 (1) 成立，因为按照定义，我们在 Zariski 拓扑中可以作黏合。

\medskip\noindent
为证明条件 (2)，设 $f : X \to Y$ 是 $S$ 上有限表现、满、平坦且固有的态射，
并且 $Y$ 仿射。由于 $Y, X, X \times_Y X$ 都拟紧且拟分离，
可以使用命题陈述中给出的纤维范畴描述。显然只须证明
$$
\textit{Vect}(Y)
\longrightarrow
\textit{Vect}(X)
\times_{\textit{Vect}(X \times_Y X)}
\textit{Vect}(X)
$$
是等价（因为这将蕴含余极限的相同结论）。这立即由有限局部自由模的
fppf 下降得到；见《下降》，命题
\ref{descent-proposition-fpqc-descent-quasi-coherent} 和
引理 \ref{descent-lemma-finite-locally-free-descends}。

\medskip\noindent
条件 (3) 正是 引理 \ref{flat-lemma-vector-bundle-I} 和
\ref{flat-lemma-vector-bundle-II} 的内容。
\end{proof}

\begin{lemma}
\label{flat-lemma-trivial-fibres-dvr}
设 $f : X \to S$ 是具有几何连通纤维的固有态射，其中 $S$ 是离散赋值环的谱。
以 $\eta \in S$ 表示一般点，并以 $X_n \subset X$ 表示由 $S$ 上某个
素元的第 $n$ 次幂截出的闭子概形。则存在整数 $n$，使得如下断言成立：
任意有限局部自由 $\mathcal{O}_X$-模 $\mathcal{E}$，只要
$\mathcal{E}|_{X_\eta}$ 和 $\mathcal{E}|_{X_n}$ 都自由，它就是自由的。
% P05-EN-CH38-0093: source 12490 writes the phrasal verb “cutout” as one word.
\end{lemma}

\begin{proof}
首先归约到 $X \to S$ 有截面的情形。记 $S = \Spec(A)$。
选取 $X_\eta$ 的闭点 $\xi$。选取离散赋值环扩张 $A \subset B$，
使得 $B$ 的分式域是 $\kappa(\xi)$。由 Krull--Akizuki
（《代数》，引理 \ref{algebra-lemma-integral-closure-Dedekind}）
以及 $\kappa(\xi)$ 是 $A$ 的分式域的有限扩张，这是可以做到的。
由固有性的赋值判据（《态射》，引理
\ref{morphisms-lemma-characterize-proper}），得到 $B$-值点
$\tau : \Spec(B) \to X$，它诱导截面 $\sigma : \Spec(B) \to X_B$。
对有限局部自由 $\mathcal{O}_X$-模 $\mathcal{E}$，令 $\mathcal{E}_B$
为它在基变换 $X_B$ 上的拉回。由平坦基变换（《概形上同调》，
引理 \ref{coherent-lemma-flat-base-change-cohomology}），有
$H^0(X_B, \mathcal{E}_B) = H^0(X, \mathcal{E}) \otimes_A B$。
因此，若 $\mathcal{E}_B$ 是秩 $r$ 的自由模，则
$H^0(X, \mathcal{E})$ 中的截面生成自由 $B$-模
$\tau^*\mathcal{E} = \sigma^*\mathcal{E}_B$。特别地，可以找到
$\mathcal{E}$ 的 $r$ 个整体截面 $s_1, \ldots, s_r$，它们生成
$\tau^*\mathcal{E}$。于是
$$
s_1, \ldots, s_r :
\mathcal{O}_X^{\oplus r}
\longrightarrow
\mathcal{E}
$$
是秩 $r$ 的有限局部自由 $\mathcal{O}_X$-模之间的映射；它在 $X_B$
上的拉回是自由 $\mathcal{O}_{X_B}$-模之间的映射，并且在每个纤维的
某一点处限制为同构。取行列式，得到函数
$g \in \Gamma(X_\eta, \mathcal{O}_{X_B})$，它在每个纤维的某一点处可逆。
由于纤维固有且连通，$g$ 必须可逆（细节从略；提示：使用《簇》，
引理 \ref{varieties-lemma-proper-geometrically-reduced-global-sections}）。
因此只须对基变换 $X_B \to \Spec(B)$ 证明本引理。
% P05-EN-CH38-0094: source 12530 places an \mathcal{O}_{X_B}-section in \Gamma(X_\eta, -), although the determinant is taken on X_B.

\medskip\noindent
假设已有截面 $\sigma : S \to X$。令 $\mathcal{E}$ 是有限局部自由
$\mathcal{O}_X$-模，并假设它在一般纤维和 $X_n$ 上自由
（稍后选取 $n$）。选取同构
$\sigma^*\mathcal{E} = \mathcal{O}_S^{\oplus r}$。考虑 $D(A)$ 中的映射
$$
K = R\Gamma(X, \mathcal{E}) \longrightarrow
R\Gamma(S, \sigma^*\mathcal{E}) = A^{\oplus r}
$$
同上论证可知，$\mathcal{E}$ 自由，当且仅当所诱导的映射
$H^0(K) = H^0(X, \mathcal{E}) \to A^{\oplus r}$ 是满射。

\medskip\noindent
令 $L = R\Gamma(X, \mathcal{O}_X^{\oplus r})$，并注意相应的映射
$L \to A^{\oplus r}$ 具有所需性质。由平坦基变换以及
$\mathcal{E}$ 在一般纤维上自由的假设，有
$K \otimes_A Q(A) \cong L \otimes_A Q(A)$。令 $\pi \in A$ 是素元。
注意
$$
K \otimes_A^\mathbf{L} A/\pi^m A =
R\Gamma(X, \mathcal{E} \xrightarrow{\pi^m} \mathcal{E})
$$
对 $L$ 亦然。以 $\mathcal{E}_{tors} \subset \mathcal{E}$ 表示由支集在
特殊纤维上的截面组成的凝聚子层，对其他 $\mathcal{O}_X$-模也作同样记号。
选取 $k > 0$，使得
$(\mathcal{O}_X)_{tors} \to \mathcal{O}_X/\pi^k \mathcal{O}_X$
是单射（《概形上同调》，引理
\ref{coherent-lemma-Artin-Rees}）。由于 $\mathcal{E}$ 局部自由，可知
$\mathcal{E}_{tors} \subset \mathcal{E}/\pi^k\mathcal{E}$。
于是对 $n \geq m + k$，在 $D(\mathcal{O}_X)$ 中有同构
\begin{align*}
(\mathcal{E} \xrightarrow{\pi^m} \mathcal{E})
& \cong
(\mathcal{E}/\pi^k\mathcal{E} \xrightarrow{\pi^m}
\mathcal{E}/\pi^{k + m}\mathcal{E}) \\
& \cong
(\mathcal{O}_X^{\oplus r}/\pi^k\mathcal{O}_X^{\oplus r} \xrightarrow{\pi^m}
\mathcal{O}_X^{\oplus r}/\pi^{k + m}\mathcal{O}_X^{\oplus r}) \\
& \cong
(\mathcal{O}_X^{\oplus r} \xrightarrow{\pi^m} \mathcal{O}_X^{\oplus r})
\end{align*}
这在 $D(A)$ 中确定一个同构
$$
K \otimes_A^\mathbf{L} A/\pi^m A \cong L \otimes_A^\mathbf{L} A/\pi^m A
$$
（当 $n \geq m + k$ 时成立）。注意，这些同构与沿 $\sigma$ 拉回相容，
因而特别可知
$K \otimes_A^\mathbf{L} A/\pi^m A \to (A/\pi^m A)^{\oplus r}$
在 $0$ 次上同调模上诱导满射（因为对 $L$ 而言这是成立的）。
% P05-EN-CH38-0095: source 12588 has the article error “an surjection”.
由于 $A$ 是离散赋值环，在 $D(A)$ 中有
$$
K \cong \bigoplus H^i(K)[-i]
\quad\text{且}\quad
L \cong
\bigoplus H^i(L)[-i]
$$
参见《代数进阶》，例
\ref{more-algebra-example-finite-injective-finite-global-dimension}。
上同调群 $H^i(K) = H^i(X, \mathcal{E})$ 和
$H^i(L) = H^i(X, \mathcal{O}_X)^{\oplus r}$
由《概形上同调》，引理
\ref{coherent-lemma-proper-over-affine-cohomology-finite} 是有限 $A$-模。
由《代数进阶》，引理
\ref{more-algebra-lemma-generalized-valuation-ring-modules}，这些模都是
循环模的直和。我们已经看到，$H^i(K)$ 与 $H^i(L)$ 的自由部分具有
相同的秩 $\beta_i$。接着注意
$$
H^i(L \otimes_A^\mathbf{L} A/\pi^m A) =
H^i(L)/\pi^m H^i(L) \oplus H^{i + 1}(L)[\pi^m]
$$
对 $K$ 亦然。令 $e$ 是满足如下条件的最大整数：对某个 $i$，
$A/\pi^eA$ 是 $H^i(X, \mathcal{O}_X)$（等价地，$H^i(L)$）的直和项。
取 $m = e + 1$，则 $H^i(L \otimes_A^\mathbf{L} A/\pi^m A)$ 是
$\beta_i$ 个 $A/\pi^m A$ 与另外若干循环模的直和，其中每个循环模都被
$\pi^e$ 零化。对 $K$ 作同样论证，并利用同构
% P05-EN-CH38-0096: source 12613--12615 does not define e when all relevant cohomology modules are torsion-free.
$K \otimes_A^\mathbf{L} A/\pi^m A \cong L \otimes_A^\mathbf{L} A/\pi^m A$
可知，$H^i(K)$ 不可能有形如 $A/\pi^l A$、其中 $l > e$ 的循环直和项。
（这还蕴含 $K$ 与 $L$ 作为 $D(A)$ 的对象同构，但我们不需要这一点。）
于是映射
$$
H^0(K \otimes^\mathbf{L}_A A/\pi^{e + 1} A) =
H^0(K)/\pi^{e + 1}H^0(K) \oplus H^1(K)[\pi^{e + 1}]
\longrightarrow
(A/\pi^{e + 1} A)^{\oplus r}
$$
为满射的唯一可能，是它在第一个直和项上为满射。这正是所要证明的。
（更准确地说，若有截面 $\sigma$，则引理陈述中的整数 $n$ 应等于
$k + e + 1$，其中 $k$ 与 $e$ 如上，并且只依赖于 $X$。）
\end{proof}

\begin{lemma}
\label{flat-lemma-trivial-over-dvrs}
设 $f : X \to S$ 是概形态射，$\mathcal{E}$ 是有限局部自由
$\mathcal{O}_X$-模。假设
\begin{enumerate}
\item $f$ 平坦且固有，并且 $\mathcal{O}_S = f_*\mathcal{O}_X$；
\item $S$ 是正规 Noether 概形；
\item 对每个余维 $1$ 的点 $s \in S$，$\mathcal{E}$ 在
$X \times_S \Spec(\mathcal{O}_{S, s})$ 上的拉回是自由的。
\end{enumerate}
则 $\mathcal{E}$ 同构于某个有限局部自由 $\mathcal{O}_S$-模的拉回。
\end{lemma}

\begin{proof}
我们证明典范映射
$$
\Phi : f^*f_*\mathcal{E} \longrightarrow \mathcal{E}
$$
是同构。由平坦基变换（《概形上同调》，引理
\ref{coherent-lemma-flat-base-change-cohomology}）以及假设 (1) 和 (3)，
对每个余维 $1$ 的点 $s \in S$，该映射在
$X \times_S \Spec(\mathcal{O}_{S, s})$ 上的拉回是同构。
由《因子》，引理 \ref{divisors-lemma-check-isomorphism-via-depth-and-ass}，
只须证明：对任意映到余维 $\geq 2$ 的点 $s \in S$ 的点 $x \in X$，有
$\text{depth}((f^*f_*\mathcal{E})_x) \geq 2$。由于 $f$ 平坦且
$(f^*f_*\mathcal{E})_x = (f_*\mathcal{E})_s \otimes_{\mathcal{O}_{S, s}}
\mathcal{O}_{X, x}$，只须证明 $\text{depth}((f_*\mathcal{E})_s) \geq 2$；见
《代数》，引理 \ref{algebra-lemma-apply-grothendieck}。
由于 $S$ 是正规 Noether 概形且 $\dim(\mathcal{O}_{S, s}) \geq 2$，
有 $\text{depth}(\mathcal{O}_{S, s}) \geq 2$；见《性质》，引理
\ref{properties-lemma-criterion-normal}。于是所需结论由《因子》，
引理 \ref{divisors-lemma-depth-pushforward} 得到。
\end{proof}

\noindent
利用上述结果，可以证明如下惊人的断言。

\begin{theorem}
\label{flat-theorem-pullback-trivial-fibres}
设 $p$ 是素数，$Y$ 是 $\mathbf{F}_p$ 上拟紧且拟分离的概形。
设 $f : X \to Y$ 是具有几何连通纤维的有限表现满固有态射。则函子
$$
\colim_F \textit{Vect}(Y) \longrightarrow \colim_F \textit{Vect}(X)
$$
全忠实，其本质像描述如下。令 $\mathcal{E}$ 是有限局部自由
$\mathcal{O}_X$-模。假设对所有 $y \in Y$，存在整数
$n_y, r_y \geq 0$，使得
% P05-EN-CH38-0097: source 12697 has singular “there exists integers”.
$$
F^{n_y, *}\mathcal{E}|_{X_{y, red}}
\cong
\mathcal{O}_{X_{y, red}}^{\oplus r_y}
$$
则对某个 $n \geq 0$，第 $n$ 次 Frobenius 幂拉回
$F^{n, *}\mathcal{E}$ 是某个有限局部自由 $\mathcal{O}_Y$-模的拉回。
\end{theorem}

\begin{proof}
证明全忠实性。由于 $Y$ 上的向量丛局部平凡，这归结为证明
% P05-EN-CH38-0098: source 12710 closes “vector bundles” into “vectorbundles”.
$$
\colim_F \Gamma(Y, \mathcal{O}_Y)
\longrightarrow
\colim_F \Gamma(X, \mathcal{O}_X)
$$
是双射。由于 $\{X \to Y\}$ 是 h 覆盖，只要证明两个映射
$$
\colim_F \Gamma(X, \mathcal{O}_X)
\longrightarrow
\colim_F \Gamma(X \times_Y X, \mathcal{O}_{X \times_Y X})
$$
相等，即由 引理 \ref{flat-lemma-h-sheaf-colim-F} 得到该结论。
令 $g \in \Gamma(X, \mathcal{O}_X)$，以 $g_1$ 和 $g_2$ 表示 $g$
在 $X \times_Y X$ 上的两个拉回。由于 $X_{y, red}$ 几何连通，
$H^0(X_{y, red}, \mathcal{O}_{X_{y, red}})$ 是 $\kappa(y)$ 的纯不可分扩张；
见《簇》，引理
\ref{varieties-lemma-proper-geometrically-reduced-global-sections}。
因此对某个 $p$ 的幂 $q$（它可以依赖于 $y$），
$g^q|_{X_{y, red}}$ 来自 $\kappa(y)$ 中的元素。于是对
$(X \times_Y X)_y = X_y \times_y X_y$ 的任意点，$g_1^q$ 和 $g_2^q$
在其剩余域中的像相同。因此 $g_1 - g_2$ 在
$(X \times_Y X)_{red}$ 上限制为零。故对某个 $n$ 有
$(g_1 - g_2)^n = 0$；按需要，可将它取为 $p$ 的幂。

\medskip\noindent
本质像的描述。令 $\mathcal{E}$ 如该命题陈述中所述。我们首先归约到
Noether 情形。
% P05-EN-CH38-0099: source 12741--12742 calls the enclosing theorem a proposition.

\medskip\noindent
令 $y \in Y$ 是一点，并将其视为从剩余域的谱到 $Y$ 的态射
$y \to Y$。可以将 $y \to Y$ 写成有限表现态射 $Y_i \to Y$ 的滤过极限，
其中 $Y_i$ 仿射。（最好自行证明这一点；它也可形式地由《极限》，
引理 \ref{limits-lemma-relative-approximation} 和
\ref{limits-lemma-limit-affine} 推出。）对每个 $i$，令
$Z_i = Y_i \times_Y X$。则 $X_y = \lim Z_i$ 且
$X_{y, red} = \lim Z_{i, red}$。由《极限》，引理
\ref{limits-lemma-descend-modules-finite-presentation}，可找到 $i$，使得
$F^{n_y, *}\mathcal{E}|_{Z_{i, red}} \cong
\mathcal{O}_{Z_{i, red}}^{\oplus r_y}$.
固定该 $i$。有 $Z_{i, red} = \lim Z_{i, j}$，其中
$Z_{i, j} \to Z_i$ 是有限表现的增厚（《极限》，引理
\ref{limits-lemma-closed-is-limit-closed-and-finite-presentation}）。
使用与此前相同的引理，可找到 $j$，使得
$F^{n_y, *}\mathcal{E}|_{Z_{i, j}} \cong \mathcal{O}_{Z_{i, j}}^{\oplus r_y}$.
由此可知，对每个 $y \in Y$，存在像包含 $y$ 的有限表现态射
$Y_y \to Y$ 以及增厚 $Z_y \to Y_y \times_Y X$，使得
$F^{n_y, *}\mathcal{E}|_{Z_y} \cong \mathcal{O}_{Z_y}^{\oplus r_y}$.
注意，$Y_y \to Y$ 的像是可构造的（《态射》，引理
\ref{morphisms-lemma-chevalley}）。由于 $Y$ 在可构造拓扑中拟紧
（《拓扑》，引理 \ref{topology-lemma-constructible-hausdorff-quasi-compact}
和 《性质》，引理
\ref{properties-lemma-quasi-compact-quasi-separated-spectral}），
存在有限多个有限表现态射
$$
Y_1 \to Y,\ Y_2 \to Y,\ \ldots,\ Y_N \to Y
$$
使得集合论意义下 $Y = \bigcup \Im(Y_a \to Y)$，并且对每个
$a \in \{1, \ldots, N\}$，存在整数 $n_a, r_a \geq 0$ 和有限表现的增厚
$Z_a \subset Y_a \times_Y X$，使得
$F^{n_a, *}\mathcal{E}|_{Z_a} \cong \mathcal{O}_{Z_a}^{\oplus r_a}$.

\medskip\noindent
按此方式表述后，该条件可下降到一个绝对 Noether 逼近。强烈建议读者跳过
本段。首先，将 $Y = \lim_{i \in I} Y_i$ 写成 $\mathbf{F}_p$ 上有限型概形
的余滤过极限，其中转移态射仿射（《极限》，引理
\ref{limits-lemma-relative-approximation}）。其次，可以假设已有固有态射
$f_i : X_i \to Y_i$，其到 $Y$ 的基变换恢复 $f : X \to Y$；见《极限》，
引理 \ref{limits-lemma-descend-finite-presentation}。增大 $i$ 后，可以假设
存在有限局部自由 $\mathcal{O}_{X_i}$-模 $\mathcal{E}_i$，它在 $X$ 上的
拉回同构于 $\mathcal{E}$；见《极限》，引理
\ref{limits-lemma-descend-invertible-modules}。
选取 $0 \in I$，以 $E \subset Y_0$ 表示 $f_0$ 的几何纤维连通的可构造子集；
见《态射进阶》，引理
\ref{more-morphisms-lemma-nr-geom-connected-components-constructible}。
于是 $Y \to Y_0$ 映入 $E$；见《态射进阶》，引理
\ref{more-morphisms-lemma-base-change-fibres-geometrically-connected}。
故当 $i \gg 0$ 时，$Y_i \to Y_0$ 映入 $E$；见《极限》，引理
\ref{limits-lemma-limit-contained-in-constructible}。因此当 $i \gg 0$ 时，
$f_i$ 的纤维几何连通。由《极限》，引理
\ref{limits-lemma-descend-finite-presentation}，对充分大的 $i$，可以找到
有限型态射 $Y_{i, a} \to Y_i$，其到 $Y$ 的基变换恢复
$Y_a \to Y$，其中 $a  \in \{1, \ldots, N\}$。必要时增大 $i$ 后，
可以找到增厚 $Z_{i, a} \subset Y_{i, a} \times_{Y_i} X_i$，其到
$Y_a \times_Y X$ 的基变换恢复 $Z_a$（使用与此前相同的参考结果，并结合
《极限》，引理 \ref{limits-lemma-descend-closed-immersion-finite-presentation}
和 \ref{limits-lemma-descend-surjective}）。由于
$Z_a = \lim Z_{i, a}$，再增大 $i$ 后，可以假设
$F^{n_a, *}\mathcal{E}_i|_{Z_{i, a}} \cong
\mathcal{O}_{Z_{i, a}}^{\oplus r_a}$；见《极限》，引理
\ref{limits-lemma-descend-modules-finite-presentation}。最后，再增大一次 $i$，
由《极限》，引理 \ref{limits-lemma-descend-surjective}，可以假设
$\coprod Y_{i, a} \to Y_i$ 是满态射。此时所有假设对
$X_i \to Y_i$ 和 $\mathcal{E}_i$ 都成立，因此只须对
$X_i \to Y_i$ 和 $\mathcal{E}_i$ 证明该结果。
% P05-EN-CH38-0100: source 12825 omits “the” in “prove result”.

\medskip\noindent
假设 $Y$ 在 $\mathbf{F}_p$ 上有限型。为证明该结果，我们对
$\dim(Y)$ 作归纳。我们要找到 $\colim_F \textit{Vect}(Y)$ 的一个对象，
其拉回是由 $\mathcal{E}$ 确定的 $\colim_F \textit{Vect}(X)$ 对象。
由已经证明的全忠实性以及 命题
\ref{flat-proposition-h-descent-vector-bundles-p}，用一个 h 覆盖的成员替换
$Y$ 并对 $X$ 作相应基变换后，只须构造 $\mathcal{E}$ 的下降。
这意味着，只要不增加 $Y$ 的维数，就可以用一个在 $Y$ 上固有且满的
概形替换 $Y$。若 $T \subset T'$ 是 $\mathbf{F}_p$ 上有限型概形的增厚，则
$\colim_F \textit{Vect}(T) = \colim_F \textit{Vect}(T')$
，因为 $\{T \to T'\}$ 是满足 $T \times_{T'} T = T$ 的 h 覆盖。
若 $T' \to T$ 是 $\mathbf{F}_p$ 上有限型概形之间的泛同胚，则
$\colim_F \textit{Vect}(T) = \colim_F \textit{Vect}(T')$
，因为 $\{T \to T'\}$ 是 h 覆盖，且对角态射
$T \subset T \times_{T'} T$ 是增厚。
% P05-EN-CH38-0101: source 12845--12849 reverses T and T' in both the h-cover arrow and the diagonal fibre product for T' -> T.

\medskip\noindent
利用上述一般说明，我们用其约化替换 $X$，并假设 $X$ 既约。
考虑 Stein 分解 $X \to Y' \to Y$；见《态射进阶》，定理
\ref{more-morphisms-theorem-stein-factorization-Noetherian}。
则 $Y' \to Y$ 是 $\mathbf{F}_p$ 上有限型概形之间的泛同胚。
由上述说明，可以用 $Y'$ 替换 $Y$。因此可以假设
$f_*\mathcal{O}_X = \mathcal{O}_Y$，且 $Y$ 既约。这就归结到下一段
讨论的情形。

\medskip\noindent
假设 $Y$ 既约，并且在 $Y$ 的某个稠密开子概形上有
$f_*\mathcal{O}_X = \mathcal{O}_Y$。则 $X \to Y$ 在某个稠密开子概形
$V \subset Y$ 上平坦；见《态射》，命题
\ref{morphisms-proposition-generic-flatness-reduced}。
由 引理 \ref{flat-lemma-flat-after-blowing-up}，存在 $V$-容许爆破
$Y' \to Y$，使 $X$ 的严格变换 $X'$ 在 $Y'$ 上平坦。由于 $Y$ 与 $Y'$
有公共稠密开子概形，$\dim(Y') = \dim(Y)$。由《态射进阶》，引理
\ref{more-morphisms-lemma-proper-flat-nr-geom-conn-comps-lower-semicontinuous}
以及 $V \subset Y'$ 稠密这一事实，$f' : X' \to Y'$ 的所有纤维都几何连通。
仍有 $(f'_*\mathcal{O}_{X'})|_V = \mathcal{O}_V$。写成
$$
Y' \times_Y X = X' \cup E \times_Y X
$$
其中 $E \subset Y'$ 是该爆破的例外除子。由上述一般说明，只须对
$Y' \times_Y X \to Y'$ 以及 $\mathcal{E}$ 在 $Y' \times_Y X$ 上的限制
证明存在性。设我们找到了 $\colim_F \textit{Vect}(Y')$ 中的对象
$\xi'$，它拉回到 $\mathcal{E}$ 在 $X'$ 上的限制（视为该余极限范畴的对象）。
对 $\dim(Y)$ 作归纳，可找到 $\colim_F \textit{Vect}(E)$ 中的对象
$\xi''$，它拉回到 $\mathcal{E}$ 在 $E \times_Y X$ 上的限制。
于是全忠实性确定唯一同构 $\xi'|_E \to \xi''$，它与给定的、到
$\mathcal{E}$ 在 $E \times_{Y'} X'$ 上限制的等同相容。由于
% P05-EN-CH38-0102: source 12892 misspells “full faithfulness” as “fully faithfullness”.
$$
\{E \times_Y X \to Y' \times_Y X, X' \to Y' \times_Y X\}
$$
是由一对闭浸入给出的 h 覆盖，且
$$
(E \times_Y X) \times_{(Y' \times_Y X)} X' = E \times_{Y'} X'
$$
由此可知 $\xi'$ 拉回到 $\mathcal{E}$ 在 $Y' \times_Y X$ 上的限制。
因此只须找到 $\xi'$，这就归结到下一段讨论的情形。

\medskip\noindent
假设 $Y$ 既约，$f$ 平坦，并且在 $Y$ 的某个稠密开子概形上有
$f_*\mathcal{O}_X = \mathcal{O}_Y$。在此情形下，考虑正规化
$Y^\nu \to Y$（《态射》，第 \ref{morphisms-section-normalization} 节）。这是在某个稠密开集上为同构的
有限满态射（《态射》，引理
\ref{morphisms-lemma-nagata-normalization} 和
\ref{morphisms-lemma-ubiquity-nagata}）。因此由上述一般说明，可以用
$Y^\nu$ 替换 $Y$，并用 $Y^\nu \times_Y X$ 替换 $X$。替换后有
$\mathcal{O}_Y = f_*\mathcal{O}_X$（因为在此情形下 Stein 分解必为同构；
略去一个小细节）。

\medskip\noindent
假设 $Y$ 是正规 Noether 概形，$f$ 平坦，且
$f_*\mathcal{O}_X = \mathcal{O}_Y$。用适当的 Frobenius 幂拉回替换
$\mathcal{E}$ 后，可以假设 $\mathcal{E}$ 在 $f$ 位于 $Y$ 各不可约分支
一般点处的概形论纤维上平凡（因为 $\colim_F \textit{Vect}(-)$ 在
泛同胚上是等价；见上文）。与上述论证（归约到 Noether 情形时）类似，
可知存在稠密开子概形 $V \subset Y$，使得
$\mathcal{E}|_{f^{-1}(V)}$ 自由。令 $Z \subset Y$ 是一个闭子概形，
使集合论意义下 $Y = V \amalg Z$。令 $z_1, \ldots, z_t \in Z$ 是 $Z$
中余维 $1$ 的不可约分支的一般点。则 $A_i = \mathcal{O}_{Y, z_i}$
是离散赋值环。令 $n_i$ 是 引理 \ref{flat-lemma-trivial-fibres-dvr}
对 $A_i$ 上概形 $X_{A_i}$ 所给出的整数。用适当的 Frobenius 幂拉回替换
$\mathcal{E}$ 后，可以假设 $\mathcal{E}$ 在
$X_{A_i/\mathfrak m_i^{n_i}}$ 上自由（同样因为余极限范畴在泛同胚下不变；
见上文）。于是 引理 \ref{flat-lemma-trivial-fibres-dvr} 表明
$\mathcal{E}$ 在 $X_{A_i}$ 上自由。最后应用
引理 \ref{flat-lemma-trivial-over-dvrs} 即得结论。
\end{proof}










\section{复形的爆破}
\label{flat-section-blowup-complexes}

\noindent
本节寻找导出范畴的完美对象在爆破后的正规形。

\begin{lemma}
\label{flat-lemma-fitting-ideals-complex}
设 $X$ 是概形，$E \in D(\mathcal{O}_X)$ 是伪凝聚的。对每个
$p, k \in \mathbf{Z}$，存在有限型拟凝聚理想层
$\text{Fit}_{p, k}(E) \subset \mathcal{O}_X$，具有如下性质：
若 $U \subset X$ 是开子集，且 $E|_U$ 同构于
% P05-EN-CH38-0103: source 12965 has the article error “an finite type”.
$$
\ldots \to
\mathcal{O}_U^{\oplus n_{b - 2}}
\xrightarrow{d_{b - 2}}
\mathcal{O}_U^{\oplus n_{b - 1}}
\xrightarrow{d_{b - 1}}
\mathcal{O}_U^{\oplus n_b} \to 0 \to \ldots
$$
则限制 $\text{Fit}_{p, k}(E)|_U$ 由 $d_p$ 的矩阵中如下阶数的
子式生成：
$$
- k + n_{p + 1} - n_{p + 2} + \ldots + (-1)^{b - p + 1} n_b
$$
约定：若 $r \leq 0$，则由 $r \times r$ 子式生成的理想是
$\mathcal{O}_U$；若 $r > \min(n_p, n_{p + 1})$，则由
$r \times r$ 子式生成的理想为零。
\end{lemma}

\begin{proof}
注意，$E$ 在 $X$ 上局部具有引理中所述的形式；见《代数进阶》，
第 \ref{more-algebra-section-pseudo-coherent} 节，《上同调》，
第 \ref{cohomology-section-pseudo-coherent} 节，以及
《概形的导出范畴》，第 \ref{perfect-section-spell-out} 节。
因此只须证明该子式理想与所选代表无关。为此，只须在局部环中核验。
在局部环 $(R, \mathfrak m, \kappa)$ 上，考虑上有界复形
$$
F^\bullet :
\ldots \to
R^{\oplus n_{b - 2}}
\xrightarrow{d_{b - 2}}
R^{\oplus n_{b - 1}}
\xrightarrow{d_{b - 1}}
R^{\oplus n_b} \to 0 \to \ldots
$$
以 $\text{Fit}_{k, p}(F^\bullet) \subset R$ 表示由 $d_p$ 的矩阵中
阶数为 $k - n_{p + 1} + n_{p + 2} - \ldots + (-1)^{b - p} n_b$ 的
子式生成的理想。假设 $F^\bullet$ 某个微分的某个矩阵系数可逆。
选取最大的整数 $i$，使 $d_i$ 有可逆的矩阵系数。由《代数》，引理
\ref{algebra-lemma-add-trivial-complex}，复形 $F^\bullet$ 同构于如下直和：
一个仅在次数 $i$ 和 $i + 1$ 有非零项的平凡复形
$\ldots \to 0 \to R \to R \to 0 \to \ldots$，以及一个复形
$(F')^\bullet$。留给读者验证
% P05-EN-CH38-0104: source 13005--13006 swaps the Fit indices and reverses the stated minor-size formula.
$\text{Fit}_{p, k}(F^\bullet) = \text{Fit}_{p, k}((F')^\bullet)$;
；这里用到了子式阶数的公式。若 $(F')^\bullet$ 还有某个微分具有可逆的
矩阵系数，就再次作同样处理，依此类推。
继续这一过程，最终得到复形 $(F^\infty)^\bullet$，其所有微分矩阵的系数
都在 $\mathfrak m$ 中。这里可能需要作无穷多步，但对任意截断，
只有有限多步会影响次数 $\geq $ 该截断的复形。因此，“极限”
$(F^\infty)^\bullet$ 是有限自由模的良定义上有界复形，并配有到起始复形的
拟同构 $(F^\infty)^\bullet \to F^\bullet$，且
% P05-EN-CH38-0105: source 13023 leaves “the cutoff” outside a malformed math fragment “degrees $\geq $”.
$\text{Fit}_{p, k}(F^\bullet) = \text{Fit}_{p, k}((F^\infty)^\bullet)$.
由《代数进阶》，引理
\ref{more-algebra-lemma-lift-pseudo-coherent-from-residue-field}，
复形 $(F^\infty)^\bullet$ 在同构意义下唯一，证明完成。
\end{proof}

\begin{lemma}
\label{flat-lemma-blowup-complex}
设 $X$ 是概形，$E \in D(\mathcal{O}_X)$ 是完美的。设
$U \subset X$ 是概形论稠密的开子概形，使得对所有
$i \in \mathbf{Z}$，$H^i(E|_U)$ 都是常秩 $r_i$ 的有限局部自由模。
则存在 $U$-容许爆破 $b : X' \to X$，使得对所有
$i \in \mathbf{Z}$，$H^i(Lb^*E)$ 都是 Tor 维数 $\leq 1$ 的完美
$\mathcal{O}_{X'}$-模。
% P05-EN-CH38-0106: source 13038, 13084, and 13129 use “scheme theoretically” without the compound-adverb hyphen.
\end{lemma}

\begin{proof}
我们在仿射局部构造并研究该爆破。具体而言，设 $V \subset X$ 是仿射开子概形，
使得 $E|_V$ 可由复形
$$
\mathcal{O}_V^{\oplus n_a} \xrightarrow{d_a}
\ldots \xrightarrow{d_{b - 1}} \mathcal{O}_V^{\oplus n_b}
$$
表示。令 $k_ i = r_{i + 1} - r_{i + 2} + \ldots + (-1)^{b - i + 1}r_b$。
略去的一个计算表明，$d_i$ 在 $U \cap V$ 上的秩是
% P05-EN-CH38-0107: source 13054 has a stray space in the subscript token k_ i.
% P05-EN-CH38-0108: source 13055 has singular “A computation” with bare “show”.
$$
\rho_i = - k_i + n_{i + 1} - n_{i + 2} + \ldots + (-1)^{b - i + 1}n_b
$$
这里的含义是 $d_i$ 的余核是秩 $n_{i + 1} - \rho_i$ 的有限局部自由模。
令 $\mathcal{I}_i \subset \mathcal{O}_V$ 是由 $d_i$ 的矩阵中
$\rho_i \times \rho_i$ 子式生成的理想。

\medskip\noindent
一方面，与 引理 \ref{flat-lemma-fitting-ideals-complex} 比较可知，
理想 $\mathcal{I}_i$ 对应于整体理想 $\text{Fit}_{i, k_i}(E)$；已经证明
后者与表示 $E|_V$ 的复形的选择无关。另一方面，$\mathcal{I}_i$ 是
$\Coker(d_i)$ 的第 $(n_{i + 1} - \rho_i)$ 个 Fitting 理想。请记住这一点。

\medskip\noindent
令 $b : X' \to X$ 是沿理想 $\text{Fit}_{i, k_i}(E)$ 之积的爆破；
这是有意义的，因为在 $X$ 上局部，这些理想中几乎全都等于单位理想
（见上文）。该爆破控制沿理想 $\text{Fit}_{i, k_i}(E)$ 的各个爆破
$b_i : X'_i \to X$；见《因子》，引理
\ref{divisors-lemma-blowing-up-two-ideals}。由《因子》，引理
\ref{divisors-lemma-blowup-fitting-ideal}，每个 $b_i$ 都是 $U$-容许爆破。
因此 $b$ 是 $U$-容许爆破（略去一个小细节；比较 《因子》，引理
\ref{divisors-lemma-dominate-admissible-blowups} 的证明）。最后，$U$
仍是 $X'$ 的概形论稠密开子概形。因此用 $X'$ 替换 $X$ 后，便处于下一段
讨论的情形。

\medskip\noindent
假设对所有 $i$，$\text{Fit}_{i, k_i}(E)$ 都是可逆理想。选取仿射开集
$V$ 以及如上表示 $E|_V$ 的有限自由模复形。由《因子》，引理
\ref{divisors-lemma-blowup-fitting-ideal}，$\Coker(d_i)$ 的 Tor 维数
$\leq 1$。因此，作为从有限局部自由模到 Tor 维数 $\leq 1$ 的有限表现模
之映射的核，$\Im(d_i)$ 是有限局部自由的。于是 $\Ker(d_i)$ 也有限局部自由
（同一论证）。故短正合列
$$
0 \to \Im(d_{i - 1}) \to \Ker(d_i) \to H^i(E)|_V \to 0
$$
给出所需结论，证明完成。
\end{proof}

\begin{lemma}
\label{flat-lemma-blowup-complex-integral}
设 $X$ 是整概形，$E \in D(\mathcal{O}_X)$ 是完美的。则存在非空开集
$U \subset X$，使得对所有 $i \in \mathbf{Z}$，$H^i(E|_U)$ 都是
常秩 $r_i$ 的有限局部自由模；并且存在 $U$-容许爆破
$b : X' \to X$，使得对所有 $i \in \mathbf{Z}$，$H^i(Lb^*E)$ 都是
Tor 维数 $\leq 1$ 的完美 $\mathcal{O}_{X'}$-模。
\end{lemma}

\begin{proof}
强烈建议读者自行证明 $U$ 的存在性。令 $\eta \in X$ 是一般点。
$E$ 在 $\eta$ 上的限制在 $D(\kappa(\eta))$ 中同构于一个微分全为零的
有限维向量空间有限复形 $V^\bullet$。令
$r_i = \dim_{\kappa(\eta)} V^i$。由各项为
$\mathcal{O}_X^{\oplus r_i}$、微分全为零的复形表示的
$D(\mathcal{O}_X)$ 中完美对象 $E'$，拉回到 $\eta$ 后与 $E$ 同构。
因此由《概形的导出范畴》，引理
\ref{perfect-lemma-descend-relatively-perfect}，存在 $\eta$ 的开邻域 $U$，
使 $E|_U$ 与 $E'|_U$ 同构。这证明第一个断言。第二个断言由第一个断言和
引理 \ref{flat-lemma-blowup-complex} 推出，因为整概形 $X$ 中任意非空开集都
概形论稠密。
\end{proof}

\begin{remark}
\label{flat-remark-when-you-have-a-complex}
设 $X$ 是概形，$E \in D(\mathcal{O}_X)$ 是完美对象，并且对所有
$i \in \mathbf{Z}$，$H^i(E)$ 都是 Tor 维数 $\leq 1$ 的完美
$\mathcal{O}_X$-模。该性质有时可将关于 $E$ 的问题归结为关于
$H^i(E)$ 的问题。例如，设
$$
\mathcal{E}^a \xrightarrow{d^a} \ldots
\xrightarrow{d^{b - 2}} \mathcal{E}^{b - 1}
\xrightarrow{d^{b - 1}} \mathcal{E}^b
$$
是表示 $E$ 的有限局部自由 $\mathcal{O}_X$-模有界复形。则对所有 $i$，
$\Im(d^i)$ 和 $\Ker(d^i)$ 都是有限局部自由 $\mathcal{O}_X$-模。
事实上，归纳假设对所有大于 $i$ 的指标均已知该结论。首先利用短正合列
$$
0 \to \Im(d^i) \to \Ker(d^{i + 1}) \to H^{i + 1}(E) \to 0
$$
以及 $H^{i + 1}(E)$ 是 Tor 维数 $\leq 1$ 的完美模这一假设，可知
$\Im(d^i)$ 有限局部自由。对短正合列
$$
0 \to \Ker(d^i) \to \mathcal{E}^i \to \Im(d^i) \to 0
$$
再次使用相同论证，得到 $\Ker(d^i)$ 有限局部自由。因此，特异三角
$$
\tau_{\leq k - 1}E \to \tau_{\leq k}E \to H^k(E)[-k] \to
(\tau_{\leq k - 1}E)[1]
$$
由如下有限局部自由模有界复形的短正合列表示：
$$
\begin{matrix}
& &
& &
& &
0 \\
& &
& &
& &
\downarrow \\
\mathcal{E}^a & \to &
\ldots & \to &
\mathcal{E}^{k - 2} & \to &
\Ker(d^{k - 1}) \\
\downarrow & &
& &
\downarrow & &
\downarrow \\
\mathcal{E}^a & \to &
\ldots & \to &
\mathcal{E}^{k - 2} & \to &
\mathcal{E}^{k - 1} & \to &
\Ker(d^k) \\
& &
& &
& &
\downarrow & &
\downarrow \\
& &
& &
& &
\Im(d^{k - 1}) & \to &
\Ker(d^k) \\
& &
& &
& &
\downarrow \\
& &
& &
& &
0
\end{matrix}
$$
这里各行是复形，图中省略了“显然”的零项。
\end{remark}





\section{完美模的爆破}
\label{flat-section-blowup-perfect-modules}

\noindent
本节试图寻找 Tor 维数 $\leq 1$ 的完美模在爆破后的正规形。
我们只取得了部分成功。

\begin{lemma}
\label{flat-lemma-blowup-pd1-derived}
设 $X$ 是概形，$\mathcal{F}$ 是 Tor 维数 $\leq 1$ 的完美
$\mathcal{O}_X$-模。对任意爆破 $b : X' \to X$，有
$Lb^*\mathcal{F} = b^*\mathcal{F}$，并且 $b^*\mathcal{F}$ 是
Tor 维数 $\leq 1$ 的完美 $\mathcal{O}_X$-模。
% P05-EN-CH38-0109: source 13228 and 13256 retain \mathcal{O}_X although the modules live on X'.
\end{lemma}

\begin{proof}
可以假设 $X = \Spec(A)$ 仿射，并且对应于 $\mathcal{F}$ 的 $A$-模 $M$
有一个表现
$$
0 \to A^{\oplus m} \to A^{\oplus n} \to M \to 0
$$
设 $I \subset A$ 是理想且 $a \in I$。回顾仿射爆破代数
$A[\frac{I}{a}]$ 是 $A_a$ 的子环。由于局部化正合，
$A_a^{\oplus m} \to A_a^{\oplus n}$ 是单射。因此
$A[\frac{I}{a}]^{\oplus m} \to A[\frac{I}{a}]^{\oplus n}$
也是单射。这证明了引理。
\end{proof}

\begin{lemma}
\label{flat-lemma-blowup-pd1}
设 $X$ 是概形，$\mathcal{F}$ 是 Tor 维数 $\leq 1$ 的完美
$\mathcal{O}_X$-模。设 $U \subset X$ 是概形论稠密的开集，使得
$\mathcal{F}|_U$ 是常秩 $r$ 的有限局部自由模。则存在 $U$-容许爆破
$b : X' \to X$，以及典范短正合列
% P05-EN-CH38-0110: source 13248 and 13370 add two further unhyphenated “scheme theoretically dense” loci.
$$
0 \to \mathcal{K} \to b^*\mathcal{F} \to \mathcal{Q} \to 0
$$
其中 $\mathcal{Q}$ 是秩 $r$ 的有限局部自由模，而 $\mathcal{K}$ 是
Tor 维数 $\leq 1$ 的完美 $\mathcal{O}_X$-模，且其在 $U$ 上的限制为零。
\end{lemma}

\begin{proof}
结合 《因子》，引理 \ref{divisors-lemma-blowup-fitting-ideal} 和
引理 \ref{flat-lemma-blowup-pd1-derived} 即得。
\end{proof}

\begin{lemma}
\label{flat-lemma-canonical-blowup-torsion-pd1}
设 $X$ 是概形，$\mathcal{F}$ 是 Tor 维数 $\leq 1$ 的完美
$\mathcal{O}_X$-模。设 $U \subset X$ 是满足 $\mathcal{F}|_U = 0$
的开集。则存在 $U$-容许爆破
$$
b : X' \to X
$$
使得 $\mathcal{F}' = b^*\mathcal{F}$ 配有两个典范的局部有限滤过
$$
0 = F^0 \subset F^1 \subset F^2 \subset \ldots \subset \mathcal{F}'
\quad\text{且}\quad
\mathcal{F}' = F_1 \supset F_2 \supset F_3 \supset \ldots \supset 0
$$
并且对每个 $n \geq 1$，存在有效 Cartier 除子 $D_n \subset X'$，使得
$$
F^i/F^{i - 1}
\quad\text{且}\quad
F_i/F_{i + 1}
$$
在 $D_i$ 上是秩 $i$ 的有限局部自由模。
\end{lemma}

\begin{proof}
选取仿射开集 $V \subset X$，使得存在某个 $n$ 和某个矩阵 $A$，以及表现
$$
0 \to \mathcal{O}_V^{\oplus n} \xrightarrow{A} \mathcal{O}_V^{\oplus n} \to
\mathcal{F} \to 0
$$
我们将沿所有 $k \geq 0$ 的 Fitting 理想
$\text{Fit}_k(\mathcal{F})$ 之积作爆破。这是有意义的，因为在上述仿射
情形中，对 $k > n$ 有
$\text{Fit}_k(\mathcal{F})|_V = \mathcal{O}_V$。显然，这是
$U$-容许爆破。由《因子》，引理
\ref{divisors-lemma-blowing-up-two-ideals}，在 $X'$ 上各理想
$\text{Fit}_k(\mathcal{F})$ 都可逆。因此归结到下一段讨论的情形。

\medskip\noindent
假设对 $k \geq 0$，$\text{Fit}_k(\mathcal{F})$ 都是可逆理想。
若 $E_k \subset X$ 是由 $\text{Fit}_k(\mathcal{F})$ 定义的有效
Cartier 除子（$k \geq 0$），则引理陈述中的有效 Cartier 除子 $D_k$ 将满足
$$
E_k = D_{k + 1} + 2 D_{k + 2} + 3 D_{k + 3} + \ldots
$$
这是有意义的，因为族 $D_k$ 将是局部有限的。此外，该等式唯一确定
有效 Cartier 除子 $D_k$，因此只须在局部构造 $D_k$。

\medskip\noindent
选取仿射开集 $V \subset X$ 以及如上的 $\mathcal{F}|_V$ 表现。
我们对表现中的整数 $n$ 作归纳，构造这些除子和滤过。对 $k > n$，令
$D_k|_V = \emptyset$，并令 $D_n|V = E_{n - 1}|_V$。
% P05-EN-CH38-0111: source 13323 omits the restriction subscript in D_n|V.
缩小 $V$ 后，可以假设 $\text{Fit}_{n - 1}(\mathcal{F})|_V$
由单个非零因子 $f \in \Gamma(V, \mathcal{O}_V)$ 生成。
由于 $\text{Fit}_{n - 1}(\mathcal{F})|_V$ 是由 $A$ 的各项生成的理想，
存在 $\Gamma(V, \mathcal{O}_V)$ 中的矩阵 $A'$，使得 $A = fA'$。
由短正合列在 $V$ 上定义 $\mathcal{F}'$：
$$
0 \to \mathcal{O}_V^{\oplus n} \xrightarrow{A'} \mathcal{O}_V^{\oplus n} \to
\mathcal{F}' \to 0
$$
由于 $A'$ 的各项在 $\Gamma(V, \mathcal{O}_V)$ 中生成单位理想，
$\mathcal{F}'$ 在 $V$ 上局部有一个将 $n$ 减少 $1$ 的表现；见《代数》，
引理 \ref{algebra-lemma-add-trivial-complex}。还要注意
$f^{n - k}\text{Fit}_k(\mathcal{F}') = \text{Fit}_k(\mathcal{F})|_V$
对 $k = 0, \ldots, n$ 成立。因此对所有 $k$，
$\text{Fit}_k(\mathcal{F}')$ 都是可逆理想。由归纳可知，存在有效
Cartier 除子 $D'_k \subset V$，使得 $\mathcal{F}'$ 具有引理陈述中的
两个典范滤过。然后对 $k = 1, \ldots, n - 1$，令
$D_k|_V = D'_k$。注意，在这些选择下，上述有效 Cartier 除子间的等式成立。
最后构造滤过。事实上，有短正合列
$$
0 \to
\mathcal{O}_{D_n \cap V}^{\oplus n} \to
\mathcal{F} \to \mathcal{F}' \to 0
\quad\text{且}\quad
0 \to
\mathcal{F}' \to \mathcal{F} \to
\mathcal{O}_{D_n \cap V}^{\oplus n} \to 0
$$
它们来自 $A$ 的两个分解 $A = A'f = f A'$。这些序列是典范的，
因为第一个序列中的子模是
$\Ker(f : \mathcal{F} \to \mathcal{F})$，第二个序列中的商模是
$\Coker(f : \mathcal{F} \to \mathcal{F})$。
\end{proof}

\begin{lemma}
\label{flat-lemma-blowup-map-pd1}
设 $X$ 是概形。设 $\varphi : \mathcal{F} \to \mathcal{G}$ 是
Tor 维数 $\leq 1$ 的完美 $\mathcal{O}_X$-模之间的同态。
设 $U \subset X$ 是概形论稠密的开集，使得
$\mathcal{F}|_U = 0$ 且 $\mathcal{G}|_U = 0$。则存在 $U$-容许爆破
$b : X' \to X$，使得 $b^*\varphi$ 的核、像和余核都是
Tor 维数 $\leq 1$ 的完美 $\mathcal{O}_{X'}$-模。
% P05-EN-CH38-0112: source 13369 misspells “homomorphism” as “homorphism”.
\end{lemma}

\begin{proof}
这些假设表明 $D(\mathcal{O}_X)$ 的对象
$(\mathcal{F} \to \mathcal{G})$ 是完美的。因此由 引理
\ref{flat-lemma-blowup-complex} 和 \ref{flat-lemma-blowup-pd1-derived}
（后者用来考察拉回后的复形形状），得到一个对余核和核都适用的
$U$-容许爆破。像是到余核之映射的核，因而也是 Tor 维数
$\leq 1$ 的完美模。
\end{proof}











\section{Berthelot 与 Ogus 引入的算子}
\label{flat-section-eta}

\noindent
请先阅读 《上同调》，第 \ref{cohomology-section-eta} 节。

\medskip\noindent
设 $X$ 是概形，$D \subset X$ 是有效 Cartier 除子。令
$\mathcal{I} = \mathcal{I}_D \subset \mathcal{O}_X$ 是 $D$ 的理想层；
见《因子》，第 \ref{divisors-section-effective-Cartier-invertible} 节。显然，可以将
《上同调》，第 \ref{cohomology-section-eta} 节的讨论应用于
$X$ 和 $\mathcal{I}$。

\begin{lemma}
\label{flat-lemma-eta-stalks}
设 $X$ 是概形，$D \subset X$ 是理想层为
$\mathcal{I} \subset \mathcal{O}_X$ 的有效 Cartier 除子。设
$\mathcal{F}^\bullet$ 是拟凝聚 $\mathcal{O}_X$-模复形，并且对所有 $i$，
$\mathcal{F}^i$ 都无 $\mathcal{I}$-挠。则
$\eta_\mathcal{I}\mathcal{F}^\bullet$ 是拟凝聚 $\mathcal{O}_X$-模复形。
此外，若 $U = \Spec(A) \subset X$ 是仿射开集且 $D \cap U = V(f)$，
则 $\eta_f(\mathcal{F}^\bullet(U))$ 典范同构于
$(\eta_\mathcal{I}\mathcal{F}^\bullet)(U)$。
\end{lemma}

\begin{proof}
从略。
\end{proof}

\begin{lemma}
\label{flat-lemma-Leta}
设 $X$ 是概形，$D \subset X$ 是理想层为
$\mathcal{I} \subset \mathcal{O}_X$ 的有效 Cartier 除子。《上同调》，
引理 \ref{cohomology-lemma-Leta} 的函子
$L\eta_\mathcal{I} : D(\mathcal{O}_X) \to D(\mathcal{O}_X)$
将 $D_\QCoh(\mathcal{O}_X)$ 映入自身。此外，若
$X = \Spec(A)$ 仿射且 $D = V(f)$，则 《代数进阶》，引理
\ref{more-algebra-lemma-Leta} 所定义的 $D(A)$ 上函子 $L\eta_f$，
与 $D_\QCoh(\mathcal{O}_X)$ 上函子 $L\eta_\mathcal{I}$，经由
《概形的导出范畴》，引理
\ref{perfect-lemma-affine-compare-bounded} 的等价相对应。
\end{lemma}

\begin{proof}
从略。
\end{proof}





\section{复形的爆破，II}
\label{flat-section-blowup-complexes-II}

\noindent
本节的材料将用于在 第 \ref{flat-section-blowup-complexes-III} 节
中构造 Macpherson 图构造的一个版本。
 
\begin{situation}
\label{flat-situation-complex-and-divisor}
这里 $X$ 是概形，$D \subset X$ 是理想层为
$\mathcal{I} \subset \mathcal{O}_X$ 的有效 Cartier 除子，而
$M$ 是 $D(\mathcal{O}_X)$ 的完美对象。
\end{situation}

\noindent
令 $(X, D, M)$ 是 情形 \ref{flat-situation-complex-and-divisor}
中的三元组。考虑满足如下条件的仿射开集
$U = \Spec(A) \subset X$：
\begin{enumerate}
\item 对某个非零因子 $f \in A$，有 $D \cap U = V(f)$；
\item 存在表示 $M|_U$ 的有限自由 $A$-模有界复形 $M^\bullet$
（经由《概形的导出范畴》，引理
\ref{perfect-lemma-affine-compare-bounded} 的等价）。
\end{enumerate}
称 $(U, A, f, M^\bullet)$ 是 $(X, D, M)$ 的一个{\it 仿射图}。
考虑《代数进阶》，第 \ref{more-algebra-section-perfect-eta} 节中定义的理想
$I_i(M^\bullet, f) \subset A$。若对每个 $x \in D$，存在满足
$x \in U$ 的仿射图 $(U, A, f, M^\bullet)$，使得对所有 $i \in \mathbf{Z}$，
$I_i(M^\bullet, f)$ 都是主理想，则称 $(X, S, M)$ 是{\it 良三元组}。
% P05-EN-CH38-0113: source 13482 changes D to the undefined S in the good-triple notation.

\begin{lemma}
\label{flat-lemma-pullback-triple-ideals-good}
在 情形 \ref{flat-situation-complex-and-divisor} 中，设 $h : Y \to X$
是概形态射，使得 $D$ 的拉回 $E = h^{-1}D$ 有定义（《因子》，
定义 \ref{divisors-definition-pullback-effective-Cartier-divisor}）。
设 $(U, A, f, M^\bullet)$ 是 $(X, D, M)$ 的仿射图，
$V = \Spec(B) \subset Y$ 是满足 $h(V) \subset U$ 的仿射开集。
以 $g \in B$ 表示 $f \in A$ 的像。则
% P05-EN-CH38-0114: source 13493--13494 uses “Let ... is” twice instead of “Let ... be”.
\begin{enumerate}
\item $(V, B, g, M^\bullet \otimes_A B)$ 是 $(Y, E, Lh^*M)$ 的仿射图；
\item 在 $B$ 中有 $I_i(M^\bullet, f)B = I_i(M^\bullet \otimes_A B, g)$；
\item 若 $(X, D, M)$ 是良三元组，则 $(Y, E, Lh^*M)$ 是良三元组。
\end{enumerate}
\end{lemma}

\begin{proof}
第一个断言由如下观察得到：$g$ 是 $B$ 中定义
$E \cap V \subset V$ 的非零因子；而 $M^\bullet \otimes_A B$ 表示
$M^\bullet \otimes_A^\mathbf{L} B$，因而由《概形的导出范畴》，
引理 \ref{perfect-lemma-quasi-coherence-pullback}，它表示 $M$ 在 $V$ 上的
拉回。(2) 由 (1) 和 《代数进阶》，引理
\ref{more-algebra-lemma-eta-base-change-pre} 推出。再结合 《代数进阶》，
引理 \ref{more-algebra-lemma-eta-base-change-pre}，可知引理的第二个断言成立。
% P05-EN-CH38-0115: source 13512--13516 cites the same lemma twice, concludes (2) twice, and never proves item (3).
\end{proof}

\begin{lemma}
\label{flat-lemma-good-triple}
设 $X, D, \mathcal{I}, M$ 如 情形
\ref{flat-situation-complex-and-divisor} 中所述。若 $(X, D, M)$ 是良三元组，
则 $L\eta_\mathcal{I}M$ 是 $D(\mathcal{O}_X)$ 的完美对象。
\end{lemma}

\begin{proof}
这是 《代数进阶》，引理
\ref{more-algebra-lemma-eta-locally-free} 的翻译。使用
引理 \ref{flat-lemma-Leta} 完成该翻译。
\end{proof}

\begin{lemma}
\label{flat-lemma-good-triple-bdd-loc-free}
设 $X, D, \mathcal{I}, M$ 如 情形
\ref{flat-situation-complex-and-divisor} 中所述。假设 $(X, D, M)$ 是良三元组。
若存在表示 $M$ 的有限局部自由 $\mathcal{O}_X$-模局部有界复形
$\mathcal{M}^\bullet$，则存在表示 $L\eta_\mathcal{I}M$ 的有限局部自由
$\mathcal{O}_{X'}$-模局部有界复形 $\mathcal{Q}^\bullet$。
% P05-EN-CH38-0116: source 13542 uses the undefined X'; the eta object still lives on X.
\end{lemma}

\begin{proof}
由《上同调》，引理 \ref{cohomology-lemma-Leta}，复形
$\mathcal{Q}^\bullet = \eta_\mathcal{I}\mathcal{M}^\bullet$
表示 $L\eta_\mathcal{I}M$。为核验该复形局部有界且由有限局部自由模组成，
可以在仿射局部工作，此时有界性显然。选取 $(X, D, M)$ 的仿射图
$(U, A, f, M^\bullet)$，使各理想 $I_i(M^\bullet, f)$ 都是主理想，
并使对每个 $i$，$\mathcal{M}^i|_U$ 都有限自由。由 $(X, D, M)$ 是良三元组
的假设，可以作此选择。写成
$N^i = \Gamma(U, \mathcal{M}^i|_U)$，得到有限自由 $A$-模有界复形
$N^\bullet$；它与复形 $M^\bullet$ 表示 $D(A)$ 中的同一对象
（由《概形的导出范畴》，引理
\ref{perfect-lemma-affine-compare-bounded}）。由《代数进阶》，引理
\ref{more-algebra-lemma-ideal-well-defined}，对所有 $i$，
$I_i(N^\bullet, f)$ 都是主理想。因此复形 $\eta_fN^\bullet$ 是有限局部自由
$A$-模有界复形。由 引理 \ref{flat-lemma-eta-stalks}，
$\mathcal{Q}^i|_U$ 是对应于 $\eta_fN^i$ 的拟凝聚
$\mathcal{O}_U$-模，结论得证。
\end{proof}

\begin{lemma}
\label{flat-lemma-complex-and-divisor-eta-pull}
在 情形 \ref{flat-situation-complex-and-divisor} 中，设 $h : Y \to X$
是概形态射，使拉回 $E = h^{-1}D$ 有定义。若 $(X, D, M)$ 是良三元组，则
$$
Lh^*(L\eta_\mathcal{I}M) = L\eta_\mathcal{J}(Lh^*M)
$$
在 $D(\mathcal{O}_Y)$ 中成立，其中 $\mathcal{J}$ 是 $E$ 的理想层。
\end{lemma}

\begin{proof}
这是 《代数进阶》，引理
\ref{more-algebra-lemma-eta-base-change} 的翻译。使用 引理
\ref{flat-lemma-eta-stalks} 和 \ref{flat-lemma-Leta} 完成该翻译。
\end{proof}

\begin{lemma}
\label{flat-lemma-complex-and-divisor-blowup-pre}
在情形 \ref{flat-situation-complex-and-divisor} 中，
存在唯一态射 $b : X' \to X$，使得
\begin{enumerate}
\item 拉回 $D' = b^{-1}D$ 有定义，且 $(X', D', M')$ 是良三元组，
其中 $M' = Lb^*M$；
\item 对任意概形态射 $h : Y \to X$，若拉回 $E = h^{-1}D$ 有定义且
$(Y, E, Lh^*M)$ 是良三元组，则 $h$ 唯一地通过 $b$ 分解。
\end{enumerate}
此外，对任意仿射图 $(U, A, f, M^\bullet)$，限制
$b^{-1}(U) \to U$ 是沿理想 $I_i(M^\bullet, f)$ 之积的爆破；而对任意
拟紧开集 $W \subset X$，限制
$b|_{b^{-1}(W)} : b^{-1}(W) \to W$ 是 $W \setminus D$-容许爆破。
\end{lemma}

\begin{proof}
证明的思路只是：对任意仿射图 $(U, A, f, M^\bullet)$，在局部沿积理想
$I_i(M^\bullet, f)$ 爆破 $X$。《代数进阶》，第 \ref{more-algebra-section-perfect-eta} 节 开头的几个引理表明，这是良定义的。
泛性质 (2) 随即由爆破的泛性质得到。以下给出细节。

\medskip\noindent
令 $U, A, f, M^\bullet$ 是 $(X, D, M)$ 的仿射图。除有限多个以外，
理想 $I_i(M^\bullet, f)$ 都等于 $A$，故可以考虑
$$
I = \prod\nolimits_i I_i(M^\bullet, f)
$$
这是 $A$ 的有限生成理想。以
$$
b_U : U' \to U
$$
表示 $U$ 沿 $I$ 的爆破。由《因子》，引理
\ref{divisors-lemma-blow-up-pullback-effective-Cartier}，
$b_U^{-1}(U \cap D)$ 有定义。回顾 $f^{r_i} \in I_i(M^\bullet, f)$，
故 $b_U$ 是 $(U \setminus D)$-容许爆破。由《因子》，引理
\ref{divisors-lemma-blowing-up-two-ideals}，对每个 $i$，态射 $b_U$
分解为 $U' \to U'_i \to U$，其中 $U'_i \to U$ 是沿
$I_i(M^\bullet, f)$ 的爆破，而 $U' \to U'_i$ 是另一个爆破。
因此 $I_i(M^\bullet, f)$ 在 $U'$ 上的拉回
$I_i(M^\bullet, f)\mathcal{O}_{U'}$ 是可逆理想层；见《因子》，引理
\ref{divisors-lemma-blow-up-pullback-effective-Cartier} 和
\ref{divisors-lemma-blowing-up-gives-effective-Cartier-divisor}。
故 $(U', b^{-1}D, Lb^*M|_U)$ 是良三元组；关于理想 $I_i(-,-)$
在拉回下的行为，见 引理 \ref{flat-lemma-pullback-triple-ideals-good}。
% P05-EN-CH38-0117: source 13638 uses global b inside the local b_U construction before b exists.
最后，我们断言 $b_U : U' \to U$ 具有引理陈述 (2) 中的泛性质。
事实上，设 $h : Y \to U$ 是概形态射，使得拉回
$E = h^{-1}(D \cap U)$ 有定义且 $(Y, E, Lh^*M)$ 是良三元组。
则 $Y$ 由仿射图 $(V, B, g, N^\bullet)$ 覆盖，并且对每个 $i$，
$I_i(N^\bullet, g)$ 都是可逆理想。由于 $g$ 与 $f$ 在 $B$ 中的像
都截出有效 Cartier 除子 $E \cap V$，二者相差一个单位。因此由
《代数进阶》，引理 \ref{more-algebra-lemma-eta-change-unit}，
可以假设 $g$ 是 $f$ 的像。再由《代数进阶》，引理
\ref{more-algebra-lemma-ideal-well-defined}，$I_i(N^\bullet, g)$
作为 $B$-模同构于 $I_i(M^\bullet \otimes_A B, g)$。因此
$I_i(M^\bullet \otimes_A B, g) = I_i(M^\bullet, f)B$
（引理 \ref{flat-lemma-pullback-triple-ideals-good}）是可逆 $B$-模。
故理想 $IB$ 可逆，从而 $I\mathcal{O}_Y$ 可逆。于是由《因子》，引理
\ref{divisors-lemma-universal-property-blowing-up}，得到 $h$ 通过
$b_U$ 的唯一分解。

\medskip\noindent
令 $\mathcal{B}$ 是如下仿射开集 $U \subset X$ 的集合：存在
$(X, D, M)$ 的仿射图 $(U, A, f, M^\bullet)$。则 $\mathcal{B}$
是 $X$ 的拓扑的一组基；细节从略。对 $U \in \mathcal{B}$，有上述构造的
态射 $b_U : U' \to U$，它在 $U$ 上满足该泛性质。若
$U_1 \subset U_2 \subset X$ 都属于 $\mathcal{B}$，则由泛性质，
$b_{U_1} : U'_1 \to U_1$ 典范同构于
$$
b_{U_2}|_{b_{U_2}^{-1}(U_1)} : b_{U_2}^{-1}(U_1) \longrightarrow U_1
$$
换言之，得到 $U_1$ 上同构
$U'_1 \to b_{U_2}^{-1}(U_1)$。这些同构满足上链条件（仍由泛性质），
故由《构造》，引理 \ref{constructions-lemma-relative-glueing}，
得到态射 $b : X' \to X$，其在 $\mathcal{B}$ 中每个 $U$ 上的限制同构于
$U' \to U$。于是态射 $b : X' \to X$ 满足引理陈述中的性质 (1) 和 (2)，
因为这些性质可在局部核验（细节从略）。

\medskip\noindent
还须证明引理的最后一个断言。设 $W \subset X$ 是拟紧开集。选取有限覆盖
$W = U_1 \cup \ldots \cup U_T$，使得对每个 $1 \leq t \leq T$，
存在仿射图 $(U_t, A_t, f_t, M_t^\bullet)$。下文将使用如下事实：
对任意仿射开集 $V = \Spec(B) \subset U_t \cap U_{t'}$，
(a) $f_t$ 与 $f_{t'}$ 在 $B$ 中的像相差一个单位；且
(b) 复形 $M_t^\bullet \otimes_A B$ 与 $M_{t'} \otimes_A B$
定义 $D(B)$ 中同构的对象。对 $i \in \mathbf{Z}$，令
% P05-EN-CH38-0119: source 13691 tensors both chart complexes over undefined A instead of A_t and A_{t'}.
$$
N_i = \max\nolimits_{t = 1, \ldots, T}
\left(\sum\nolimits_{j \geq i} (-1)^{j - i} rk(M_t^j)\right)
$$
则 $N_t - \sum\nolimits_{j \geq i} (-1)^{j - i} rk(M_t^j) \geq 0$，
因而可以考虑理想
% P05-EN-CH38-0118: source 13697 uses N_t although the maximum was defined as N_i.
$$
I_{t, i} =
f_t^{N_i - \sum\nolimits_{j \geq i} (-1)^{j - i} rk(M_t^j)}
I_i(M_t^\bullet, f_t)
\subset A_t
$$
由《代数进阶》，引理 \ref{more-algebra-lemma-eta-change-unit} 和
\ref{more-algebra-lemma-ideal-well-defined}，理想 $I_{t, i}$ 黏合成
拟凝聚有限型理想 $\mathcal{I}_i \subset \mathcal{O}_W$。此外，除有限多个
以外，这些理想都等于 $\mathcal{O}_W$。显然，上述构造的态射
$X' \to X$ 在 $W$ 上限制为沿理想 $\mathcal{I}_i$ 之积的爆破。
证明完成。
\end{proof}

\begin{lemma}
\label{flat-lemma-complex-and-divisor-blowup}
在 情形 \ref{flat-situation-complex-and-divisor} 中，令 $b : X' \to X$
是 引理 \ref{flat-lemma-complex-and-divisor-blowup-pre} 的态射。考虑理想层为
$\mathcal{I}' \subset \mathcal{O}_{X'}$ 的有效 Cartier 除子
$D' = b^{-1}D$。则 $Q = L\eta_{\mathcal{I}'}Lb^*M$ 是
$D(\mathcal{O}_{X'})$ 的完美对象。
\end{lemma}

\begin{proof}
由 引理 \ref{flat-lemma-complex-and-divisor-blowup-pre} 和
\ref{flat-lemma-good-triple} 得到。
\end{proof}

\begin{lemma}
\label{flat-lemma-complex-and-divisor-blowup-base-change}
在 情形 \ref{flat-situation-complex-and-divisor} 中，设 $h : Y \to X$
是概形态射，使拉回 $E = h^{-1}D$ 有定义。令 $b : X' \to X$ 和
$c : Y' \to Y$ 分别是 引理
\ref{flat-lemma-complex-and-divisor-blowup-pre} 对 $D \subset X$ 与 $M$、
以及对 $E \subset Y$ 与 $Lh^*M$ 所构造的态射。则 $Y'$ 是 $Y$ 关于
$b : X' \to X$ 的严格变换（精确表述见证明），并且
$$
L\eta_{\mathcal{J}'}L(h \circ c)^*M = L(Y' \to X')^*Q
$$
其中 $Q = L\eta_{\mathcal{I}'}Lb^*M$ 如
引理 \ref{flat-lemma-complex-and-divisor-blowup} 中所述。特别地，若
$(Y, E, Lh^*M)$ 是良三元组，且 $k : Y \to X'$ 是满足
$h = b \circ k$ 的唯一态射，则
$L\eta_\mathcal{J}Lh^*M = Lk^*Q$。
\end{lemma}

\begin{proof}
记 $E' = c^{-1}E$。则 $(Y', E', L(h \circ c)^*M)$ 是良三元组。
故由 引理 \ref{flat-lemma-complex-and-divisor-blowup-pre} 的泛性质，
存在唯一态射
$$
h' : Y' \longrightarrow X'
$$
使得 $b \circ h' = h \circ c$。特别地，存在态射
$(h', c) : Y' \to X' \times_X Y$。我们断言：给定拟紧开集
$W \subset X$，使 $b^{-1}(W) \to W$ 是爆破，该态射将 $Y'|_W$
等同于 $Y_W$ 关于 $b^{-1}(W) \to W$ 的严格变换。该断言是在 $W$
上的局部问题，因此可以在仿射图上证明。下一段即作此事。

\medskip\noindent
令 $(U, A, f, M^\bullet)$ 是 $(X, D, M)$ 的仿射图。回顾
引理 \ref{flat-lemma-complex-and-divisor-blowup} 的证明，$b : X' \to X$
在 $U$ 上的限制是 $U = \Spec(A)$ 沿理想 $I_i(M^\bullet, f)$ 之积的爆破。
% P05-EN-CH38-0120: source 13767 cites the later consequence lemma, while the local blowup construction is in lemma-complex-and-divisor-blowup-pre.
若 $V = \Spec(B) \subset Y$ 是满足 $h(V) \subset U$ 的任意仿射开集，
则 $(V, B, g, M^\bullet \otimes_A B)$ 是 $(Y, E, Lh^*M)$ 的仿射图，
其中 $g \in B$ 是 $f$ 的像；见
引理 \ref{flat-lemma-pullback-triple-ideals-good}。因此 $c : Y' \to Y$
在 $V$ 上的限制是沿理想 $I_i(M^\bullet, f)B$ 之积的爆破；换言之，
$c : Y' \to Y$ 在 $h^{-1}(U)$ 上是 $h^{-1}(U)$ 沿理想
$\prod I_i(M^\bullet, f) \mathcal{O}_{h^{-1}(U)}$ 的爆破。由于对严格变换
也有相同描述，关于严格变换的断言成立。

\medskip\noindent
由此，等式
$L\eta_{\mathcal{J}'}L(h \circ c)^*M = L(Y' \to X')^*Q$
由 引理 \ref{flat-lemma-complex-and-divisor-eta-pull} 得到。
最后一个断言是其特例（即 $c = \text{id}_Y$ 且 $k = h'$ 的情形）。
\end{proof}

\begin{lemma}
\label{flat-lemma-complex-and-divisor-blowup-good}
在 情形 \ref{flat-situation-complex-and-divisor} 中，令 $W \subset X$
是 $M$ 的上同调层在其上局部自由的最大开子概形。则 引理
\ref{flat-lemma-complex-and-divisor-blowup-pre} 的态射 $b : X' \to X$
在 $W$ 上是同构。
\end{lemma}

\begin{proof}
这是因为，对任意满足 $U \subset W$ 的仿射图
$(U, A, f, M^\bullet)$，由《代数进阶》，引理
\ref{more-algebra-lemma-eta-cohomology-free}，各
$I_i(M^\bullet, f)$ 在局部由 $f$ 的某个幂生成。由于 $f$ 是非零因子，
爆破 $b^{-1}(U) \to U$ 是同构。
\end{proof}

\begin{lemma}
\label{flat-lemma-complex-and-divisor-blowup-good-T}
设 $X, D, \mathcal{I}, M$ 如 情形
\ref{flat-situation-complex-and-divisor} 中所述。若 $(X, D, M)$ 是良三元组，
则存在有限表现的闭浸入
$$
i : T \longrightarrow D
$$
该态射具有如下性质：
\begin{enumerate}
\item $T$ 在概形论意义下包含 $D \cap W$，其中 $W \subset X$
是 $M$ 的上同调层在其上局部自由的最大开集；
\item $Li^*L\eta_\mathcal{I}M$ 的上同调层局部自由；
\item 对任意像满足 $x = i(t) \in W$ 的点 $t \in T$，
$H^i(M)_x$ 作为 $\mathcal{O}_{X, x}$-模的秩，等于
$H^i(Li^*L\eta_\mathcal{I}M)_t$ 作为 $\mathcal{O}_{T, t}$-模的秩。
\end{enumerate}
% P05-EN-CH38-0121: source 13820, 13870, 13898, 13949, and 14056 add five further unhyphenated “scheme theoretically” loci.
\end{lemma}

\begin{proof}
令 $(U, A, f, M^\bullet)$ 是 $(X, D, M)$ 的仿射图，使得对所有
$i \in \mathbf{Z}$，$I_i(M^\bullet, f)$ 都是主理想。将
$T \cap U \subset D \cap U$ 定义为由理想
$$
J(M^\bullet, f) = \sum J_i(M^\bullet, f) \subset A/fA
$$
定义的闭子概形；该理想在《代数进阶》，引理
\ref{more-algebra-lemma-eta-vanishing-beta-plus-pre} 和
\ref{more-algebra-lemma-eta-vanishing-beta-plus} 中研究。按第二个引理的
表述，$T \cap U \to D \cap U$ 由其中研究的环同态
$A/fA \to C$ 给出。由于 $(X, D, M)$ 是良三元组，可以用这种形式的
仿射图覆盖 $X$；由两个引理中的第一个，这些构造可以黏合。因此得到闭子概形
$T \subset D$，它在上述良仿射图上由理想 $J(M^\bullet, f)$ 给出。
性质 (1) 和 (2) 随即由第二个引理得到；细节从略。给读者一个有用的小观察：
由《代数进阶》，引理 \ref{more-algebra-lemma-eta-locally-free}，
$\eta_fM^\bullet$ 是局部自由模复形，故
$Li^*L\eta_\mathcal{I}M|_{T \cap U}$ 由 $C$-模复形
$\eta_fM^\bullet \otimes_A C$ 表示。关于秩的断言 (3) 由《上同调》，
引理 \ref{cohomology-lemma-eta-cohomology-locally-free} 得到。
\end{proof}

\begin{lemma}
\label{flat-lemma-complex-and-divisor-blowup-T}
在 情形 \ref{flat-situation-complex-and-divisor} 中，令 $b : X' \to X$
和 $D'$ 如 引理 \ref{flat-lemma-complex-and-divisor-blowup-pre} 中所述，
并令 $Q = L\eta_{\mathcal{I}'}Lb^*M$ 如
引理 \ref{flat-lemma-complex-and-divisor-blowup} 中所述。令 $W \subset X$
是 $M$ 的上同调模在其上局部自由的最大开集。则存在有限表现的闭浸入
$i : T \to D'$，使得
\begin{enumerate}
\item 在概形论意义下 $D' \cap b^{-1}(W) \subset T$；
\item $Li^*Q$ 具有局部自由上同调层；
\item 对映到 $w \in W$ 的 $t \in T$，$H^i(Li^*Q)_t$ 作为
$\mathcal{O}_{T, t}$-模的秩等于 $H^i(M)_x$ 作为
$\mathcal{O}_{X, x}$-模的秩。
\end{enumerate}
% P05-EN-CH38-0122: source 13872 introduces w but item (3) uses undefined x in both H^i(M)_x and O_{X,x}.
\end{lemma}

\begin{proof}
引理 \ref{flat-lemma-complex-and-divisor-blowup-good} 表明 $b$ 在 $W$ 上
是同构。因此 $b^{-1}(W) \subset X'$ 包含于 $Lb^*M$ 具有局部自由
上同调层的最大开集 $W' \subset X'$。于是本引理的各个断言立即由
引理 \ref{flat-lemma-complex-and-divisor-blowup-good-T} 对
$(X', D', Lb^*M)$ 的应用以及陈述中提到的其他引理得到。
\end{proof}

\begin{lemma}
\label{flat-lemma-complex-and-divisor-blowup-T-ranks}
在 情形 \ref{flat-situation-complex-and-divisor} 中，令 $b : X' \to X$、
$D' \subset X'$ 和 $Q$ 如 引理
\ref{flat-lemma-complex-and-divisor-blowup} 中所述。令
$\rho = (\rho_i)_{i \in \mathbf{Z}}$ 是整数族。令
$W(\rho) \subset X$ 是满足如下性质的最大开子概形：对所有 $i$，
$H^i(M)$ 都是秩 $\rho_i$ 的局部自由模。令 $i : T \to D'$ 如
引理 \ref{flat-lemma-complex-and-divisor-blowup-T} 中所述。则存在既开又闭的
子概形 $T(\rho) \subset T$，它在概形论意义下包含
$D' \cap b^{-1}(W(\rho))$，并且对所有 $i$，
$H^i(Li^*Q|_{T(\rho)})$ 都是秩 $\rho_i$ 的局部自由模。
\end{lemma}

\begin{proof}
令 $T(\rho) \subset T$ 是如下既开又闭的子概形：对所有 $i$，
$H^i(Li^*Q)$ 的秩都是 $\rho_i$。则该断言立即由 引理
\ref{flat-lemma-complex-and-divisor-blowup-T} 中关于上同调模秩的断言得到。
\end{proof}

\begin{lemma}
\label{flat-lemma-complex-and-divisor-blowup-T-represented-bdd-loc-free}
在 情形 \ref{flat-situation-complex-and-divisor} 中，令 $b : X' \to X$、
$D' \subset X'$ 和 $Q$ 如 引理
\ref{flat-lemma-complex-and-divisor-blowup} 中所述。若存在表示 $M$ 的
有限局部自由 $\mathcal{O}_X$-模局部有界复形 $\mathcal{M}^\bullet$，
则存在表示 $Q$ 的有限局部自由 $\mathcal{O}_{X'}$-模局部有界复形
$\mathcal{Q}^\bullet$。
\end{lemma}

\begin{proof}
回顾 $Q = L\eta_{\mathcal{I}'}Lb^*M$，其中 $\mathcal{I}'$ 是有效
Cartier 除子 $D'$ 的理想层。有限局部自由 $\mathcal{O}_{X'}$-模
局部有界复形 $(\mathcal{M}')^\bullet = b^*\mathcal{M}^\bullet$
表示 $Lb^*M$。因此本引理由
引理 \ref{flat-lemma-good-triple-bdd-loc-free} 得到。
\end{proof}

\begin{lemma}
\label{flat-lemma-complex-and-divisor-derived}
设 $X$ 是概形，$D \subset X$ 是有效 Cartier 除子，且
$M \in D(\mathcal{O}_X)$ 是完美对象。令 $W \subset X$ 是上同调层
$H^i(M)$ 在其上局部自由的最大开集。存在固有态射
$b : X' \longrightarrow X$ 和 $D(\mathcal{O}_{X'})$ 中的对象 $Q$，
具有如下性质：
\begin{enumerate}
\item $b : X' \to X$ 在 $X \setminus D$ 上是同构；
\item $b : X' \to X$ 在 $W$ 上是同构；
\item $D' = b^{-1}D$ 是有效 Cartier 除子；
\item $Q = L\eta_{\mathcal{I}'}Lb^*M$，其中 $\mathcal{I}'$ 是 $D'$ 的理想层；
\item $Q$ 是 $D(\mathcal{O}_{X'})$ 的完美对象；
\item 存在有限表现的闭浸入 $i : T \to D'$，使得
\begin{enumerate}
\item 在概形论意义下 $D' \cap b^{-1}(W) \subset T$；
\item $Li^*Q$ 具有有限局部自由上同调层；
\item 对像为 $w \in W$ 的 $t \in T$，$H^i(Li^*Q)_t$ 作为
$\mathcal{O}_{T, t}$-模的秩等于 $H^i(M)_x$ 作为
$\mathcal{O}_{X, x}$-模的秩；
\end{enumerate}
\item 对任意 $(X, D, M)$ 的仿射图 $(U, A, f, M^\bullet)$，
$b$ 在 $U$ 上的限制是 $U = \Spec(A)$ 沿理想
$I = \prod I_i(M^\bullet, f)$ 的爆破；
\item 对任意 $(X', D', Lb^*N)$ 的仿射图 $(V, B, g, N^\bullet)$，
若 $I_i(N^\bullet, g)$ 是主理想，则
% P05-EN-CH38-0123: source 13958 uses Lb^*N in the triple although N^\bullet is the chart complex and the pulled-back object is Lb^*M.
\begin{enumerate}
\item $Q|_V$ 对应于 $\eta_gN^\bullet$；
\item $T \subset V \cap D'$ 对应于 《代数进阶》，引理
\ref{more-algebra-lemma-eta-vanishing-beta-plus} 中研究的理想
$J(N^\bullet, g) = \sum J_i(N^\bullet, g) \subset B/gB$
。
% P05-EN-CH38-0128: source 13962 places all of T inside an arbitrary chart V; the local closed subscheme should be T \cap V.
\end{enumerate}
\item 若 $M$ 可由有限局部自由 $\mathcal{O}_X$-模局部有界复形表示，
则 $Q$ 可由有限局部自由 $\mathcal{O}_{X'}$-模有界复形表示。
% P05-EN-CH38-0124: source 13967--13970 strengthens “locally bounded” to “bounded”, unlike the lemma it collects.
\end{enumerate}
\end{lemma}

\begin{proof}
该断言汇集了 引理 \ref{flat-lemma-pullback-triple-ideals-good}、
\ref{flat-lemma-good-triple},
\ref{flat-lemma-complex-and-divisor-eta-pull},
\ref{flat-lemma-complex-and-divisor-blowup-pre},
\ref{flat-lemma-complex-and-divisor-blowup},
\ref{flat-lemma-complex-and-divisor-blowup-base-change},
\ref{flat-lemma-complex-and-divisor-blowup-good},
\ref{flat-lemma-complex-and-divisor-blowup-good-T},
\ref{flat-lemma-complex-and-divisor-blowup-T} 以及
\ref{flat-lemma-complex-and-divisor-blowup-T-represented-bdd-loc-free}
中得到的信息。
\end{proof}










\section{复形的爆破，III}
\label{flat-section-blowup-complexes-III}

\noindent
本节给出 \cite[Section 18.1]{F} 中 Macpherson 图构造版本的一个
“代数版本”。
% P05-EN-CH38-0127: source 14001 repeats “version” in “an algebra version of the version”.

\medskip\noindent
设 $X$ 是概形，$E$ 是 $D(\mathcal{O}_X)$ 的完美对象。令
$U \subset X$ 是使得 $E|_U$ 具有局部自由上同调层的最大开子概形。

\medskip\noindent
考虑交换图
$$
\xymatrix{
\mathbf{A}^1_X \ar[r] \ar[rd] &
\mathbf{P}^1_X \ar[d]^p &
(\mathbf{P}^1_X)_\infty \ar[l] \ar[ld] \\
& X \ar@/_1em/[ur]_\infty
}
$$
回顾 $\mathbf{A}^1 = D_+(T_0)$ 是 $\mathbf{P}^1$ 的第一个标准仿射开集，
而 $\infty = V_+(T_0)$ 是其补有效 Cartier 除子；上图是这些概形到 $X$
的拉回。注意，$\infty : X \to (\mathbf{P}^1_X)_\infty$ 是同构。于是
$$
(\mathbf{P}^1_X, (\mathbf{P}^1_X)_\infty, Lp^*E)
$$
是第 \ref{flat-section-blowup-complexes-II} 节情形
\ref{flat-situation-complex-and-divisor} 中的三元组。令
$$
b : W \longrightarrow \mathbf{P}^1_X\quad\text{且}\quad
W_\infty = b^{-1}((\mathbf{P}^1_X)_\infty)
$$
是引理 \ref{flat-lemma-complex-and-divisor-blowup-pre} 从该三元组出发构造的
爆破和有效 Cartier 除子。另记
$$
Q = L\eta_{\mathcal{I}} Lb^*M = L\eta_\mathcal{I} L(p \circ b)^*E
$$
为 引理 \ref{flat-lemma-complex-and-divisor-blowup} 中考虑的
$D(\mathcal{O}_W)$ 的完美对象。这里
$\mathcal{I} \subset \mathcal{O}_W$ 是 $W_\infty$ 的理想层。
% P05-EN-CH38-0125: source 14040 uses undefined M in Lb^*M; the triple contains Lp^*E.

\begin{lemma}
\label{flat-lemma-graph-construction}
上述构造具有如下性质：
\begin{enumerate}
\item $b$ 在 $\mathbf{P}^1_U \cup \mathbf{A}^1_X$ 上是同构；
\item $Q$ 在 $\mathbf{A}^1_X$ 上的限制等于 $E$ 的拉回；
\item 存在有限表现的闭浸入 $i : T \to W_\infty$，使得在概形论意义下
$(W_\infty \to X)^{-1}U \subset T$，且 $Li^*Q$ 具有局部自由上同调层；
\item 对像为 $u \in U$ 的 $t \in T$，$H^i(Li^*Q)_t$ 作为
$\mathcal{O}_{T, t}$-模的秩等于 $H^i(M)_u$ 作为
$\mathcal{O}_{U, u}$-模的秩；
\item 若 $E$ 可由有限局部自由 $\mathcal{O}_X$-模局部有界复形表示，
则 $Q$ 可由有限局部自由 $\mathcal{O}_W$-模局部有界复形表示。
\end{enumerate}
% P05-EN-CH38-0126: source 14060 uses undefined M in the rank comparison; the section's object is E.
\end{lemma}

\begin{proof}
这立即由第 \ref{flat-section-blowup-complexes-II} 节中的结果得到；
汇集全部所需内容的断言见
引理 \ref{flat-lemma-complex-and-divisor-derived}。
\end{proof}
